% !TeX root = ../Book1.tex
% 纲要标题：# 第三编　域扩张与 Galois 理论
% 纲要位置：计划纲要.md，第 498 行。

\BookPart{域扩张与 Galois 理论}{b1:part:03}
\BookPartGuide
解一个多项式方程，常常意味着把原来的数域扩大。加入一个根以后，其他根是否也随之出现，新增的根之间有哪些代数关系，哪些变化仍然保持原有运算？这些问题把第二编的不可约性推进为域扩张的结构理论。Galois 理论所研究的自同构，正是保持基域而重新安排这些根的方式。

\ProseParagraphBreak
\chapref{b1:ch:14}至\chapref{b1:ch:16}建立扩张次数、分裂域、嵌入与可分性。扩张次数是向量空间维数，超越次数则衡量代数独立的参数数目；两者不能互相代替。\chapref{b1:ch:17}用这些工具构造和计算有限域，\chapref{b1:ch:18}证明嵌入独立性，并由乘法算子定义迹与范数。

\ProseParagraphBreak
\chapref{b1:ch:19}由嵌入独立性证明不动域的次数公式，建立有限 Galois 对应。\chapref{b1:ch:20}和\chapref{b1:ch:21}随后研究循环扩张、根式可解性、分圆域与尺规作图；\chapref{b1:ch:22}则由有限层上的相容自同构进入无限 Galois 理论。

\begin{center}
\begin{tikzcd}[column sep=large,row sep=large]
\begin{gathered}\text{第14--16章}\\\text{域扩张与可分性}\end{gathered}\arrow[r]\arrow[dr]
&\begin{gathered}\text{第18.1、19.1--2节}\\\text{有限Galois对应}\end{gathered}\arrow[r]\arrow[dr]
&\begin{gathered}\text{第20--21章}\\\text{循环扩张与根式}\end{gathered}\\
&\begin{gathered}\text{第17章}\\\text{有限域与计算}\end{gathered}\arrow[r,dashed]
&\begin{gathered}\text{第22章}\\\text{无限Galois理论}\end{gathered}
\end{tikzcd}
\end{center}

实线表示主要依赖，虚线表示有限域为无限理论提供典型例子。

\ProseParagraphBreak
不同问题所需的结果仍须分别核对：
\begin{itemize}
\item 有限 Galois 对应使用\cref{b1:thm:embedding-extension,b1:thm:normal-embedding-criterion,b1:thm:separable-embeddings,b1:cor:embedding-independence}，据此证明\cref{b1:thm:artin-fixed-field,b1:thm:finite-galois-correspondence}；迹、范数的算术应用可以暂缓。
\item 特征零的根式可解性除有限对应与可解群外，还用范数的共轭元公式、\cref{b1:thm:hilbert90-multiplicative,b1:thm:kummer-cyclic}及必要的单位根知识；一般 Kummer 对应不是其前提。一般五次及更高次方程的不可解性还需证明独立参数的一般方程具有 \(S_n\) 群，并用到 \(n\geq5\) 时 \(S_n\) 不可解；具体多项式的群则可另由\cref{b1:lem:reduction-cycle-type}求出循环型。
\item 尺规作图的二次扩张塔判别不依赖 Hilbert 第 90 定理；正多边形的完整判据另用分圆域和有限 Galois 对应。
\item 无限 Galois 理论使用一般嵌入延拓和有限对应，并在\secref{b1:sec:2201}、\secref{b1:sec:2202}建立相容自同构、拓扑与紧性；有限域用于例子，根式与尺规作图可以暂缓。
\end{itemize}

本编使用前两编的群同态、正规子群、可解群、商环和多项式不可约性，也使用本科线性代数的基、维数、行列式与迹。基域、特征、单位根以及无限自同构群的连续性与闭性，分别决定相应结论的适用范围。

% !TeX root = ../../Book1.tex
% 纲要标题：## 第14章　域扩张、代数性与超越性
% 纲要位置：计划纲要.md，第 500 行。

\chapter{域扩张、代数性与超越性}
\label{b1:ch:14}
\BookChapterNumber{b1:ch:14}

\begin{chapterabstract}
添入代数元与添入超越元会产生不同性质的域扩张。本章用最小多项式和向量空间维数研究代数扩张，证明次数的塔式公式；再以代数独立性、超越基与函数域描述自由参数的数量。
\end{chapterabstract}

\begin{learninggoals}
\begin{itemize}
\item 用最小多项式辨认单扩张，从不可约多项式构造含根的扩张，并用塔式公式计算次数。
\item 区分有限、有限生成与代数扩张，说明有限个代数元为何生成有限扩张。
\item 使用交换引理解释超越次数为何与超越基的选择无关。
\item 从具体函数域的生成元和关系中辨认独立参数及有限代数部分。
\end{itemize}
\end{learninggoals}

从 \(\mathbb Q\) 添入 \(\sqrt2\) 后，\((1+\sqrt2)^{-1}=\sqrt2-1\)，一切四则运算都可以整理成 \(a+b\sqrt2\) 的形式，其中 \(a,b\in\mathbb Q\)。若添入的是一个超越元 \(t\)，其各次幂却不能这样用有限个固定元素表示。两种添元过程的区别，不能仅从“加入一个新符号”看出，而要考察新元满足怎样的多项式关系，以及所得域作为原域上的向量空间有多大。扩张次数和代数独立性将分别量出这些信息。

\ProseParagraphBreak

本章需要域上多项式的带余除法、不可约性、分式域，以及线性代数中的基和维数。\cref{b1:prop:field-poly-division}及\cref{b1:thm:gauss-descent}提供主要的多项式工具。一般超越基的存在性使用 Zorn 引理，有限生成扩张中的参数选择则只用有限次选取与交换。代数闭包的存在性将在下一章证明。

\ProseParagraphBreak
有限 Galois 理论主要依赖\secref{b1:sec:1401}和\secref{b1:sec:1402}的代数扩张与次数计算。\secref{b1:sec:1403}和\secref{b1:sec:1404}的超越基、超越次数与函数域不属于下一章的前提；它们用于后续独立参数的一般方程与代数几何问题。

\ProseParagraphBreak
域扩张与代数独立性的另一种组织方式，可见 Roman 的《Field Theory》\cite[Chapter 2; Sections 4.2--4.3]{Roman2006FieldTheory}。若要补习这里使用的基、维数与线性映射，可读席南华的《基础代数》第二卷\cite[第1--2章]{Xi2018BasicAlgebraII}。对函数域的几何意义感兴趣的读者，还可在读完本章后参阅 Shafarevich 的《Basic Algebraic Geometry》第一卷\cite[Chapter 1, Sections 1.3, 3.1--3.3]{Shafarevich2013BasicAGI}，了解有理函数与函数域的关系；这属于进一步的代数几何学习。

% !TeX root = ../../Book1.tex
% 纲要标题：### 14.1 域扩张及其次数
% 纲要位置：计划纲要.md，第 502 行。

\section{域扩张及其次数}
\label{b1:sec:1401}


% 纲要内容（仅作写作计划，尚非教材正文）：
%
% * 域扩张、素域与特征
% * 有限扩张与塔式公式
% * 单扩张及最小多项式
%

把多项式的一个根添入原来的域，并非只多出一个数：加、减、乘、除会随之产生更多元素。扩张次数用线性代数的语言衡量这种变化。这里的域均为交换域，子域与原域具有同一个单位；基与维数沿用线性代数中的通常约定。

\subsection{域扩张、素域与特征}
\label{b1:subsec:1401:01}

\begin{definition}\label{b1:def:field-extension}
设 \(K,L\) 为域。如果 \(K\subseteq L\)，并且 \(K\) 的加法与乘法都是 \(L\) 中相应运算的限制，则称 \(L\) 为 \(K\) 的一个\term[yu-kuozhang]{域扩张}[field extension]，记为 \(L/K\)。给定域扩张 \(L/K\)，如果子域 \(E\subseteq L\) 满足 \(K\subseteq E\)，则称 \(E\) 为 \(L/K\) 的\term[zhongjian-yu]{中间域}[intermediate field]。给定子集 \(S\subseteq L\)，域 \(L\) 中包含 \(K\cup S\) 的最小子域记为 \(K(S)\)；它等于所有包含 \(K\cup S\) 的子域之交。
\end{definition}
\symidx{\(L/K\)：域扩张}
\symidx{\(K(S)\)：由 \(K\) 与 \(S\) 生成的子域}
记号 \(L/K\) 不表示商集。若原始数据是域同态 \(\iota:K\rightarrow L\)，则它必为单射：若非零 \(a\) 在核中，则 \(\iota(1)=\iota(a)\iota(a^{-1})=0\)，与 \(\iota(1)=1\neq0\) 矛盾。因而可把 \(K\) 与像 \(\iota(K)\) 认同，再使用包含的写法；认同时仍须记住具体嵌入。

\begin{proposition}\label{b1:prop:prime-field}
每个域 \(K\) 的特征为零或一个素数 \(p\)。它包含唯一的最小子域，称为\term[su-yu]{素域}[prime field]；特征零时同构于 \(\mathbb Q\)，特征 \(p\) 时同构于 \(\mathbb F_p=\mathbb Z/p\mathbb Z\)。
\end{proposition}
\begin{proof}
映射 \(\mathbb Z\rightarrow K\)、\(n\mapsto n1_K\) 的核为 \(m\mathbb Z\)，其中取 \(m\ge0\)。若 \(m=1\)，便有 \(1_K=0\)，与域非零矛盾，所以 \(m=0\) 或 \(m>1\)。若 \(m>1\) 是合数，可写成 \(m=ab\)，其中 \(1<a,b<m\)。此时 \(a,b\) 均不在核中，所以 \(a1_K\) 与 \(b1_K\) 都非零，但 \((a1_K)(b1_K)=m1_K=0\)，与域无零因子矛盾。因此 \(m=0\) 或 \(m=p\) 为素数。在后一情形，像同构于域 \(\mathbb F_p\)，每个子域必须包含它。在前一情形，像是 \(\mathbb Z\) 的副本，取其非零元素的逆元便得到 \(\mathbb Q\) 的副本；映射 \(a/b\mapsto(a1_K)(b1_K)^{-1}\) 良定义且单射。任一子域必须包含这些元素，所以所得子域最小，唯一性也随之成立。
\end{proof}
例如 \(\mathbb Q\subseteq\mathbb R\subseteq\mathbb C\) 是域扩张塔；\(\mathbb Z\subseteq\mathbb Q\) 则不是域扩张，因为 \(\mathbb Z\) 不是域。域的特征不会随扩张改变，因为单位及其整数倍均未改变。

\subsection{有限扩张与塔式公式}
\label{b1:subsec:1401:02}

设 \(L/K\) 为域扩张。以 \(L\) 的加法为加法，以域 \(L\) 中的乘法 \(K\times L\rightarrow L\)、\((c,\ell)\mapsto c\ell\) 为标量乘法，\(L\) 成为 \(K\) 向量空间\footnote{“向量空间”也称“线性空间”；参见席南华《基础代数》第二卷\cite[定义1.1, p. 1]{Xi2018BasicAlgebraII}。}。这里 \(c\in K\) 本来也是 \(L\) 的元素，所以标量乘法就是在 \(L\) 中作乘法；分配律、标量乘法结合律以及 \(1\ell=\ell\) 都继承自域 \(L\) 的运算。把这个向量空间的维数称为扩张 \(L/K\) 的\term[kuozhang-cishu]{扩张次数}[degree of an extension]，记为 \([L:K]\)。如果 \([L:K]\) 有限，则称 \(L/K\) 为\term[youxian-kuozhang]{有限扩张}[finite extension]；否则称为无限扩张。这里的有限指维数，不指底集合，例如 \(\mathbb Q(\sqrt2)/\mathbb Q\) 的次数为二，但两个域都是无限集。
\symidx{\([L:K]\)：域扩张的次数}

\ProseParagraphBreak
例如在域扩张塔 \(K\subseteq E\subseteq L\) 中，若 \(E/K\) 的一组基为 \(\{1,u\}\)，\(L/E\) 的一组基为 \(\{1,v,v^2\}\)，则任意 \(z\in L\) 唯一写成 \(e_0+e_1v+e_2v^2\)，其中 \(e_i\in E\)。再将每个 \(e_i\) 唯一写成 \(a_i+b_i u\)，其中 \(a_i,b_i\in K\)，便得到
\[
z=a_0+b_0u+a_1v+b_1uv+a_2v^2+b_2uv^2.
\]
两层坐标都唯一，所以 \(\{1,u,v,uv,v^2,uv^2\}\) 是 \(L/K\) 的一组基：三个 \(E\) 坐标各需两个 \(K\) 坐标，总共需 \(3\cdot2=6\) 个 \(K\) 坐标。下面的公式把这种逐层展开推广到任意两组基。

\ProseParagraphBreak
有限 Galois 理论的次数计算只涉及有限维时的整数公式。一般扩张也允许无限基，此时乘积基的大小用基数乘法表示。\secref{b1:sec:1403}比较无限超越基的大小时，还将使用无限基数的计数与比较。

\begin{theorem}[塔式公式]\label{b1:thm:tower-law}
设 \(K\subseteq E\subseteq L\) 为域扩张塔。若 \((u_i)_{i\in I}\) 为 \(E\) 在 \(K\) 上的一组向量空间基，\((v_j)_{j\in J}\) 为 \(L\) 在 \(E\) 上的一组向量空间基，则 \((u_iv_j)_{(i,j)\in I\times J}\) 为 \(L\) 在 \(K\) 上的一组向量空间基。因此
\[
[L:K]=[L:E]\,[E:K],
\]
无限维时右边理解为基数乘法。特别地，\(L/K\) 有限当且仅当 \(L/E\) 与 \(E/K\) 都有限。
\end{theorem}
\begin{proof}
任意 \(z\in L\) 是有限多个 \(v_j\) 的 \(E\) 线性组合；把这些系数分别写成有限多个 \(u_i\) 的 \(K\) 线性组合，即得 \(z\) 是有限多个 \(u_iv_j\) 的 \(K\) 线性组合。若有限和 \(\sum_{i,j}a_{ij}u_iv_j=0\)，其中 \(a_{ij}\in K\)，先按 \(j\) 收集。\(v_j\) 在 \(E\) 上线性无关，故每个 \(\sum_i a_{ij}u_i=0\)；再用 \(u_i\) 在 \(K\) 上线性无关，得所有 \(a_{ij}=0\)。这证明所列乘积是基。两个非零基数的乘积有限，恰好在二者均有限时发生。
\end{proof}
这一结果称为\term[tashi-gongshi]{塔式公式}[tower law]。扩张次数为一恰好等价于两域相等：若 \([L:K]=1\)，非零向量 \(1\) 本身就是 \(L/K\) 的一组基，因此 \(L=K\cdot1=K\)；反过来，\(K/K\) 以 \(1\) 为基，次数为一。有限扩张中，中间域的次数必整除总次数。因此，素数次数的扩张没有严格处于两端之间的中间域；这个判断不用知道扩张的具体生成元。

\subsection{单扩张及最小多项式}
\label{b1:subsec:1401:03}

\begin{definition}\label{b1:def:algebraic-element}
设 \(L/K\) 为域扩张，\(\alpha\in L\)。如果存在非零多项式 \(f\in K[x]\) 使 \(f(\alpha)=0\)，则称 \(\alpha\) 为 \(K\) 上的\term[daishu-yuan]{代数元}[algebraic element]，也说它在 \(K\) 上代数；否则称其为 \(K\) 上的\term[chaoyue-yuan]{超越元}[transcendental element]，也说它在 \(K\) 上超越。如果一个域扩张由单个元素 \(\alpha\) 生成，即形如 \(K(\alpha)/K\)，则称它为\term[dan-kuozhang]{单扩张}[simple extension]。
\end{definition}

求值同态
\[
\operatorname{ev}_{\alpha}:K[x]\longrightarrow L,\qquad
f\longmapsto f(\alpha)
\]
的像记为 \(K[\alpha]\)，即 \(\alpha\) 的全部 \(K\) 系数多项式值；它的核则收集了 \(\alpha\) 满足的全部多项式关系。代数情形下，这个核非零；由于 \(K[x]\) 是主理想整环，这个关系理想有唯一的首一生成元。
\symidx{\(\operatorname{ev}_{\alpha}\)：多项式的代入同态}
\symidx{\(K[\alpha]\)：元素 \(\alpha\) 的 \(K\) 系数多项式值组成的子环}

\begin{proposition}\label{b1:prop:minimal-polynomial-degree}
设 \(L/K\) 为域扩张，\(\alpha\in L\)。若 \(\alpha\) 在 \(K\) 上代数，则存在唯一首一不可约多项式 \(m_{\alpha,K}\in K[x]\)，使对每个 \(f\in K[x]\) 都有
\[
f(\alpha)=0\quad\Longleftrightarrow\quad m_{\alpha,K}\mid f.
\]
它称为 \(\alpha\) 在 \(K\) 上的\term[zuixiao-duoxiangshi]{最小多项式}[minimal polynomial]。若其次数为 \(d\)，则
\[
K(\alpha)=K[\alpha]\cong K[x]/(m_{\alpha,K}),\qquad
\bigl[K(\alpha):K\bigr]=d,
\]
且 \(1,\alpha,\ldots,\alpha^{d-1}\) 为一组基。若 \(\alpha\) 超越，则代入 \(x\mapsto\alpha\) 给出 \(K(\alpha)\cong K(x)\)，其中 \(K(x)\) 是多项式环 \(K[x]\) 的分式域。
\end{proposition}
\begin{proof}
代数情形下，求关系理想的首一生成元可以从次数最小的关系入手。在零化 \(\alpha\) 的非零多项式中取次数最小者，并除以首项系数，得到首一 \(m\)。对任意 \(f\) 作\cref{b1:prop:field-poly-division}中的带余除法 \(f=qm+r\)、\(\deg r<\deg m\)。若 \(f(\alpha)=0\)，则 \(r(\alpha)=0\)，最小性迫使 \(r=0\)。若 \(m=gh\) 而两因子均为正次数，则域中 \(g(\alpha)h(\alpha)=0\) 迫使一个因子零化 \(\alpha\)，再次矛盾。故 \(m\) 不可约；另一个满足条件的首一多项式与它相互整除，必相等。

代入同态 \(K[x]\rightarrow L\) 的像是 \(K[\alpha]\)，核为 \((m)\)，因此商环与 \(K[\alpha]\) 同构。要说明这个像已是域，目标是把每个非零 \(g(\alpha)\) 的逆也写成 \(\alpha\) 的多项式值。若 \(g(\alpha)\neq0\)，则 \(m\nmid g\)，由不可约性知 \(\gcd(m,g)=1\)。多项式 Euclid 算法给出 \(um+vg=1\)，代入后得 \(v(\alpha)g(\alpha)=1\)，恰好找到了这样的逆。所以 \(K[\alpha]\) 已是域，等于 \(K(\alpha)\)。带余除法给出所列幂的张成性；次数小于 \(d\) 的多项式不能零化 \(\alpha\)，给出线性无关性。

超越情形下代入同态的核为零，故 \(K[\alpha]\cong K[x]\)。取分式就把它延拓为 \(K(x)\rightarrow L\)，像恰是包含 \(K,\alpha\) 的最小子域，亦即 \(K(\alpha)\)。
\end{proof}
\symidx{\(m_{\alpha,K}\)：\(\alpha\) 在 \(K\) 上的最小多项式}
代数情形下，\(K[\alpha]=K(\alpha)\)，每个非零多项式值的逆仍是多项式值，取逆不增加元素。超越情形下，\(K[\alpha]\cong K[x]\) 只是多项式环，而 \(K(\alpha)\cong K(x)\) 允许多项式比；例如 \(1/\alpha\) 不能写成 \(\alpha\) 的多项式，否则清除分母就会得到非零多项式关系。

\ProseParagraphBreak
反过来，设 \(K\) 为域，\(q\in K[x]\) 为首一不可约多项式。要构造一个根，可以把 \(q(x)=0\) 作为新关系，取商环
\[
L=K[x]/(q),\qquad \alpha=x+(q).
\]
由于 \(K[x]\) 是主理想整环，\cref{b1:prop:pid-irreducible-prime,b1:prop:prime-max-quotient}说明 \(L\) 是域。常数映射 \(K\rightarrow L\) 是单射，因为正次数的 \(q\) 不可能整除非零常数。将 \(K\) 与其像识别后，在商环中直接计算
\[
q(\alpha)=q(x)+(q)=(q)=0_L.
\]
理想 \((q)\) 中的多项式被变为零，故 \(x\) 的剩余类确实成为 \(q\) 的根。又有 \(L=K[\alpha]=K(\alpha)\)；由\cref{b1:prop:minimal-polynomial-degree}，\(\alpha\) 在 \(K\) 上的最小多项式就是 \(q\)。这便从给定多项式构造了含有其根的域扩张。

\ProseParagraphBreak

例如在 \(\mathbb F_2[x]\) 中，二次多项式 \(q=x^2+x+1\) 在 \(0\) 与 \(1\) 处的值都为 \(1\)，故不可约。取 \(L=\mathbb F_2[x]/(q)\)、\(\alpha=x+(q)\)。按 \(q\) 带余除法，\(L\) 的四个元素为 \(0,1,\alpha,\alpha+1\)。这里特征为二，所以 \(-1=1\)、\(-\alpha=\alpha\)；从 \(q(\alpha)=0\) 得
\[
\alpha^2=-\alpha-1=\alpha+1,\qquad \alpha(\alpha+1)=1.
\]
所以 \(1,\alpha\) 是 \(L\) 在 \(\mathbb F_2\) 上的一组基，第二个等式也给出 \(\alpha^{-1}=\alpha+1\)。

\ProseParagraphBreak

若根已经位于某个扩张域中，同一结论则用于辨认该扩张。例如 \(x^2-2\) 在 \(\mathbb Q[x]\) 中不可约，因此 \(\mathbb Q(\sqrt2)\) 的元素唯一写成 \(a+b\sqrt2\)。其中非零元素的逆为
\[
(a+b\sqrt2)^{-1}=\frac{a-b\sqrt2}{a^2-2b^2}.
\]
分母若为零，且 \(b\neq0\)，就会得到有理数 \(a/b\) 的平方为二；若 \(b=0\)，则 \(a=0\)。所以对非零元素，所写分母确实非零。最小多项式依赖基域：\(\sqrt2\) 在 \(\mathbb Q(\sqrt2)\) 上的最小多项式已变为 \(x-\sqrt2\)。

\begin{proposition}\label{b1:prop:biquadratic-example}
域 \(L=\mathbb Q(\sqrt2,\sqrt3)\) 在 \(\mathbb Q\) 上的次数为四，基为 \(1,\sqrt2,\sqrt3,\sqrt6\)。元素 \(\theta=\sqrt2+\sqrt3\) 生成 \(L\)，其最小多项式为 \(x^4-10x^2+1\)。
\end{proposition}
\begin{proof}
若 \(\sqrt3=a+b\sqrt2\)、\(a,b\in\mathbb Q\)，平方后比较 \(1,\sqrt2\) 的系数得 \(ab=0\)、\(a^2+2b^2=3\)。\(b=0\) 会使 \(3\) 为有理平方，\(a=0\) 会使 \(3/2\) 为有理平方，均不可能，因为有理平方的素因子指数都为偶数。因此 \(x^2-3\) 在 \(\mathbb Q(\sqrt2)\) 上无根，二次塔的两层次数都是二。由\cref{b1:thm:tower-law}，总次数为四，所列四个乘积构成基。

令 \(\theta=\sqrt2+\sqrt3\)，先计算
\[
(\sqrt2+\sqrt3)(\sqrt3-\sqrt2)=3-2=1.
\]
所以 \(\theta^{-1}=\sqrt3-\sqrt2\)，并且
\[
\sqrt2=\frac{\theta-\theta^{-1}}2,\qquad
\sqrt3=\frac{\theta+\theta^{-1}}2.
\]
两个平方根均属于 \(\mathbb Q(\theta)\)，因而 \(L=\mathbb Q(\theta)\)。再由 \(\theta^2=5+2\sqrt6\) 得 \(P(\theta)=0\)，其中 \(P=x^4-10x^2+1\)。由 \([\mathbb Q(\theta):\mathbb Q]=4\) 及\cref{b1:prop:minimal-polynomial-degree}，最小多项式 \(m_{\theta,\mathbb Q}\) 的次数为四，而且它整除 \(P\)。两者次数相同且都首一，故 \(m_{\theta,\mathbb Q}=P\)。
\end{proof}

% !TeX root = ../../Book1.tex
% 纲要标题：### 14.2 代数扩张
% 纲要位置：计划纲要.md，第 508 行。

\section{代数扩张}
\label{b1:sec:1402}


% 纲要内容（仅作写作计划，尚非教材正文）：
%
% * 代数元构成子域
% * 有限扩张蕴含代数扩张
% * 代数扩张不必有限
%

两个分别满足多项式方程的元素相加、相乘后，是否仍满足多项式方程？直接消元往往使计算膨胀；有限维向量空间提供了更短的解释。本节先把有限生成的代数数据放进一个有限扩张，再由此处理一般代数扩张。

\subsection{代数元构成的子域}
\label{b1:subsec:1402:01}

逐个添入代数元时，下一步的基域已经包含前面添入的元素。扩大基域不会使原来的零化方程失效，最小多项式的次数因而只能下降；这使各步次数可以用最初基域上的次数统一估计。

\begin{lemma}\label{b1:lem:finite-algebraic-generators}
设 \(L/K\) 为域扩张，\(r\geq0\) 为整数，\(\alpha_1,\ldots,\alpha_r\in L\)。若这些元素都在 \(K\) 上代数，则 \(K(\alpha_1,\ldots,\alpha_r)/K\) 有限。更具体地，若 \(d_i=\deg m_{\alpha_i,K}\)，其次数至多为 \(d_1\cdots d_r\)，其中空乘积约定为一。
\end{lemma}
\begin{proof}
令 \(K_0=K\)、\(K_i=K(\alpha_1,\ldots,\alpha_i)\)。对每个 \(1\leq i\leq r\)，多项式 \(m_{\alpha_i,K}\) 也属于 \(K_{i-1}[x]\)，且仍然零化 \(\alpha_i\)。由\cref{b1:prop:minimal-polynomial-degree}，\(m_{\alpha_i,K_{i-1}}\) 整除这个多项式，整除关系在 \(K_{i-1}[x]\) 中理解。因此
\[
[K_i:K_{i-1}]
=\deg m_{\alpha_i,K_{i-1}}
\leq\deg m_{\alpha_i,K}=d_i.
\]
由\cref{b1:thm:tower-law}沿这个有限塔相乘，得所述估计。\(r=0\) 时扩张就是 \(K/K\)，次数为一。
\end{proof}

在域扩张 \(E/K\) 中取 \(a\in E\)。若 \(1,a,a^2,\ldots\) 的每个有限初段都线性无关，\(E\) 的维数就不可能有限。因此有限维会强迫这串幂满足某个有限的非平凡线性关系，而关系的系数正好给出零化 \(a\) 的多项式。

\begin{lemma}\label{b1:lem:finite-dimensional-algebraic}
设 \(E/K\) 为域扩张。若 \([E:K]=n<\infty\)，则每个 \(a\in E\) 都在 \(K\) 上代数，而且 \(\deg m_{a,K}\leq n\)。
\end{lemma}
\begin{proof}
在 \(n\) 维 \(K\) 向量空间 \(E\) 中，\(n+1\) 个向量 \(1,a,\ldots,a^n\) 线性相关。一个不全为零的相关系数列给出次数至多 \(n\) 的非零多项式 \(f\)，且 \(f(a)=0\)。最小多项式整除 \(f\)，故次数不超过 \(n\)。
\end{proof}

\begin{proposition}\label{b1:prop:algebraic-elements-subfield}
设 \(L/K\) 为域扩张。则 \(L\) 中在 \(K\) 上代数的全部元素组成一个包含 \(K\) 的子域。
\end{proposition}
\begin{proof}
每个 \(c\in K\) 被 \(x-c\) 零化。若 \(a,b\) 代数，\cref{b1:lem:finite-algebraic-generators}说明 \(E=K(a,b)\) 有限。\(a-b,ab\) 及 \(a\neq0\) 时的 \(a^{-1}\) 均在 \(E\) 内，由\cref{b1:lem:finite-dimensional-algebraic}也在 \(K\) 上代数。这正是包含 \(K\) 的子域所需的封闭性。
\end{proof}
这个证明给出方程存在性的结构解释：\(a-b,ab\) 以及非零 \(a\) 的逆元都落在同一个有限扩张里，所以各自满足某个非零的 \(K\) 系数多项式方程。证明不必从 \(a,b\) 的方程直接算出这些新方程。

\ProseParagraphBreak
例如所有实代数数构成 \(\mathbb R\) 的子域。对于 \(\theta=\sqrt2+\sqrt3\)，有限扩张 \(\mathbb Q(\sqrt2,\sqrt3)\) 先保证方程存在；若要找到具体方程，再用 \(\theta^2=5+2\sqrt6\) 平方消元，得到 \(\theta^4-10\theta^2+1=0\)。\cref{b1:prop:biquadratic-example}还说明这个四次关系就是最小多项式。

\subsection{代数扩张及其传递性}
\label{b1:subsec:1402:02}

\begin{definition}\label{b1:def:algebraic-extension}
设 \(L/K\) 为域扩张。如果 \(L\) 的每个元素都在 \(K\) 上代数，则称 \(L/K\) 为\term[daishu-kuozhang]{代数扩张}[algebraic extension]；否则称其为超越扩张。如果存在有限集 \(S\subseteq L\) 使 \(L=K(S)\)，则称 \(L/K\) 为\term[youxian-shengcheng-kuozhang]{有限生成扩张}[finitely generated extension]。
\end{definition}
由\cref{b1:lem:finite-dimensional-algebraic}，有限扩张必代数。它也必有限生成：若 \(v_1,\ldots,v_n\) 是 \(L\) 在 \(K\) 上的一组有限线性基，那么包含 \(K\) 及这些基向量的子域，必包含全部 \(K\) 线性组合；这些线性组合就是 \(L\) 的所有元素，所以 \(K(v_1,\ldots,v_n)=L\)。

\ProseParagraphBreak
反过来，若 \(L/K\) 既代数又有限生成，选取有限个域生成元后，每个生成元都代数，\cref{b1:lem:finite-algebraic-generators}就说明 \(L/K\) 有限。因此
\[
\text{有限扩张}
\quad\Longleftrightarrow\quad
\text{有限生成的代数扩张}.
\]
右端的两个条件必须同时成立；本节最后的例子将分别检验只保留其中一个条件时会发生什么。

\begin{theorem}\label{b1:thm:algebraic-tower}
设 \(K\subseteq E\subseteq L\) 为域扩张塔。则 \(L/K\) 代数当且仅当 \(E/K\) 与 \(L/E\) 都代数。
\end{theorem}
\begin{proof}
关键是反向蕴含：\(E/K\) 可能无限，不能直接把整个 \(E\) 当作有限的一层。但零化一个固定元素 \(a\) 的多项式只有有限多个系数，因此可以先把这些系数放进有限中间域 \(E_0\)，再使用 \(K\subseteq E_0\subseteq E_0(a)\) 这个两层有限塔。

若 \(L/K\) 代数，限制到 \(E\) 即知 \(E/K\) 代数；把零化元素的 \(K\) 系数多项式视为 \(E\) 系数多项式，即知 \(L/E\) 代数。

反过来，设 \(E/K\) 与 \(L/E\) 都代数。给定 \(a\in L\)，取零化它的非零多项式 \(f=c_0+c_1x+\cdots+c_dx^d\in E[x]\)，并令 \(E_0=K(c_0,\ldots,c_d)\)。因每个 \(c_i\) 在 \(K\) 上代数，\cref{b1:lem:finite-algebraic-generators}说明 \(E_0/K\) 有限。原来的 \(f\) 已经属于 \(E_0[x]\)，且仍是非零多项式，所以 \(a\) 在 \(E_0\) 上代数，由\cref{b1:prop:minimal-polynomial-degree}知 \(E_0(a)/E_0\) 也有限。由\cref{b1:thm:tower-law}，\(E_0(a)/K\) 有限，再由\cref{b1:lem:finite-dimensional-algebraic}知 \(a\) 在 \(K\) 上代数。任意 \(a\) 均如此，故 \(L/K\) 代数。
\end{proof}
这里的 \(E_0\) 可以随 \(a\) 及所选方程而改变，不要求所有元素所需的系数同时落在一个固定有限子域中。这使论证也适用于 \(E/K\) 无限的情形。

\subsection{无限代数扩张}
\label{b1:subsec:1402:03}

\begin{proposition}\label{b1:prop:infinite-algebraic-root-tower}
对整数 \(n\geq1\)，在 \(\mathbb R\) 中取正实根 \(a_n=2^{1/2^n}\)，令
\[
K_n=\mathbb Q(a_n),\qquad L=\bigcup_{n\geq1}K_n.
\]
则 \(L/\mathbb Q\) 是无限代数扩张，而且不是有限生成扩张。
\end{proposition}
\begin{proof}
因为 \(a_n=a_{n+1}^2\)，这些域递增。任意两个元素落在同一个 \(K_n\) 中，所以其差、积及非零元素的逆也落在该域中；因此 \(L\) 是域。每个 \(a_n\) 满足 \(x^{2^n}-2=0\)，故 \(K_n/\mathbb Q\) 有限；每个 \(L\) 中的元素属于一个这样的有限扩张，因此在 \(\mathbb Q\) 上代数。

另一方面，\(x^{2^n}-2\) 对素数二满足\cref{b1:thm:eisenstein}的条件，故 \([K_n:\mathbb Q]=2^n\)。如果 \([L:\mathbb Q]\) 有限，\cref{b1:thm:tower-law}就会使它至少为每个 \(2^n\)，不可能。

若 \(L/\mathbb Q\) 有限生成，每个生成元又都在 \(\mathbb Q\) 上代数，\cref{b1:lem:finite-algebraic-generators}便会使 \(L/\mathbb Q\) 有限，与刚才的结论矛盾。因此它也不是有限生成扩张。
\end{proof}
这里每个 \(K_n/\mathbb Q\) 都有限，每个 \(a\in L\) 也都落入某个有限层；但层数可以随 \(a\) 改变，全体元素不能同时落入一个固定有限中间域。

\ProseParagraphBreak
单独的有限生成或代数条件都不能推出有限。若 \(t\) 是 \(K\) 上的不定元，\(K(t)/K\) 只需 \(t\) 一个元素生成，却不是代数扩张，也不是有限扩张，因为 \(1,t,t^2,\ldots\) 线性无关。上述根塔 \(L/\mathbb Q\) 则代数，却既非有限也非有限生成。这两个反例分别说明，不能从“有限生成的代数扩张”中去掉代数或有限生成的条件。

\ProseParagraphBreak
有限中间域未必排成一条升链，但它们可以组成有向族。这里“有向”的含义是：任取有限多个有限中间域，都能找到一个仍然有限的中间域共同包含它们。这样，原来分别位于不同有限层的元素就能移到同一层中计算。

\begin{proposition}\label{b1:prop:algebraic-directed-union}
设 \(L/K\) 为代数域扩张。则 \(L\) 是其有限中间扩张的并；任意有限多个这样的中间域都包含于另一个有限中间扩张。
\end{proposition}
\begin{proof}
每个 \(a\in L\) 属于有限扩张 \(K(a)\)。若 \(E_1,\ldots,E_s\) 为有限中间域，分别选取它们在 \(K\) 上的有限基 \(B_1,\ldots,B_s\)，并令 \(E=K(B_1\cup\cdots\cup B_s)\)。合在一起的基向量是有限多个代数元，故由\cref{b1:lem:finite-algebraic-generators}，\(E/K\) 仍有限。对每个 \(i\)，\(E\) 包含 \(K\) 和 \(B_i\)，所以包含 \(B_i\) 的全部 \(K\) 线性组合，也就是整个 \(E_i\)。而所用生成元都在 \(L\) 中，故 \(E\subseteq L\)，这给出共同的有限中间域。
\end{proof}
例如计算分别来自有限中间域 \(E_1,E_2\) 的两个元素的差或积时，先选共同包含这两个域的有限中间域 \(E\)，运算就能在 \(E\) 内完成；非零元素的逆也留在该域中。因此每个只涉及有限多个元素的计算都可移到某个有限层。处理整个扩张的自同构或子域结构时，则必须另外检查各有限层之间的相容性；有限层的结论不能不经说明便推到无限层。

% !TeX root = ../../Book1.tex
% 纲要标题：### 14.3 超越扩张
% 纲要位置：计划纲要.md，第 514 行。

\section{超越扩张}
\label{b1:sec:1403}


% 纲要内容（仅作写作计划，尚非教材正文）：
%
% * 代数独立与超越基
% * 交换引理及超越次数的不变性
% * 有限生成扩张与纯超越扩张
% * 超越次数的塔式公式
%

一个超越元没有单变量多项式关系，几个超越元之间却仍可能有关系。例如 \(t\) 与 \(t^2\) 分别在 \(K\) 上超越，但满足 \(Y-X^2=0\)。要衡量扩张中真正独立的参数数目，必须同时考察有限多个元素的多项式关系。

\ProseParagraphBreak
有限生成扩张可从已给生成元中逐个选取独立参数。一般超越基的存在性则使用 Zorn 引理；比较无限超越基的大小，还需对有限支集作基数计数。

\subsection{代数独立与超越基}
\label{b1:subsec:1403:01}

\begin{definition}\label{b1:def:transcendence-basis}
设 \(L/K\) 为域扩张，\(T\subseteq L\)。如果对任意整数 \(r\geq1\)、任意互异的 \(t_1,\ldots,t_r\in T\) 以及任意 \(f\in K[X_1,\ldots,X_r]\)，等式 \(f(t_1,\ldots,t_r)=0\) 都蕴含 \(f=0\)，则称 \(T\) 在 \(K\) 上\term[daishu-duli]{代数独立}[algebraic independence]。如果 \(T\) 在 \(K\) 上代数独立，并且 \(L/K(T)\) 是代数扩张，则称 \(T\) 为 \(L/K\) 的\term[chaoyue-ji]{超越基}[transcendence basis]。空集视为代数独立，并约定 \(K(\varnothing)=K\)。
\end{definition}

\begin{lemma}\label{b1:lem:independence-adjoining}
设 \(L/K\) 为域扩张，\(T\subseteq L\) 在 \(K\) 上代数独立，且 \(a\in L\setminus T\)。则 \(T\cup\{a\}\) 在 \(K\) 上代数独立，当且仅当 \(a\) 在 \(K(T)\) 上超越。因此，\(L/K\) 的超越基恰是 \(L\) 中按包含关系极大的 \(K\) 上代数独立子集。
\end{lemma}
\begin{proof}
\(K(T)\) 是由 \(T\) 中有限多个元素的多项式比组成的域。若 \(a\) 满足 \(K(T)\) 系数的多项式方程，清除有限多个分母，就得到 \(a\) 与 \(T\) 中有限多个元素之间的非零多项式关系；多项式仍非零，是因为 \(T\) 代数独立。反过来，把一个这样的关系按 \(a\) 的幂收集，至少一个系数是 \(T\) 中元素的非零多项式值，因而非零；这便给出 \(a\) 在 \(K(T)\) 上的代数方程。极大性的等价刻画随即得到。
\end{proof}

对于有限生成扩张，可以直接选择超越基。若 \(L=K(s_1,\ldots,s_n)\)，依次检查这些生成元；只在 \(s_i\) 对已选元素生成的域仍超越时，才把它选入集合 \(T\)。由\cref{b1:lem:independence-adjoining}，\(T\) 始终代数独立。每个未选入的 \(s_i\) 在当时的参数域上代数，因而在最终的 \(K(T)\) 上仍代数；由\cref{b1:lem:finite-algebraic-generators}，它们有限个生成的扩张 \(L/K(T)\) 也代数，所以 \(T\) 是 \(L/K\) 的超越基。一般扩张未必有有限生成集可供逐项检查，下面的存在定理因而采用 Zorn 引理。

\begin{theorem}\label{b1:thm:transcendence-basis}
每个域扩张 \(L/K\) 都有超越基。更一般地，任意在 \(K\) 上代数独立的子集 \(A\subseteq L\) 都可扩充成 \(L/K\) 的超越基。
\end{theorem}
\begin{proof}
考虑包含 \(A\) 的代数独立子集，按包含排序。它包含 \(A\)，因而非空。任一非空链的并仍代数独立：一个多项式关系只涉及有限多个元素，它们已经同属链中的某一个成员。这个并仍包含 \(A\)，故是该链的上界；空链则以 \(A\) 为上界。于是 Zorn 引理（\cref{b1:thm:A-zorn}）给出极大元 \(T\)。由\cref{b1:lem:independence-adjoining}，每个不在 \(T\) 中的元素都在 \(K(T)\) 上代数；\(T\) 中元素本来就在该域中，所以 \(L/K(T)\) 代数。
\end{proof}

\subsection{交换引理及超越次数的不变性}
\label{b1:subsec:1403:02}

\begin{lemma}[交换引理]\label{b1:lem:transcendence-exchange}
设 \(L/F\) 为域扩张，\(x,y\in L\)。若 \(x\) 在 \(F(y)\) 上代数而在 \(F\) 上超越，则 \(y\) 在 \(F(x)\) 上代数。
\end{lemma}
\begin{proof}
把 \(x\) 的一个 \(F(y)\) 系数方程清除分母，得到 \(P(x,y)=0\)，其中 \(P(X,Y)\in F[X,Y]\) 非零。把它按 \(Y\) 展开为 \(\sum_j q_j(X)Y^j\)。由于 \(x\) 在 \(F\) 上超越，每个非零 \(q_j\) 都满足 \(q_j(x)\neq0\)，所以 \(P(x,Y)\) 是 \(F(x)[Y]\) 的非零多项式。它以 \(y\) 为根，故结论成立。
\end{proof}
该结论称为\term[jiaohuan-yinli]{交换引理}[exchange lemma]。它允许用一个新的独立参数替换旧参数，但替换时必须保证新参数在其余已保留参数上超越。

\begin{lemma}\label{b1:lem:finite-transcendence-bound}
设 \(L/K\) 为域扩张，\(b_1,\ldots,b_m\in L\)。若 \(L/K(b_1,\ldots,b_m)\) 代数，则 \(L\) 中任意在 \(K\) 上代数独立的子集至多含 \(m\) 个元素。这里不要求 \(b_1,\ldots,b_m\) 本身代数独立。
\end{lemma}
\begin{proof}
反设 \(a_1,\ldots,a_{m+1}\in L\) 在 \(K\) 上代数独立。设已选取其中的 \(a_1,\ldots,a_r\)，令 \(F=K(a_1,\ldots,a_r)\)，并以 \(D\) 表示尚未替换的 \(b_i\) 构成的集合。归纳假设是 \(\lvert D\rvert\leq m-r\) 且 \(L/E\) 代数，其中 \(E=F(D)\)；\(r=0\) 时这正是题设。令 \(a=a_{r+1}\)。它在 \(F\) 上超越，却在 \(E\) 上代数，因此 \(D\) 非空。在 \(D\) 中取一个包含极小的非空子集 \(C\)，使 \(a\) 在 \(F(C)\) 上代数。任取 \(b\in C\)，若删去 \(b\) 后 \(a\) 仍代数，就与 \(C\) 的极小性矛盾。因此 \(a\) 在 \(F\bigl(C\setminus\{b\}\bigr)\) 上仍超越；这正是交换引理所需的假设。

对底域 \(F\bigl(C\setminus\{b\}\bigr)\) 应用\cref{b1:lem:transcendence-exchange}，得 \(b\) 在 \(F\bigl(C\setminus\{b\},a\bigr)\) 上代数。令
\[
E'=F\bigl(D\setminus\{b\},a\bigr),\qquad
M=F(D,a)=E(a)=E'(b).
\]
由于 \(a\) 在 \(E\) 上代数，\(M/E\) 代数。交换所得的方程在更大的系数域 \(E'\) 上仍有效，故 \(b\) 在 \(E'\) 上代数，即 \(M/E'\) 代数。又因 \(E\subseteq M\subseteq L\)，归纳假设给出 \(L/M\) 代数。旧参数域 \(E\) 与新参数域 \(E'\) 未必互相包含，不能把二者直接排成一条域塔；共同的上层域 \(M\) 给出了如下包含关系：
\[
\begin{tikzcd}[row sep=large, column sep=large]
& L & \\
& M \arrow[u, hook, "\text{代数}"] & \\
E \arrow[ur, hook, "\text{代数}"] && E' \arrow[ul, hook, "\text{代数}"']
\end{tikzcd}
\]
沿右侧的 \(E'\subseteq M\subseteq L\) 应用\cref{b1:thm:algebraic-tower}，得到 \(L/E'\) 代数。这样便用 \(a\) 替换了 \(b\)，同时保持了归纳所需的代数性。每次替换使 \(D\) 少一个元素，故至多 \(m\) 次后它为空。此时 \(L/F\) 代数，但 \(a_{r+1}\) 在 \(F\) 上超越，矛盾。
\end{proof}

\begin{theorem}\label{b1:thm:transcendence-cardinality}
设 \(L/K\) 为域扩张。任意两个 \(L/K\) 的超越基都具有相同基数；这个共同基数称为扩张 \(L/K\) 的\term[chaoyue-cishu]{超越次数}[transcendence degree]，记为 \(\operatorname{trdeg}_K L\)。
\end{theorem}
\begin{proof}
设 \(A,B\) 为超越基。若其中一个有限，不妨设 \(B\) 有限。\cref{b1:lem:finite-transcendence-bound}给出 \(\lvert A\rvert\leq\lvert B\rvert\)，于是 \(A\) 也有限；交换二者再用该引理即得等号。

现在设二者均无限。每个 \(a\in A\) 在 \(K(B)\) 上代数，其一个代数方程只使用 \(B\) 中有限多个元素的有理函数，故可选有限子集 \(B_a\subseteq B\)，使 \(a\) 在 \(K(B_a)\) 上代数。对任意有限集 \(C\subseteq B\)，令
\[
A_C=\{\,a\in A\mid B_a=C\,\},\qquad M_C=K(C)(A_C).
\]
\(A_C\) 的每个元素都在 \(K(C)\) 上代数；由\cref{b1:prop:algebraic-elements-subfield}，\(M_C/K(C)\) 也代数。把\cref{b1:lem:finite-transcendence-bound}用于 \(M_C/K\)，以 \(C\) 为至多 \(\lvert C\rvert\) 个参数，便得到 \(\lvert A_C\rvert\leq\lvert C\rvert\)，因为 \(A_C\subseteq A\) 仍在 \(K\) 上代数独立。

由\cref{b1:prop:finite-support-cardinality}，无限集 \(B\) 的有限子集总数为 \(\lvert B\rvert\)，而每个 \(A_C\) 都有限。因此对应 \(a\mapsto B_a\) 给出 \(\lvert A\rvert\leq\lvert B\rvert\)。交换 \(A,B\) 得反向不等式，再由\cref{b1:prop:cantor-bernstein}得到等号。为各 \(a\) 同时选取有限支集 \(B_a\)，以及这里的无限基数计数，均采用\cref{b1:def:A-choice}中的选择公理；它与选取超越基时使用的选择原则一致。
\end{proof}
\symidx{\(\operatorname{trdeg}_K L\)：扩张的超越次数}
最后一段不能仅用有限次交换代替：无限基的大小是基数问题，有限个参数的局部估计必须结合有限支集的计数，才能得到所需结论。

\subsection{有限生成扩张与纯超越扩张}
\label{b1:subsec:1403:03}

\begin{definition}\label{b1:def:purely-transcendental}
设 \(L/K\) 为域扩张。如果存在 \(K\) 上代数独立的子集 \(T\subseteq L\)，使 \(L=K(T)\)，则称 \(L/K\) 为\term[chun-chaoyue-kuozhang]{纯超越扩张}[purely transcendental extension]。
\end{definition}
对有限集 \(T=\{t_1,\ldots,t_r\}\)，代入给出 \(K(T)\cong K(X_1,\ldots,X_r)\)。一般超越扩张未必纯超越；超越基只保证以参数域 \(K(T)\) 为新基域后，\(L/K(T)\) 是代数扩张。

\begin{proposition}\label{b1:prop:finitely-generated-field-structure}
设 \(L/K\) 为有限生成域扩张。则存在有限超越基 \(T\subseteq L\)，使 \(L/K(T)\) 有限。因而每个有限生成扩张都可以分成一个有限参数的纯超越扩张和其后的有限代数扩张。
\end{proposition}
\begin{proof}
从有限生成集 \(S\) 中按上文的方法逐个选取，得到超越基 \(T\subseteq S\)。每个 \(s\in S\setminus T\) 在 \(K(T)\) 上代数；\(S\setminus T\) 有限，故\cref{b1:lem:finite-algebraic-generators}说明 \(L=K(T)\bigl(S\setminus T\bigr)\) 在 \(K(T)\) 上有限。
\end{proof}

\begin{proposition}[超越次数的塔式公式]\label{b1:prop:transcendence-tower}
设 \(K\subseteq E\subseteq L\) 为域扩张塔。则下式中的加法为基数加法：
\begin{equation}\label{b1:eq:transcendence-tower}
\operatorname{trdeg}_K L=\operatorname{trdeg}_K E+\operatorname{trdeg}_E L.
\end{equation}
\end{proposition}
\begin{proof}
取 \(E/K\) 的超越基 \(T\) 及 \(L/E\) 的超越基 \(U\)。\(T\subseteq E\)，而 \(U\) 的任何元素都不属于 \(E\)，否则它在 \(E\) 上代数；故 \(T\cap U=\varnothing\)。

先证 \(T\cup U\) 在 \(K\) 上代数独立。只涉及 \(T\) 或只涉及 \(U\) 的关系已由各自的独立性排除。一般关系涉及有限多个 \(t_1,\ldots,t_r\in T\) 及 \(u_1,\ldots,u_s\in U\)，按 \(u_j\) 的单项式收集后可写成
\[
\sum_{\beta}P_\beta(t_1,\ldots,t_r)
u_1^{\beta_1}\cdots u_s^{\beta_s}=0,
\qquad P_\beta\in K[X_1,\ldots,X_r],
\]
其中 \(\beta\in\mathbb Z_{\geq0}^s\)，只有有限多个 \(P_\beta\) 非零。各系数 \(P_\beta(t_1,\ldots,t_r)\) 属于 \(E\)；由 \(U\) 在 \(E\) 上代数独立，它们全为零。再由 \(T\) 在 \(K\) 上代数独立，每个 \(P_\beta\) 都是零多项式，故原关系只能是零关系。

又因 \(T\) 是 \(E/K\) 的超越基，\(E/K(T)\) 代数。\(E\) 中每个元素在更大的 \(K(T,U)\) 上仍代数，所以 \(E(U)=K(T,U)(E)\) 在 \(K(T,U)\) 上代数。\(L/E(U)\) 也因 \(U\) 是 \(L/E\) 的超越基而代数；由\cref{b1:thm:algebraic-tower}，\(L/K(T,U)\) 代数。因此 \(T\cup U\) 是 \(L/K\) 的超越基。由于 \(T\cap U=\varnothing\)，其基数为 \(\lvert T\rvert+\lvert U\rvert\)，即得所述等式。
\end{proof}

有限代数扩张的次数沿塔相乘；超越次数沿塔按基数相加。

% !TeX root = ../../Book1.tex
% 纲要标题：### 14.4 函数域与超越次数
% 纲要位置：计划纲要.md，第 520 行。

\section{函数域与超越次数}
\label{b1:sec:1404}


% 纲要内容（仅作写作计划，尚非教材正文）：
%
% * 有理函数域与代数函数域
% * 次数与超越次数的不同角色；剩余有限次数随超越基的选择而变
% * 第三卷维数理论的域论背景
%
% ---
%

从多项式环取分式域，就得到可以进行除法的函数表达式。变量之间若有代数关系，取分式以后关系仍然存在。本节用几个可直接计算的例子说明：扩张次数描述代数方程的复杂度，超越次数则记录独立参数的数目。

\subsection{有理函数域与代数函数域}
\label{b1:subsec:1404:01}

\begin{definition}\label{b1:def:function-field}
设 \(L/k\) 为域扩张。如果 \(L/k\) 有限生成，则称 \(L\) 为 \(k\) 上的\term[hanshu-yu]{函数域}[function field]。本书允许函数域的超越次数为零，因而有限代数扩张也属于函数域。设 \(r\geq0\) 为整数；如果 \(t_1,\ldots,t_r\in L\) 在 \(k\) 上代数独立，则称 \(k(t_1,\ldots,t_r)\) 为 \(k\) 上的\term[youli-hanshu-yu]{有理函数域}[rational function field]；有理函数域的有限代数扩张通常称为\term[daishu-hanshu-yu]{代数函数域}[algebraic function field]。若 \(L\) 中所有在 \(k\) 上代数的元素都属于 \(k\)，则称 \(k\) 在 \(L\) 中相对代数闭。本定义不要求这一条件；它也不同于要求 \(k\) 自身是代数闭域。
\end{definition}
由\cref{b1:prop:finitely-generated-field-structure}，每个函数域都是某个有理函数域的有限扩张。例如 \(\mathbb Q(\sqrt2)/\mathbb Q\) 的超越次数为零，仍符合本书的函数域定义；此时超越基为空，参数域就是 \(\mathbb Q\)。

\ProseParagraphBreak
相对代数闭只比较 \(k\) 与当前的扩域 \(L\)。例如设 \(t\) 在 \(\mathbb Q\) 上超越，则 \(\mathbb Q\) 在 \(\mathbb Q(t)\) 中相对代数闭：其中在 \(\mathbb Q\) 上代数的有理函数只有有理常数，\cref{b1:ex:14-10}给出了这一事实的核验。但 \(\mathbb Q\) 本身不是代数闭域，因为 \(x^2-2\) 没有有理根。前者只问所给扩域是否带来新的代数元，后者要求 \(k[x]\) 中每个非常数多项式都有 \(k\) 中的根。

\ProseParagraphBreak
若整环 \(A\) 包含 \(k\)，且 \(A=k[a_1,\ldots,a_n]\)，其分式域是函数域。反过来，若 \(L=k(a_1,\ldots,a_n)\)，先取 \(L\) 中的子环 \(A=k[a_1,\ldots,a_n]\)。它位于域 \(L\) 中，因而是整环；由\cref{b1:prop:fraction-in-ambient-field}，\(\operatorname{Frac}(A)\) 正是 \(L\) 中包含 \(A\) 的最小子域，即
\[
L=k(a_1,\ldots,a_n)=\operatorname{Frac}\bigl(k[a_1,\ldots,a_n]\bigr).
\]
写 \(k[a_1,\ldots,a_n]\) 时，元素由生成元的 \(k\) 系数多项式给出；写 \(k(a_1,\ldots,a_n)\) 时，则还允许非零分母。二者未必相等：若 \(t\) 超越，\(1/t\in k(t)\)，却不在 \(k[t]\) 中，因为多项式 \(tq(t)\) 的常数项为零，不可能等于 \(1\)。

\ProseParagraphBreak
例如，设 \(k\) 为域、\(t\) 在 \(k\) 上超越，则 \(k[t^2,t^3]\) 的分式域仍为 \(k(t)\)，因为 \(t=t^3/t^2\)。所以不同的环可能有同一个函数域。还可考虑
\[
L=k(t)(u),\qquad u^2=t^3-t.
\]
多项式 \(x^2-(t^3-t)\) 在 \(k(t)[x]\) 中不可约。为证明这一点，假设 \(t^3-t=(A/B)^2\)，其中 \(A,B\in k[t]\) 均非零，并已将分式约成 \(\gcd(A,B)=1\) 的形式。清除分母后有
\[
A^2=(t^3-t)B^2=t(t-1)(t+1)B^2.
\]
在 \(k[t]\) 的唯一分解中，\(t\) 在 \(t^3-t\) 中恰出现一次。因此右边的因子 \(t\) 出现次数为一加上它在 \(B\) 中出现次数的两倍，是奇数；左边的出现次数则是它在 \(A\) 中出现次数的两倍，是偶数，矛盾。所以 \(t^3-t\) 不是 \(k(t)\) 中的平方。二次多项式无根便不可约，故可用不可约多项式的商域构造 \(L\)，并有 \(\bigl[L:k(t)\bigr]=2\)。

\ProseParagraphBreak
这给出一个带有两个生成元 \(t,u\)、一个关系的函数域，但不能将超越次数机械地算作“生成元数减方程数”。方程可以重复或由其他方程推出，生成元也可以冗余；独立参数的数目须由代数独立性判断。这里 \(t\) 超越，而 \(u\) 在 \(k(t)\) 上代数，才是超越次数为一的理由。刚才算出的次数二是相对于参数域 \(k(t)\) 的次数，并未判断是否存在另一参数 \(s\) 使 \(L=k(s)\)。

\subsection{次数与超越次数}
\label{b1:subsec:1404:02}

设 \(t\) 在 \(k\) 上超越，\(n\geq1\)。则 \(t^n\) 仍超越，否则将零化 \(t^n\) 的非零多项式中的变量换成 \(x^n\)，会得到零化 \(t\) 的非零多项式。并且
\begin{equation}\label{b1:eq:rational-power-degree}
\bigl[k(t):k(t^n)\bigr]=n.
\end{equation}
事实上，\(t\) 被 \(x^n-t^n\in k(t^n)[x]\) 零化，所以次数至多为 \(n\)。若 \(1,t,\ldots,t^{n-1}\) 存在非平凡的 \(k(t^n)\) 线性关系，清除分母得到
\[
\sum_{i=0}^{n-1}p_i(t^n)t^i=0,\qquad p_i\in k[x].
\]
清除分母后至少有一个 \(p_i\ne0\)。写 \(p_i(X)=\sum_jc_{ij}X^j\)，则上式变为
\[
\sum_{i=0}^{n-1}\sum_j c_{ij}t^{nj+i}=0.
\]
第 \(i\) 组中的幂指数模 \(n\) 余 \(i\)；在 \(0\le i<n\) 的范围内，不同的 \((i,j)\) 给出不同的指数 \(nj+i\)，所以不同组没有同次幂，不能相互抵消。\(t\) 超越使每个系数 \(c_{ij}=0\)，进而所有 \(p_i=0\)，与关系非平凡矛盾，证明了线性无关性。次数因而恰为 \(n\)，而 \(x^n-t^n\) 是 \(n\) 次首一零化多项式；由\cref{b1:prop:minimal-polynomial-degree}，它正是 \(t\) 在 \(k(t^n)\) 上的最小多项式，在 \(k(t^n)[x]\) 中不可约。

\ProseParagraphBreak
固定同一个纯超越扩张 \(L=k(t)/k\)。由上面的证明，\(\{t\}\)、\(\{t^2\}\) 和 \(\{t^n\}\) 都是它的超越基：这些元素在 \(k\) 上超越，且 \(t\) 在各自生成的参数域上代数。采用超越基 \(\{t\}\) 时，参数域为 \(k(t)=L\)，剩余次数 \([L:k(t)]=1\)；采用 \(\{t^2\}\) 时，参数域为 \(k(t^2)\)，剩余次数 \([L:k(t^2)]=2\)；一般地，对 \(n\ge1\)，采用 \(\{t^n\}\) 时，参数域为 \(k(t^n)\)，剩余次数 \([L:k(t^n)]=n\)。

\ProseParagraphBreak
扩张 \(L/k\) 的超越次数始终为一，剩余有限次数却随参数域而变。特别地，\([L:k(t^2)]=2\) 不能用来断定 \(L/k\) 不是纯超越扩张；\(L=k(t)\) 已经给出了反例。

\ProseParagraphBreak
当 \(\operatorname{char}k=p\) 且 \(n=p\) 时，次数仍为 \(p\)，\(x^p-t^p\) 在 \(k(t^p)[x]\) 中仍是不可约的最小多项式。把系数域扩大到 \(k(t)\) 后，才有分解 \(x^p-t^p=(x-t)^p\)，其中 \(t\) 已经属于新的系数域。不可约性必须注明系数域；这个分解也说明它只有一个不同的根，根的个数与扩张次数须加以区分，\chapref{b1:ch:16}将解释其差别。

\ProseParagraphBreak
上一小节的 \(L=k(t)(u)\) 同样满足 \(\operatorname{trdeg}_k L=1\)，因为 \(L/k(t)\) 代数，而 \(\{t\}\) 仍独立，故 \(\{t\}\) 是 \(L/k\) 的超越基。关系 \(u^2=t^3-t\) 表明 \(u\) 在这个参数域上代数，所以两个生成元只有一个独立参数。

\subsection{维数理论的域论背景}
\label{b1:subsec:1404:03}

把 \(t_1,\ldots,t_r\) 当作自由参数时，多项式环 \(k[t_1,\ldots,t_r]\) 没有参数之间的多项式关系，其分式域的超越次数就是 \(r\)。现在取 \(k[X_1,\ldots,X_n]\) 的素理想 \(\mathfrak p\)，令
\[
A=k[X_1,\ldots,X_n]/\mathfrak p.
\]
由\cref{b1:prop:prime-max-quotient}，素理想条件保证 \(A\) 是整环，因而可以取分式域。非零常数是多项式环的单位，不能属于真理想 \(\mathfrak p\)，所以 \(k\) 自然嵌入 \(A\)。剩余类 \(x_i\) 生成函数域 \(\operatorname{Frac}(A)=k(x_1,\ldots,x_n)\)。从这组生成元中抽出的超越基刻画其中相互独立的参数；其余生成元在这些参数上代数。

\ProseParagraphBreak
例如 \(k[X,Y]/(Y-X^2)\cong k[t]\)，同构由 \(X\mapsto t,Y\mapsto t^2\) 给出：用 \(Y=X^2\) 可把每个剩余类化为 \(X\) 的多项式，而这样的多项式映到零只能本来为零。商环中 \(Y\) 的剩余类已由 \(X\) 的剩余类平方得到，因此 \(Y\) 是这组生成元中的冗余项，可以消去。所得函数域 \(k(t)\) 只有一个独立参数。一般的方程也可能冗余，商环可能不再是整环，所以增加一个方程后，独立参数的变化要根据具体理想判断。

\ProseParagraphBreak
第三卷第5章将建立整扩张与 Noether 正规化，第9章将定义 Krull 维数\footnote{克鲁尔（Wolfgang Krull，1899-1971），德国数学家，发展交换环与理想理论，定义了以其名字命名的维数。}，并在域上有限生成整环的情形证明维数等于函数域的超越次数。这里尚未定义那种环的维数，因而只把超越次数视为已建立的域论不变量。


\begin{chaptersummary}
最小多项式把单个代数元的全部多项式关系压缩成一个首一不可约多项式，其次数等于相应单扩张的维数。域扩张塔中的基可以逐层相乘，得到次数的乘法公式。有限个代数元总落在有限扩张中，因此代数元对域运算封闭，代数性也沿扩张塔传递；无限代数扩张则可看作有限中间扩张的有向并。

\ProseParagraphBreak
超越基保留相互独立的参数，使剩余扩张变为代数扩张。交换引理与有限支集计数保证其基数不随选择改变。有限生成扩张因此分成纯超越部分和有限代数部分；超越次数由原扩张决定，剩余有限次数却可能随参数的选择而改变。

\ProseParagraphBreak
本章主要研究加入给定元素后的次数与代数性。下一章转向同时容纳一个多项式全部根的扩张，讨论这种扩张的构造及其同构意义下的唯一性。
\end{chaptersummary}

\BookExercises

\begin{exercise}\label{b1:ex:14-1}
证明域 \(K\) 的每个自同构都逐点固定其素域。若 \(K\) 有限，证明其元素个数必为 \(p^n\)，其中 \(p=\operatorname{char}K\)、\(n\geq1\)。
\end{exercise}
\begin{solution}
自同构固定 \(1\)，因而固定所有整数倍 \(m1\)。特征 \(p\) 时这些元素已组成素域；特征零时还保持逆元，所以固定所有 \((a1)(b1)^{-1}\)，即固定 \(\mathbb Q\) 的副本。有限域不可能含无限素域 \(\mathbb Q\)，故特征为素数 \(p\)。把 \(K\) 看成 \(\mathbb F_p\) 向量空间，因其底集有限可取有限基，设长度为 \(n\)。每个元素与唯一的 \(n\) 元坐标组对应，每个坐标有 \(p\) 种可能，故元素总数为 \(p^n\)。
\end{solution}

\begin{exercise}\label{b1:ex:14-2}
设 \(a^3=2\)。在 \(\mathbb Q(a)\) 中把 \((1+a)^{-1}\) 写成 \(1,a,a^2\) 的有理线性组合，并求 \(1+a\) 的最小多项式。
\end{exercise}
\begin{solution}
\((1+a)(1-a+a^2)=1+a^3=3\)，故逆元为 \((1-a+a^2)/3\)。若 \(b=1+a\)，则 \((b-1)^3-2=0\)，即 \(b^3-3b^2+3b-3=0\)。\(x^3-2\) 由 Eisenstein 判据不可约，平移 \(x\mapsto x-1\) 是 \(\mathbb Q[x]\) 的自同构，保持不可约性。因此 \(x^3-3x^2+3x-3\) 就是 \(b\) 的最小多项式。
\end{solution}

\begin{exercise}\label{b1:ex:14-3}
证明 \(\left[\mathbb Q(\sqrt2,\sqrt5):\mathbb Q\right]=4\)，并给出一组基。对每个 \(c\in\mathbb Q\)，求 \(\theta_c=\sqrt2+c\sqrt5\) 的最小多项式，并判断它何时生成整个 \(\mathbb Q(\sqrt2,\sqrt5)\)。
\end{exercise}
\begin{solution}
若 \(\sqrt5=a+b\sqrt2\)、\(a,b\in\mathbb Q\)，平方后比较 \(1,\sqrt2\) 的系数得 \(2ab=0\)、\(a^2+2b^2=5\)。\(b=0\) 将使 \(5\) 为有理平方，\(a=0\) 将使 \(5/2\) 为有理平方，二者均由素因数指数的奇偶性排除。所以 \(\sqrt5\notin\mathbb Q(\sqrt2)\)，二次塔的总次数为四，乘积基为 \(1,\sqrt2,\sqrt5,\sqrt{10}\)。

对任意 \(c\in\mathbb Q\)，令 \(\theta_c=\sqrt2+c\sqrt5\)。由 \(\theta_c^2=2+5c^2+2c\sqrt{10}\) 移项再平方，得 \(\bigl(\theta_c^2-(2+5c^2)\bigr)^2=40c^2\)。展开并整理常数项，得到首一关系
\[
P_c(\theta_c)=0,\qquad
P_c(X)=X^4-(4+10c^2)X^2+(2-5c^2)^2.
\]
当 \(c=0\) 时，\(P_0(X)=(X^2-2)^2\)，而 \(\theta_0=\sqrt2\) 的最小多项式是 \(X^2-2\)，它只生成二次子域。若 \(c\neq0\)，则 \(2-5c^2\neq0\)，因为 \(2/5\) 不是有理数的平方。由
\[
\theta_c(\sqrt2-c\sqrt5)=2-5c^2
\]
可知 \(\sqrt2-c\sqrt5=(2-5c^2)\theta_c^{-1}\in\mathbb Q(\theta_c)\)。将它与 \(\theta_c\) 相加、相减，再除以非零的 \(2\) 与 \(2c\)，便分别得到 \(\sqrt2\) 和 \(\sqrt5\)。故 \(\mathbb Q(\theta_c)=\mathbb Q(\sqrt2,\sqrt5)\)，其次数为四。由\cref{b1:prop:minimal-polynomial-degree}，\(\theta_c\) 的最小多项式次数也为四；它整除首一四次零化多项式 \(P_c\)，所以与 \(P_c\) 相等。这里确定最小多项式的依据，还包括 \(\theta_c\) 生成整个四次扩张。
\end{solution}

\begin{exercise}\label{b1:ex:14-4}
设 \(E/K\)、\(F/K\) 是同一扩张域中的有限扩张，次数互素。证明 \(E\cap F=K\)。反向命题是否成立？
\end{exercise}
\begin{solution}
由\cref{b1:thm:tower-law}，\([E\cap F:K]\) 同时整除 \([E:K]\)、\([F:K]\)，故为一。反向不成立：取 \(E=\mathbb Q(\sqrt2)\)、\(F=\mathbb Q(\sqrt3)\)。两者次数都为二；若交不为 \(\mathbb Q\)，二次塔公式迫使 \(E=F\)，但\cref{b1:prop:biquadratic-example}中 \(\sqrt3\notin E\) 的计算排除了这一点。所以交为 \(\mathbb Q\)，次数却不互素。

塔式公式给出了中间域次数必须满足的整除约束，但这些整数约束并不完全决定子域的位置关系。这里互素足以迫使交为基域，交为基域却不能反推出次数互素。
\end{solution}

\begin{exercise}\label{b1:ex:14-5}
设 \([L:K]=n<\infty\)，\(f\in K[x]\) 不可约且次数 \(d\) 与 \(n\) 互素。若 \(d>1\)，证明 \(f\) 在 \(L\) 中没有根；进一步判断：在上述次数互素的假设下，\(f\) 是否仍在 \(L[x]\) 中不可约？
\end{exercise}
\begin{solution}
若 \(a\in L\) 是根，则其最小多项式为 \(f\) 的首一化，故 \(\bigl[K(a):K\bigr]=d\)。塔式公式迫使 \(d\mid n\)，与互素且 \(d>1\) 矛盾。

判断 \(f\) 在 \(L[x]\) 中是否不可约，需要比较其根在 \(K\) 与 \(L\) 上的次数。关键是同一个大域可以按两条域塔计算次数：一条经过 \(L\)，另一条经过 \(K(a)\)。在包含 \(L\) 的扩张中取 \(f\) 的一个根 \(a\)，令 \(m=\bigl[L(a):L\bigr]\leq d\)。塔式公式给出
\[
mn=\bigl[L(a):K\bigr]=\bigl[L(a):K(a)\bigr]\,d.
\]
所以 \(d\mid mn\)，结合 \(\gcd(d,n)=1\) 得 \(d\mid m\)。于是 \(m=d\)，\(a\) 在 \(L\) 上的最小多项式仍有次数 \(d\)，只能等于 \(f\) 的首一化，故 \(f\) 仍不可约。这里添根可直接在 \(L[x]\) 中选 \(f\) 的一个不可约因子并取商域。
\end{solution}

\begin{exercise}\label{b1:ex:14-6}
设域 \(L\) 是两个子域 \(E,F\) 的并。证明 \(L=E\) 或 \(L=F\)。
\end{exercise}
\begin{solution}
若两者都是真子域，则既有 \(a\in E\setminus F\)，也有 \(b\in F\setminus E\)；否则一者包含另一者，并域就是后者。元素 \(a+b\) 属于 \(E\) 时，由减去 \(a\) 得 \(b\in E\)，矛盾；属于 \(F\) 时，由减去 \(b\) 得 \(a\in F\)，也矛盾。故不可能以两个真子域覆盖 \(L\)。
\end{solution}

\begin{exercise}\label{b1:ex:14-7}
设 \(K\subseteq E\subseteq L\)，且 \(E/K\) 为代数扩张。
\begin{enumerate}
\item 仅用代数性的传递性证明：对 \(a\in L\)，它在 \(E\) 上代数当且仅当在 \(K\) 上代数。
\item 利用超越次数的塔式公式证明 \(\operatorname{trdeg}_K L=\operatorname{trdeg}_E L\)。
\end{enumerate}
\end{exercise}
\begin{solution}
1.\ 在 \(K\) 上代数的方程也可视为 \(E\) 系数方程。反过来，若 \(a\) 在 \(E\) 上代数，则 \(E(a)/E\) 代数，结合 \(E/K\) 代数，由\cref{b1:thm:algebraic-tower}知 \(a\) 在 \(K\) 上代数。

2.\ \(E/K\) 为代数扩张，故 \(\operatorname{trdeg}_K E=0\)。将此代入\eqref{b1:eq:transcendence-tower}，得到 \(\operatorname{trdeg}_K L=\operatorname{trdeg}_E L\)。
\end{solution}

\begin{exercise}\label{b1:ex:14-8}
设 \(x,y\) 在 \(k\) 上代数独立。证明 \(x,x+y\) 也代数独立，并给出 \(k(x,y)\) 的另一组不是 \(k\)-线性坐标变换所得的超越基。
\end{exercise}
\begin{solution}
变量替换 \((X,Y)\mapsto(X,X+Y)\) 是多项式环的自同构，逆替换为 \((X,Y)\mapsto(X,Y-X)\)。所以非零多项式 \(P\) 代入 \((x,x+y)\) 后若为零，将产生 \(x,y\) 的非零关系，矛盾。另取 \(x,xy\)；它们生成的域包含 \(y=(xy)/x\)，故等于 \(k(x,y)\)。若二者相关，由 \(x\) 在 \(k\) 上超越及\cref{b1:lem:independence-adjoining}，\(xy\) 在 \(k(x)\) 上代数；于是 \(k(x,y)=k(x,xy)\) 在 \(k(x)\) 上代数，\cref{b1:lem:finite-transcendence-bound}便使该域的超越次数至多为一，与 \(x,y\) 独立矛盾。因此也是超越基。\(k\)-线性坐标变换所得的坐标必须形如 \(ax+by\)、\(a,b\in k\)；若 \(xy=ax+by\)，就会给出非零关系 \(XY-aX-bY=0\)。所以 \(xy\) 不是这样的线性式，所取超越基不由 \(k\)-线性坐标变换得到。
\end{solution}

\begin{exercise}\label{b1:ex:14-9}
设 \(x,y\) 在 \(k\) 上代数独立，令 \(s=x+y,t=xy\)。证明 \(s,t\) 代数独立，并证明 \(\bigl[k(x,y):k(s,t)\bigr]=2\)。
\end{exercise}
\begin{solution}
\(x,y\) 都是 \(X^2-sX+t\) 的根，且 \(y=s-x\)，故总次数至多为二。如果 \(s,t\) 相关，从两元素中可抽出的超越基至多有一个；\(k(x,y)/k(s,t)\) 代数于是迫使其超越次数至多为一，与 \(x,y\) 独立矛盾。因此 \(s,t\) 独立。

交换变量 \(x,y\) 给出 \(k(x,y)\) 的非恒等 \(k\)-自同构，它固定 \(s,t\) 而不固定 \(x\)。固定生成元的自同构会固定它们的一切有理表达式，因而逐点固定 \(k(s,t)\)。若 \(x\) 属于这个域，也必须被固定，矛盾。所以 \(x\notin k(s,t)\)，次数不能为一，故恰为二。这里用的是一个一般原则：属于被逐点固定的子域的元素，必被相应自同构固定。这一论证也适用于特征二。
\end{solution}

\begin{exercise}\label{b1:ex:14-10}
设 \(t\) 在 \(k\) 上超越，\(r\in k(t)\setminus k\)。证明 \(t\) 在 \(k(r)\) 上代数，进而证明 \(k(t)\) 中在 \(k\) 上代数的元素恰为 \(k\)。
\end{exercise}
\begin{solution}
写 \(r=A(t)/B(t)\)，其中 \(B\neq0\)。多项式 \(A(X)-r\cdot B(X)\in k(r)[X]\) 非零：若其系数全为零，取 \(B\) 的非零系数 \(b_i\) 就会得到 \(r=a_i/b_i\in k\)，矛盾。它以 \(t\) 为根，故 \(t\) 在 \(k(r)\) 上代数。若 \(r\) 还在 \(k\) 上代数，代数性的塔传递性会迫使 \(t\) 在 \(k\) 上代数，矛盾。所以一切非常数有理函数都超越，题目结论成立。

用\cref{b1:def:function-field}中的术语，这说明 \(k\) 在 \(k(t)\) 中相对代数闭；这个性质不要求 \(k\) 自身是代数闭域。
\end{solution}

\begin{exercise}\label{b1:ex:14-11}
设 \(t\) 超越，\(a,b,c,d\in k\) 且 \(ad-bc\neq0\)。证明 \(k(t)=k\bigl((at+b)/(ct+d)\bigr)\)，并说明当 \(ad-bc=0\) 且分母非零时发生什么。
\end{exercise}
\begin{solution}
令 \(u=(at+b)/(ct+d)\)。分母非零，因为 \(c,d\) 不同时为零且 \(t\) 超越。由 \((cu-a)t=b-du\) 得 \(t=(b-du)/(cu-a)\)。若 \(cu-a=0\)，则同时 \(b-du=0\)，推出 \(ad=bc\)，与假设矛盾，所以反解有效，两个域相等。若行列式为零且分母非零，两行 \((a,b),(c,d)\) 成比例，所给比值为 \(k\) 中常数，生成域退化为 \(k\)。
\end{solution}

\begin{exercise}\label{b1:ex:14-12}
设 \(T\) 是 \(k\) 上可数无限的代数独立集合，并将其不重复地枚举为 \(T=\{t_1,t_2,\ldots\}\)。证明 \(k(T)=\bigcup_{n\geq1}k(t_1,\ldots,t_n)\) 不是 \(k\) 上有限生成的扩张。
\end{exercise}
\begin{solution}
每个有理表达式只使用有限多个变量，故所给并等于 \(k(T)\)。若该域由有限多个元素生成，把这些元素的表达式所涉及变量合并，便有某个 \(N\) 使全部生成元都在 \(k(t_1,\ldots,t_N)\) 中，于是整个域都包含于该子域。但 \(t_{N+1}\) 在此子域上超越，特别地不在其中，矛盾。

关键是有限个域生成元总共只使用有限多个变量。无限超越基的基数证明也使用了同一有限支集思想：一个代数方程的有限多个系数，只涉及有限多个参数；这里则是有限多个有理表达式的变量范围可以同时限制在一个有限集合内。
\end{solution}

\begin{exercise}\label{b1:ex:14-13}
证明 \(A=k[X,Y]/(XY-1)\) 是整环，且其分式域同构于 \(k(t)\)。求该函数域在 \(k\) 上的超越次数。
\end{exercise}
\begin{solution}
映射 \(k[X,Y]\rightarrow k[t,t^{-1}]\)、\(X\mapsto t,Y\mapsto t^{-1}\) 满射。模 \(XY-1\)，每个混合单项式可以消去一个 \(XY\)，所以每个剩余类可写成 \(\sum_{i\geq0}a_iX^i+\sum_{j\geq1}b_jY^j\) 的有限和。其像为不同幂次的 Laurent 多项式，等于零当且仅当全部系数为零。因此核恰为 \((XY-1)\)，商环同构于 \(k[t,t^{-1}]\)，是整环；其分式域为 \(k(t)\)，超越次数为一。
\end{solution}

\begin{exercise}\label{b1:ex:14-14}
设整环 \(A\) 包含域 \(k\)，并由 \(k\) 与有限多个元素作为环生成，\(S\subseteq A\setminus\{0\}\) 为乘法集。证明 \(A\) 与 \(S^{-1}A\) 的分式域自然同构。由此说明只允许更多分母，不改变这里的超越次数。
\end{exercise}
\begin{solution}
由于 \(A\) 是整环，可把 \(S^{-1}A\) 视为 \(\operatorname{Frac}(A)\) 中由指定分式组成的子环。其分式域因此包含于 \(\operatorname{Frac}(A)\)。反方向，\(A\subseteq S^{-1}A\)，所以 \(A\) 的每个非零元素在 \(\operatorname{Frac}(S^{-1}A)\) 中可逆，\(A\) 的全部分式都属于后者。两个域由此相同，作为 \(k\) 扩张的超越次数当然相同。这个结论比较的是分式域，局部化前后的环未必同构。

分式域相同的证明只用了 \(A\) 是整环和 \(S\) 不含零。\(A\) 在 \(k\) 上有限生成的假设，保证其分式域是本节意义下的函数域，并未用于分式域相等本身。
\end{solution}

% !TeX root = ../../Book1.tex
% 纲要标题：## 第15章　分裂域、代数闭包与正规扩张
% 纲要位置：计划纲要.md，第 528 行。

\chapter{分裂域、代数闭包与\mbox{正规扩张}}
\label{b1:ch:15}
\BookChapterNumber{b1:ch:15}

\begin{chapterabstract}
分裂域把一个多项式的全部根置于同一最小扩张中，代数闭包则容纳所有代数方程的根。本章构造这两类域并证明相应的唯一性，研究嵌入延拓，最后用共轭根刻画正规扩张。
\end{chapterabstract}

\begin{learninggoals}
\begin{itemize}
\item 构造并计算低次多项式的分裂域，区分同构类型唯一与同构映射唯一。
\item 解释代数闭包的存在性构造，以及 Zorn 引理在极大理想和部分嵌入中的作用。
\item 用分裂条件、嵌入条件和根的生成条件判断正规性。
\item 构造正规闭包，把共轭根之间的同构延拓到更大的代数扩张。
\end{itemize}
\end{learninggoals}

多项式 \(X^3-2\) 有实根 \(\sqrt[3]{2}\)，但域 \(\mathbb Q(\sqrt[3]{2})\) 仍没有它的另外两个复根。若只为方程找到一个根，便无法同时研究三个根之间的代数关系；继续添入三次单位根后，才得到包含全部根的最小扩张。不同的添根次序又可能给出外观不同的域。分裂域的构造与唯一性要回答的，正是如何在不依赖添根次序的情况下讨论方程的所有根。

\ProseParagraphBreak

本章使用\cref{b1:prop:minimal-polynomial-degree}的单扩张结构及\cref{b1:thm:algebraic-tower}的代数性传递。分裂域只需有限次添根；代数闭包与一般嵌入延拓使用 Zorn 引理。实数和复数仅用于具体低次算例；分裂域的存在性证明适用于任意特征。

\ProseParagraphBreak
分裂域、嵌入和正规性可结合 Lang 的《Algebra》\cite[Chapter V, \S\S 2--4]{Lang2002}阅读；Roman 的《Field Theory》\cite[Sections 2.7--2.9, 3.1--3.2]{Roman2006FieldTheory}也可用于比较逐次添根与嵌入延拓。比较不同教材的证明时，应特别核对它们处理的是有限扩张，还是任意代数扩张，以及何处需要选择原理。

% !TeX root = ../../Book1.tex
% 纲要标题：### 15.1 分裂域
% 纲要位置：计划纲要.md，第 530 行。

\section{分裂域}
\label{b1:sec:1501}


% 纲要内容（仅作写作计划，尚非教材正文）：
%
% * 多项式的分裂域构造
% * 嵌入延拓与分裂域的唯一性；系数变换与根的像须保持已有关系
% * 具体低次多项式的分裂域
%

添入一个根使某个不可约因子变为一次因子，却未必得到其他根。例如实数 \(\sqrt[3]2\) 不会使 \(x^3-2\) 在实数域中完全分裂。若要同时研究全部根之间的对称性，应选取由全部根生成的域；它避免了任意附加与原方程无关的元素。

\subsection{多项式的分裂域构造}
\label{b1:subsec:1501:01}

\begin{definition}\label{b1:def:splitting-field}
设 \(K\) 为域，\(n\geq1\) 为整数，\(f\in K[x]\) 为 \(n\) 次多项式。如果域扩张 \(L/K\) 中存在 \(a\in K^\times\) 及 \(\alpha_1,\ldots,\alpha_n\in L\)，使
\[
f=a\prod_{i=1}^n(x-\alpha_i),\qquad L=K(\alpha_1,\ldots,\alpha_n),
\]
则称 \(L\) 为 \(f\) 在 \(K\) 上的\term[fenlie-yu]{分裂域}[splitting field]。非零常数多项式的分裂域约定为 \(K\)。
\end{definition}
第一个条件要求全部根都已进入 \(L\)，第二个条件要求 \(L\) 恰由这些根生成。例如 \(X^2-2\) 在 \(\mathbb C\) 中分裂，但 \(\mathbb C\) 不是它在 \(\mathbb Q\) 上的分裂域；两个根只生成 \(\mathbb Q(\sqrt2)\)。根按重数列出，重复列出一个根并不改变生成的域。

\begin{proposition}\label{b1:prop:splitting-field-exists}
设 \(K\) 为域。每个非零多项式 \(f\in K[x]\) 都有分裂域，且分裂域为 \(K\) 的有限扩张。若 \(\deg f=n\geq1\)，则该扩张的次数至多为 \(n!\)。
\end{proposition}
\begin{proof}
若 \(f\) 非常数，取一个不可约因子 \(q\)。域 \(K_1=K[x]/(q)\) 包含元素 \(\alpha_1=x+(q)\)，它满足 \(q(\alpha_1)=0\)，从而也是 \(f\) 的根。这里商环为域可由多项式 Bézout 关系直接验证：非零剩余类的代表元与 \(q\) 互素，故有逆元。由因式定理，在 \(K_1[x]\) 中写 \(f=(x-\alpha_1)f_1\)。

在 \(K_1\) 上重复此构造，为次数 \(n-1\) 的 \(f_1\) 添入一个根。每除去一个一次因子，剩余多项式的次数就减少一，所以至多 \(n\) 步便得到由全部这些根生成、使 \(f\) 分裂的域。每步所选不可约因子的次数不超过当时剩余多项式的次数，故各步扩张次数分别至多为 \(n,n-1,\ldots,1\)。由\cref{b1:thm:tower-law}知总次数至多为 \(n!\)。若某个根已经在当前域内，就可继续除去相应的一次因子而不扩大域；重根只需按其重数继续除去，不要求反复添入同一个元素。
\end{proof}
这是一个存在性构造。实际计算不必机械地逐根相添；根之间已知的关系常能显著减少步骤，也能把 \(n!\) 这个粗上界缩小。

\subsection{嵌入延拓与分裂域的唯一性}
\label{b1:subsec:1501:02}

若将基域元素按 \(\sigma\) 变换，那么关系 \(m(\alpha)=0\) 的系数也随之变换。因此 \(\alpha\) 的像必须满足 \(\sigma(m)(\beta)=0\)，不能仅凭 \(m(\beta)=0\) 选择它。下面的引理说明，这一条件也足以得到单扩张上的延拓。

\begin{lemma}\label{b1:lem:simple-embedding-extension}
设 \(L/K\) 为域扩张，\(\alpha\in L\) 在 \(K\) 上代数，最小多项式为 \(m\in K[x]\)。设 \(\sigma:K\rightarrow K'\) 为域同构，\(\Omega/K'\) 为域扩张，并记 \(\sigma(m)\in K'[x]\) 为对 \(m\) 的系数施行 \(\sigma\) 所得的多项式。若预先选定 \(\sigma(m)\) 在 \(\Omega\) 中的一个根 \(\beta\)，则存在唯一的域嵌入 \(\widetilde\sigma_\beta:K(\alpha)\rightarrow\Omega\)，使
\[
\widetilde\sigma_\beta|_K=\sigma,\qquad \widetilde\sigma_\beta(\alpha)=\beta.
\]
反过来，\(\sigma\) 的每个此类延拓都把 \(\alpha\) 送到 \(\sigma(m)\) 的一个根。因此延拓与 \(\sigma(m)\) 在 \(\Omega\) 中的不同根一一对应。
\end{lemma}
\begin{proof}
若 \(\tau:K(\alpha)\rightarrow\Omega\) 延拓 \(\sigma\)，令 \(\beta=\tau(\alpha)\)。把 \(m(\alpha)=0\) 经 \(\tau\) 映过去即得 \(\sigma(m)(\beta)=0\)。反过来，给定这样的根 \(\beta\)，系数同构把不可约 \(m\) 变为不可约 \(\sigma(m)\)。由\cref{b1:prop:minimal-polynomial-degree}，\(K(\alpha)=K[\alpha]\)，故可定义
\[
\widetilde\sigma_\beta:\sum_i a_i\alpha^i\longmapsto\sum_i\sigma(a_i)\beta^i.
\]
若两个多项式表达式在 \(\alpha\) 处相等，其差 \(h\in K[x]\) 属于代入映射的核 \((m)\)，即 \(h=qm\)。变换系数后有 \(\sigma(h)=\sigma(q)\sigma(m)\)，在 \(\beta\) 处仍为零，所以像与表达式的选择无关。该映射保持加法、乘法和单位，并在域之间单射；其像为 \(K'(\beta)\)。每个元素都是 \(\alpha\) 的多项式，故在已选定 \(\beta\) 后，延拓只能是 \(\widetilde\sigma_\beta\)。
\end{proof}

例如取 \(K=\mathbb Q(\sqrt2)\)、\(\alpha=\sqrt[4]2>0\)，并在 \(\mathbb C\) 中计算。由\cref{b1:thm:eisenstein}，\(X^4-2\) 在 \(\mathbb Q\) 上不可约，故 \([\mathbb Q(\alpha):\mathbb Q]=4\)。又有 \(\sqrt2=\alpha^2\)，塔式公式给出 \([K(\alpha):K]=2\)，所以 \(\alpha\) 在 \(K\) 上的最小多项式为
\[
m=X^2-\sqrt2.
\]
令 \(\sigma:K\rightarrow K\) 把 \(\sqrt2\) 送到 \(-\sqrt2\)，则
\[
\sigma(m)=X^2+\sqrt2,
\qquad \sigma(m)\;\text{的根为}\;\pm i\alpha.
\]
因此 \(\sigma\) 从 \(K(\alpha)\) 到 \(\mathbb C\) 的两个嵌入延拓必须分别把 \(\alpha\) 送到 \(i\alpha\) 与 \(-i\alpha\)。若送到 \(\pm\alpha\)，平方仍为 \(\sqrt2\)，就不满足系数已经变换后的关系。这里定义域为实域 \(K(\alpha)=\mathbb Q(\alpha)\)，目标域为 \(\mathbb C\)，两个延拓的像域均为 \(K(i\alpha)\)，含有非实元素。嵌入的定义域、目标域与像域是不同的对象。

\ProseParagraphBreak
即使基域保持不变，也不能独立指定各生成元的共轭根。例如 \(\alpha\) 与 \(b=\alpha^2\) 在 \(\mathbb Q\) 上分别满足 \(X^4-2\) 与 \(X^2-2\)。若一个 \(\mathbb Q\) 同态固定 \(\alpha\)，便必须把 \(b\) 送到 \(\alpha^2=b\)，不能另选 \(-b\)。逐次延拓时要使用当前中间域上的最小多项式：在 \(\mathbb Q(\alpha)\) 上，\(b\) 的最小多项式已是 \(X-\alpha^2\)，没有另选符号的自由。

\begin{proposition}\label{b1:prop:splitting-field-unique}
设 \(f\in K[x]\) 为非零多项式，\(L/K\) 是 \(f\) 的分裂域，\(\sigma:K\rightarrow K'\) 是域同构，而 \(L'/K'\) 是 \(\sigma(f)\) 的分裂域。则 \(\sigma\) 可以延拓为域同构 \(L\rightarrow L'\)。特别地，同一多项式的两个分裂域在保持基域的同构意义下唯一，但这样的同构映射未必唯一。
\end{proposition}
\begin{proof}
先逐根延拓基域同构以取得嵌入，再利用目标域恰由全部根生成来证明它满射。把 \(L\) 写成 \(K(\alpha_1,\ldots,\alpha_r)\)，其中所列元素是 \(f\) 的全部不同根。假设已在 \(E=K(\alpha_1,\ldots,\alpha_{i-1})\) 上构造嵌入 \(\tau:E\rightarrow L'\)。因为 \(f(\alpha_i)=0\)，\cref{b1:prop:minimal-polynomial-degree}说明 \(\alpha_i\) 在 \(E\) 上的最小多项式整除 \(f\)。所以 \(\tau(m_{\alpha_i,E})\) 整除 \(\sigma(f)\)，后者在 \(L'\) 中分裂。于是前者在 \(L'\) 中有根，由\cref{b1:lem:simple-embedding-extension}可把 \(\tau\) 延拓到 \(E(\alpha_i)\)。有限次延拓给出嵌入 \(L\rightarrow L'\)。

在 \(L[x]\) 中，\(f\) 的一次分解经此嵌入变为 \(\sigma(f)\) 的一次分解。因此 \(\sigma(f)\) 的全部根都在像域中，像域又包含 \(K'\)。\(L'\) 由这些根生成，故像恰为 \(L'\)，嵌入为同构。
\end{proof}
唯一的是同构类型，不是同构本身。例如 \(\mathbb Q(\sqrt2)\) 到自身有保持 \(\mathbb Q\) 的两个同构，分别把 \(\sqrt2\) 送到 \(\sqrt2\) 与 \(-\sqrt2\)。

\subsection{具体低次多项式的分裂域}
\label{b1:subsec:1501:03}

在 \(\mathbb Q\) 上，\(x^2-2\) 的两个根是 \(\pm\sqrt2\)，分裂域为 \(\mathbb Q(\sqrt2)\)，次数二。对 \(x^4-2\)，令 \(a=\sqrt[4]2>0\)，其根为 \(a,-a,ia,-ia\)，故分裂域为 \(\mathbb Q(a,i)\)。Eisenstein 判据给出 \(\bigl[\mathbb Q(a):\mathbb Q\bigr]=4\)，而 \(\mathbb Q(a)\subseteq\mathbb R\) 不含 \(i\)，所以再添 \(i\) 的次数为二，总次数为八。

\ProseParagraphBreak
对 \(x^3-2\)，令 \(a=\sqrt[3]2>0\)，\(\zeta=(-1+\sqrt{-3})/2\)。其三个根为 \(a,a\zeta,a\zeta^2\)，故分裂域为 \(\mathbb Q(a,\zeta)\)。\(x^3-2\) 在 \(\mathbb Q\) 上不可约，而 \(\zeta\) 不实且满足 \(x^2+x+1=0\)，所以
\begin{equation}\label{b1:eq:cubic-splitting-degree}
\begin{aligned}
\bigl[\mathbb Q(a,\zeta):\mathbb Q\bigr]
&=\bigl[\mathbb Q(a,\zeta):\mathbb Q(a)\bigr]\,\bigl[\mathbb Q(a):\mathbb Q\bigr]\\
&=2\cdot3=6.
\end{aligned}
\end{equation}
根 \(a\) 生成的三次扩张并未包含另外两个根；分裂域把它们一起纳入，才为后续置换全部根提供合适的域。添一个根的次数由它的最小多项式控制，分裂域的次数还反映其余根需要补充多少代数信息。

\ProseParagraphBreak
由单扩张嵌入引理，\(a\mapsto a\zeta\) 给出 \(\mathbb Q(a)\) 到 \(\mathbb C\) 的 \(\mathbb Q\) 嵌入，其像为 \(\mathbb Q(a\zeta)\)。这个像含有非实元素，而 \(\mathbb Q(a)\) 是实域，所以两个子域不同，尽管它们保持 \(\mathbb Q\) 地同构。同构描述代数结构相同，并不意味着在共同环境中占据同一个子域。

% !TeX root = ../../Book1.tex
% 纲要标题：### 15.2 代数闭包
% 纲要位置：计划纲要.md，第 536 行。

\section{代数闭包}
\label{b1:sec:1502}


% 纲要内容（仅作写作计划，尚非教材正文）：
%
% * 代数闭包的存在；有限多条根关系的相容性与递增并域
% * 一般嵌入延拓与代数闭包在基域同构意义下的唯一性
% * 选择原理在证明中的位置
%

一个分裂域只处理指定多项式。若要在同一个域中讨论基域上的所有代数方程，可先构造代数扩张，使 \(K[x]\) 中的每个非常数多项式至少有一个根。但得到新域 \(F_1\) 后，\(F_1[x]\) 中还可能出现系数不全属于 \(K\) 的多项式，刚才的构造并未为这些方程安排根。因此对新域继续添根，得到递增的域列 \(F_0=K\subseteq F_1\subseteq F_2\subseteq\cdots\)。利用每个多项式只有有限多个系数，可以证明这些域的并处理了其中全部的多项式方程。

\subsection{代数闭包的存在性}
\label{b1:subsec:1502:01}

\begin{definition}\label{b1:def:algebraic-closure}
设 \(\overline K/K\) 为域扩张。如果 \(\overline K/K\) 是代数扩张，并且 \(\overline K\) 是代数闭域，即 \(\overline K[x]\) 中每个非常数多项式都在 \(\overline K\) 中有根，则称 \(\overline K\) 为 \(K\) 的\term[daishu-bibao]{代数闭包}[algebraic closure]。
\end{definition}
\symidx{\(\overline K\)：\(K\) 的一个代数闭包}

要使域 \(F\) 上的每个非常数多项式都有根，可以为每个首一非常数多项式 \(f\) 设置一个形式符号 \(X_f\)，再在多项式环中施加关系 \(f(X_f)=0\)。若这些关系在一个非零商环中同时成立，\(X_f\) 的像便成为 \(f\) 的根。下面还要使这个商环成为域，并保证所得扩张代数。

\begin{lemma}\label{b1:lem:adjoin-all-roots}
对任意域 \(F\)，存在代数扩张 \(F'/F\)，使 \(F[x]\) 中每个非常数多项式在 \(F'\) 中至少有一个根。
\end{lemma}
\begin{proof}
为每个首一非常数多项式 \(f\in F[x]\) 设置一个不定元 \(X_f\)，令 \(R=F[X_f\mid f]\)。每个环元素只涉及有限多个不定元。令 \(I\) 为全部 \(f(X_f)\) 生成的理想；在 \(R/I\) 中，这些关系便同时成立。

我们先证 \(I\neq R\)。关键是任何有限多个这样的关系都能在某个扩域中同时实现。如果 \(1\in I\)，按生成理想的定义，就有有限表达式
\[
1=\sum_{i=1}^r h_i f_i(X_{f_i}),\qquad h_i\in R,
\]
其中可将重复的 \(f_i\) 合并。由\cref{b1:prop:splitting-field-exists}逐次为 \(f_1,\ldots,f_r\) 添入根，得到某个扩张域 \(E/F\)，在其中同时选定 \(f_i\) 的根 \(a_i\)。把 \(X_{f_i}\) 送到 \(a_i\)，其余不定元送到零，就得到保持 \(F\) 的环同态 \(R\rightarrow E\)。余变量取零只是为了给整个多项式环定义同态，并不要求这次赋值满足它们各自对应的添根关系；上式只需所列的有限条关系成立。上式的右边映为零，左边仍为一，矛盾。因此有限子系统的相容性排除了 \(1\in I\)。

由\cref{b1:prop:maximal-ideal-exists}，真理想 \(I\) 包含于一个极大理想 \(\mathfrak m\)。令 \(F'=R/\mathfrak m\)，则它是域；\(F\cap\mathfrak m=0\)，因为非零常数可逆，所以 \(F\) 嵌入 \(F'\)。记 \(X_f\) 的像为 \(\overline X_f\)，则 \(f(\overline X_f)=0\)，所以每个 \(\overline X_f\) 都在 \(F\) 上代数。任意 \(a\in F'\) 都是有限多个 \(\overline X_{f_1},\ldots,\overline X_{f_s}\) 的多项式表达，因而属于
\[
F(\overline X_{f_1},\ldots,\overline X_{f_s}).
\]
由\cref{b1:lem:finite-algebraic-generators}，这个中间域在 \(F\) 上有限，再由\cref{b1:lem:finite-dimensional-algebraic}知 \(a\) 在 \(F\) 上代数。因此 \(F'/F\) 是代数扩张。任意非常数多项式除以首项系数便成为首一多项式，根不改变，故所有非常数多项式都有根。
\end{proof}

\begin{theorem}\label{b1:thm:algebraic-closure}
每个域 \(K\) 都存在代数闭包。
\end{theorem}
\begin{proof}
从 \(F_0=K\) 出发，反复应用\cref{b1:lem:adjoin-all-roots}构造代数扩张 \(F_{n+1}/F_n\)，使 \(F_n\) 上每个非常数多项式在 \(F_{n+1}\) 中有根。将各步嵌入认同为包含，令 \(\Omega=\bigcup_{n\geq0}F_n\)。任意两个元素都落在某个共同的 \(F_n\) 中，其差、积以及非零元素的逆仍在该域内，因此 \(\Omega\) 是域。反复应用\cref{b1:thm:algebraic-tower}可知每个 \(F_n/K\) 代数。任意 \(a\in\Omega\) 属于某个 \(F_n\)，所以在 \(K\) 上代数；这证明了 \(\Omega/K\) 代数。

若 \(g=\sum_{i=0}^d c_ix^i\in\Omega[x]\) 非常数，每个系数 \(c_i\) 属于某个 \(F_{n_i}\)。系数只有有限多个，取 \(n=\max\{n_0,\ldots,n_d\}\)，并利用域列递增，就有所有 \(c_i\in F_n\)。所以 \(g\in F_n[x]\)，在 \(F_{n+1}\subseteq\Omega\) 中有根。故 \(\Omega\) 代数闭，满足代数闭包的两项要求。
\end{proof}
这项构造不要求基域可数。每一步同时为该域的全部多项式安排不定元；迭代后的有限系数论证直接证明并域代数闭，不需要先判定第一步所得域是否已代数闭。“可数次迭代”也不意味着所得代数闭包是可数集。

\subsection{一般嵌入延拓与代数闭包的唯一性}
\label{b1:subsec:1502:02}

\begin{theorem}[嵌入延拓]\label{b1:thm:embedding-extension}
设 \(L/K\) 为任意代数扩张，\(\Omega\) 为代数闭域，\(\sigma:K\rightarrow\Omega\) 为域嵌入。则存在嵌入 \(\widetilde\sigma:L\rightarrow\Omega\) 延拓 \(\sigma\)。
\end{theorem}
\begin{proof}
只要一个部分嵌入的定义域还没有覆盖 \(L\)，就尝试再添入一个元素并延拓。Zorn 引理给出极大部分嵌入，而代数性将保证其定义域不能遗漏任何元素。

令 \(\mathcal P\) 为所有 \((E,\tau)\) 组成的集合，其中 \(E\) 是 \(L/K\) 的中间域，\(\tau:E\hookrightarrow\Omega\) 是满足 \(\tau|_K=\sigma\) 的域嵌入。规定 \((E,\tau)\le(E',\tau')\) 当且仅当 \(E\subseteq E'\) 且 \(\tau'|_E=\tau\)。初始对 \((K,\sigma)\) 使 \(\mathcal P\) 非空。

对一条非空链 \(\mathcal C\)，将其定义域取并，得到中间域 \(E_{\mathcal C}\)。链中任意两个定义域按包含关系可比，较大域上的映射限制到较小域后就是原映射，所以两个映射在重叠处绝不冲突。于是可定义 \(\tau_{\mathcal C}(a)=\tau(a)\)，其中任选一个包含 \(a\) 的链成员 \((E,\tau)\)，所得值与选择无关。任意有限多个元素都可放进同一个链成员的定义域，在那里验证 \(\tau_{\mathcal C}\) 保持域运算并且单射，就知它是延拓 \(\sigma\) 的嵌入。故 \((E_{\mathcal C},\tau_{\mathcal C})\) 给出上界；空链以 \((K,\sigma)\) 为上界。Zorn 引理（\cref{b1:thm:A-zorn}）给出极大对 \((E,\tau)\)。

若 \(E\neq L\)，取 \(a\in L\setminus E\)。因为 \(L/K\) 代数，\(a\) 满足某个非零的 \(K\) 系数多项式方程。\(K\subseteq E\)，所以同一个方程也是非零的 \(E\) 系数方程，\(a\) 因而在 \(E\) 上代数。其最小多项式经系数同构 \(\tau:E\rightarrow\tau(E)\) 变为 \(\tau(E)\) 上的不可约多项式。\(\Omega\) 代数闭，故该多项式在 \(\Omega\) 中有根。\cref{b1:lem:simple-embedding-extension}便把 \(\tau\) 延拓到 \(E(a)\)，与极大性矛盾。所以 \(E=L\)，得到所需延拓。
\end{proof}

\begin{proposition}\label{b1:prop:algebraic-closure-unique}
若 \(\Omega/K\)、\(\Omega'/K\) 都是代数闭包，则存在保持 \(K\) 的同构 \(\Omega\rightarrow\Omega'\)。
\end{proposition}
\begin{proof}
对代数扩张 \(\Omega/K\) 和嵌入 \(K\hookrightarrow\Omega'\) 应用\cref{b1:thm:embedding-extension}，得到保持 \(K\) 的嵌入 \(\tau:\Omega\rightarrow\Omega'\)。为了证明像域 \(\tau(\Omega)\) 代数闭，任取其中的非常数多项式
\[
g(x)=\sum_{i=0}^d\tau(a_i)x^i,
\qquad a_i\in\Omega,
\]
将系数拉回，得到 \(g^\tau(x)=\sum_{i=0}^d a_ix^i\in\Omega[x]\)。因 \(\tau\) 单射，这个多项式仍非常数。\(\Omega\) 代数闭，故有 \(c\in\Omega\) 使 \(g^\tau(c)=0\)。映回去便有
\[
g(\tau(c))=\tau\bigl(g^\tau(c)\bigr)=0,
\]
所以 \(g\) 在像域中有根。

任意 \(b\in\Omega'\) 在 \(K\) 上代数。由于 \(\tau\) 保持 \(K\)，有 \(K\subseteq\tau(\Omega)\)；原来零化 \(b\) 的非零 \(K\) 系数多项式仍是非零的 \(\tau(\Omega)\) 系数多项式，所以扩大基域以后 \(b\) 仍代数。其在像域上的不可约最小多项式只能为一次，因为像域代数闭。于是 \(b\in\tau(\Omega)\)，所以 \(\tau\) 满射。这里利用了“代数闭域没有非平凡代数扩张”，而非错误地认为无限集上的单射自动为满射。
\end{proof}
同构通常并不唯一；选择不同共轭根会给出不同延拓。后文写“固定一个代数闭包”时，是选定共同环境，后续的嵌入与根都在这个选定的域中讨论。

\subsection{证明中的选择原理}
\label{b1:subsec:1502:03}

\cref{b1:lem:adjoin-all-roots}由真理想取得极大理想，调用了先前用 Zorn 引理证明的极大理想存在性；\cref{b1:thm:embedding-extension}则直接对部分嵌入使用 Zorn 引理。

\ProseParagraphBreak
对链取并时必须检查两件不同的事：定义域之并对域运算封闭；映射在重叠处给出同一个值。只有在偏序要求“延拓”而不只是“定义域更大”时，第二件事才得到保证。极大部分嵌入的论证还依赖 \(L/K\) 代数：若剩余元素超越，不能再以最小多项式及其根作为延拓依据。

\ProseParagraphBreak
本书在这些一般存在性论证中采用通常的选择公理。算法性问题另作处理：抽象地得到一个极大理想或代数闭包，并未给出可计算的元素编码、多项式分解程序或根的选择算法。具体基域上的有效计算，须另加适当表示方法。

% !TeX root = ../../Book1.tex
% 纲要标题：### 15.3 正规扩张
% 纲要位置：计划纲要.md，第 542 行。

\section{正规扩张}
\label{b1:sec:1503}


% 纲要内容（仅作写作计划，尚非教材正文）：
%
% * 分裂条件与嵌入条件；代数扩张到自身的嵌入自动为自同构
% * 正规闭包；有限生成元最小多项式的分裂域与全部嵌入像的生成域
% * 正规性对中间域和塔的行为
%

分裂域包含一个方程的所有根。正规扩张把这一要求推广为：凡是不可约方程已有一个根落入扩张，其余根也必须一并落入。用嵌入来表述时，这意味着改变根的选择不会把整个扩张带到另一个子域。本节的扩张均为代数扩张。

\ProseParagraphBreak
给定代数扩张 \(L/K\)，由\cref{b1:thm:algebraic-closure}可取 \(L\) 的代数闭包 \(\Omega\)。因为 \(\Omega/L\) 与 \(L/K\) 都代数，\cref{b1:thm:algebraic-tower}说明 \(\Omega/K\) 代数；\(\Omega\) 又是代数闭域，所以它同时是包含指定 \(L\) 的 \(K\) 的代数闭包。下文的根与嵌入均在这样选定的 \(\Omega\) 中讨论。

\subsection{分裂条件与嵌入条件}
\label{b1:subsec:1503:01}

\begin{definition}\label{b1:def:normal-extension}
设 \(L/K\) 为代数扩张。如果每个在 \(L\) 中有根的 \(K\) 系数不可约多项式都在 \(L\) 中分裂，则称 \(L/K\) 为\term[zhenggui-kuozhang]{正规扩张}[normal extension]。
\end{definition}
这里的分裂允许重根：只要求多项式成为一次因子的乘积，不要求这些因子互不相同。正规性关心所有根是否留在域内；不同根的个数是下一章可分性所研究的问题。

\begin{proposition}\label{b1:prop:algebraic-self-embedding}
设 \(L/K\) 为代数扩张，\(\sigma:L\rightarrow L\) 为逐点固定 \(K\) 的域嵌入。则 \(\sigma\) 为 \(L\) 的 \(K\) 自同构，无须假设 \(L/K\) 正规。
\end{proposition}
\begin{proof}
即使 \(L/K\) 无限，每个元素仍受一个有限次数的方程约束。给定 \(a\in L\)，令 \(R\) 为 \(m_{a,K}\) 在 \(L\) 中的不同根组成的集合。它有限且包含 \(a\)；\(\sigma\) 固定多项式的系数，故将 \(R\) 单射地映入自身。有限集上的单射必为双射，所以存在 \(b\in R\) 使 \(\sigma(b)=a\)。每个 \(a\in L\) 都有原像，故 \(\sigma\) 满射。
\end{proof}
本命题的目标域已经是 \(L\)，所以只须解决满射性。正规性的嵌入判据还须保证：当目标扩大为代数闭包时，像不会离开 \(L\)。

\begin{theorem}\label{b1:thm:normal-embedding-criterion}
设 \(L/K\) 为代数扩张，\(\Omega\) 为包含 \(L\) 的 \(K\) 的代数闭包。以下条件等价：
\begin{enumerate}
\item \(L/K\) 正规。
\item 每个保持 \(K\) 的嵌入 \(\sigma:L\rightarrow\Omega\) 都满足 \(\sigma(L)=L\)。
\item 存在一族 \(K[x]\) 中的非零多项式，使 \(L\) 在 \(K\) 上由它们在 \(\Omega\) 中的全部根生成。
\end{enumerate}
有限扩张时，这些条件还等价于 \(L\) 是某一个非零 \(K\) 系数多项式的分裂域。
\end{theorem}
\begin{proof}
先设正规。对任意 \(a\in L\)，其最小多项式在 \(L\) 中分裂，\(\sigma(a)\) 也是该多项式的根，所以 \(\sigma(L)\subseteq L\)。将目标限制为 \(L\) 后，\cref{b1:prop:algebraic-self-embedding}给出满射性，故 \(\sigma(L)=L\)。这里适用的是逐元素的有限根集论证，既不要求 \(L/K\) 有限，也没有把一般无限集上的单射当作满射。

设第二个条件成立。若不可约 \(f\in K[x]\) 在 \(L\) 中有根 \(a\)，要证明它分裂，就须把任意另一个根 \(b\in\Omega\) 也纳入 \(L\)。\cref{b1:lem:simple-embedding-extension}给出 \(K(a)\rightarrow K(b)\)、\(a\mapsto b\) 的 \(K\) 同构，再由\cref{b1:thm:embedding-extension}延拓为 \(L\rightarrow\Omega\)。其像按假设为 \(L\)，故 \(b\in L\)。因此全部根属于 \(L\)，证明正规性。

正规时，取 \(L\) 中所有元素在 \(K\) 上的最小多项式；这些非零多项式的全部根都在 \(L\) 中，并在 \(K\) 上生成 \(L\)，得到第三个条件。反过来，若 \(L\) 在 \(K\) 上由一族非零多项式在 \(\Omega\) 中的全部根生成，则每个 \(K\) 嵌入把每个有限根集单射地映入自身，因而置换该根集。生成域遂被映到自身且映满，满足第二个条件。

若 \(L/K\) 有限，取有限生成元 \(a_1,\ldots,a_r\)，令 \(f=\prod_{i=1}^r m_{a_i,K}\)。正规性使各最小多项式都在 \(L\) 中分裂，故全部根生成的域包含于 \(L\)；这个域又包含各 \(a_i\)，因而包含 \(L\)。所以 \(L\) 正是 \(f\) 的分裂域；反方向已由第三个条件得到。
\end{proof}
第三个条件必须取所选多项式的全部根。仅由其中一些根生成的域可以是任意代数扩张，未必正规。例如 \(\mathbb Q(\sqrt2)/\mathbb Q\) 正规；\(\mathbb Q(\sqrt[3]2)/\mathbb Q\) 不正规，因为 \(x^3-2\) 的另两个根不实。此处三次扩张和六次分裂域之间的差别，正是正规性所检测的差别。

\subsection{正规闭包}
\label{b1:subsec:1503:02}

\begin{proposition}\label{b1:prop:normal-closure}
设 \(L/K\) 为代数扩张，\(\Omega\) 为包含 \(L\) 的 \(K\) 的代数闭包。则在所选 \(\Omega\) 中，按子域的包含关系存在包含 \(L\) 的最小正规扩张 \(N/K\)。它由所有 \(a\in L\) 的最小多项式在 \(\Omega\) 中的全部根生成，称为 \(L/K\) 在 \(\Omega\) 中的\term[zhenggui-bibao]{正规闭包}[normal closure]。它也有描述
\[
N=K\left(\bigcup_{\sigma\in\operatorname{Hom}_K(L,\Omega)}\sigma(L)\right),
\]
其中 \(\operatorname{Hom}_K(L,\Omega)\) 表示逐点固定 \(K\) 的嵌入 \(L\rightarrow\Omega\) 全体。若 \(L/K\) 有限且 \(L=K(a_1,\ldots,a_r)\)，则 \(N\) 正是 \(\prod_{i=1}^r m_{a_i,K}\) 在 \(\Omega\) 中的分裂域，特别地 \(N/K\) 有限。
\end{proposition}
\begin{proof}
按所述根集生成 \(N\)，\cref{b1:thm:normal-embedding-criterion}的第三个条件直接保证 \(N/K\) 正规，且每个 \(a\in L\) 都在生成根集中，故 \(L\subseteq N\)。若正规扩张 \(M/K\)、\(L\subseteq M\subseteq\Omega\)，则每个 \(a\) 的最小多项式在 \(M\) 中分裂，全部生成根都属于 \(M\)，故 \(N\subseteq M\)。

记所有像域生成的域为 \(M_0\)。每个 \(K\) 嵌入将 \(a\in L\) 送到 \(m_{a,K}\) 的一个根，故 \(\sigma(L)\subseteq N\)，从而 \(M_0\subseteq N\)。反过来，任取 \(m_{a,K}\) 在 \(\Omega\) 中的根 \(\beta\)。\cref{b1:lem:simple-embedding-extension}给出 \(a\mapsto\beta\) 的 \(K(a)\rightarrow\Omega\) 嵌入，再由\cref{b1:thm:embedding-extension}延拓到 \(L\)。于是 \(\beta\) 属于某个像域，故 \(N\) 的全部生成根都属于 \(M_0\)。这证明 \(N=M_0\)：正规闭包把 \(L\) 在共同代数闭包中的所有 \(K\) 同构像一起纳入。

若 \(L=K(a_1,\ldots,a_r)\) 为有限扩张，取 \(m_{a_i,K}\) 的乘积在 \(\Omega\) 内的分裂域 \(N_0\)。它有限、正规并包含 \(L\)，故 \(N\subseteq N_0\)。另一方面，\(N/K\) 正规且包含各 \(a_i\)，所以包含各 \(m_{a_i,K}\) 的全部根，故 \(N_0\subseteq N\)。因此 \(N=N_0\) 为有限扩张；这个具体描述不依赖生成集的选择。
\end{proof}
例如 \(\mathbb Q(\sqrt[3]2)\) 的正规闭包就是 \(\mathbb Q(\sqrt[3]2,\zeta)\)，其中 \(\zeta^2+\zeta+1=0\)、\(\zeta\neq1\)。这里先固定共同代数闭包，才可以谈“最小子域”；若不固定环境，结论应改写成保持基域和已给扩张的同构意义下唯一。

\subsection{正规性与中间域、扩张塔}
\label{b1:subsec:1503:03}

\begin{proposition}\label{b1:prop:normal-tower-behavior}
设 \(K\subseteq E\subseteq L\) 为代数扩张塔。若 \(L/K\) 正规，则 \(L/E\) 正规。在同一代数闭包内，任意一族正规 \(K\) 扩张的交以及它们生成的合成域都对 \(K\) 正规。但 \(L/K\) 正规不保证 \(E/K\) 正规，\(E/K\) 与 \(L/E\) 都正规也不保证 \(L/K\) 正规。
\end{proposition}
\begin{proof}
每个 \(E\) 嵌入 \(L\rightarrow\Omega\) 也是 \(K\) 嵌入，故由\cref{b1:thm:normal-embedding-criterion}保持 \(L\)，从而 \(L/E\) 正规。

若 \(a\) 属于一族正规域之交，其在 \(K\) 上的最小多项式在每个域中分裂，所以全部根属于交域，证明交域正规。对于合成域，每个被合成的域由某族 \(K\) 系数多项式的根生成；把各族合并即可满足\cref{b1:thm:normal-embedding-criterion}的第三个条件。

第一个反例取 \(L=\mathbb Q(\sqrt[3]2,\zeta)\)、\(E=\mathbb Q(\sqrt[3]2)\)、\(K=\mathbb Q\)，其正规与非正规已在本节说明。第二个取
\[
K=\mathbb Q,\quad E=\mathbb Q(\sqrt2),\quad L=\mathbb Q(\sqrt[4]2).
\]
\(E\) 是 \(x^2-2\) 的分裂域，而 \(L/E\) 是 \(x^2-\sqrt2\) 的分裂域，故两层正规。但 \(x^4-2\) 在 \(\mathbb Q\) 上不可约，其非实根不在实域 \(L\) 中，故总扩张不正规。
\end{proof}
交和合成是在同一环境中组合若干扩张，域塔则是在一个扩张内改变基域；前两种运算保持正规性，并不意味着正规性沿塔传递。将来用子群描述中间域时，这些不同的行为会分别对应不同的群论条件。

% !TeX root = ../../Book1.tex
% 纲要标题：### 15.4 嵌入的延拓与共轭
% 纲要位置：计划纲要.md，第 548 行。

\section{嵌入的延拓与共轭}
\label{b1:sec:1504}


% 纲要内容（仅作写作计划，尚非教材正文）：
%
% * 中间域嵌入的延拓与自同构；代数闭包上的自同构延拓
% * 多项式根的共轭性
% * 为自同构计数准备工具
%
% ---
%

比较两个根时，我们通常不只关心它们满足同一个方程，还关心替换这个根能否同时替换域中的其他元素。嵌入延拓定理保证代数关系可相容地延续；正规性决定延拓后的像是否仍是原域。本节把这两种作用分清，并整理随后计数所需的形式。

\subsection{中间域嵌入的延拓与自同构}
\label{b1:subsec:1504:01}

\begin{corollary}\label{b1:cor:intermediate-embedding-extension}
设 \(K\subseteq E\subseteq L\) 为域扩张塔，并且 \(L/K\) 代数；设 \(\Omega\) 为包含 \(L\) 的 \(K\) 的代数闭包。每个 \(K\) 嵌入 \(\tau:E\rightarrow\Omega\) 均能延拓为 \(K\) 嵌入 \(L\rightarrow\Omega\)。若 \(L/K\) 正规，将该延拓的目标限制为 \(L\) 后，它为 \(L\) 的 \(K\) 自同构。此外，\(\tau\) 总能延拓为 \(\Omega\) 的 \(K\) 自同构。
\end{corollary}
\begin{proof}
\(L/E\) 代数，对它与嵌入 \(\tau\) 使用\cref{b1:thm:embedding-extension}，便得到 \(\widetilde\tau:L\rightarrow\Omega\)。由于它延拓 \(\tau\)，在 \(K\) 上仍为恒等；当 \(L/K\) 正规时，\cref{b1:thm:normal-embedding-criterion}给出 \(\widetilde\tau(L)=L\)。所以将目标限制为 \(L\) 后，才得到 \(L\) 的自同构。

\(\Omega/K\) 代数，故 \(\Omega/E\) 也代数。再次应用嵌入延拓定理，可将 \(\tau\) 延拓为 \(\Omega\rightarrow\Omega\) 的 \(K\) 嵌入。\cref{b1:prop:algebraic-self-embedding}保证它满射，因而是 \(\Omega\) 的 \(K\) 自同构。
\end{proof}
上述结论也允许 \(E/K\) 无限；在代数闭包中，局部选定的共轭根可以相容地延拓到整个域的自同构。

\subsection{多项式根的共轭性}
\label{b1:subsec:1504:02}

\begin{definition}\label{b1:def:field-conjugate-elements}
设 \(\Omega\) 为域 \(K\) 的代数闭包，\(a,b\in\Omega\) 都在 \(K\) 上代数。如果 \(a,b\) 具有相同的 \(K\) 上最小多项式，则称它们为 \(K\) 上的\term[gonge-yuan]{共轭元}[conjugate elements]，也说它们在 \(K\) 上共轭。
\end{definition}
这等价于 \(a,b\) 是同一个 \(K[x]\) 中首一不可约多项式的根。这里的共轭是域论术语，其定义与群中的 \(gxg^{-1}\) 不同；下面通过域的自同构解释它。

\begin{proposition}\label{b1:prop:conjugates-automorphism}
设 \(\Omega\) 为域 \(K\) 的代数闭包，\(a,b\in\Omega\) 都在 \(K\) 上代数。则 \(a,b\) 在 \(K\) 上共轭，当且仅当存在 \(\Omega\) 的 \(K\) 自同构把 \(a\) 送到 \(b\)。若还有中间域 \(K\subseteq L\subseteq\Omega\) 包含 \(a,b\)，且 \(L/K\) 正规，则也存在 \(L\) 的 \(K\) 自同构实现这一替换。
\end{proposition}
\begin{proof}
同一个不可约多项式的两个根由\cref{b1:lem:simple-embedding-extension}给出 \(K(a)\rightarrow K(b)\) 的 \(K\) 同构，先在单扩张上实现所需替换。再应用\cref{b1:cor:intermediate-embedding-extension}，把它延拓为 \(\Omega\) 的 \(K\) 自同构；若 \(L/K\) 正规，也可将它延拓为 \(L\) 的 \(K\) 自同构。反过来，若自同构送 \(a\) 到 \(b\)，则把 \(m_{a,K}(a)=0\) 映过去，得到 \(m_{a,K}(b)=0\)。由于它首一不可约，必为 \(b\) 的最小多项式。
\end{proof}
共轭性依赖基域。例如 \(\sqrt2\) 与 \(-\sqrt2\) 在 \(\mathbb Q\) 上的最小多项式都是 \(X^2-2\)，但在 \(K=\mathbb Q(\sqrt2)\) 上分别为 \(X-\sqrt2\) 和 \(X+\sqrt2\)，不再相同。扩大基域后，可用的系数增加，原来不能区分的根可能被区分。即使在正特征中出现重根，命题谈的也是不同的元素；重数不会产生新的嵌入。

\subsection{嵌入个数与扩张次数}
\label{b1:subsec:1504:03}

保持 \(K\) 的域同态 \(L\rightarrow\Omega\) 全体记为 \(\operatorname{Hom}_K(L,\Omega)\)。域同态必单射，所以此处它就是嵌入集合。保持 \(K\) 的 \(L\) 的自同构全体记为 \(\operatorname{Aut}_K(L)\)，复合使它成为群。
\symidx{\(\operatorname{Hom}_K(L,\Omega)\)：保持 \(K\) 的域嵌入集合}
\symidx{\(\operatorname{Aut}_K(L)\)：保持 \(K\) 的域自同构群}

\begin{proposition}\label{b1:prop:embedding-degree-bound}
设 \(L/K\) 为有限域扩张，\(\Omega\) 为包含 \(L\) 的 \(K\) 的代数闭包。则
\[
1\leq\#\operatorname{Hom}_K(L,\Omega)\leq[L:K].
\]
若 \(L/K\) 还正规，则这些嵌入恰为 \(\operatorname{Aut}_K(L)\) 的元素。
\end{proposition}
\begin{proof}
\cref{b1:thm:embedding-extension}将 \(\operatorname{id}_K\) 延拓到 \(L\)，所以至少有一个嵌入。为证明上界，写 \(L=K(a_1,\ldots,a_r)\)，令 \(E_0=K\)、\(E_i=K(a_1,\ldots,a_i)\)。给定 \(E_{i-1}\) 的一个嵌入，它到 \(E_i\) 的延拓由 \(a_i\) 的最小多项式经系数变换后的不同根决定。该多项式的次数为 \([E_i:E_{i-1}]\)，故\cref{b1:lem:simple-embedding-extension}给出的延拓数至多为这个数。每个 \(E_i\) 的嵌入都限制为一个 \(E_{i-1}\) 的嵌入，不同限制所对应的延拓彼此不同，因此
\[
\#\operatorname{Hom}_K(E_i,\Omega)
\leq [E_i:E_{i-1}]\,\#\operatorname{Hom}_K(E_{i-1},\Omega).
\]
从 \(E_0\) 的唯一嵌入出发，逐层相乘并用\cref{b1:thm:tower-law}，便得 \(\#\operatorname{Hom}_K(L,\Omega)\leq[L:K]\)。正规情形的断言正是\cref{b1:thm:normal-embedding-criterion}：每个嵌入的像都等于 \(L\)。
\end{proof}
这个上界可能严格。取素数 \(p\)，令 \(t\) 在 \(\mathbb F_p\) 上超越，并置
\[
K=\mathbb F_p(t^p),\qquad L=\mathbb F_p(t).
\]
取包含 \(L\) 的 \(K\) 的代数闭包 \(\Omega\)。由\eqref{b1:eq:rational-power-degree}，\([L:K]=p\)，所以 \(t\) 在 \(K\) 上的最小多项式是 \(X^p-t^p\)。在特征 \(p\) 下，
\[
X^p-t^p=(X-t)^p.
\]
它在 \(\Omega\) 中只有一个不同的根 \(t\)。每个 \(K\) 嵌入 \(L\rightarrow\Omega\) 都须把 \(t\) 送到这个根，因而只能是恒等映射。于是
\[
\#\operatorname{Hom}_K(L,\Omega)=|\operatorname{Aut}_K(L)|=1<p=[L:K].
\]
同时，\(L\) 是 \(X^p-t^p\) 在 \(K\) 上的分裂域，故 \(L/K\) 正规。这个例子表明，正规性保证上述嵌入的像仍为 \(L\)，却不保证嵌入数达到扩张次数；重根的重数不会产生新的嵌入。下一章将证明：对有限扩张，如果没有不可分现象，嵌入数便等于扩张次数；若同时正规，这个数就是自同构群的阶。


\begin{samepage}
本章各类嵌入结论的条件列于\cref{b1:tab:embedding-normality-count}；达到次数上界的条件将在下一章讨论。
\begin{table}[H]
\centering
\small
\caption{嵌入延拓、正规性与计数的条件比较}\label{b1:tab:embedding-normality-count}
\begin{tabularx}{\linewidth}{@{}p{0.27\linewidth}X@{}}
\toprule
所处理的对象 & 假设与结论 \\
\midrule
单扩张 \(K(\alpha)/K\) & 给定域同构 \(\sigma:K\rightarrow K'\) 及扩张域 \(\Omega/K'\)，并选定 \(\sigma(m_{\alpha,K})\) 在 \(\Omega\) 中的根 \(\beta\) 后，\cref{b1:lem:simple-embedding-extension}给出唯一满足 \(\alpha\mapsto\beta\) 的延拓。 \\
一般代数扩张 \(L/K\) & 目标域 \(\Omega\) 代数闭时，给定的基域嵌入 \(K\hookrightarrow\Omega\) 可延拓为 \(L\hookrightarrow\Omega\)；延拓通常不唯一（\cref{b1:thm:embedding-extension}）。 \\
正规扩张 \(L/K\) & 在包含 \(L\) 的 \(K\) 的代数闭包 \(\Omega\) 中，每个逐点固定 \(K\) 的嵌入 \(\tau:L\hookrightarrow\Omega\) 都满足 \(\tau(L)=L\)（\cref{b1:thm:normal-embedding-criterion}）。 \\
有限扩张 \(L/K\) & 若 \([L:K]<\infty\)，则 \(L\) 到包含它的 \(K\) 的代数闭包 \(\Omega\) 中的 \(K\) 嵌入数至多为 \([L:K]\)（\cref{b1:prop:embedding-degree-bound}）。 \\
\bottomrule
\end{tabularx}
\end{table}
\end{samepage}

\begin{chaptersummary}
添入一个代数元时，嵌入的延拓由最小多项式的一个根决定。逐次使用这一事实，可构造分裂域并证明其在基域同构意义下唯一。把部分嵌入按延拓排序，则能将同样的想法推广到任意代数扩张；其中的极大性论证必须同时保证定义域的封闭性和映射的相容性。

\ProseParagraphBreak
设 \(L/K\) 为代数扩张，\(\Omega\) 为包含 \(L\) 的 \(K\) 的代数闭包。正规性可由一族非零 \(K\) 系数多项式的全部根生成来刻画；当 \(L/K\) 正规时，每个逐点固定 \(K\) 的嵌入 \(\sigma:L\hookrightarrow\Omega\) 都满足 \(\sigma(L)=L\)。这个判据也适用于无限代数扩张。正规性可以沿塔向上层限制，也对交和合成保持，但既不自动向中间域下降，也不沿两层正规扩张传递。

\ProseParagraphBreak
若 \([L:K]<\infty\)，则逐点固定 \(K\) 的嵌入 \(L\hookrightarrow\Omega\) 的数目至多为 \([L:K]\)；何时达到这个上界，将由下一章的可分性理论说明。
\end{chaptersummary}

\BookExercises

\begin{exercise}\label{b1:ex:15-1}
求 \(x^4+4\) 在 \(\mathbb Q\) 上的分裂域及其次数。
\end{exercise}
\begin{solution}
在 \(\mathbb C\) 中分解为 \((x^2-2x+2)(x^2+2x+2)\)，四个根为 \(1\pm i,-1\pm i\)。它们都属于 \(\mathbb Q(i)\)，而其中的根 \(1+i\) 减去 \(1\) 就得到 \(i\)，所以它们在 \(\mathbb C\) 中生成的子域恰为 \(\mathbb Q(i)\)。因此 \(\mathbb Q(i)\) 是一个分裂域；任意其他分裂域都与它保持 \(\mathbb Q\) 地同构。\(x^2+1\) 在 \(\mathbb Q\) 上无根，故次数为二。
\end{solution}

\begin{exercise}\label{b1:ex:15-2}
设 \(a=\sqrt[4]3>0\)，\(b=a^2=\sqrt3\)。
\begin{enumerate}
\item 求 \(x^4-3\) 在 \(\mathbb Q\) 上的分裂域，并列出 \(\mathbb Q(a)\) 到 \(\mathbb C\) 的全部 \(\mathbb Q\) 嵌入。
\item 分别求 \(a,b\) 在 \(\mathbb Q\) 上的最小多项式。虽然 \(a\mapsto a\)、\(b\mapsto-b\) 分别把它们送到各自最小多项式的根，证明这两个指定不能同时由一个 \(\mathbb Q\) 域同态实现，并说明在 \(\mathbb Q(a)\) 上为何不能再独立选择 \(b\) 的像。
\end{enumerate}
\end{exercise}
\begin{solution}
1.\ 根为 \(a,-a,ia,-ia\)，分裂域为 \(\mathbb Q(a,i)\)。\(x^4-3\) 对素数三满足 Eisenstein 判据，故 \(\bigl[\mathbb Q(a):\mathbb Q\bigr]=4\)。该域实，添入 \(i\) 再给出二次扩张，所以分裂域次数为八。由\cref{b1:lem:simple-embedding-extension}，\(\mathbb Q(a)\) 的嵌入恰由把 \(a\) 送到上述四个根给出；其中只有像为 \(a,-a\) 的两个嵌入把 \(\mathbb Q(a)\) 映入自身。

2.\ \(a,b\) 在 \(\mathbb Q\) 上的最小多项式分别为 \(X^4-3\) 与 \(X^2-3\)。所指定的两个像各自都是对应多项式的根，但若同态固定 \(a\)，就必须把 \(b=a^2\) 送到 \(a^2=b\)，不能送到 \(-b\)。在当前中间域 \(\mathbb Q(a)\) 上，\(b\) 的最小多项式是 \(X-a^2\)，所以逐次延拓时也没有另选其符号的自由。
\end{solution}

\begin{exercise}\label{b1:ex:15-3}
设 \(f\in K[x]\) 非零，\(r\geq1\) 为整数。证明 \(f\) 与 \(f^r\) 在同一个 \(K\) 的代数闭包中的分裂域相同。再取非零 \(g\in K[x]\)，并在同一个代数闭包中取 \(f,g\) 的分裂域；说明乘积 \(fg\) 的分裂域如何由它们得到。
\end{exercise}
\begin{solution}
\(f\) 与 \(f^r\) 有相同的不同根，只是重数改变，故由根生成的域相同。若 \(f,g\) 的分裂域分别为 \(L_f,L_g\)，乘积 \(fg\) 的根集是两个根集的并，故其分裂域是包含 \(L_f,L_g\) 的最小子域，即合成域 \(L_fL_g\)。这也包括常数多项式的情形。
\end{solution}

\begin{exercise}\label{b1:ex:15-4}
在 \(\mathbb F_2[x]\) 中考虑 \(f=x^2+x+1\)。不用有限域分类，直接构造其分裂域，并写出保持 \(\mathbb F_2\) 的两个自同构。
\end{exercise}
\begin{solution}
\(f(0)=f(1)=1\)，所以二次多项式不可约，\(L=\mathbb F_2[x]/(f)\) 是域。记 \(a=x+(f)\)，则 \(a^2+a+1=0\)，又 \(f(a+1)=f(a)=0\)，故它的两个根为 \(a,a+1\)，均在 \(L\) 中。由 \(L=\mathbb F_2(a)\) 知它就是分裂域。两个自同构分别为 \(a\mapsto a\) 与 \(a\mapsto a+1\)；后一映射重复两次恢复 \(a\)，也可直接验证为自同构。
\end{solution}

\begin{exercise}\label{b1:ex:15-5}
证明每个代数闭域都是无限域。若 \(K\) 可数，证明其任意代数闭包也可数。
\end{exercise}
\begin{solution}
若域 \(F\) 有限，则 \(1+\prod_{a\in F}(x-a)\) 为非常数多项式，在每个 \(a\in F\) 处值为一，故没有根，\(F\) 不是代数闭域。若 \(K\) 可数，\(K[x]\) 中非零多项式只有可数多个，每个在代数闭包中只有有限多个不同根，而代数闭包的每个元素都属于某个这样的有限根集。因此代数闭包至多可数；由刚证的无限性，它恰为可数无限集。
\end{solution}

\begin{exercise}\label{b1:ex:15-6}
设 \(\Omega\) 是代数闭域，\(K\subseteq\Omega\)，不要求 \(\Omega/K\) 代数。证明 \(\Omega\) 中在 \(K\) 上代数的元素组成 \(K\) 的一个代数闭包。
\end{exercise}
\begin{solution}
记该集合为 \(A\)。\cref{b1:prop:algebraic-elements-subfield}说明它是子域，且 \(A/K\) 按定义代数。给定非常数 \(f\in A[x]\)，在 \(\Omega\) 中取根 \(a\)。其有限多个系数在 \(K\) 上生成有限扩张 \(E\)，而 \(a\) 在 \(E\) 上代数。\cref{b1:thm:algebraic-tower}给出 \(a\) 在 \(K\) 上代数，故 \(a\in A\)。于是 \(A\) 代数闭，同时满足代数闭包的两个条件。
\end{solution}

\begin{exercise}\label{b1:ex:15-7}
设 \(L/K\) 代数，\(\Omega\) 为 \(K\) 的代数闭包。证明每个 \(K\) 嵌入 \(L\rightarrow\Omega\) 的像都是一个与 \(L\) 同构的中间域，但这些像未必相同。给出具体例子。
\end{exercise}
\begin{solution}
域嵌入保持运算与逆元，故像是包含 \(K\) 的子域，映射本身给出到像的 \(K\) 同构。取 \(L=\mathbb Q(a)\)、\(a=\sqrt[3]2>0\)，并在一个包含三个根的代数闭包中取非实根 \(a\zeta\)。嵌入 \(a\mapsto a\zeta\) 的像为 \(\mathbb Q(a\zeta)\)，包含非实元素；原来的 \(L\) 全部实，故二者不相同。它们尽管嵌入在同一环境中位置不同，抽象的 \(\mathbb Q\) 域结构仍同构。
\end{solution}

\begin{exercise}\label{b1:ex:15-8}
设 \(L/K\) 正规，\(E/K\) 是同一代数闭包中的任意代数扩张。证明合成域 \(LE/E\) 正规。若 \(L/K\) 有限，证明 \([LE:E]\leq[L:K]\)。
\end{exercise}
\begin{solution}
\(L\) 由某族 \(K\) 系数多项式的全部根生成，故 \(LE\) 由 \(E\) 与同一族根生成；把系数视为 \(E\) 中元素，\cref{b1:thm:normal-embedding-criterion}即说明 \(LE/E\) 正规。有限情形下，取 \(L/K\) 的基 \(u_1,\ldots,u_n\)。它们的 \(E\) 线性生成空间包含 \(E,L\)，且对乘法封闭，因为 \(u_iu_j\) 可用 \(K\) 系数写成该基的线性组合。非零元素在这个有限维空间上的乘法为单射，因而满射，所以它也有逆元。因此该空间是域，恰等于 \(LE\)，从而 \([LE:E]\leq n\)。
\end{solution}

\begin{exercise}\label{b1:ex:15-9}
证明每个二次扩张都正规，不论其特征如何。
\end{exercise}
\begin{solution}
设 \([L:K]=2\)，取 \(a\in L\setminus K\)。塔式公式给出 \(L=K(a)\)，其最小多项式写成 \(x^2+bx+c\)。已知一个根 \(a\)，另一个根为 \(-b-a\)，也在 \(L\) 中，包括二者相等的情形。因此 \(L\) 是该多项式的分裂域，由\cref{b1:thm:normal-embedding-criterion}正规。证明没有除以二，所以也适用于特征二。
\end{solution}

\begin{exercise}\label{b1:ex:15-10}
设 \(\sigma:L\rightarrow L\) 为保持 \(K\) 的域嵌入，\(L/K\) 代数，但未必正规。证明 \(\sigma\) 必为自同构。给出删去代数性后的反例。
\end{exercise}
\begin{solution}
给定 \(a\in L\)，令 \(R\) 为 \(m_{a,K}\) 在 \(L\) 中的不同根组成的有限非空集合。\(\sigma\) 固定系数，因此将 \(R\) 单射地映入自身，必为双射，故 \(a\) 有 \(R\) 中原像。对每个 \(a\) 都如此，所以 \(\sigma\) 满射。反例是 \(k(t)\rightarrow k(t)\)、\(t\mapsto t^2\)：超越性保证这是嵌入，但像为 \(k(t^2)\)，由\eqref{b1:eq:rational-power-degree}知它是次数二的真子域。
\end{solution}

\begin{exercise}\label{b1:ex:15-11}
设 \(E/K\) 代数，\(\Omega/K\) 为代数闭包且 \(E\subseteq\Omega\)。证明 \(\Omega\) 也是 \(E\) 的代数闭包，并说明 \(E\) 的任意 \(K\) 嵌入到 \(\Omega\) 能否延拓为 \(\Omega\) 的自同构。
\end{exercise}
\begin{solution}
\(\Omega\) 本来代数闭；其中元素在 \(K\) 上的代数方程也可视为 \(E\) 系数方程，所以 \(\Omega/E\) 代数。因此它是 \(E\) 的代数闭包。给定嵌入，用\cref{b1:thm:embedding-extension}延拓到 \(\Omega\rightarrow\Omega\)，所得映射仍固定 \(K\)。它是代数扩张 \(\Omega/K\) 到自身的嵌入，由上一题的有限根集论证必满射，故为自同构；也可直接用\cref{b1:cor:intermediate-embedding-extension}之后的结论。
\end{solution}

\begin{exercise}\label{b1:ex:15-12}
令 \(a=\sqrt[3]2>0\)，\(\zeta^2+\zeta+1=0\)、\(\zeta\neq1\)，\(L=\mathbb Q(a,\zeta)\)。证明复共轭在 \(L\) 上给出自同构，并求它的固定元素。
\end{exercise}
\begin{solution}
共轭固定 \(a\)，并把 \(\zeta\) 送到 \(\zeta^2=-1-\zeta\)，故保持生成域 \(L\)，其平方为恒等。\eqref{b1:eq:cubic-splitting-degree}给出 \(\bigl[L:\mathbb Q(a)\bigr]=2\)，所以元素唯一写成 \(u+v\zeta\)、\(u,v\in\mathbb Q(a)\)。其共轭为 \(u-v-v\zeta\)，与原式相等便由系数比较得 \(v=-v\)，故 \(v=0\)。所以固定元素恰为 \(\mathbb Q(a)\)。
\end{solution}

\begin{exercise}\label{b1:ex:15-13}
设 \(L/K\) 为有限扩张，整数 \(r\geq1\)，且 \(L=K(a_1,\ldots,a_r)\)。固定包含 \(L\) 的 \(K\) 的代数闭包 \(\Omega\)，并在其中记 \(L\) 的正规闭包为 \(N\)。证明 \(N\) 是 \(\prod_{i=1}^r m_{a_i,K}\) 在 \(\Omega\) 中的分裂域。进一步证明
\[
N=K\left(\bigcup_{\sigma\in\operatorname{Hom}_K(L,\Omega)}\sigma(L)\right),
\]
并由此说明 \(N\) 不依赖生成集的选择。
\end{exercise}
\begin{solution}
令 \(N_0\) 为题中多项式在 \(\Omega\) 中的分裂域。它正规并包含全部 \(a_i\)，因而包含 \(L\)。若 \(M/K\) 是 \(\Omega\) 中包含 \(L\) 的正规扩张，每个 \(m_{a_i,K}\) 都在 \(M\) 中分裂，所以 \(N_0\subseteq M\)。由\cref{b1:prop:normal-closure}的最小性，\(N=N_0\)。

记 \(M_0=K\bigl(\bigcup_{\sigma\in\operatorname{Hom}_K(L,\Omega)}\sigma(L)\bigr)\)。对每个 \(K\) 嵌入 \(\sigma:L\hookrightarrow\Omega\)，\(\sigma(a_i)\) 是 \(m_{a_i,K}\) 的根，故 \(\sigma(L)\subseteq N\)，进而 \(M_0\subseteq N\)。反过来，任取 \(m_{a_i,K}\) 在 \(\Omega\) 中的根 \(\beta\)。由\cref{b1:lem:simple-embedding-extension}，映射 \(a_i\mapsto\beta\) 给出 \(K(a_i)\hookrightarrow\Omega\) 的 \(K\) 嵌入；再由\cref{b1:thm:embedding-extension}将它延拓为 \(\sigma:L\hookrightarrow\Omega\)。于是 \(\beta\in\sigma(L)\subseteq M_0\)，所有这些最小多项式的根都属于 \(M_0\)，故 \(N\subseteq M_0\)。两边相等，而右边只取决于 \(L,K,\Omega\)，不取决于所选生成元。
\end{solution}

\begin{exercise}\label{b1:ex:15-14}
设 \(K\subseteq E\subseteq L\)，\(L/K\) 正规。证明：\(E/K\) 正规，当且仅当 \(L\) 的每个 \(K\) 自同构都将 \(E\) 映到自身。
\end{exercise}
\begin{solution}
若 \(E/K\) 正规，任意自同构的限制都是 \(E\) 到共同代数闭包的 \(K\) 嵌入，由\cref{b1:thm:normal-embedding-criterion}其像为 \(E\)。反过来，给定 \(E\) 到代数闭包的任意 \(K\) 嵌入，\cref{b1:cor:intermediate-embedding-extension}将它延拓到 \(L\) 的 \(K\) 自同构；按假设它把 \(E\) 映到自身。因此所有这样的嵌入都保持 \(E\)，再次用正规性判据得 \(E/K\) 正规。这里无需假定扩张可分。
\end{solution}

% !TeX root = ../../Book1.tex
% 纲要标题：## 第16章　可分性与本原元素
% 纲要位置：计划纲要.md，第 556 行。

\chapter{可分性与本原元素}
\label{b1:ch:16}
\BookChapterNumber{b1:ch:16}

\begin{chapterabstract}
正特征中，多项式的次数可能大于不同根的个数，扩张次数也可能大于可选的嵌入数。本章以形式导数和可分性整理这种差别，讨论纯不可分扩张与完美域，并证明有限可分扩张的本原元素定理。
\end{chapterabstract}

\begin{learninggoals}
\begin{itemize}
\item 用形式导数与最大公因子判断重根，并解释正特征中不可约但不可分的现象。
\item 计算可分次数和不可分次数，区分可分性、正规性与纯不可分性。
\item 使用 Frobenius 判据判断完美域，构造不完美函数域及其纯不可分扩张。
\item 寻找本原元素，说明有限性与可分性在单生成结论中的作用。
\end{itemize}
\end{learninggoals}

设 \(p\) 为素数。在有理函数域 \(\mathbb F_p(t)\) 上，\(X^p-t\) 的次数为 \(p\)，进入含有 \(t^{1/p}\) 的扩张后却只有一个不同的根，因为 \((X-t^{1/p})^p=X^p-t\)。上一章的单扩张嵌入引理说明，若 \(a\) 在 \(K\) 上代数，\(\Omega\) 是包含 \(a\) 的 \(K\) 的代数闭包，则每个 \(K\) 嵌入 \(K(a)\rightarrow\Omega\) 由 \(a\) 的像唯一决定，而且这个像可以是最小多项式的任一个不同根。因此
\[
 \#\operatorname{Hom}_K(K(a),\Omega)
 =\#\{\,m_{a,K}\text{ 在 }\Omega\text{ 中的不同根}\,\}.
\]
重根合并了可供选择的像，多项式次数便可能大于嵌入个数。找出何时两者相等，须从重根与形式导数入手；进一步的问题是，在已有可分元时，能否用一个元素生成整个有限扩张。

\ProseParagraphBreak

本章用到\cref{b1:prop:multiple-root}的重根判别、\cref{b1:lem:simple-embedding-extension}的单扩张嵌入，以及\cref{b1:thm:embedding-extension}的一般延拓定理。扩张次数与嵌入个数都依赖所取基域，比较两者时须固定同一个基域。有限域的存在与分类仍留在下一章；本章只从一个已经给定的有限域出发证明所需性质。

\ProseParagraphBreak
可分性与嵌入的联系可参阅 Roman 的《Field Theory》\cite[Sections 3.1--3.6]{Roman2006FieldTheory}。Matsumura 的《Commutative Ring Theory》\cite[Sections 25--26]{Matsumura1989CommutativeRingTheory}进一步讨论微分与可分性的关系；其第 26 节还借助张量积与约化环，把可分性用于非代数域扩张及交换代数。在代数域扩张上，这与本章的定义一致。\(L/K\) 的可分性考察 \(L\) 中元素在 \(K\) 上的最小多项式；第16.3节的完美性则要求 \(L[x]\) 中所有不可约多项式可分，是对域 \(L\) 自身的要求。即使 \(K\) 不完美，平凡扩张 \(K/K\) 也可分，所以扩张可分不意味着扩张域本身完美。理解 Matsumura 的该节还需用到模、导子与张量积。

% !TeX root = ../../Book1.tex
% 纲要标题：### 16.1 可分多项式
% 纲要位置：计划纲要.md，第 558 行。

\section{可分多项式}
\label{b1:sec:1601}


% 纲要内容（仅作写作计划，尚非教材正文）：
%
% * 重根、形式导数与最大公因子
% * 可分元与可分扩张
% * 特征 p 的不可分现象；区分不同根数与根的共同重数
%

扩张次数是相应向量空间的维数。上一章的\cref{b1:lem:simple-embedding-extension}说明：给定基域嵌入，若目标域包含系数变换后的最小多项式的全部根，则单扩张的延拓数等于这些根的不同个数。它何时等于扩张次数，取决于最小多项式是否有重根。特征零时不可约性已经排除了重根；正特征中导数可能恒为零，这会产生新的现象。本节始终在代数闭包中比较根。

\subsection{重根、形式导数与最大公因子}
\label{b1:subsec:1601:01}

\begin{definition}\label{b1:def:separable-polynomial}
设 \(K\) 为域，\(f\in K[x]\) 为非零多项式。如果 \(f\) 在 \(K\) 的一个代数闭包中没有重根，则称 \(f\) 为\term[kefen-duoxiangshi]{可分多项式}[separable polynomial]。这一条件与代数闭包的选择无关；非零常数多项式按此约定也可分。
\end{definition}
次数为 \(n\) 的非零多项式在代数闭包中按重数共有 \(n\) 个根，所以可分性等价于它有 \(\deg f\) 个不同根；非零常数多项式的次数和根数都为零。由\cref{b1:prop:algebraic-closure-unique}，任意两个 \(K\) 的代数闭包之间存在保持 \(K\) 的同构。将根的分解沿此同构搬过去，各根及其重数都被保持，这就说明定义与代数闭包的选择无关。

\ProseParagraphBreak
\cref{b1:prop:multiple-root}的重根判别给出 \(f\) 可分当且仅当 \(\gcd(f,f')=1\)。虽然重根是在代数闭包中检查的，最大公因子却只需在 \(K[x]\) 中计算。若它为一，就有 \(u,v\in K[x]\) 使
\[
uf+vf'=1.
\]
这条等式进入任意扩域后仍成立；若 \(f,f'\) 有公共根，代入就会得到 \(0=1\)。反之，若 \(K[x]\) 中的最大公因子非常数，它在代数闭包中必有根；由于它同时整除 \(f,f'\)，这个根便是二者的公共根。

\begin{proposition}\label{b1:prop:irreducible-separable-derivative}
设 \(K\) 为域，\(f\in K[x]\) 为不可约多项式。则 \(f\) 可分当且仅当 \(f'\neq0\)。因此特征零域上的所有不可约多项式均可分。
\end{proposition}
\begin{proof}
记首一最大公因子为 \(d=\gcd(f,f')\)。它整除 \(f\)，而 \(f\) 不可约，所以 \(d\) 只能为一或与 \(f\) 相伴；这正是不可约性在证明中的作用。若 \(f'\neq0\)，则 \(\deg f'<\deg f\)，第二种可能会要求 \(f\mid f'\)，与次数不等式矛盾。因此 \(d=1\)，\(f\) 可分。若 \(f'=0\)，最大公因子就是首一化的 \(f\)，非常数，故不可分。特征零时次数为 \(n\geq1\)、首项系数为 \(a_n\) 的多项式，其导数首项系数 \(na_n\neq0\)，所以不可约多项式总属于前一种情形。
\end{proof}
特征零保证的是不可约多项式全部可分；一般多项式的不可约分解中仍可能重复出现同一个因子，例如 \((x-1)^2\)。不同的不可约因子互素，由 Bézout 关系知它们没有公共根，所以在特征零域上，一个非零多项式可分恰好等价于其不可约分解中没有重复因子。

\subsection{可分元与可分扩张}
\label{b1:subsec:1601:02}

\begin{definition}\label{b1:def:separable-extension}
设 \(L/K\) 为域扩张，\(a\in L\) 在 \(K\) 上代数。如果最小多项式 \(m_{a,K}\) 可分，则称 \(a\) 为 \(K\) 上的\term[kefen-yuan]{可分元}[separable element]，也说它在 \(K\) 上可分。如果 \(L/K\) 是代数扩张，并且 \(L\) 的每个元素都在 \(K\) 上可分，则称 \(L/K\) 为\term[kefen-kuozhang]{可分扩张}[separable extension]。
\end{definition}
若 \(a\) 被一个可分 \(K\) 系数多项式 \(f\) 零化，其最小多项式 \(m_{a,K}\) 整除 \(f\)。若这个因子有重根，该根在 \(f\) 中也至少出现两次，便与 \(f\) 可分矛盾。因此 \(m_{a,K}\) 可分，\(a\) 也可分。

\ProseParagraphBreak
若扩大基域为中间域 \(K\subseteq E\subseteq L\)，原来的关系 \(m_{a,K}(a)=0\) 仍成立，所以 \(m_{a,E}\) 在 \(E[x]\) 中整除 \(m_{a,K}\)。增加可用的系数可能让原最小多项式分解，并使新的最小多项式次数降低；但新的多项式只是原多项式的因子，不能产生新的重根。故 \(a\) 在 \(K\) 上可分就仍在 \(E\) 上可分。反方向一般不成立：在 \(K(a)\) 上，\(a\) 的最小多项式总是一次的，即使它在 \(K\) 上不可分。

\ProseParagraphBreak
例如 \(\mathbb Q(\sqrt2)/\mathbb Q\) 既可分又正规：特征零给出可分性，而它是 \(x^2-2\) 的分裂域。\(\mathbb Q(\sqrt[3]2)/\mathbb Q\) 同样可分，却只包含 \(x^3-2\) 的实根，因而不正规。可分性检查每个最小多项式的根是否相异；正规性检查这些根是否全部留在所讨论的域里。正特征中还会出现正规但不可分以及两者都不满足的扩张，下面的算例将直接检验它们。

\subsection{特征 \texorpdfstring{\(p\)}{p} 的不可分现象}
\label{b1:subsec:1601:03}

设 \(\operatorname{char}K=p>0\)。由导数公式，\(f'=0\) 当且仅当 \(f\) 的所有非零幂指数均被 \(p\) 整除，即 \(f(x)=g(x^p)\)。这并不必然使 \(f\) 在 \(K[x]\) 中可约，因为系数未必有 \(p\) 次根。

\begin{proposition}\label{b1:prop:inseparable-polynomial-shape}
设 \(K\) 为特征 \(p>0\) 的域，\(f\in K[x]\) 为首一不可约多项式。则存在非负整数 \(r\) 及首一不可约可分多项式 \(g\in K[x]\)，使
\[
f(x)=g(x^{p^r}).
\]
\(f\) 的每个不同根的重数均为 \(p^r\)，不同根数为 \(\deg g\)。若 \(a\) 是 \(f\) 的根，则 \(g\) 是 \(a^{p^r}\) 在 \(K\) 上的最小多项式，因而 \(a^{p^r}\) 在 \(K\) 上可分。
\end{proposition}
\begin{proof}
只要导数为零，就把所有非零项的正指数同时除以 \(p\)。多项式仍非常数而次数严格下降，所以有限步后得到导数非零的 \(g\)。这里 \(r\) 是使 \(f\) 的所有非零项指数都被 \(p^r\) 整除的最大整数：所得 \(g\) 至少有一个非零项的指数不被 \(p\) 整除，否则还可以再除一次。若 \(g\) 可约，把分解中的变量换成 \(x^{p^r}\) 就使 \(f\) 可约，矛盾。故 \(g\) 不可约，由\cref{b1:prop:irreducible-separable-derivative}可分。

在代数闭包中，设 \(g\) 的互异根为 \(b_1,\ldots,b_s\)。每个 \(b_i\) 有唯一 \(p^r\) 次根 \(a_i\)：存在性来自代数闭性；若 \(u^{p^r}=v^{p^r}\)，则 \((u-v)^{p^r}=0\)，而域没有非零幂零元，所以 \(u=v\)。这正是幂映射 \(u\mapsto u^p\) 及其 \(r\) 次迭代的单射性。于是
\[
f(x)=\prod_{i=1}^s(x^{p^r}-b_i)=\prod_{i=1}^s(x-a_i)^{p^r}.
\]
不同 \(b_i\) 对应不同 \(a_i\)，故不同根数为 \(s=\deg g\)，每个根的重数为 \(p^r\)，并且
\[
\deg f=(\deg g)p^r.
\]
因此 \(g\) 的次数记录了全部不同根的个数，而 \(p^r\) 记录它们共有的重数。又因为 \(a^{p^r}\) 是首一不可约多项式 \(g\) 的根，\(g\) 就是它在 \(K\) 上的最小多项式，且 \(a^{p^r}\) 在 \(K\) 上可分。
\end{proof}
在 \(K=\mathbb F_p(t)\) 上，\(x^p-t\) 不可约：它在唯一分解整环 \(\mathbb F_p[t]\) 上对素元 \(t\) 满足 Eisenstein 判据，再由 Gauss 分解定理得到分式域上的不可约性。在一个代数闭包中取 \(u\) 使 \(u^p=t\)，并令 \(L=K(u)\)。此时 \([L:K]=p\)，但在 \(L[x]\) 中有
\[
x^p-t=(x-u)^p.
\]
这个等式给出一个重数为 \(p\) 的根，故 \(u\) 在 \(K\) 上不可分，\(L/K\) 也不可分；同时 \(L\) 已是该多项式的分裂域，由\cref{b1:thm:normal-embedding-criterion}知 \(L/K\) 正规。不可约性是在 \(K[x]\) 中成立的，进入 \(L[x]\) 后同一个多项式已经分解，所以不可约与否必须指明系数域。

\ProseParagraphBreak
正规性和可分性也可以同时失败。取 \(K=\mathbb F_2(t)\)，在代数闭包中取 \(a^6=t\)，令 \(L=K(a)=\mathbb F_2(a)\)。多项式 \(x^6-t\) 同样对素元 \(t\) 满足 Eisenstein 判据，故在 \(K[x]\) 中不可约；其导数为零，所以 \(L/K\) 不可分。在代数闭包中有
\[
x^6-t=(x^3-a^3)^2
=(x-a)^2(x^2+ax+a^2)^2.
\]
若它在 \(L\) 中完全分裂，\(L\) 就包含一个满足 \(\zeta^2+\zeta+1=0\) 的三次单位根 \(\zeta\)。可是 \(a\) 超越于 \(\mathbb F_2\)；若写 \(\zeta=P(a)/Q(a)\)，其中非零 \(P,Q\in\mathbb F_2[x]\)，就会得到 \(P^2+PQ+Q^2=0\)。当 \(\deg P\ne\deg Q\) 时，平方次数较高的一项无法消去；当次数相同时，最高项系数为 \(1+1+1=1\)，仍不能消去。因此这样的 \(\zeta\) 不在 \(L\) 中，\(x^6-t\) 不在 \(L\) 中分裂，\(L/K\) 也不正规。

% !TeX root = ../../Book1.tex
% 纲要标题：### 16.2 可分次数
% 纲要位置：计划纲要.md，第 564 行。

\section{可分次数}
\label{b1:sec:1602}


% 纲要内容（仅作写作计划，尚非教材正文）：
%
% * 嵌入个数与可分次数
% * 可分次数的塔式公式
% * 纯不可分扩张及其正规性；在混合次数算例中计算最大可分中间域
%

计数嵌入时，生成元的像必须是对应最小多项式的根，并且逐层选择要保持此前已经确定的映射。上一章已给出延拓的存在性与次数上界；现在把重根造成的差额单独记下来。

\subsection{嵌入个数与可分次数}
\label{b1:subsec:1602:01}

\begin{definition}\label{b1:def:separable-degree}
设 \(L/K\) 为有限域扩张，\(\Omega\) 为包含 \(L\) 的 \(K\) 的代数闭包。把整数
\[
[L:K]_{\mathrm s}\coloneqq\#\operatorname{Hom}_K(L,\Omega)
\]
称为扩张 \(L/K\) 的\term[kefen-cishu]{可分次数}[separable degree]。代数闭包之间的 \(K\) 同构通过复合给出嵌入集合的双射，所以该数不依赖所选代数闭包。
\end{definition}
\symidx{\([L:K]_{\mathrm s}\)：可分次数}

对单扩张 \(L=K(a)\)，\cref{b1:lem:simple-embedding-extension}把每个 \(K\) 嵌入与 \(a\) 的一个共轭根对应起来，因而
\[
[K(a):K]_{\mathrm s}
=\#\{\,b\in\Omega\mid m_{a,K}(b)=0\,\}
\leq[K(a):K]=\deg m_{a,K}.
\]
总次数是向量空间维数，可分次数是不同嵌入的个数，而一个根出现多次并不会给出多个嵌入。在特征 \(p>0\) 时，若 \(m_{a,K}(x)=g(x^{p^r})\)，其中 \(g\) 不可约且可分，\cref{b1:prop:inseparable-polynomial-shape}给出
\[
[K(a):K]=p^r\deg g,\qquad
[K(a):K]_{\mathrm s}=\deg g,\qquad
\frac{[K(a):K]}{[K(a):K]_{\mathrm s}}=p^r.
\]
这里第三个数字正是各根共有的重数，记录总次数中没有转化为不同嵌入的因子。对一般有限扩张，这个比值也总是正整数；\cref{b1:prop:separable-inseparable-decomposition}将证明这一点，并把它定义为不可分次数。

\begin{theorem}\label{b1:thm:separable-embeddings}
有限扩张 \(L/K\) 满足 \([L:K]_{\mathrm s}\leq[L:K]\)，等号成立当且仅当 \(L/K\) 可分。若 \(L\) 由有限多个在 \(K\) 上可分的元素生成，则 \(L/K\) 可分。
\end{theorem}
\begin{proof}
固定包含 \(L\) 的 \(K\) 的代数闭包 \(\Omega\)。上界已由\cref{b1:prop:embedding-degree-bound}证明。假设 \(L=K(a_1,\ldots,a_r)\)，每个 \(a_i\) 在 \(K\) 上可分。令 \(E_0=K\)、\(E_i=K(a_1,\ldots,a_i)\)。每个 \(E_i\) 的嵌入都由 \(E_{i-1}\) 的一个嵌入延拓而来；若在同一层每个已有嵌入都有相同数目的延拓，则该层的嵌入总数就是原总数乘以这个数。下面要证明，第 \(i\) 层的延拓数恰为 \([E_i:E_{i-1}]\)。

记 \(a_i\) 在 \(E_{i-1}\) 上的最小多项式为 \(m_i\)。它整除 \(a_i\) 在 \(K\) 上的可分最小多项式，故仍可分。于是存在 \(u_i,v_i\in E_{i-1}[x]\)，使
\[
u_i m_i+v_i m_i'=1.
\]
给定 \(K\) 嵌入 \(\tau:E_{i-1}\rightarrow\Omega\)，逐项变换系数与形式求导相容，即 \(\tau(m_i')=\tau(m_i)'\)。把 \(\tau\) 施于上式的系数，得到
\[
\tau(u_i)\tau(m_i)+\tau(v_i)\tau(m_i)'=1.
\]
因此 \(\tau(m_i)\) 与其导数互素，仍没有重根；系数嵌入不把非零首项系数变成零，所以其次数仍为 \(\deg m_i=[E_i:E_{i-1}]\)。在代数闭包 \(\Omega\) 中，它恰有这么多个不同根。由\cref{b1:lem:simple-embedding-extension}，\(\tau\) 也恰有这么多个延拓。不同的限制对应互不相交的延拓集合，从 \(E_0=K\) 的唯一 \(K\) 嵌入开始逐层计数，得
\[
\#\operatorname{Hom}_K(L,\Omega)
=\prod_{i=1}^r[E_i:E_{i-1}]=[L:K].
\]
特别地，若 \(L/K\) 可分，取任何有限生成集即得到等号。

反过来，设嵌入总数已等于 \([L:K]\)。任取 \(a\in L\)，记其最小多项式的不同根数为 \(s\)。从 \(K(a)\) 到 \(\Omega\) 恰有 \(s\) 个 \(K\) 嵌入。对其中任意一个嵌入，把 \(L/K(a)\) 写成逐次单扩张，每层延拓的选择数至多为相应最小多项式的次数，所以其延拓总数至多为 \(\bigl[L:K(a)\bigr]\)。将所有 \(L\) 的嵌入按它们在 \(K(a)\) 上的限制分组，便得第一个不等号；第二个不等号则来自不同根数不超过多项式次数：
\[
[L:K]\leq s\bigl[L:K(a)\bigr]\leq\bigl[K(a):K\bigr]\,\bigl[L:K(a)\bigr]=[L:K].
\]
首项等于嵌入总数是此方向的假设，末项等于总次数是扩张次数的塔式公式。首尾相等迫使所有不等号均为等号，所以 \(s=\bigl[K(a):K\bigr]\)，即 \(a\) 的最小多项式没有重根。任意 \(a\) 都可分，故扩张可分；这也说明前面由可分生成元得到的等号足以推出整个扩张可分。
\end{proof}

可分扩张的定义要求每个元素都可分，这个定理却把有限生成情形的检查缩小到生成元：只需验证 \(a_1,\ldots,a_r\) 的最小多项式没有重根，就能保证 \(K(a_1,\ldots,a_r)\) 中的所有元素可分。特别是，从 \(a,b\) 各自的最小多项式看不出 \(a-b,ab\) 及非零 \(a\) 的逆元的可分性时，可以把它们放入共同的有限扩张 \(K(a,b)\)。这与\cref{b1:prop:algebraic-elements-subfield}用有限维数代替显式消元的做法相同。

\begin{proposition}\label{b1:prop:separable-subfield}
任意扩张 \(L/K\) 中在 \(K\) 上代数且可分的元素构成一个子域。
\end{proposition}
\begin{proof}
若 \(a,b\) 都可分，\cref{b1:thm:separable-embeddings}说明有限扩张 \(K(a,b)/K\) 可分。其元素 \(a-b,ab\) 以及 \(a\neq0\) 时的 \(a^{-1}\) 因而均可分。底域元素的最小多项式为一次，所以也包含在这个集合中。
\end{proof}

\subsection{可分次数的塔式公式}
\label{b1:subsec:1602:02}

\begin{theorem}\label{b1:thm:separable-degree-tower}
设 \(K\subseteq E\subseteq L\) 为代数扩张塔。则 \(L/K\) 可分当且仅当 \(E/K\) 与 \(L/E\) 都可分。若进一步有 \([L:K]<\infty\)，则可分次数满足
\[
[L:K]_{\mathrm s}=[L:E]_{\mathrm s}[E:K]_{\mathrm s}.
\]
\end{theorem}
\begin{proof}
先设 \([L:K]<\infty\)，固定包含 \(L\) 的 \(K\) 的代数闭包 \(\Omega\)；它也是 \(E\) 的代数闭包。把 \(L\) 的 \(K\) 嵌入按它们在 \(E\) 上的限制分组，一共有 \([E:K]_{\mathrm s}\) 种可能的限制。只要证明每组都恰有 \([L:E]_{\mathrm s}\) 个元素，便能相乘得到总数。因此考虑限制映射
\[
\operatorname{res}:\operatorname{Hom}_K(L,\Omega)
\longrightarrow\operatorname{Hom}_K(E,\Omega),
\qquad \sigma\longmapsto\sigma|_E.
\]
对每个 \(\tau\in\operatorname{Hom}_K(E,\Omega)\)，我们想把固定 \(E\) 的嵌入 \(\rho:L\rightarrow\Omega\) 变为在 \(E\) 上等于 \(\tau\) 的嵌入。\(\rho\) 的值在 \(\Omega\) 中，未必属于 \(E\)，所以只定义在 \(E\) 上的 \(\tau\) 不能直接与它复合。为此，把\cref{b1:cor:intermediate-embedding-extension}用于塔 \(K\subseteq E\subseteq\Omega\)，将 \(\tau\) 延拓为整个 \(\Omega\) 的一个 \(K\) 自同构 \(\widehat\tau\)。若 \(\rho\in\operatorname{Hom}_E(L,\Omega)\)，则 \(\widehat\tau\circ\rho\) 固定 \(K\)，并且对每个 \(e\in E\) 有
\[
(\widehat\tau\circ\rho)(e)=\widehat\tau(e)=\tau(e).
\]
所以它属于纤维 \(\operatorname{res}^{-1}(\tau)\)。反过来，若 \(\psi\in\operatorname{res}^{-1}(\tau)\)，则对每个 \(e\in E\) 有
\[
(\widehat\tau^{-1}\circ\psi)(e)
=\widehat\tau^{-1}(\tau(e))=e,
\]
故 \(\widehat\tau^{-1}\circ\psi\in\operatorname{Hom}_E(L,\Omega)\)。这两个复合互为逆映射，给出了 \(\operatorname{Hom}_E(L,\Omega)\) 与该纤维之间的双射。计数只要求存在一个双射，并不要求双射唯一：\(\widehat\tau\) 未必唯一，但任一选择都给出上述双射，足以说明纤维的大小等于同一个数 \([L:E]_{\mathrm s}\)。目标集合有 \([E:K]_{\mathrm s}\) 个元素，得到乘法公式。

有限情形下，每层可分次数至多等于该层次数。将刚证明的公式与\cref{b1:thm:tower-law}比较，两个有限正数的乘积取满值当且仅当每个因子取满值，再由\cref{b1:thm:separable-embeddings}得到可分性的塔判据。

一般情形，若 \(L/K\) 可分，限制到 \(E\) 以及扩大最小多项式的系数域，分别说明 \(E/K\)、\(L/E\) 可分。反过来，假设这两层都可分。即使 \(E/K\) 无限，一个具体多项式也只涉及有限多个系数，因而可以像\cref{b1:thm:algebraic-tower}的证明一样，把问题放入有限塔。给定 \(a\in L\)，写出其在 \(E\) 上的可分最小多项式
\[
f(x)=x^n+c_{n-1}x^{n-1}+\cdots+c_0,
\qquad E_0=K(c_0,\ldots,c_{n-1}).
\]
每个 \(c_i\) 都在 \(K\) 上代数且可分，所以\cref{b1:thm:separable-embeddings}说明 \(E_0/K\) 有限可分。

因为 \(f\) 在 \(E[x]\) 中不可约且可分，\(f'\neq0\)。其系数已经全在 \(E_0\) 中，缩小系数域不改变求导公式，故这个非零多项式 \(f'\) 也属于 \(E_0[x]\)。若 \(f\) 在 \(E_0[x]\) 中可约，同一分解会使它在 \(E[x]\) 中可约，矛盾；所以 \(f\) 在 \(E_0[x]\) 中仍不可约。再由\cref{b1:prop:irreducible-separable-derivative}，\(f\) 在 \(E_0[x]\) 中可分，故 \(E_0(a)/E_0\) 有限可分。应用已经证明的有限塔判据，\(E_0(a)/K\) 可分，于是 \(a\) 在 \(K\) 上可分。任意 \(a\) 均如此，结论成立。
\end{proof}
例如特征零的有限扩张中可分次数总等于次数；在 \(\mathbb F_p(t^{1/p})/\mathbb F_p(t)\) 中，可分次数为一而次数为 \(p\)。两者反映的是不同根的选择数与向量空间大小之间的差别。

\subsection{纯不可分扩张}
\label{b1:subsec:1602:03}

\begin{definition}\label{b1:def:purely-inseparable}
设 \(L/K\) 为代数扩张。如果 \(\operatorname{char}K=p>0\)，并且对每个 \(a\in L\) 都存在整数 \(r=r(a)\geq0\) 使 \(a^{p^r}\in K\)，则称 \(L/K\) 为\term[chun-bukefen-kuozhang]{纯不可分扩张}[purely inseparable extension]。在特征零情形，只把平凡扩张 \(K/K\) 称为纯不可分扩张。对任意在 \(K\) 上代数的元素 \(a\)，若 \(K(a)/K\) 纯不可分，就称 \(a\) 在 \(K\) 上纯不可分。
\end{definition}

这里的量词是对每个 \(a\in L\) 存在一个可以依赖 \(a\) 的指数 \(r=r(a)\)，没有要求一开始就用同一个指数处理所有元素。特征零时，不可约多项式全部可分；若它只有一个不同根，次数只能为一，根已经在 \(K\) 中。因此只有平凡扩张符合这种行为，这也解释了定义中的特征零约定。

\begin{proposition}\label{b1:prop:purely-inseparable-criteria}
设 \(K\) 为域，\(\Omega\) 为 \(K\) 的代数闭包，\(a\in\Omega\) 在 \(K\) 上代数。则 \(a\) 在 \(K\) 上纯不可分，当且仅当其最小多项式在 \(\Omega\) 中只有一个不同根。当这两个等价条件成立且 \(\operatorname{char}K=p>0\) 时，该最小多项式形如 \(x^{p^r}-c\)，其中 \(r\geq0\) 为整数、\(c\in K\)。因此，对任意代数扩张 \(L/K\)，固定一个包含 \(L\) 的 \(K\) 的代数闭包以后，扩张 \(L/K\) 纯不可分，当且仅当从 \(L\) 到该代数闭包的 \(K\) 嵌入只有一个。
\end{proposition}
\begin{proof}
先设 \(\operatorname{char}K=p>0\)。若 \(a^{p^r}=c\in K\)，则 \(x^{p^r}-c=(x-a)^{p^r}\)，其任何非常数因子都只有根 \(a\)。反过来，由\cref{b1:prop:inseparable-polynomial-shape}，只有一个不同根的不可约多项式的可分部分必须为一次，故它就是 \(x^{p^r}-c\)，并有 \(a^{p^r}=c\in K\)。这也保证整个 \(K(a)/K\) 纯不可分：其中每个元素都可写成 \(a\) 的 \(K\) 系数有理式，取 \(p^r\) 次幂后各系数与 \(a^{p^r}\) 都在 \(K\) 中，分母仍非零，因而所得值在 \(K\) 中。特征零时不可约多项式可分，只有一个根就只能为一次，等价于 \(a\in K\)。

固定自然包含 \(\iota:L\hookrightarrow\Omega\)。纯不可分时，每个 \(a\in L\) 的像必须是其最小多项式的根，而唯一的根就是 \(a\) 自身，故
\[
\operatorname{Hom}_K(L,\Omega)=\{\iota\}.
\]
若扩张不是纯不可分，某个 \(a\) 的最小多项式 \(m\) 有不同根 \(b\neq a\)。两次代入分别给出
\[
K(a)\cong_K K[x]/(m)\cong_K K(b),
\]
从而得到把 \(a\) 送到 \(b\) 的 \(K\) 嵌入 \(K(a)\rightarrow\Omega\)。再由\cref{b1:thm:embedding-extension}把这个嵌入延拓到 \(L\)，便得到不同于 \(\iota\) 的嵌入。
\end{proof}

特别地，若 \(\operatorname{char}K=p>0\)，且 \(m_{a,K}(x)=x^{p^r}-c\)，则在代数闭包中
\[
m_{a,K}(x)=(x-a)^{p^r},\qquad
[K(a):K]=p^r,\qquad [K(a):K]_{\mathrm s}=1.
\]
只有一个嵌入可选，总次数中的其余因子完全来自这一个根的重数。

\begin{corollary}\label{b1:cor:purely-inseparable-normal}
设 \(L/K\) 为纯不可分代数扩张，不要求有限。则 \(L/K\) 正规，且 \(\operatorname{Aut}_K(L)=\{\operatorname{id}_L\}\)。
\end{corollary}
\begin{proof}
固定包含 \(L\) 的 \(K\) 的代数闭包 \(\Omega\)。若 \(K[x]\) 中的不可约多项式在 \(L\) 中有根 \(a\)，其首一化就是 \(a\) 的最小多项式。由\cref{b1:prop:purely-inseparable-criteria}，它在 \(\Omega\) 中只有根 \(a\)，因此已经在 \(L[x]\) 中完全分裂。由正规扩张的定义，\(L/K\) 正规。每个 \(K\) 自同构同样必须把 \(a\in L\) 送到其最小多项式的根，所以固定每个元素，只能为 \(\operatorname{id}_L\)。
\end{proof}

\begin{proposition}\label{b1:prop:separable-inseparable-decomposition}
设 \(L/K\) 为有限域扩张，令 \(E\subseteq L\) 为其中所有在 \(K\) 上可分的元素组成的子域。则 \(E\) 是 \(L/K\) 的最大可分中间域：\(E/K\) 可分，且任何中间域 \(M\) 若 \(M/K\) 可分就有 \(M\subseteq E\)。此外，\(L/E\) 纯不可分，而且
\[
[L:K]_{\mathrm s}=[E:K].
\]
由扩张次数的塔式公式，商
\[
[L:K]_{\mathrm i}\coloneqq\frac{[L:K]}{[L:K]_{\mathrm s}}
=\frac{[L:K]}{[E:K]}=[L:E]
\]
是正整数，称为\term[bukefen-cishu]{不可分次数}[inseparable degree]；当 \(\operatorname{char}K=p>0\) 时它为 \(p\) 的幂，特征零时为一。
\end{proposition}
\begin{proof}
\cref{b1:prop:separable-subfield}保证 \(E\) 是域，它按定义对 \(K\) 可分。若中间域 \(M/K\) 可分，其每个元素都属于 \(E\)，所以 \(M\subseteq E\)。

正特征时，要让 \(a\in L\) 落到可分元素组成的域 \(E\) 中，可以从它的最小多项式出发。由\cref{b1:prop:inseparable-polynomial-shape}写成 \(m_{a,K}(x)=g(x^{p^r})\)，其中 \(g\) 不可约且可分。令 \(b=a^{p^r}\)，则 \(g(b)=0\)，所以 \(g\) 是 \(b\) 在 \(K\) 上的最小多项式，\(b\) 在 \(K\) 上可分。这正说明 \(a^{p^r}\in E\)，于是每个 \(a\) 在 \(E\) 上纯不可分，\(L/E\) 纯不可分。特征零时所有元素都可分，故 \(E=L\)。

由\cref{b1:prop:purely-inseparable-criteria}，\([L:E]_{\mathrm s}=1\)；又因 \(E/K\) 可分，\([E:K]_{\mathrm s}=[E:K]\)。可分次数塔式公式遂给出
\[
[L:K]_{\mathrm s}=[L:E]_{\mathrm s}[E:K]_{\mathrm s}=[E:K].
\]
再用总次数的塔式公式，商 \([L:K]/[L:K]_{\mathrm s}=[L:E]\) 为正整数。

若 \(\operatorname{char}K=p>0\)，取有限生成集 \(L=E(a_1,\ldots,a_m)\)，令 \(F_0=E\)、\(F_i=E(a_1,\ldots,a_i)\)。每个 \(a_i\) 的某个 \(p\) 次幂迭代已经在 \(E\subseteq F_{i-1}\) 中，所以 \(a_i\) 在 \(F_{i-1}\) 上纯不可分。由\cref{b1:prop:purely-inseparable-criteria}，其最小多项式形如 \(x^{p^{s_i}}-c_i\)，其中 \(s_i\geq0\)、\(c_i\in F_{i-1}\)。于是
\[
[F_i:F_{i-1}]=p^{s_i},\qquad
[L:E]=\prod_{i=1}^m p^{s_i}=p^{s_1+\cdots+s_m}.
\]
特征零时 \(E=L\)，不可分次数为一。
\end{proof}
\symidx{\([L:K]_{\mathrm i}\)：不可分次数}

上述结论把三种次数联系在同一域塔 \(K\subseteq E\subseteq L\) 中：下层 \(E/K\) 可分，其次数是 \([L:K]_{\mathrm s}\)；上层 \(L/E\) 纯不可分，其次数是 \([L:K]_{\mathrm i}\)，而
\[
[L:K]=[E:K][L:E]=[L:K]_{\mathrm s}[L:K]_{\mathrm i}.
\]
域 \(E\) 由 \(L\) 中所有 \(K\) 可分元素唯一确定。这里得到的是一层可分、上一层纯不可分的域塔，并没有一般声称还存在纯不可分中间域 \(P/K\)，使 \(L=EP\)。

\begin{corollary}\label{b1:cor:simple-extension-separable-part}
设 \(K\) 为特征 \(p>0\) 的域，\(L=K(a)\) 为其有限扩张。若 \(a\) 在 \(K\) 上的最小多项式为 \(m_{a,K}(x)=g(x^{p^r})\)，其中 \(r\geq0\) 为整数，\(g\in K[x]\) 首一、不可约且可分，则 \(L/K\) 的最大可分中间域为 \(E=K(a^{p^r})\)，并且
\[
[E:K]=\deg g,\qquad [L:E]=p^r.
\]
\end{corollary}
\begin{proof}
令 \(b=a^{p^r}\)。因为 \(g(b)=0\) 且 \(g\) 不可约，它就是 \(b\) 在 \(K\) 上的最小多项式。因此 \(K(b)/K\) 可分，且 \([K(b):K]=\deg g\)。由\cref{b1:prop:inseparable-polynomial-shape}，\(m_{a,K}\) 有 \(\deg g\) 个不同根；单扩张的嵌入对应于这些根，所以 \([L:K]_{\mathrm s}=\deg g\)。\cref{b1:prop:separable-inseparable-decomposition}中的最大可分中间域包含 \(K(b)\)，且次数同为 \(\deg g\)，故二者相等。最后由 \([L:K]=p^r\deg g\) 及扩张次数的塔式公式得到 \([L:E]=p^r\)。
\end{proof}

取奇素数 \(p\)，令 \(u\) 在 \(\mathbb F_p\) 上超越，并设 \(t=u^{2p}\)、\(K=\mathbb F_p(t)\)、\(L=\mathbb F_p(u)\)。多项式 \(x^{2p}-t\) 在 \(\mathbb F_p[t][x]\) 中对素元 \(t\) 满足 Eisenstein 判据，因而在 \(K[x]\) 中不可约，是 \(u\) 在 \(K\) 上的最小多项式。它又可写为
\[
x^{2p}-t=g(x^p),\qquad g(y)=y^2-t.
\]
\(g\) 同样满足 Eisenstein 判据，且因 \(p\) 为奇数而有 \(g'(y)=2y\neq0\)，所以不可约且可分。在代数闭包中，原多项式为 \((x-u)^p(x+u)^p\)，只有 \(u,-u\) 两个不同根，每个根重数为 \(p\)。

\ProseParagraphBreak

令 \(v=u^p\)，则 \(v^2=t\)，上述推论给出最大可分中间域
\[
E=K(v)=\mathbb F_p(u^p),\qquad
[E:K]=2,\qquad [L:E]=p.
\]
下层 \(E/K\) 的可分性来自最小多项式 \(y^2-t\) 的两个不同根；上层满足 \(u^p=v\)，是纯不可分扩张。于是 \([L:K]=2p\)、\([L:K]_{\mathrm s}=2\)、\([L:K]_{\mathrm i}=p\)：两个嵌入的选择来自下层，而上层的 \(p\) 倍次数全来自重数。

% !TeX root = ../../Book1.tex
% 纲要标题：### 16.3 完美域
% 纲要位置：计划纲要.md，第 570 行。

\section{完美域}
\label{b1:sec:1603}


% 纲要内容（仅作写作计划，尚非教材正文）：
%
% * 特征零的域与有限域
% * Frobenius 映射的判据
% * 不完美函数域的例子；一般完美闭包、最小性与反复 p 次根的逐层构造
%

不可分性在正特征中可能出现，却并非每个正特征域都会出现。对特征为素数 \(p\) 的域，决定性的区别是域中的每个元素能否在原域内开 \(p\) 次方。有限域的有限性恰好保证这一点；有理函数域则提供反例。

\subsection{特征零的域与有限域}
\label{b1:subsec:1603:01}

\begin{definition}\label{b1:def:perfect-field}
设 \(K\) 为域。如果对每个不可约多项式 \(f\in K[x]\)，\(f\) 都可分，则称 \(K\) 为\term[wanmei-yu]{完美域}[perfect field]。等价地，对每个代数域扩张 \(L/K\)，\(L/K\) 都可分。
\end{definition}
第一种说法逐个检查 \(K\) 系数不可约多项式，第二种说法则量化所有代数扩张，包括无限代数扩张。二者通过最小多项式联系起来：若所有不可约多项式可分，任一代数扩张 \(L/K\) 中每个元素 \(a\) 的最小多项式 \(m_{a,K}\) 就可分，所以整个扩张可分。反过来，对任一首一不可约多项式 \(f\in K[x]\)，在代数闭包中添入一个根 \(a\)，便得到代数扩张 \(K(a)/K\)；它可分就使 \(m_{a,K}=f\) 可分。任意不可约多项式与其首一化相伴，故也可分。由\cref{b1:prop:irreducible-separable-derivative}，特征零的域全都完美。因此 \(\mathbb Q,\mathbb R,\mathbb C\) 及它们的任意扩张域都是完美域。

\ProseParagraphBreak
设 \(K\) 为特征 \(p>0\) 的域。二项式系数 \(\binom pi\) 对 \(0<i<p\) 均被 \(p\) 整除，故 \((a+b)^p=a^p+b^p\)；又有 \((ab)^p=a^pb^p\) 和 \(1^p=1\)。因此映射
\[
F_K:K\longrightarrow K,\qquad a\longmapsto a^p
\]
同时保持加法、乘法和单位，是域同态，称为域 \(K\) 的\term[frobenius-yingshe]{Frobenius 映射}[Frobenius map]。它单射，因为 \(a^p=b^p\) 蕴含 \((a-b)^p=0\)，而域是整环，没有非零幂零元，故 \(a-b=0\)。若 \(K\) 有限，单射必满射；下面的判据遂说明有限域都是完美域。
\symidx{\(F_K\)：特征 \(p\) 域上的 Frobenius 映射}

\subsection{Frobenius 映射的判据}
\label{b1:subsec:1603:02}

\begin{proposition}\label{b1:prop:perfect-frobenius}
设 \(K\) 为特征 \(p>0\) 的域，\(F_K:a\mapsto a^p\) 为其 Frobenius 映射。则 \(K\) 完美，当且仅当 \(F_K\) 满射，即每个 \(c\in K\) 都是 \(K\) 中元素的 \(p\) 次幂。
\end{proposition}
\begin{proof}
先设 Frobenius 满射。关键是：导数为零使多项式只含 \(p\) 的倍数次幂，满射性又使这些项的系数都能开 \(p\) 次方，两者合在一起便使整个多项式成为一个 \(p\) 次幂。具体地，若不可约多项式 \(f\) 不可分，\cref{b1:prop:irreducible-separable-derivative}给出 \(f'=0\)，故 \(f=\sum_i a_ix^{pi}\)。为每个系数取 \(b_i\in K\) 使 \(b_i^p=a_i\)，便有
\[
f=\left(\sum_i b_ix^i\right)^p.
\]
因为 \(f\) 非常数，括号中的多项式也非常数；\(p\ge2\)，所以右边给出两个非常数因子的乘积，与不可约性矛盾。故 \(K\) 完美。

反过来，若 Frobenius 不满射，取 \(c\in K\) 不是 \(p\) 次幂，在代数闭包中取 \(a^p=c\)。由于 \(a^p\in K\)，\(a\) 纯不可分；由\cref{b1:prop:purely-inseparable-criteria}，其最小多项式 \(m_{a,K}\) 形如 \(x^{p^r}-d\)，其中 \(r\ge0\)、\(d\in K\)。另一方面，\(m_{a,K}\) 整除零化 \(a\) 的多项式 \(x^p-c\)，故
\[
\deg m_{a,K}=p^r\le p.
\]
所以 \(r\) 只能为零或一，次数只能为一或 \(p\)。若为一次，则 \(a\in K\)，使 \(c=a^p\) 成为 \(K\) 中的 \(p\) 次幂，与选择矛盾。因此次数为 \(p\)，而 \(m_{a,K}\) 与 \(x^p-c\) 都首一且前者整除后者，必有 \(m_{a,K}=x^p-c\)。这个多项式不可约，又因 \(x^p-c=(x-a)^p\) 而不可分，故 \(K\) 不完美。
\end{proof}
若 \(K\) 是有限域，其 Frobenius 映射是有限集合 \(K\) 到自身的单射，故必为满射；由\cref{b1:prop:perfect-frobenius}，\(K\) 完美。这里的有限性正用于把已经证明的单射性变成开 \(p\) 次方所需的满射性。

\begin{corollary}\label{b1:cor:algebraic-perfect-extension}
设 \(K\) 为完美域，\(E/K\) 为代数扩张，不要求有限。则 \(E\) 仍是完美域。
\end{corollary}
\begin{proof}
给定任意不可约多项式 \(f\in E[x]\)，在 \(E\) 的代数闭包中添入它的一个根 \(a\)。因 \(E/K\) 与 \(E(a)/E\) 都代数，\cref{b1:thm:algebraic-tower}给出 \(E(a)/K\) 代数。\(K\) 完美，故 \(a\) 在 \(K\) 上可分，即 \(m_{a,K}\) 无重根。把它视为 \(E\) 系数多项式后仍零化 \(a\)，所以 \(m_{a,E}\mid m_{a,K}\)，从而 \(m_{a,E}\) 也无重根。最后，\(f\) 是 \(a\) 在 \(E\) 上的不可约零化多项式，故与 \(m_{a,E}\) 相伴；首一化后两者相同。因此 \(f\) 可分，证明 \(E\) 完美。
\end{proof}

当 \(K\) 不完美时，可以在代数闭包中反复加入缺失的 \(p\) 次根。Frobenius 的单射性使这些根没有选择上的歧义；把所有这样的元素收齐，便得到包含 \(K\) 的最小完美子域。

\begin{proposition}\label{b1:prop:perfect-closure}
设 \(K\) 为特征 \(p>0\) 的域，固定 \(K\) 的代数闭包 \(\Omega\)，令
\[
K^{\mathrm{perf}}=\{\,a\in\Omega\mid
\exists r\in\mathbb Z_{\ge0},\ a^{p^r}\in K\,\}.
\]
则 \(K^{\mathrm{perf}}\) 是 \(\Omega\) 中包含 \(K\) 的最小完美子域，且 \(K^{\mathrm{perf}}/K\) 是纯不可分代数扩张。这个子域称为 \(K\) 在 \(\Omega\) 中的\term[wanmei-bibao]{完美闭包}[perfect closure]。
\end{proposition}
\symidx{\(K^{\mathrm{perf}}\)：在选定代数闭包中的完美闭包}
\begin{proof}
取 \(r=0\) 可见 \(K\subseteq K^{\mathrm{perf}}\)。给定 \(a,b\in K^{\mathrm{perf}}\)，把各自的指数提高到同一个 \(r\)，便有 \(a^{p^r},b^{p^r}\in K\)。利用 Frobenius 的加法与乘法相容性，得到
\[
(a-b)^{p^r}=a^{p^r}-b^{p^r}\in K,
\qquad (ab)^{p^r}=a^{p^r}b^{p^r}\in K.
\]
若 \(a\ne0\)，还有 \((a^{-1})^{p^r}=(a^{p^r})^{-1}\in K\)。因此这个集合是子域。每个 \(a\) 满足某个 \(a^{p^r}=c\in K\)，因而被 \(x^{p^r}-c\) 零化，故它在 \(K\) 上代数且纯不可分。

对任意 \(a\in K^{\mathrm{perf}}\)，在 \(\Omega\) 中取 \(b^p=a\)。若 \(a^{p^r}\in K\)，则 \(b^{p^{r+1}}\in K\)，所以 \(b\in K^{\mathrm{perf}}\)。这说明 Frobenius 在该子域上满射，由\cref{b1:prop:perfect-frobenius}，它完美。

最后，设 \(M\) 是 \(\Omega\) 中包含 \(K\) 的完美子域。给定 \(a\in K^{\mathrm{perf}}\)，取 \(a^{p^r}=c\in K\)。在 \(M\) 中反复利用 Frobenius 满射，可找到 \(d\in M\) 使 \(d^{p^r}=c\)。在 \(\Omega\) 中，\((a-d)^{p^r}=0\) 迫使 \(a=d\)，故 \(a\in M\)。于是 \(K^{\mathrm{perf}}\subseteq M\)，证明最小性。
\end{proof}

\subsection{不完美函数域的例子}
\label{b1:subsec:1603:03}

设 \(p\) 为素数，\(t\) 在 \(\mathbb F_p\) 上超越，取 \(K=\mathbb F_p(t)\)。以 \(K^p=\{\,a^p\mid a\in K\,\}\) 记 Frobenius 的像；一般的特征 \(p\) 域完美，当且仅当 \(K^p=K\)。任取 \(A,B\in\mathbb F_p[x]\) 且 \(B\neq0\)，Frobenius 将 \(A(t)/B(t)\) 送到 \(A(t^p)/B(t^p)\)，因为系数在 \(\mathbb F_p\) 中。反过来，每个 \(A(t^p)/B(t^p)\) 都是 \(A(t)/B(t)\) 的 \(p\) 次幂，故
\[
K^p=\mathbb F_p(t^p).
\]
由\eqref{b1:eq:rational-power-degree}，
\[
\bigl[\mathbb F_p(t):\mathbb F_p(t^p)\bigr]=p,
\]
所以 \(K^p\subsetneq K\)，\(K\) 不完美。直接观察也能看出 \(t\) 没有 \(p\) 次根：有理函数的 \(p\) 次幂中，素因子 \(t\) 的指数必被 \(p\) 整除，而 \(t\) 的指数为一。

\ProseParagraphBreak
这说明“完美域的代数扩张仍完美”不能删去“代数”：\(\mathbb F_p\) 完美，但纯超越扩张 \(\mathbb F_p(t)\) 不完美。这个函数域的完美闭包还可逐层写出。固定 \(K\) 的代数闭包 \(\Omega\)，对每个 \(n\ge0\)，在其中取 \(t\) 的唯一 \(p^n\) 次根 \(t^{1/p^n}\)，则
\[
K^{\mathrm{perf}}=\bigcup_{n\geq0}\mathbb F_p(t^{1/p^n})\subseteq\Omega.
\]
这些域逐层包含，因为 \((t^{1/p^{n+1}})^p=t^{1/p^n}\)。为核对这个等式，任取第 \(n\) 层的元素 \(z=A(t^{1/p^n})/B(t^{1/p^n})\)，其中 \(A,B\in\mathbb F_p[x]\)、\(B\neq0\)，便有 \(z^{p^n}=A(t)/B(t)\in K\)，故右侧并域包含在完美闭包中。反过来，若 \(a^{p^n}=A(t)/B(t)\in K\)，那么第 \(n\) 层中的 \(A(t^{1/p^n})/B(t^{1/p^n})\) 是它的 \(p^n\) 次根；由根的唯一性，这个元素就是 \(a\)。因此并域也包含整个完美闭包。\cref{b1:prop:perfect-closure}已说明这个并域完美且对 \(K\) 纯不可分。它补齐的是反复开 \(p\) 次方得到的根：任一层中元素的 \(p\) 次根在下一层中，再迭代便得到它的所有 \(p^r\) 次根。

\ProseParagraphBreak
这样的完美域仍未必代数闭。例如取素数 \(\ell\ne p\)。\(x^\ell-t\) 在 \(\mathbb F_p[t]\) 上对素元 \(t\) 满足 Eisenstein 判据，故在 \(K[x]\) 中不可约；其导数为 \(\ell x^{\ell-1}\ne0\)，所以可分。若它的某个根 \(a\) 属于 \(K^{\mathrm{perf}}\)，则 \(a\) 必在某一层 \(K_n=\mathbb F_p(t^{1/p^n})\) 中。\(t^{1/p^n}\) 仍在 \(\mathbb F_p\) 上超越，应用\eqref{b1:eq:rational-power-degree}得到 \([K_n:K]=p^n\)，而不可约性给出 \([K(a):K]=\ell\)。塔式公式会迫使 \(\ell\mid p^n\)，与 \(\ell\ne p\) 矛盾。因此这个可分多项式的根并未被加入；完美性要求代数扩张可分，不要求每个代数多项式都已在域内分裂。

% !TeX root = ../../Book1.tex
% 纲要标题：### 16.4 本原元素定理
% 纲要位置：计划纲要.md，第 576 行。

\section{本原元素定理}
\label{b1:sec:1604}


% 纲要内容（仅作写作计划，尚非教材正文）：
%
% * 先证明域的有限乘法子群为循环群；第17.2节复用此引理，不以本原元素定理反向证明循环性
% * 完整陈述本原元素定理，并连续证明无限与有限两种基域情形；解释区分全部嵌入的目标与坏参数的来源
% * 证明中的有限性与可分性
% * 非单生成有限扩张的明确反例；由指数余数对证明双变量单项式基
%
% ---
%

一个有限扩张总能由有限多个元素生成，能否进一步只用一个？可分性给出肯定答案。无限基域允许从线性组合中避开有限多个坏参数，从而把不同嵌入区分开；有限基域没有这样的参数选择余地，却能利用扩张域乘法群的循环性。两种办法分别把嵌入计数和乘法生成转化为域的单生成性。

\subsection{有限乘法子群的循环性}
\label{b1:subsec:1604:04}

\begin{lemma}\label{b1:lem:finite-multiplicative-cyclic}
任意域 \(F\) 的乘法群 \(F^\times\) 的有限子群都是循环群。
\end{lemma}
\begin{proof}
设有限子群为 \(G\)。平凡群可直接取生成元 \(1\)，以下设 \(G\ne\{1\}\)，令 \(m\) 为所有元素阶的最小公倍数。每个元素都是 \(x^m-1\) 的根，所以根数上界给出 \(\lvert G\rvert\le m\)。因此只要在 \(G\) 中构造一个阶为 \(m\) 的元素，它的 \(m\) 个幂就会占满整个群。

对每个素数幂 \(\ell^e\) 恰好整除 \(m\)，至少有一个元素 \(g_\ell\) 的阶中含有完整的 \(\ell^e\) 因子：否则所有元素阶中 \(\ell\) 的指数都小于 \(e\)，其最小公倍数也不可能含有 \(\ell^e\)。设该阶为 \(\ell^e d\)，其中 \(\ell\nmid d\)，则 \(h_\ell=g_\ell^d\) 的阶为 \(\ell^e\)：\(h_\ell^r=1\) 当且仅当 \(\ell^e d\mid dr\)，即 \(\ell^e\mid r\)。

\(G\) 交换，且各 \(h_\ell\) 的阶两两互素，所以其乘积 \(h\) 的阶为 \(m\)。为核对这一点，若两个交换元素的阶为互素的 \(a,b\)，且 \((xy)^r=1\)，则 \(x^r=y^{-r}\) 位于 \(\langle x\rangle\cap\langle y\rangle\)。交中元素的阶同时整除 \(a,b\)，故只有单位，从而 \(a\mid r,b\mid r\)，即 \(ab\mid r\)；反向 \((xy)^{ab}=1\) 直接成立。有限次应用此结论便得 \(h\) 的阶为 \(m\)。

由\cref{b1:prop:factor-root}得到的上界 \(\lvert G\rvert\le m\)，与 \(h\) 的 \(m\) 个不同幂给出的下界相等。因此 \(G=\langle h\rangle\)。
\end{proof}

\subsection{有限可分扩张的单生成性}
\label{b1:subsec:1604:01}

\begin{theorem}[本原元素]\label{b1:thm:primitive-element}
每个有限可分扩张 \(L/K\) 都是单扩张。如果 \(\theta\in L\) 满足 \(L=K(\theta)\)，则称 \(\theta\) 为该扩张的\term[benyuan-yuansu]{本原元素}[primitive element]。
\end{theorem}
\begin{proof}
先考虑无限基域 \(K\)。固定包含 \(L\) 的代数闭包 \(\Omega\)，由\cref{b1:thm:separable-embeddings}，\(L\) 有恰好 \(n=[L:K]\) 个 \(K\) 嵌入 \(\sigma_1,\ldots,\sigma_n\)。我们要找到一个元素，使这些嵌入在它上面取值两两不同；这样一个元素的最小多项式就必须容纳 \(n\) 个不同根，其次数才能达到整个扩张的次数。

先设 \(L=K(a,b)\)，考察 \(a+cb\)、\(c\in K\)。对于 \(i\neq j\)，若 \(\sigma_i(b)=\sigma_j(b)\)，必有 \(\sigma_i(a)\neq\sigma_j(a)\)，否则它们在两个生成元上相同，因而为同一个映射。此时
\[
 \sigma_i(a+cb)-\sigma_j(a+cb)=\sigma_i(a)-\sigma_j(a)\ne0,
\]
所以每个 \(c\in K\) 都把这一对嵌入区分开。若 \(\sigma_i(b)\neq\sigma_j(b)\)，等式
\[
\sigma_i(a)+c\cdot\sigma_i(b)=\sigma_j(a)+c\cdot\sigma_j(b)
\]
只可能在
\[
 c=\frac{\sigma_j(a)-\sigma_i(a)}{\sigma_i(b)-\sigma_j(b)}
\]
时成立。这个值若不在 \(K\) 中，就不排除任何系数；若在 \(K\) 中，也只排除一个。嵌入对只有有限多个，坏系数也只有有限多个。这里才用到 \(K\) 无限：这些值不能耗尽 \(K\)，所以可以同时避开它们。

令 \(\theta=a+cb\)。每个 \(\sigma_i\) 固定最小多项式的系数，所以
\[
 m_{\theta,K}(\sigma_i(\theta))
 =\sigma_i\bigl(m_{\theta,K}(\theta)\bigr)=0.
\]
这些值两两不同，故 \(m_{\theta,K}\) 至少有 \(n\) 个不同根。由次数与根数的关系及塔式公式，
\[
n\leq\bigl[K(\theta):K\bigr]\leq[L:K]=n.
\]
所以 \(\bigl[L:K(\theta)\bigr]=1\)，即 \(L=K(\theta)\)。一般情形取有限生成集 \(a_1,\ldots,a_r\)。合并前两个后得到 \(\theta_2\)，满足 \(K(\theta_2)=K(a_1,a_2)\)。因为 \(\theta_2\in L\) 且 \(L/K\) 可分，\(\theta_2\) 仍在 \(K\) 上可分；于是 \(K(\theta_2,a_3)/K\) 仍是有限可分扩张，可以再次应用二生成元情形。这样逐次合并，便得到生成整个 \(L\) 的一个元素。

若基域有限，有限多个坏值有可能占满整个基域，因而不能再靠上面的避让选出系数。记 \(\lvert K\rvert=q\)、\([L:K]=n\)。选一组 \(K\) 基，每个元素由 \(n\) 个 \(K\) 坐标唯一确定，所以 \(\lvert L\rvert=q^n\)，是有限域。\cref{b1:lem:finite-multiplicative-cyclic}说明 \(L^\times=\langle\theta\rangle\)。域 \(K(\theta)\) 包含 \(\theta\) 的全部整数幂，因而包含 \(L\) 的每个非零元素；它也包含零，故等于 \(L\)。这实际上证明了有限域的任意有限扩张都单生成，论证本身不需要可分性。有限域的完美性已由\cref{b1:prop:perfect-frobenius}得到，所以这些扩张也都可分。
\end{proof}
这一结果称为\term[benyuan-yuansu-dingli]{本原元素定理}[primitive element theorem]。有限域乘法群的生成元一定是域扩张的本原元素，反之未必；前者必须用乘法的幂得到所有非零元素，后者允许加法、减法和除法。

\ProseParagraphBreak
例如 \(X^2+1\) 在 \(\mathbb F_3\) 中无根，故 \(F=\mathbb F_3[X]/(X^2+1)\) 是九元素域。令 \(a\) 为 \(X\) 的剩余类，则 \(F=\mathbb F_3(a)\)，但 \(a^2=-1\ne1\)、\(a^4=1\)，所以 \(a\) 的乘法阶为四，不能生成八阶的 \(F^\times\)。

\begin{corollary}\label{b1:prop:primitive-element-infinite}
设 \(K\) 为无限域，\(L/K\) 为有限可分扩张，\(a,b\in L\) 且 \(L=K(a,b)\)。记 \(n=[L:K]\)。则使 \(L\ne K(a+cb)\) 的系数 \(c\in K\) 至多有 \(\binom{n}{2}\) 个。
\end{corollary}
\begin{proof}
取包含 \(L\) 的 \(K\) 的代数闭包 \(\Omega\)，沿用\cref{b1:thm:primitive-element}对无限基域的证明。对 \(n\) 个不同的 \(K\) 嵌入 \(L\rightarrow\Omega\) 中的每一无序对，等式 \(\sigma_i(a+cb)=\sigma_j(a+cb)\) 或无解，或在 \(K\) 中至多有一个解。这样的嵌入对共 \(\binom{n}{2}\) 个；避开这些解时，全部嵌入在 \(a+cb\) 上取值不同，前一定理的次数论证便给出 \(L=K(a+cb)\)。
\end{proof}

例如 \(\theta=\sqrt2+\sqrt3\) 生成 \(\mathbb Q(\sqrt2,\sqrt3)\)。这不只由存在性定理得出，还可直接恢复两个根：\(\theta^{-1}=\sqrt3-\sqrt2\)，从而
\[
\sqrt2=\frac{\theta-\theta^{-1}}2,\qquad
\sqrt3=\frac{\theta+\theta^{-1}}2.
\]

\subsection{有限性与可分性的作用}
\label{b1:subsec:1604:02}

在无限基域的证明中，有限性保证嵌入数有限，因而只有有限多个嵌入对和坏参数；可分性则保证嵌入数恰为 \([L:K]\)。所以，把这些嵌入在一个元素上的值全部分开，才能迫使该元素的最小多项式次数达到 \([L:K]\)。若嵌入数小于次数，这一办法只给出较弱的次数下界，不能推出该元素生成整个扩张。

\ProseParagraphBreak
无限次代数扩张不可能单生成：若 \(L/K\) 是这样的扩张且 \(L=K(a)\)，则 \(a\) 在 \(K\) 上代数，由\cref{b1:prop:minimal-polynomial-degree}知 \([L:K]=\deg m_{a,K}<\infty\)，矛盾。因此，即使扩张可分，也不能省去本原元素定理的有限性假设。例如，\cref{b1:prop:infinite-algebraic-root-tower}中的 \(\bigcup_{n\geq1}\mathbb Q(2^{1/2^n})/\mathbb Q\) 是无限次代数扩张，并因特征为零而可分。

\ProseParagraphBreak
另一方面，可分性是充分条件，并非单生成的必要条件。\(\mathbb F_p(t^{1/p})/\mathbb F_p(t)\) 本来就是单扩张，却纯不可分。因此有限可分推出单生成，单生成却未必可分；而下面的例子表明，有限纯不可分扩张也未必单生成。

\subsection{非单生成有限扩张}
\label{b1:subsec:1604:03}

\begin{proposition}\label{b1:prop:non-simple-finite-extension}
设 \(p\) 为素数，\(u,v\) 在 \(\mathbb F_p\) 上代数独立，并令
\[
L=\mathbb F_p(u,v),\qquad K=\mathbb F_p(u^p,v^p).
\]
则 \([L:K]=p^2\)，但对每个 \(a\in L\) 都有 \(\bigl[K(a):K\bigr]\leq p\)。特别地，这个有限扩张不是单扩张。
\end{proposition}
\begin{proof}
因为 \(u^p,v^p\in K\)，高次幂的指数都能模 \(p\) 降低，所以自然考察 \(p\times p\) 个单项式
\[
 \{\,u^iv^j:0\le i,j<p\,\}.
\]
\(u,v\) 分别满足 \(x^p-u^p\)、\(x^p-v^p\)，所以总次数至多 \(p^2\)。要证明这些单项式是一组基，只需证明它们在 \(K\) 上线性无关。若有关系，清除系数中 \(u^p,v^p\) 的有理函数分母，得到
\[
 \sum_{0\le i,j<p}P_{ij}(u^p,v^p)u^iv^j=0,
 \qquad P_{ij}\in\mathbb F_p[U,V].
\]
展开第 \((i,j)\) 项后，每个单项式的指数对形如 \((pm+i,pn+j)\)。不同 \((i,j)\) 给出不同的模 \(p\) 余数对，所以这些项所含的单项式集合互不相交。\(u,v\) 在 \(\mathbb F_p\) 上代数独立，迫使各 \(P_{ij}\) 的全部系数为零，进而迫使原关系的系数全为零。于是扩张次数恰为 \(p^2\)，这些单项式组成一组 \(K\) 基。

对任意 \(a=A(u,v)/B(u,v)\)，其中 \(A,B\in\mathbb F_p[U,V]\)、\(B(u,v)\ne0\)，特征 \(p\) 使取幂时的交叉项消失，且系数满足 \(c^p=c\)，所以
\[
 a^p=\frac{A(u,v)^p}{B(u,v)^p}
     =\frac{A(u^p,v^p)}{B(u^p,v^p)}\in K.
\]
分母仍非零，因为域中非零元素的幂非零。因此 \(a\) 被次数为 \(p\) 的多项式 \(x^p-a^p\in K[x]\) 零化，给出 \(\bigl[K(a):K\bigr]\leq p<p^2\)。因而任何单个 \(a\) 都不能生成 \(L\)。
\end{proof}
这里所说的两个 \(p\) 次根方向，是指 \(u\) 与 \(v\) 各带来一个次数为 \(p\) 的步骤，总次数为 \(p^2\)，而每个元素单独只能带来至多 \(p\) 的次数。它不是说 \(u,v\) 在 \(K\) 上代数独立；二者在 \(K\) 上都已代数。


\begin{chaptersummary}
不可约多项式有重根，当且仅当其导数为零。在特征 \(p\) 中剥去变量的 \(p\) 次幂，可得到可分的不可约部分；这个过程同时记录不同根数和各根重数。有限扩张的可分次数等于嵌入个数，沿塔相乘；它达到总次数，恰好等价于扩张可分。对有限扩张 \(L/K\)，\(L\) 中在 \(K\) 上可分的元素组成最大可分中间域 \(E\)，而上层扩张 \(L/E\) 纯不可分。

\ProseParagraphBreak
特征零域和有限域都完美。正特征的一般域是否完美，则由 Frobenius 是否满射决定。有限可分扩张中的多个生成元可以合并成一个：无限基域靠避开有限多个坏参数，有限基域靠乘法群的循环性。纯不可分扩张可能单生成，也可能必须使用多个生成元，不能用一个笼统的正特征结论代替具体检查。
\end{chaptersummary}

\BookExercises

\begin{exercise}\label{b1:ex:16-1}
设 \(K\) 的特征为 \(p>0\)。证明 \(x^{p^r}-x-c\) 对每个 \(r\geq1,c\in K\) 都可分。说明这是否意味着它总不可约。
\end{exercise}
\begin{solution}
其形式导数为 \(-1\)，故与原多项式最大公因子为一，由重根判据可分。但不可约性不由此保证，例如取 \(c=0\)，则 \(0\) 是根，次数 \(p^r>1\) 时已有一次因子 \(x\)。可分多项式可以有多个不可约因子；准确地说，在其不可约分解中，每个因子都须可分，且每个因子只出现一次。
\end{solution}

\begin{exercise}\label{b1:ex:16-2}
在 \(K=\mathbb F_3(t)\) 上，证明 \(f=x^6+t\) 不可约，并求其不同根数和每个根的重数。
\end{exercise}
\begin{solution}
把 \(f\) 视为 \(\mathbb F_3[t][x]\) 中多项式，它对素元 \(t\) 满足 Eisenstein 判据，所以在分式域 \(K\) 上不可约。写成 \(g(x^3)\)，其中 \(g(y)=y^2+t\)。若 \(g(y)=r(y)s(y)\)，两个因子都为正次数，就会有 \(f(x)=r(x^3)s(x^3)\) 的非平凡分解，与不可约性矛盾。因此 \(g\) 不可约，又有 \(g'=2y\neq0\)，所以可分。由\cref{b1:prop:inseparable-polynomial-shape}，\(f\) 有两个不同根，每个重数为三。
\end{solution}

\begin{exercise}\label{b1:ex:16-3}
设 \(L/K\) 有限，\(\operatorname{char}K=p>0\)，且 \(p\nmid[L:K]\)。证明 \(L/K\) 可分。说明逆命题是否成立。
\end{exercise}
\begin{solution}
\cref{b1:prop:separable-inseparable-decomposition}说明不可分次数为 \(p\) 的幂，并整除总次数，所以只能为一。于是可分次数等于总次数，由\cref{b1:thm:separable-embeddings}可分。逆命题不成立：在 \(\mathbb F_2\) 上，\(x^2+x+1\) 无根故不可约，导数为一故可分。添入一个根所得扩张次数为二，恰被特征二整除。
\end{solution}

\begin{exercise}\label{b1:ex:16-4}
设 \(p\) 为素数，\(r\ge1\)、\(n>1\) 为整数且 \(p\nmid n\)，\(u\) 在 \(\mathbb F_p\) 上超越。令
\[
 K=\mathbb F_p(u^{np^r}),\quad L=\mathbb F_p(u),\quad
 E=\mathbb F_p(u^{p^r}),\quad F=\mathbb F_p(u^n).
\]
计算两条域塔 \(K\subseteq E\subseteq L\)、\(K\subseteq F\subseteq L\) 各层的次数，并判定各层的可分性或纯不可分性。求 \(L/K\) 的三种次数及最大可分中间域，证明 \(E\cap F=K\)、\(EF=L\)。
\end{exercise}
\begin{solution}
记 \(m=p^r\)。由\eqref{b1:eq:rational-power-degree}，两条塔给出
\[
 [E:K]=n,\quad [L:E]=m,\qquad
 [F:K]=m,\quad [L:F]=n.
\]
\(u^m\) 在 \(K\) 上的最小多项式为 \(X^n-u^{mn}\)，\(u\) 在 \(F\) 上的最小多项式为 \(X^n-u^n\)：这些多项式零化相应元素，且次数与已算出的扩张次数相同。因 \(p\nmid n\) 且常数项非零，它们与导数互素，所以 \(E/K\)、\(L/F\) 可分。又有 \(u^m\in E\)、\((u^n)^m\in K\)，对相应有理函数取 \(m\) 次幂也落入基域，故 \(L/E\)、\(F/K\) 纯不可分。于是可分次数塔式公式给出
\[
 [L:K]=nm,\qquad [L:K]_{\mathrm s}=n,\qquad
 [L:K]_{\mathrm i}=m.
\]
最大可分中间域包含 \(E\)，且在 \(K\) 上的次数也是 \(n\)，所以恰为 \(E\)。

\([E\cap F:K]\) 同时整除 \([E:K]=n\) 与 \([F:K]=m\)，两者互素，故交域的次数为一，即 \(E\cap F=K\)。取整数 \(A,B\) 使 \(Am+Bn=1\)，则
\[
 u=(u^m)^A(u^n)^B\in EF.
\]
负指数在域中合法，因为这些元素非零。因此 \(EF=L\)。本例既有可分层在下、纯不可分层在上的域塔，也有顺序相反的另一条域塔；二者的中间域分别为 \(E\) 与 \(F\)。
\end{solution}

\begin{exercise}\label{b1:ex:16-5}
证明同时可分且纯不可分的代数扩张必为平凡扩张。推论：若 \(E/K\) 可分、\(F/K\) 纯不可分并位于同一扩张域中，则 \(E\cap F=K\)。
\end{exercise}
\begin{solution}
其中任一元素的最小多项式既无重根，又由纯不可分性只有一个不同根，所以只能为一次，元素属于 \(K\)。交域中的元素从 \(E\) 继承可分性，从 \(F\) 继承纯不可分性，故全部属于 \(K\)，反向包含原本成立。
\end{solution}

\begin{exercise}\label{b1:ex:16-6}
设 \(K\subseteq E\subseteq L\)，且 \(E/K\)、\(L/E\) 都是纯不可分扩张。证明 \(L/K\) 纯不可分。若进一步有 \(\operatorname{char}K=p>0\) 且 \([L:K]<\infty\)，证明存在整数 \(r\geq0\)，使每个 \(a\in L\) 都满足 \(a^{p^r}\in K\)。
\end{exercise}
\begin{solution}
若 \(\operatorname{char}K=0\)，两层纯不可分扩张均平凡，结论立即成立。若 \(\operatorname{char}K=p>0\)，给定 \(a\in L\)，先取 \(a^{p^m}\in E\)，再取 \((a^{p^m})^{p^n}\in K\)，便有 \(a^{p^{m+n}}\in K\)。若 \(L=K\)，第二问取 \(r=0\)。否则，在有限情形取生成元 \(a_1,\ldots,a_s\)，分别选 \(r_i\) 使 \(a_i^{p^{r_i}}\in K\)，令 \(r=\max_i r_i\)。任意元素是这些生成元的有理式；取 \(p^r\) 次幂后，系数与全部生成元的幂都在 \(K\) 中，分母仍非零，故整个值在 \(K\) 中。定义只要求对每个元素分别存在指数，这里利用有限生成性得到一个对全部元素同时适用的指数，即从 \(\forall a\,\exists r(a)\) 得到了 \(\exists r\,\forall a\)。
\end{solution}

\begin{exercise}\label{b1:ex:16-7}
若 \(a\) 在特征 \(p>0\) 的域 \(K\) 上代数且可分，证明对每个 \(r\geq1\)，有 \(K(a)=K(a^{p^r})\)。
\end{exercise}
\begin{solution}
\(a\) 被 \(x^{p^r}-a^{p^r}\) 零化，所以 \(K(a)/K(a^{p^r})\) 纯不可分。另一方面，\(a\) 在 \(K\) 上的最小多项式可分，扩大系数域后其新的最小多项式整除原多项式，故该扩张也可分。其最小多项式因此既无重根又只有一个不同根，必须为一次，得到两域相等。
\end{solution}

\begin{exercise}\label{b1:ex:16-8}
设 \(K\) 完美，\(L/K\) 代数。证明 \(L\) 的任意中间域都完美。再举一个超越扩张，说明完美性不向任意扩张域传递。
\end{exercise}
\begin{solution}
任意中间域 \(E/K\) 仍代数，\cref{b1:cor:algebraic-perfect-extension}适用，故 \(E\) 完美。反例为 \(K=\mathbb F_p\)、\(L=\mathbb F_p(t)\)：底域有限故完美，扩张域中的 \(t\) 不是 \(p\) 次幂，因为 \(t\) 的素因子指数为一，而有理函数的 \(p\) 次幂中所有指数均被 \(p\) 整除。由\cref{b1:prop:perfect-frobenius}，后者不完美。
\end{solution}

\begin{exercise}\label{b1:ex:16-9}
设 \(\operatorname{char}K=p>0\)，\(L/K\) 为纯不可分代数扩张，固定包含 \(L\) 的 \(K\) 的代数闭包 \(\Omega\)。证明在 \(\Omega\) 中有 \(K^{\mathrm{perf}}=L^{\mathrm{perf}}\)，并举例说明不能删去纯不可分假设。再取 \(K=\mathbb F_p(t)\)、\(L=\mathbb F_p(t^{1/p})\)，其中 \(t\) 在 \(\mathbb F_p\) 上超越，说明两个基域都不完美，却有同一个完美闭包。
\end{exercise}
\begin{solution}
由\cref{b1:prop:perfect-closure}的最小性，\(L^{\mathrm{perf}}\) 是包含 \(K\) 的完美子域，所以 \(K^{\mathrm{perf}}\subseteq L^{\mathrm{perf}}\)。另一方面，纯不可分性使每个 \(a\in L\) 都有某个 \(p\) 次幂的迭代落入 \(K\)，所以 \(L\subseteq K^{\mathrm{perf}}\)。后者是包含 \(L\) 的完美子域，再用最小性得到 \(L^{\mathrm{perf}}\subseteq K^{\mathrm{perf}}\)，两者相等。

取 \(K=\mathbb F_3\)、\(L=\mathbb F_3(a)\)，其中 \(a^2=-1\)，则 \([L:K]=2\)。两域都有限，因而已完美，所以各自的完美闭包就是自身，彼此不等；这说明一般代数扩张不能代替纯不可分扩张。

在题给函数域中，Frobenius 的像分别为 \(\mathbb F_p(t^p)\subsetneq K\) 与 \(\mathbb F_p(t)\subsetneq L\)，严格性均由\eqref{b1:eq:rational-power-degree}给出，所以两个域都不完美。其共同完美闭包为
\[
 \bigcup_{j\ge0}\mathbb F_p(t^{1/p^j}).
\]
这也说明完美域的子域即使只相差代数扩张，仍可能不完美。
\end{solution}

\begin{exercise}\label{b1:ex:16-10}
令 \(a=\sqrt2+\sqrt3\)、\(b=\sqrt2-\sqrt3\)。确定所有 \(c\in\mathbb Q\)，使 \(a+cb\) 不能生成 \(\mathbb Q(\sqrt2,\sqrt3)\)。
\end{exercise}
\begin{solution}
因为 \(a+cb=(1+c)\sqrt2+(1-c)\sqrt3\)，其四个可能的共轭值为
\[
\pm(1+c)\sqrt2\pm(1-c)\sqrt3.
\]
当 \(c\neq\pm1\) 时，\(1+c\) 与 \(1-c\) 均非零；\(\sqrt2\) 与 \(\sqrt3\) 在 \(\mathbb Q\) 上线性无关，故上述四个值两两不同。由\cref{b1:prop:biquadratic-example}，全扩张次数为四，于是 \(a+cb\) 生成全域。若 \(c=1\)，则 \(a+cb=2\sqrt2\)；若 \(c=-1\)，则 \(a+cb=2\sqrt3\)。两者都只生成二次子域，故所求参数恰为 \(c=\pm1\)。
\end{solution}

\begin{exercise}\label{b1:ex:16-11}
设 \(u,v\) 在 \(\mathbb F_p\) 上代数独立，\(K=\mathbb F_p(u^p,v^p)\)、\(L=\mathbb F_p(u,v)\)。证明对每个 \(c\in K\)，域 \(K(u+cv)\) 的次数为 \(p\)，且这些域两两不同。由此得到一个有限扩张有无限多个中间域的例子。
\end{exercise}
\begin{solution}
由\cref{b1:prop:non-simple-finite-extension}的单项式基，\(1,u,v\) 在 \(K\) 上线性无关，所以 \(u+cv\notin K\)。它的 \(p\) 次幂在 \(K\) 中，最小多项式的次数只能为一或 \(p\)，故恰为 \(p\)。若 \(c\neq d\) 且两个域相同，则该域同时包含 \(u+cv,u+dv\)，作差除以 \(c-d\) 得 \(v\)，进而包含 \(u\)。这使它等于 \(L\)，却分别应有次数 \(p\) 与 \(p^2\)，矛盾。\(K\) 无限，因而得到无限多个不同中间域。这里许多中间域来自两个纯不可分生成元的不同组合；整个扩张却只有包含映射这一个 \(K\) 嵌入，中间域的多少并不等于可分次数的大小。
\end{solution}

\begin{exercise}\label{b1:ex:16-12}
令 \(F=\mathbb F_3[X]/(X^2+1)\)，\(a\) 为 \(X\) 的剩余类。枚举 \(F/\mathbb F_3\) 的全部本原元素与 \(F^\times\) 的全部生成元，比较两者的数量，并求其余非零元素的乘法阶。
\end{exercise}
\begin{solution}
每个元素唯一写成 \(b+ca\)，其中 \(b,c\in\mathbb F_3\)。若 \(c\ne0\)，由 \(a=(b+ca-b)/c\) 可恢复 \(a\)，故该元素生成整个 \(F\)；若 \(c=0\)，则元素已在基域中。因此域扩张的本原元素恰为
\[
 \{\,b+ca:b\in\mathbb F_3,\ c\in\mathbb F_3^\times\,\},
\]
共有六个。当 \(b,c\) 均非零时，\(b^2=c^2=1\)，所以
\[
 (b+ca)^2=2bca,\qquad (b+ca)^4=-1,\qquad(b+ca)^8=1.
\]
这些元素的阶恰为八，正是乘法群的四个生成元 \(1+a,1-a,-1+a,-1-a\)。其余非零元素为 \(1,-1,a,-a\)，阶依次为 \(1,2,4,4\)。所以两个 \(\pm a\) 都是域本原元素，却不是乘法群生成元。
\end{solution}

\begin{exercise}\label{b1:ex:16-13}
设 \(E/K\)、\(F/K\) 为同一代数闭包中的可分代数扩张，不要求有限。证明合成域 \(EF/K\) 可分。
\end{exercise}
\begin{solution}
合成域中每个元素只使用 \(E\cup F\) 中有限多个元素的域运算，因而落在由有限多个对 \(K\) 可分的元素生成的域中。\cref{b1:thm:separable-embeddings}说明该有限扩张可分，所以所取元素在 \(K\) 上可分。对全部元素均成立，故合成域可分。
\end{solution}

\begin{exercise}\label{b1:ex:16-14}
设 \(u\) 在 \(\mathbb F_3\) 上超越，令 \(K=\mathbb F_3(u^6)\)、\(L=\mathbb F_3(u)\)。计算 \(L/K\) 的三种次数及最大可分中间域，证明该扩张正规但不可分。求全部 \(K\) 自同构，以及被这些自同构逐点固定的元素所成子域，并比较后者与最大可分中间域。
\end{exercise}
\begin{solution}
由\eqref{b1:eq:rational-power-degree}，\([L:K]=6\)，所以 \(X^6-u^6\) 是 \(u\) 在 \(K\) 上的最小多项式。特征为三，故
\[
 X^6-u^6=(X^2-u^2)^3=(X-u)^3(X+u)^3.
\]
不同根为 \(u,-u\)，所以可分次数为二、不可分次数为三。令 \(E=\mathbb F_3(u^3)\)，则 \([E:K]=2\)，且 \(u^3\) 在 \(K\) 上的最小多项式 \(Y^2-u^6\) 可分。因此 \(E\) 是最大可分中间域。上述最小多项式在 \(L\) 中分裂，且其根生成 \(L\)，由\cref{b1:thm:normal-embedding-criterion}，\(L/K\) 正规；可分次数小于总次数，故不可分。

自同构必须把 \(u\) 送到 \(u\) 或 \(-u\)，这两种代入都定义了 \(\mathbb F_3(u)\) 的自同构并固定 \(u^6\)，所以自同构群恰有两个元素。令 \(F=\mathbb F_3(u^2)\)，非恒等自同构固定 \(F\)。又有 \([L:F]=2\)，所以每个元素唯一写成 \(A+Bu\)、\(A,B\in F\)。它被 \(u\mapsto-u\) 固定，当且仅当 \(A+Bu=A-Bu\)，即 \(B=0\)。因此固定元素子域恰为 \(F\)。最后，\([F:K]=3\)、\([E:K]=2\)，二者不同；自同构数为二，也小于扩张总次数六。
\end{solution}

\begin{exercise}\label{b1:ex:16-15}
设 \(L/K\) 有限可分，\(E\) 为中间域。证明给定的每个 \(K\) 嵌入 \(E\rightarrow\Omega\) 恰有 \([L:E]\) 个延拓至 \(L\)，其中 \(\Omega\) 为共同代数闭包。
\end{exercise}
\begin{solution}
\(L/E\) 可分，由\cref{b1:thm:separable-embeddings}，固定 \(E\) 的 \(L\) 的嵌入恰有 \([L:E]\) 个。给定 \(\tau:E\rightarrow\Omega\)，将它延拓为 \(\Omega\) 的 \(K\) 自同构 \(\widehat\tau\)。复合 \(\rho\mapsto\widehat\tau\circ\rho\) 把前述嵌入双射地对应到限制为 \(\tau\) 的嵌入，逆由 \(\widehat\tau^{-1}\) 复合给出。所以每个 \(\tau\) 的延拓数都恰为 \([L:E]\)。这正是可分次数塔式证明中固定一个限制后的纤维计数；可分性把一般延拓定理保证的存在提升成了精确的延拓数。
\end{solution}

% !TeX root = ../../Book1.tex
% 纲要标题：## 第17章　有限域
% 纲要位置：计划纲要.md，第 585 行。

\chapter{有限域}
\label{b1:ch:17}
\BookChapterNumber{b1:ch:17}

\begin{chapterabstract}
有限域的阶恰为素数幂，每个素数幂都决定唯一的有限域同构类型。本章证明这些事实，并通过乘法群和 Frobenius 自同构研究子域与共轭元；不可约多项式的计数和商域运算把结构结论用于具体计算。
\end{chapterabstract}

\begin{learninggoals}
\begin{itemize}
\item 说明有限域为何恰有素数幂个元素，区分同构意义下的唯一与环境域中的唯一子域。
\item 计算乘法本原元、Frobenius 的阶、有限域的子域与自同构。
\item 证明并使用不可约多项式计数公式，熟练进行具体商域中的运算。
\item 从共轭轨道计算最小多项式，并理解有限域方法在分解和纠错中的作用。
\end{itemize}
\end{learninggoals}

在 \(\mathbb F_2[t]/(t^2+t+1)\) 中，记 \(t\) 的剩余类为 \(\alpha\)，则 \(\alpha^2=\alpha+1\)。所有元素都能写成 \(0,1,\alpha,1+\alpha\)，并且 \(\alpha^3=1\)。从一个不可约多项式可以造出四元素域；反过来，若有另一个四元素域，它是否只能是同一个域的另一种写法？考察每个元素满足的方程 \(X^4-X=0\)，便能把构造、唯一性和自同构联系起来。一般素数幂阶有限域也由这一思路控制。

\ProseParagraphBreak

本章使用有限扩张的基与塔式公式、分裂域的存在唯一性，以及特征 \(p\) 下的幂映射。乘法群循环性的主要引理已在\cref{b1:lem:finite-multiplicative-cyclic}证明。末节计算有限域扩张的迹和范数；统一定义与一般性质见下一章。

\ProseParagraphBreak
有限域的存在、子域、共轭元以及迹和范数，可与 Shoup 的《A Computational Introduction to Number Theory and Algebra》\cite[Chapter 19]{Shoup2008NumberTheoryAlgebra}对照；多项式的平方自由分解见该书\cite[Section 20.3]{Shoup2008NumberTheoryAlgebra}。求值码中带错误的多项式插值见该书\cite[Section 17.5.2]{Shoup2008NumberTheoryAlgebra}。Roman 的《Field Theory》\cite[Chapter 9; Section 10.1]{Roman2006FieldTheory}则可供从域扩张的角度继续学习。

% !TeX root = ../../Book1.tex
% 纲要标题：### 17.1 有限域的存在与唯一性
% 纲要位置：计划纲要.md，第 587 行。

\section{有限域的存在与唯一性}
\label{b1:sec:1701}

% 纲要内容（仅作写作计划，尚非教材正文）：
%
% * 素数幂阶的必要性
% * 分裂域构造有限域；由 Frobenius 迭代说明根集本身闭成子域
% * 给定阶的有限域在同构意义下唯一
%

有限域不能有任意多个元素。例如六个元素的交换环 \(\mathbb Z/6\mathbb Z\) 含有零因子，因而不是域。限制域的大小的首先是素域和向量空间维数；反过来，给定一个允许的大小，分裂域把所需的全部元素一次构造出来。

\subsection{有限域阶的素数幂条件}
\label{b1:subsec:1701:01}

设 \(F\) 是有限域。它的特征不可能为零，否则整数的像给出无限多个不同元素。其特征是某个素数 \(p\)，素域为 \(\mathbb F_p\)。从线性无关元逐次扩充，若它们尚不生成 \(F\)，就选一个不在其线性生成空间中的元素加入。每次选入的元素都不同，而 \(F\) 是有限集合，所以这个过程必定停止，所得有限线性无关集又生成 \(F\)。故 \(n=[F:\mathbb F_p]\) 是正整数。取一组基 \(e_1,\ldots,e_n\)，每个元素唯一写成
\[
a_1e_1+\cdots+a_ne_n,\qquad a_i\in\mathbb F_p.
\]
基表示的存在性与唯一性说明，\(F\) 的元素与坐标向量 \((a_1,\ldots,a_n)\in\mathbb F_p^n\) 一一对应。每个坐标有 \(p\) 种选择，所以 \(|F|=p^n\)。特别地，六元域、十二元域均不存在；八元域与九元域则尚须构造，不能由必要条件直接断言存在。

\ProseParagraphBreak
更一般地，若 \(K\subseteq F\) 是有限域，\(|K|=q\)，\([F:K]=m\)，同样的基计数给出 \(|F|=q^m\)。这里 \(q\) 本身已经是素数幂。

\subsection{有限域的分裂域构造}
\label{b1:subsec:1701:02}

\begin{theorem}\label{b1:thm:finite-field-existence}
对每个素数 \(p\) 和每个正整数 \(n\)，存在一个恰有 \(p^n\) 个元素的域。它可取为 \(x^{p^n}-x\) 在 \(\mathbb F_p\) 上的分裂域。
\end{theorem}
\begin{proof}
记 \(q=p^n\)，取 \(x^q-x\) 的一个分裂域 \(L\)，并令 \(F\) 为它在 \(L\) 中的根集。需要证明这些根在域运算下仍然是根，才能把根集看成一个域。特征 \(p\) 下二项式展开给出 \((a+b)^p=a^p+b^p\)，再次取 \(p\) 次幂便有
\[
(a+b)^{p^2}=(a^p+b^p)^p=a^{p^2}+b^{p^2}.
\]
这样迭代 \(n\) 次，得到 \((a+b)^q=a^q+b^q\)。取幂也保持乘法和单位，所以 \(a\mapsto a^q\) 是环同态；它因而保持加法逆元，给出 \((-a)^q=-a^q\)。因此若 \(a^q=a,b^q=b\)，则
\[
(a+b)^q=a+b,\qquad (ab)^q=ab,\qquad (-a)^q=-a.
\]
零和一属于 \(F\)；若 \(a\ne0\)，则 \((a^{-1})^q=(a^q)^{-1}=a^{-1}\)。所以 \(F\) 对域的运算封闭，是 \(L\) 的子域，并且包含素域 \(\mathbb F_p\)。它已经含有全部根，便含有这些根在 \(\mathbb F_p\) 上生成的域，而这个域就是分裂域 \(L\)；结合 \(F\subseteq L\)，得到 \(L=F\)。另一方面，\(q=p^n\) 在特征 \(p\) 中为零，所以 \(x^q-x\) 的导数为 \(-1\)，没有重根；它在 \(L\) 中分裂为 \(q\) 个互异的一次因子，所以 \(|F|=q\)。
\end{proof}
例如 \(x^2+x+1\) 在 \(\mathbb F_2\) 上没有根，因而不可约。在 \(\mathbb F_2[x]\) 中有
\[
x^4-x=x(x+1)(x^2+x+1).
\]
商域 \(\mathbb F_2[x]/(x^2+x+1)\) 有四个元素。记 \(x\) 的剩余类为 \(\alpha\)，则不可约二次因子的两个根为 \(\alpha,\alpha+1\)，加上零和一恰为全部四个元素。这个商域由 \(\alpha\) 生成，也是 \(x^4-x\) 的分裂域，正是上述构造在 \(q=4\) 时的一种具体表示。

\subsection{给定阶有限域的唯一性}
\label{b1:subsec:1701:03}

\begin{theorem}\label{b1:thm:finite-field-unique}
设 \(p\) 为素数，\(n\geq1\) 为整数，并令 \(q=p^n\)。任意两个具有 \(q\) 个元素的域，都在其素域 \(\mathbb F_p\) 上同构。
\end{theorem}
\begin{proof}
若 \(|F|=q=p^n\)，由已证的素数幂条件及整数素因子分解的唯一性，\(F\) 的素域是 \(\mathbb F_p\)。其乘法群 \(F^\times\) 有 \(q-1\) 个元素。由\cref{b1:thm:lagrange}，每个非零 \(a\in F\) 都满足 \(a^{q-1}=1\)，故 \(F\) 的全部 \(q\) 个元素都是 \(x^q-x\) 的根。次数为 \(q\) 的多项式不可能有更多根，因此这些元素就是它在任意扩张域中的全部根。根集既然等于整个 \(F\)，由根在素域上生成的域也等于 \(F\)，所以 \(F\) 是该多项式在 \(\mathbb F_p\) 上的分裂域。\cref{b1:prop:splitting-field-unique}保证同一个多项式的两个分裂域在保持基域的同构意义下唯一，故两个这样的域同构。
\end{proof}
以后用 \(\mathbb F_q\) 表示任一具有 \(q\) 个元素的域。
\symidx{\(\mathbb F_q\)：具有 \(q\) 个元素的有限域}
这个记号指定的是同构类型，并不指定某一组元素名称，也不指定唯一同构。例如四元域中的两个非素域元素可以交换。未选定共同的环境域时，两个这样的符号也不自带集合包含关系。需要比较多个有限域的包含关系时，本书先固定 \(\overline{\mathbb F}_p\)，再把 \(\mathbb F_{p^n}\) 取为其中 \(x^{p^n}-x\) 的根集；此时它才是环境域中的确定子域，后面的包含、交与合成均在这个意义下使用。

% !TeX root = ../../Book1.tex
% 纲要标题：### 17.2 乘法群与 Frobenius
% 纲要位置：计划纲要.md，第 593 行。

\section{乘法群与 Frobenius}
\label{b1:sec:1702}

% 纲要内容（仅作写作计划，尚非教材正文）：
%
% * 有限域乘法群的循环性
% * Frobenius 自同构
% * 有限域的子域与扩张；区分域生成元与乘法本原元，解释交域与合成域的次数
% * 乘法群循环性的主要引理已在第16.4节建立；本节用它计算有限域的生成元并组织 Frobenius 及子域的应用。
%

有限域的加法来自特征为 \(p\) 的向量空间，乘法却有更强的结构：全部非零元素可以由一个元素的幂列出。与此相配合，Frobenius 的迭代既描述共轭，也能直接辨认子域。

\subsection{有限域乘法群的循环性}
\label{b1:subsec:1702:01}

\begin{corollary}\label{b1:cor:finite-field-cyclic}
设 \(q\) 为素数幂。则有限域 \(\mathbb F_q\) 的乘法群是 \(q-1\) 阶循环群。
\end{corollary}
\begin{proof}
\cref{b1:lem:finite-multiplicative-cyclic}已经证明：任意域 \(K\) 的乘法群 \(K^\times\) 的每个有限子群都是循环群。取 \(K=\mathbb F_q\)、\(G=\mathbb F_q^\times\)，则 \(G\) 有 \(q-1\) 个元素，故 \(G\cong C_{q-1}\)。
\end{proof}
设 \(q=p^r\) 为素数幂，\(g\in\mathbb F_q^\times\)。如果 \(g\) 生成乘法群 \(\mathbb F_q^\times\)，则称 \(g\) 为有限域 \(\mathbb F_q\) 的\term[benyuan-yuan-youxianyu]{本原元}[primitive element of a finite field]。这里的意思强于“生成域扩张的元素”：乘法本原元 \(g\) 的幂给出全部非零元素，所以当然有 \(\mathbb F_p(g)=\mathbb F_q\)。域生成则只要求 \(g\) 不属于 \(\mathbb F_q\) 的任何真子域：若 \(\mathbb F_p(g)\) 是真子域，它本身就包含 \(g\)；反过来，含有 \(g\) 的任意子域都包含 \(\mathbb F_p(g)\)。这个条件并不要求 \(g\) 的乘法阶为 \(q-1\)。例如在 \(\mathbb F_9=\mathbb F_3(u)\)、\(u^2=-1\) 中，\(u\) 生成二次扩张，但其乘法阶为四，小于八。

\ProseParagraphBreak
循环群中，若 \(g\) 的阶为 \(q-1\)，则 \(g^a\) 生成整个群当且仅当 \(\gcd(a,q-1)=1\)。因此本原元的个数为 \(\varphi(q-1)\)，其中 \(\varphi(m)\) 表示模 \(m\) 与 \(m\) 互素的剩余类个数。

\subsection{Frobenius 自同构}
\label{b1:subsec:1702:02}

在特征 \(p\) 的域中，\(a\mapsto a^p\) 保持加法、乘法和单位，并且单射；有限域上的单射必为满射，所以它是自同构。对 \(\mathbb F_{q^n}/\mathbb F_q\)，其中 \(q=p^r\)，使用其第 \(r\) 次迭代
\[
\sigma:\mathbb F_{q^n}\longrightarrow\mathbb F_{q^n},\qquad a\longmapsto a^q.
\]
对每个 \(c\in\mathbb F_q\)，都有 \(c^q=c\)，所以 \(\sigma\) 逐点固定 \(\mathbb F_q\)，确是这个扩张的自同构，称为扩张 \(\mathbb F_{q^n}/\mathbb F_q\) 的 Frobenius 自同构。由于每个 \(a\in\mathbb F_{q^n}\) 都满足 \(a^{q^n}=a\)，有 \(\sigma^n=1\)。要排除 \(\sigma\) 的阶比 \(n\) 小，只须排除某个 \(0<d<n\) 满足 \(\sigma^d=1\)。后一个等式要求 \(a^{q^d}=a\) 对全部 \(q^n\) 个元素成立，即它们全是 \(x^{q^d}-x\) 的根；但其次数仅为 \(q^d<q^n\)，不可能有这么多根。因此 \(\sigma\) 的阶恰为 \(n\)。其不动元为 \(x^q-x\) 在 \(\mathbb F_{q^n}\) 中的全部根，恰组成 \(\mathbb F_q\)。

\begin{theorem}\label{b1:thm:finite-field-automorphisms}
设 \(q\) 为素数幂，\(n\geq1\) 为整数。则 \(\mathbb F_{q^n}\) 上保持 \(\mathbb F_q\) 中每个元素不变的自同构恰为 \(1,\sigma,\ldots,\sigma^{n-1}\)，其中 \(\sigma(a)=a^q\)。
\end{theorem}
\begin{proof}
由 \(\sigma\) 的阶为 \(n\)，所列 \(n\) 个自同构互异。为了证明不存在其他自同构，只需找一个生成域扩张的元素，再限制它的像。取乘法本原元 \(g\)，因为它的幂给出全部非零元素，便自动有 \(\mathbb F_{q^n}=\mathbb F_q(g)\)。基计数说明 \([\mathbb F_{q^n}:\mathbb F_q]=n\)，故\cref{b1:prop:minimal-polynomial-degree}给出 \(\deg m_{g,\mathbb F_q}=n\)。任一 \(\mathbb F_q\) 自同构 \(\tau\) 都由 \(\tau(g)\) 唯一决定；又因它固定最小多项式的系数，对 \(m_{g,\mathbb F_q}(g)=0\) 施以 \(\tau\)，得到 \(m_{g,\mathbb F_q}(\tau(g))=0\)。次数为 \(n\) 的多项式至多有 \(n\) 个根，所以这样的自同构至多有 \(n\) 个。这就是此前用生成元的共轭像计数嵌入的方法，已列出的 \(n\) 个自同构因而包含全部。
\end{proof}

\subsection{有限域的子域与扩张}
\label{b1:subsec:1702:03}

\begin{theorem}\label{b1:thm:finite-field-subfields}
设 \(q\) 为素数幂，\(m,n\) 为正整数，并在 \(\mathbb F_q\) 的一个固定代数闭包中考虑这些有限域。\(\mathbb F_{q^n}\) 中包含 \(\mathbb F_q\) 的子域恰为 \(\mathbb F_{q^d}\)，其中 \(d\mid n\)，每个这样的子域唯一。此外
\[
\mathbb F_{q^m}\cap\mathbb F_{q^n}=\mathbb F_{q^{\gcd(m,n)}},\qquad
\mathbb F_{q^m}\mathbb F_{q^n}=\mathbb F_{q^{\operatorname{lcm}(m,n)}}.
\]
\end{theorem}
\begin{proof}
若 \(E\) 是所述子域，记 \([E:\mathbb F_q]=d\)。\cref{b1:thm:tower-law}给出 \(d\mid n\)，而 \(|E|=q^d\)。对非零 \(a\in E\)，由 \(E^\times\) 的阶为 \(q^d-1\) 得 \(a^{q^d}=a\)，零也满足此式。因此 \(E\) 包含于固定代数闭包中的 \(x^{q^d}-x\) 的根集。根集恰有 \(q^d\) 个元素，与 \(E\) 的基数相同，故两者相等，这证明了子域的唯一性。反过来，若 \(d\mid n\)，任意满足 \(a^{q^d}=a\) 的元素在迭代 \(n/d\) 次后满足 \(a^{q^n}=a\)，所以代数闭包中的 \(\mathbb F_{q^d}\) 包含于 \(\mathbb F_{q^n}\)。

令 \(g=\gcd(m,n)\)，并记交域 \(E=\mathbb F_{q^m}\cap\mathbb F_{q^n}\) 在 \(\mathbb F_q\) 上的次数为 \(e\)。由两条域塔，\(e\mid m\) 且 \(e\mid n\)，所以 \(e\mid g\)。另一方面，\(g\mid m,n\) 使 \(\mathbb F_{q^g}\) 同时包含于两个域，因而包含于 \(E\)；塔式公式又给出 \(g\mid e\)。故 \(e=g\)，结合已证的子域唯一性，得到交域公式。

令 \(\ell=\operatorname{lcm}(m,n)\)，并记合成域 \(M=\mathbb F_{q^m}\mathbb F_{q^n}\) 在 \(\mathbb F_q\) 上的次数为 \(s\)。两个域都包含于 \(\mathbb F_{q^\ell}\)，所以由它们生成的 \(M\) 也包含于该域，塔式公式给出 \(s\mid\ell\)。另一方面，\(M\) 包含这两个域，故 \(m\mid s\) 且 \(n\mid s\)，于是 \(\ell\mid s\)。所以 \(s=\ell\)，再由子域唯一性得到合成域公式。
\end{proof}
子域定理把有限域的包含关系化为指数的整除关系：\(\mathbb F_{q^a}\subseteq\mathbb F_{q^b}\) 当且仅当 \(a\mid b\)，交域对应最大公因数，合成域对应最小公倍数。例如 \(\mathbb F_{64}\) 包含 \(\mathbb F_2,\mathbb F_4,\mathbb F_8,\mathbb F_{64}\)，但不包含 \(\mathbb F_{16}\)；\(\mathbb F_4\cap\mathbb F_8=\mathbb F_2\)，而 \(\mathbb F_4\mathbb F_8=\mathbb F_{64}\)。不能仅凭域的元素个数较小就断定存在包含关系。

\ProseParagraphBreak
这也让我们能直接检验域生成与乘法群生成的区别。取 \(\mathbb F_{16}^\times\) 的一个生成元 \(g\)，其阶为十五，故 \(g^3\) 的阶为五。\(\mathbb F_{16}\) 的真子域只有 \(\mathbb F_2\) 和 \(\mathbb F_4\)，它们的乘法群阶分别为一和三，都不可能含有五阶元。所以 \(g^3\) 不在任何真子域中，从而 \(\mathbb F_2(g^3)=\mathbb F_{16}\)；但 \(g^3\) 只生成五阶乘法子群，仍不是有限域的乘法本原元。

\subsection{乘法群生成元的计算}
\label{b1:subsec:1702:04}

设 \(N=q-1\)，且已分解出 \(N\) 的不同素因子 \(\ell_1,\ldots,\ell_s\)。元素 \(g\in\mathbb F_q^\times\) 是本原元当且仅当
\begin{equation}\label{b1:eq:primitive-finite-field-test}
g^{N/\ell_j}\ne1\qquad(1\leq j\leq s).
\end{equation}
事实上，其阶 \(d\) 整除 \(N\)；若 \(d<N\)，则 \(N/d>1\)，取 \(N/d\) 的任意一个素因子 \(\ell_j\)，便有 \(d\mid N/\ell_j\)，所以 \(g^{N/\ell_j}=1\)，违反检验。这个 \(\ell_j\) 必也是 \(N\) 的素因子，故检查列出的不同素因子已能发现每一种 \(d<N\) 的情形，不必再检查 \(\ell_j^2,\ell_j^3\) 等素数幂。反过来，若 \(d=N\)，因为 \(N/\ell_j<N\)，这些幂都不可能为一。若 \(N=1\)，没有素因子，唯一的非零元就是一，空的检验也与结论相符。

\ProseParagraphBreak
在 \(\mathbb F_9=\mathbb F_3(u)\)、\(u^2=2\) 中，取 \(g=1+u\)。直接计算
\[
g^2=2u,\qquad g^4=2,\qquad g^8=1.
\]
八的唯一素因子为二，而 \(g^4\ne1\)，所以 \(g\) 是本原元。这也给出一个实际搜索办法：逐个枚举非零元素，用重复平方计算上述幂，直到找到全部检验通过的元素。定理保证搜索终会成功；若不知道 \(q-1\) 的素因子，还须先处理整数分解，不能把这一步的代价忽略。

% !TeX root = ../../Book1.tex
% 纲要标题：### 17.3 不可约多项式
% 纲要位置：计划纲要.md，第 600 行。

\section{不可约多项式}
\label{b1:sec:1703}

% 纲要内容（仅作写作计划，尚非教材正文）：
%
% * 不可约多项式次数与有限域扩张
% * 不可约多项式的计数公式；从递归计数到除数关系上的 Möbius 反演
% * 有限域的显式表示
%

有限域和不可约多项式可以互相构造：不可约多项式给出商域，有限域的元素则提供不可约多项式的根。计数公式的关键不是枚举所有系数组合，而是把 \(x^{q^n}-x\) 按其不可约因子的次数分组。

\subsection{不可约多项式次数与有限域扩张}
\label{b1:subsec:1703:01}

设 \(q\) 为素数幂，\(f\in\mathbb F_q[x]\) 为首一不可约多项式，\(\deg f=d\geq1\)。令 \(A=\mathbb F_q[x]/(f)\)，记 \(\bar x\) 为 \(x\) 的剩余类。不可约性使 \(A\) 成为域；由\cref{b1:prop:minimal-polynomial-degree}，\(1,\bar x,\ldots,\bar x^{d-1}\) 是它的一组 \(\mathbb F_q\) 基，故 \(A\) 有 \(q^d\) 个元素。其乘法群的阶为 \(q^d-1\)，由\cref{b1:thm:lagrange}，每个非零元素满足 \(a^{q^d-1}=1\)，连同零便得到对所有 \(a\in A\) 都有 \(a^{q^d}=a\)。这样 \(X^{q^d}-X\) 已在 \(A\) 中找到 \(q^d\) 个不同的根，达到其次数；它的全部根就是 \(A\) 的全部元素，当然生成 \(A\)。所以 \(A\) 是这个多项式在 \(\mathbb F_q\) 上的分裂域。

\ProseParagraphBreak
在已选定的包含关系 \(\mathbb F_q\subseteq\mathbb F_{q^d}\) 下，后者也是 \(X^{q^d}-X\) 在 \(\mathbb F_q\) 上的分裂域。由\cref{b1:prop:splitting-field-unique}，存在固定 \(\mathbb F_q\) 的同构 \(\varphi:A\xrightarrow{\sim}\mathbb F_{q^d}\)。于是 \(f(\varphi(\bar x))=\varphi(f(\bar x))=0\)，即 \(f\) 在指定的 \(\mathbb F_{q^d}\) 中有根。

\ProseParagraphBreak
若 \(\alpha\) 是 \(f\) 在固定代数闭包中的一个根，则 \(f=m_{\alpha,\mathbb F_q}\)，所以 \(d=[\mathbb F_q(\alpha):\mathbb F_q]\)。给定正整数 \(n\)，根 \(\alpha\) 是否属于 \(\mathbb F_{q^n}\)，就是问它生成的 \(q^d\) 元子域是否包含于这个 \(q^n\) 元域中；\cref{b1:thm:finite-field-subfields}把这个问题化为 \(d\mid n\)。而 \(\alpha\in\mathbb F_{q^n}\) 又等价于 \(\alpha^{q^n}-\alpha=0\)，即它的最小多项式 \(f\) 整除 \(x^{q^n}-x\)。

\begin{proposition}\label{b1:prop:finite-field-factorization}
设 \(q\) 为素数幂，\(n\geq1\) 为整数。首一不可约多项式 \(f\in\mathbb F_q[x]\) 整除 \(x^{q^n}-x\)，当且仅当 \(\deg f\mid n\)。从而
\begin{equation}\label{b1:eq:finite-field-polynomial-product}
x^{q^n}-x=\prod_{d\mid n}\ \prod_{\substack{f\text{ 首一不可约}\\\deg f=d}}f(x),
\end{equation}
其中每个因子只出现一次。
\end{proposition}
\begin{proof}
在 \(\mathbb F_q\) 的固定代数闭包中取 \(f\) 的一个根 \(\alpha\)，令 \(d=\deg f\)。若 \(f\mid x^{q^n}-x\)，则 \(\alpha\in\mathbb F_{q^n}\)，其生成的 \(q^d\) 元子域由\cref{b1:thm:finite-field-subfields}满足 \(d\mid n\)。若 \(d\mid n\)，则 \(\alpha\in\mathbb F_q(\alpha)=\mathbb F_{q^d}\subseteq\mathbb F_{q^n}\)，所以 \(\alpha^{q^n}=\alpha\)，其最小多项式 \(f\) 必整除 \(x^{q^n}-x\)。最后，这个多项式的导数为 \(-1\)，所以没有重复的不可约因子。它及所选的每个不可约因子都首一，故乘积前的常数只能为一，得到\eqref{b1:eq:finite-field-polynomial-product}。
\end{proof}

\subsection{不可约多项式的计数公式}
\label{b1:subsec:1703:02}

对正整数 \(d\)，记 \(N_q(d)\) 为 \(\mathbb F_q\) 上首一 \(d\) 次不可约多项式的个数。比较\eqref{b1:eq:finite-field-polynomial-product}两端次数，得到
\[
q^n=\sum_{d\mid n}d\cdot N_q(d).
\]
\symidx{\(N_q(d)\)：有限域上首一 \(d\) 次不可约多项式的个数}
把 \(d=n\) 的一项单独取出，便得到递归式
\[
N_q(n)=\frac1n\left(q^n-\sum_{\substack{d\mid n\\d<n}}d\cdot N_q(d)\right).
\]
因此从 \(N_q(1)=q\) 开始，只要已知较小次数的计数，就能逐次求出 \(N_q(n)\)。下面的 Möbius 反演会把这个递归式改写成一个只含 \(q,n\) 的闭式。

\ProseParagraphBreak
先看 \(n=6\) 时怎样直接数根。在 \(\mathbb F_{q^6}\) 中，元素 \(a\) 对 \(\mathbb F_q\) 的次数是子域 \(\mathbb F_q(a)\) 的次数，因而必为六的除数，只能是 \(1,2,3,6\)。次数为一的元素在 \(\mathbb F_q\) 中，次数为二或三的元素分别在 \(\mathbb F_{q^2}\) 或 \(\mathbb F_{q^3}\) 中。因此，由\cref{b1:thm:finite-field-subfields}，次数小于六的元素恰好组成
\[
\mathbb F_{q^2}\cup\mathbb F_{q^3},
\qquad \mathbb F_{q^2}\cap\mathbb F_{q^3}=\mathbb F_q.
\]
从全部 \(q^6\) 个元素中删去这两个子域时，交集被删了两次，须加回一次。因此次数恰为六的元素共有 \(q^6-q^2-q^3+q\) 个。有限域上的 Frobenius 映射单射即满射，故由\cref{b1:prop:perfect-frobenius}，\(\mathbb F_q\) 完美，每个不可约多项式都可分；又由\cref{b1:prop:finite-field-factorization}，每个首一六次不可约多项式在 \(\mathbb F_{q^6}\) 中分裂，因而贡献六个不同的根。不同的不可约多项式没有公共根，因为一个根只有一个首一最小多项式，故
\[
N_q(6)=\frac{q^6-q^3-q^2+q}{6}.
\]
交替的加减修正了子域的重叠，除以六则把根的个数换成了不可约多项式的个数。

\ProseParagraphBreak
对一般的 \(n\)，次数等式把“次数恰为 \(d\)”的元素数 \(d\cdot N_q(d)\)，沿 \(d\mid n\) 累加为 \(q^n\)。要从这类累积计数恢复单个次数的计数，可以在除数关系上作反演。正整数上的\term[Mobius-hanshu]{Möbius 函数}[Möbius function]给出反演所需的系数：\footnote{默比乌斯（August Ferdinand Möbius，1790-1868），德国数学家、天文学家，提出 Möbius 反演公式，也研究拓扑中的单侧曲面。}\(\mu(1)=1\)；若正整数含有素数的平方因子，令 \(\mu(m)=0\)；若它为 \(r\) 个不同素数的乘积，令 \(\mu(m)=(-1)^r\)。
\symidx{\(\mu(m)\)：Möbius 函数}
选择不同素因子的子集给出
\[
\sum_{d\mid m}\mu(d)=
\begin{cases}1,&m=1,\\0,&m>1.\end{cases}
\]
若 \(m>1\) 有 \(r\) 个不同素因子，则上述和为 \((1-1)^r=0\)。于是
\begin{equation}\label{b1:eq:finite-field-irreducible-count}
N_q(n)=\frac1n\sum_{d\mid n}\mu(d)q^{n/d}.
\end{equation}
对每个 \(d\mid n\)，先前的次数等式给出
\[
q^{n/d}=\sum_{e\mid n/d}e\cdot N_q(e).
\]
代入后，求和的指标对 \((d,e)\) 满足 \(de\mid n\)。固定 \(e\) 时，允许的 \(d\) 恰好是 \(n/e\) 的除数，故交换有限求和得到
\[
\begin{aligned}
\sum_{d\mid n}\mu(d)q^{n/d}
&=\sum_{\substack{d,e\ge1\\de\mid n}}\mu(d)e\cdot N_q(e)\\
&=\sum_{e\mid n}e\cdot N_q(e)\sum_{d\mid n/e}\mu(d)\\
&=n\cdot N_q(n).
\end{aligned}
\]
最后一步中，内层和仅在 \(n/e=1\)，即 \(e=n\) 时为一，其余项全为零。例如 \(N_2(2)=1\)、\(N_2(3)=2\)、\(N_2(4)=3\)。

\ProseParagraphBreak
每个正次数实际上都有不可约多项式。由\cref{b1:cor:finite-field-cyclic}，可取 \(\mathbb F_{q^n}^\times\) 的一个生成元 \(g\)。子域 \(\mathbb F_q(g)\) 包含 \(g\) 的所有整数幂，也就包含 \(\mathbb F_{q^n}\) 的全部非零元素；再加上零，便有 \(\mathbb F_q(g)=\mathbb F_{q^n}\)。所以由\cref{b1:prop:minimal-polynomial-degree}，\(g\) 的最小多项式次数恰为 \(n\)。

\subsection{有限域的显式表示}
\label{b1:subsec:1703:03}

要在计算机或纸面上进行运算，须选择一个具体的不可约多项式 \(f\)，并把域写成 \(\mathbb F_q[t]/(f)\)。元素用次数小于 \(\deg f\) 的多项式表示；加法逐系数进行，乘法先相乘再除以 \(f\) 取余。对非零代表 \(a(t)\)，有 \(\gcd(a,f)=1\)，扩展 Euclid 算法给出 \(ua+vf=1\)，所以 \(u\) 的剩余类为其逆元。

\ProseParagraphBreak
若 \(f,g\in\mathbb F_q[t]\) 是不同的首一 \(d\) 次不可约多项式，它们给出的商域都作为 \(\mathbb F_q\) 扩张域同构于 \(\mathbb F_{q^d}\)。但两个商域中名为 \(\bar t\) 的元素分别满足 \(f(\bar t)=0\) 和 \(g(\bar t)=0\)，不能把它们直接对应；需要选定具体同构，才能在两种表示之间转换坐标。

\ProseParagraphBreak
例如 \(f=t^3+t+1\) 在 \(\mathbb F_2\) 上没有根，三次多项式若可约必有一次因子，故它不可约。令 \(\alpha=\bar t\)。这里特征为二，减号与加号相同，故 \(\alpha^3+\alpha+1=0\) 给出 \(\alpha^3=\alpha+1\)，于是
\[
\mathbb F_8=\{\,a+b\alpha+c\alpha^2\mid a,b,c\in\mathbb F_2\,\},
\qquad \alpha^3=\alpha+1.
\]
因为 \(\alpha(\alpha^2+1)=1\)，有 \(\alpha^{-1}=\alpha^2+1\)。又如
\[
(\alpha^2+\alpha)(\alpha+1)=\alpha^3+\alpha=1,
\]
故 \((\alpha^2+\alpha)^{-1}=\alpha+1\)。这些运算全在三维系数空间中完成；特征不同，\(\mathbb F_8\) 不可能作为子域嵌入 \(\mathbb C\)。

% !TeX root = ../../Book1.tex
% 纲要标题：### 17.4 有限域的计算与编码
% 纲要位置：计划纲要.md，第 606 行。

\section{有限域的计算与编码}
\label{b1:sec:1704}

% 纲要内容（仅作写作计划，尚非教材正文）：
%
% * 最小多项式、轨道和与轨道积；区分不同共轭元与全部嵌入像的计数
% * 多项式分解的有限域方法；不可约性检验与重复因子的重数恢复
% * Hamming 距离、求值码与唯一纠错；进一步的编码算法置于附录
%
% ---
%

Frobenius 轨道列出有限域元素的共轭元，由此可以计算最小多项式，并解释共轭元的和与积为何回到基域。下一章将从乘法线性算子的迹与行列式定义域迹、域范数，再证明它们在有限域扩张中可以用全部嵌入像的和与积计算，本节先计算这些和与积。有限域运算还可以构造纠错码：多项式在互异点处取值，根数界便控制不同编码向量之间的差别。

\subsection{最小多项式、迹与范数的计算}
\label{b1:subsec:1704:01}

\begin{proposition}[Frobenius 轨道与最小多项式]\label{b1:prop:frobenius-orbit-minimal}
设 \(q\) 为素数幂、\(n\geq1\) 为整数，\(\alpha\in\mathbb F_{q^n}\)。令 \(d\) 为使 \(\alpha^{q^d}=\alpha\) 的最小正整数。则 \(d\mid n\)，而 \(\alpha\) 在 \(\mathbb F_q\) 上的首一最小多项式为
\begin{equation}\label{b1:eq:frobenius-minimal-polynomial}
m_{\alpha,\mathbb F_q}(x)=\prod_{j=0}^{d-1}\bigl(x-\alpha^{q^j}\bigr).
\end{equation}
特别地，\([\mathbb F_q(\alpha):\mathbb F_q]=d\)，即元素的扩张次数等于其 Frobenius 轨道长度。
\end{proposition}
\begin{proof}
令 \(F:\mathbb F_{q^n}\rightarrow\mathbb F_{q^n}\) 为 \(q\) 次 Frobenius 自同构，\(F(z)=z^q\)。它满足 \(F^n=1\)，不动元恰为 \(\mathbb F_q\)。将 \(n\) 除以 \(d\)，写成 \(n=kd+r\)、\(0\leq r<d\)。由于 \(F^d(\alpha)=\alpha\)，有
\[
\alpha=F^n(\alpha)=F^r(\alpha).
\]
若 \(r>0\)，这与 \(d\) 是最小正周期矛盾，故 \(r=0\)，即 \(d\mid n\)。轨道中的 \(\alpha,\alpha^q,\ldots,\alpha^{q^{d-1}}\) 两两不同。

记右端乘积为 \(P(x)\)。对 \(P\) 的各个系数施行 \(F\)，并保持形式变量 \(x\) 不变，得到
\[
F(P)(x)=\prod_{j=0}^{d-1}\bigl(x-\alpha^{q^{j+1}}\bigr)=P(x),
\]
因为 \(F\) 只是把这 \(d\) 个根循环置换。因此每个系数 \(c\) 都满足 \(c^q=c\)，所以 \(P\in\mathbb F_q[x]\)。

另一方面，若非零 \(h\in\mathbb F_q[x]\) 满足 \(h(\alpha)=0\)，其系数都被 \(F\) 固定，故
\[
h(\alpha^q)=h(\alpha)^q=0.
\]
迭代此等式，轨道中的全部 \(d\) 个元素都是 \(h\) 的根，因而 \(\deg h\geq d\)。\(P\) 首一、以 \(\alpha\) 为根且次数为 \(d\)，于是它就是最小多项式；\cref{b1:prop:minimal-polynomial-degree}再给出所述扩张次数。
\end{proof}

在 \(\mathbb F_8=\mathbb F_2(\alpha)\)、\(\alpha^3=\alpha+1\) 中，元素在 \(\mathbb F_2\) 上的次数只能是三的除数。除 \(0,1\) 外，剩下六个元素的次数都为三，故平方映射将它们分成两个长度为三的轨道。全部轨道如下：
\begin{center}
\begin{tabular}{@{}lcl@{}}
\toprule
轨道 & 长度 & 对应的首一最小多项式 \\
\midrule
\(\{0\}\) & \(1\) & \(x\) \\
\(\{1\}\) & \(1\) & \(x+1\) \\
\(\{\alpha,\alpha^2,\alpha^4\}\) & \(3\) & \(x^3+x+1\) \\
\(\{\alpha^3,\alpha^6,\alpha^5\}\) & \(3\) & \(x^3+x^2+1\) \\
\bottomrule
\end{tabular}
\end{center}
由于 \(\alpha^3=\alpha+1\)，可核对 \(\alpha^3\) 满足 \(x^3+x^2+1\)；两个三次式在 \(\mathbb F_2\) 中都无根，故不可约。四个轨道互不相交，长度之和为 \(1+1+3+3=8\)，每个轨道对应一个不可约因子。因此
\[
x^8-x=x(x+1)(x^3+x+1)(x^3+x^2+1),
\]
并且 \(N_2(3)=2\)。\(\alpha\) 的三个轨道元素之和为零，积为 \(\alpha^7=1\)；在 \(\mathbb F_8/\mathbb F_2\) 中它们也是 \(\alpha\) 的全部嵌入像，下一章的迹与范数公式将分别给出零和一。

\ProseParagraphBreak
元素所处的扩张不能省略。例如把 \(1\) 视为 \(\mathbb F_4/\mathbb F_2\) 的元素，按扩张的两个嵌入计，和为 \(1+1=0\)、积为一；只按它的单元素轨道计则和为一。一般地，设 \(q\) 为素数幂、\(n\geq1\)，记 \(F(z)=z^q\)，若 \(\alpha\in\mathbb F_{q^n}\) 的轨道长为 \(d\)，则 \(d\mid n\)。该扩张的各个嵌入的像都是代数闭包中同一个 \(q^n\) 元子域，所以它们正是\cref{b1:thm:finite-field-automorphisms}列出的 \(n\) 个自同构 \(F^0,\ldots,F^{n-1}\)。这些嵌入在 \(\alpha\) 上的取值每经过 \(d\) 步便重复一次，所以每个轨道元素出现 \(n/d\) 次。下一章的迹公式将轨道元素的和重复 \(n/d\) 次相加，范数公式则将轨道元素的积取 \(n/d\) 次幂，不能总以最小多项式的一组不同根的和与积代替。

\subsection{有限域上的多项式分解}
\label{b1:subsec:1704:02}

设 \(q\) 为素数幂。非零 \(f\in\mathbb F_q[x]\) 若每个首一不可约因子都只出现一次，就称为无重因子多项式。\cref{b1:prop:perfect-frobenius}已说明有限域是完美域，即有限域上的不可约多项式都可分，所以对非常数 \(f\)，这等价于 \(f\) 在代数闭包中没有重根。由\cref{b1:prop:finite-field-factorization}，对这样的 \(f\) 计算
\[
\gcd\bigl(f,x^{q^d}-x\bigr)
\]
就能提取 \(f\) 中次数整除 \(d\) 的全部不可约因子。例如 \(d=6\) 时，这个最大公因式可能同时含有一次、二次、三次和六次因子，不能直接当作六次因子的乘积。幂 \(x^{q^d}\) 可以不断取 \(q\) 次幂并模 \(f\) 取余，无须展开到巨大次数。按 \(d\) 从小到大处理，从本次最大公因式中除去其中已经找出的低次因子，便得到次数恰为 \(d\) 的因子乘积。

\begin{proposition}[有限域上的不可约性检验]\label{b1:prop:finite-field-irreducibility-test}
设 \(q\) 为素数幂，\(f\in\mathbb F_q[x]\) 为次数 \(n>1\) 的首一多项式。则 \(f\) 不可约，当且仅当 \(f\mid x^{q^n}-x\)，并且对每个整除 \(n\) 的素数 \(\ell\)，都有
\[
\gcd\bigl(f,x^{q^{n/\ell}}-x\bigr)=1.
\]
\end{proposition}
\begin{proof}
若 \(f\) 不可约，\cref{b1:prop:finite-field-factorization}给出 \(f\mid x^{q^n}-x\)。对整除 \(n\) 的素数 \(\ell\)，有 \(n/\ell<n\)，故 \(n\) 不整除 \(n/\ell\)，于是 \(f\) 不整除 \(x^{q^{n/\ell}}-x\)。由 \(f\) 不可约，相应最大公因式只能为一。

反过来，\(x^{q^n}-x\) 的导数为 \(-1\)，所以条件 \(f\mid x^{q^n}-x\) 首先排除了 \(f\) 的重复因子。\cref{b1:prop:finite-field-factorization}还说明 \(f\) 的每个不可约因子 \(h\) 的次数 \(d\) 都整除 \(n\)。若 \(d<n\)，整数 \(n/d>1\) 有某个素因子 \(\ell\)，从而 \(n/d=\ell s\)，即 \(n/\ell=ds\)。因此 \(d\mid n/\ell\)，同一命题给出 \(h\mid x^{q^{n/\ell}}-x\)，与该最大公因式为一矛盾。这说明每个不可约因子的次数都等于 \(n\)；而 \(\deg f=n\)，故只有一个因子，\(f\) 不可约。
\end{proof}

处理任意非零 \(f\in\mathbb F_q[x]\) 时，先提出首项系数，对剩余的首一多项式递归；遇到常数一便停止。若 \(f'\ne0\)，记 \(h=\gcd(f,f')\)。当 \(h=1\) 时，\(f\) 已无重因子；当 \(h\ne1\) 时，须分别递归处理 \(h\) 和 \(f/h\)，二者次数都严格小于 \(\deg f\)。不能只处理商 \(f/h\)，因为它可能遗漏只在 \(h\) 中留下的不可约因子。若 \(f'=0\)，则所有非零项的指数都被特征 \(p\) 整除，可写成 \(f=\sum_i a_ix^{pi}\)。由\cref{b1:prop:perfect-frobenius}，有限域上的 Frobenius 满射，并且仍为单射，所以每个系数 \(a_i\) 有唯一的 \(p\) 次根 \(b_i\)。于是
\[
g=\sum_i b_ix^i\qquad\text{满足}\qquad f=g^p.
\]
递归处理 \(g\) 后，将所得因子的重数都乘以 \(p\)。不同分支得到相同首一不可约因子时，还须把重数相加。每次递归都降低正次数，因此过程终止；每一步保持乘积等式，所以因子及其重数均被保留。\subsecref{b1:subsec:B01:02}给出这一预处理的进一步算例。

\ProseParagraphBreak
例如在 \(\mathbb F_2[x]\) 中取 \(g=x\)、\(h=x+1\)，令 \(f=g^3h^2\)。特征为二，故
\[
f'=g^2h^2,\qquad \gcd(f,f')=g^2h^2,\qquad
\frac{f}{\gcd(f,f')}=g.
\]
商中完全没有 \(h\)，若只分解它就会漏掉这个因子。另一分支 \(g^2h^2=(gh)^2\) 的导数为零，开平方得到无重因子的 \(gh\)，再把两个一次因子的重数乘以二。与商分支的一个 \(g\) 合并后，\(g\) 的重数为 \(2+1=3\)，\(h\) 的重数为二，恢复原来的 \(g^3h^2\)。

\ProseParagraphBreak
按次数分组不一定已经得到每个不可约因子。对于同次数因子的乘积，可以枚举该次数的首一多项式并作试除；枚举有限且必终止。这说明原则上能通过有限搜索完成分解，但候选数随次数增长很快，实际计算还需要改善效率。\subsecref{b1:subsec:B01:03}进一步给出 Berlekamp 分解\footnote{伯利坎普（Elwyn Ralph Berlekamp，1940-2019），美国数学家，研究代数编码与组合博弈，提出有限域多项式的分解算法。}，用线性代数和最大公因式分开这些因子。

\subsection{编码理论的应用}
\label{b1:subsec:1704:03}

\begin{definition}\label{b1:def:hamming-distance}
设 \(q\) 为素数幂，\(N\geq1\) 为整数。对 \(v,w\in\mathbb F_q^N\)，数
\[
d_H(v,w)\coloneqq\#\{\,i\in\{1,\ldots,N\}\mid v_i\ne w_i\,\}
\]
称为 \(v,w\) 的\term[Hamming-juli]{Hamming 距离}[Hamming distance]。\footnote{汉明（Richard Wesley Hamming，1915-1998），美国数学家与计算机科学家，建立纠错码理论中的汉明码及相关距离概念。}若 \(C\subseteq\mathbb F_q^N\) 至少包含两个向量，则不同向量之间 Hamming 距离的最小值称为 \(C\) 的\term[zuixiao-juli]{最小距离}[minimum distance]，记为 \(d(C)\)。
\end{definition}
\symidx{\(d_H(v,w)\)：Hamming 距离}
\symidx{\(d(C)\)：码的最小距离}
距离为零恰好表示两个向量相同，距离也显然对称。若 \(v_i\ne w_i\)，则对任意向量 \(z\)，必有 \(v_i\ne z_i\) 或 \(z_i\ne w_i\)。对这些坐标计数，便有三角不等式
\[
d_H(v,w)\leq d_H(v,z)+d_H(z,w).
\]

\begin{samepage}
下面构造一个\term[qiuzhi-ma]{求值码}[evaluation code]。取 \(\mathbb F_5\) 中五个互异点 \(0,1,2,3,4\)。对信息 \((a,b)\in\mathbb F_5^2\)，取 \(f(x)=a+bx\in\mathbb F_5[x]\)；它可以是零多项式，否则次数不超过一。编码向量为
\[
\bigl(f(0),f(1),f(2),f(3),f(4)\bigr)\in\mathbb F_5^5.
\]
\end{samepage}
记全部编码向量组成的集合为 \(C\)，其中的向量称为码字。若两个信息对应不同的多项式 \(f,g\)，则 \(h=f-g\ne0\)，且 \(\deg h\leq1\)。在五个互异求值点中，至多一个点满足 \(h(a)=0\)，也就是 \(f(a)=g(a)\)。所以相应码字至少在四个坐标不同，即 \(d(C)\geq4\)。零多项式与 \(x\) 的取值只在点零处相同，给出距离恰为四的一对码字，因此 \(d(C)=4\)。若传输中至多改坏一个坐标，原码字便唯一：若两个码字 \(c,c'\) 都与收到的向量 \(y\) 相距至多一，三角不等式便给出 \(d_H(c,c')\leq2<4\)，与最小距离矛盾。

\ProseParagraphBreak
例如 \(f(x)=1+2x\) 给出 \((1,3,0,2,4)\)。若收到 \((1,3,1,2,4)\)，可枚举所有形如 \(a+bx\) 的多项式（包括零多项式），找出在至少四个位置吻合者，唯一得到原来的 \(f\)。也可以枚举五个求值点的所有两点子集，每次由相应的两个接收值插值得到一个次数不超过一的多项式，再检验它与全部五个接收值的吻合数。由于至多一处错误，必有选中的两个位置都正确，此时插值恢复原来的 \(f\)；包含错误位置的候选也须通过至少四处吻合的检验，不能任取两个位置就相信所得结果。前一种方法只有二十五个候选，后一种方法只需尝试十组点，两者都给出完整的有限解码步骤；更长的码需要更有效的方法。这一例子说明了多项式根数界如何转化为纠错能力。


\begin{chaptersummary}
若 \(q=p^r\)，其中 \(p\) 为素数、\(r\ge1\)，则 \(\mathbb F_q\) 是 \(x^q-x\) 在 \(\mathbb F_p\) 上的分裂域；作为 \(\mathbb F_p\) 向量空间，其维数为 \(r\)；选取一个 \(\mathbb F_p[x]\) 中的 \(r\) 次不可约多项式，还可以将它表示为相应的商域。Frobenius 把这些描述连接起来：其轨道给出最小多项式，其不动元辨认子域，其幂列出全部基域自同构。乘法本原元提供枚举非零元素的方法，但域生成元未必是乘法群生成元。不可约多项式的次数分组和 Möbius 计数，以及由根数界得到的求值码，均依赖这些可以明确检验的结构。
\end{chaptersummary}

\BookExercises

\begin{exercise}\label{b1:ex:17-cardinality}
设 \(F\) 为十六元域。确定它的特征、作为素域上向量空间的维数以及全部可能的子域大小。说明 \(\mathbb Z/16\mathbb Z\) 为何不能作为 \(F\) 的另一种表示。
\end{exercise}
\begin{solution}
由素数幂的唯一分解，\(16=2^4\)，故特征为二，维数为四。\cref{b1:thm:finite-field-subfields}说明子域对应四的正因子，大小为二、四、十六。\(\mathbb Z/16\mathbb Z\) 中 \(\bar2\ne0\)、\(\bar8\ne0\)，但 \(\bar2\bar8=0\)，存在零因子；其特征又为十六，均与域的条件不合。
\end{solution}

\begin{exercise}\label{b1:ex:17-four-field}
令 \(F=\mathbb F_2(u)\)、\(u^2+u+1=0\)。写出各非零元素的逆和全部 \(\mathbb F_2\) 自同构。
\end{exercise}
\begin{solution}
元素为 \(0,1,u,u+1\)，且 \(u(u+1)=1\)，所以一自逆，\(u\) 与 \(u+1\) 互逆。保持素域的自同构由 \(u\) 的像决定，像必须是最小多项式的根 \(u,u+1\)。这两种选择分别给出恒等和 Frobenius，因为 \(u^2=u+1\)。两者都保持定义关系，故确实是自同构。
\end{solution}

\begin{exercise}\label{b1:ex:17-embedding-count}
判断 \(\mathbb F_{16}\) 能否嵌入 \(\mathbb F_{64}\)。再求保持 \(\mathbb F_2\) 的嵌入 \(\mathbb F_8\rightarrow\mathbb F_{64}\) 有多少个。
\end{exercise}
\begin{solution}
前者不可能，因为像将是次数为四的子域，但四不整除六。记 \(E\subseteq\mathbb F_{64}\) 为唯一的八元子域，并固定一个保持 \(\mathbb F_2\) 的同构 \(\iota:\mathbb F_8\xrightarrow{\sim}E\)。任意另一个保持 \(\mathbb F_2\) 的嵌入 \(j:\mathbb F_8\rightarrow\mathbb F_{64}\) 都以 \(E\) 为像，因此
\[
\theta=\iota^{-1}\circ j\in\operatorname{Aut}_{\mathbb F_2}(\mathbb F_8),
\qquad j=\iota\circ\theta.
\]
由\cref{b1:thm:finite-field-automorphisms}，\(\theta\) 只能是恒等、平方映射 \(F:a\mapsto a^2\) 或 \(F^2\)。故全部嵌入为 \(\iota,\iota\circ F,\iota\circ F^2\)，共三个；像子域唯一不意味着到该子域的嵌入唯一。
\end{solution}

\begin{exercise}\label{b1:ex:17-subfield-lattice}
列出 \(\mathbb F_{q^{12}}\) 中包含 \(\mathbb F_q\) 的全部子域，并计算 \(\mathbb F_{q^4}\) 与 \(\mathbb F_{q^6}\) 的交域、合成域。
\end{exercise}
\begin{solution}
十二的正因子为 \(1,2,3,4,6,12\)，相应子域为这些次数的 \(\mathbb F_q\) 扩张，且每个唯一。子域的包含关系由扩张次数的整除关系确定。由\cref{b1:thm:finite-field-subfields}，交域次数为 \(\gcd(4,6)=2\)，合成域次数为 \(\operatorname{lcm}(4,6)=12\)，分别是 \(\mathbb F_{q^2}\) 和整个域。
\end{solution}

\begin{exercise}\label{b1:ex:17-frobenius-power-fixed}
设 \(\sigma(a)=a^q\) 作用在 \(\mathbb F_{q^n}\) 上。对正整数 \(r\)，直接证明 \(\sigma^r\) 的不动元构成 \(\mathbb F_{q^{\gcd(n,r)}}\)。
\end{exercise}
\begin{solution}
令 \(d=\gcd(n,r)\)，取整数 \(u,v\) 使 \(d=un+vr\)。若 \(\sigma^r(a)=a\)，则 \(\sigma^n(a)=a\) 也成立；由于 \(\sigma\) 可逆，其整数幂都可使用，故 \(\sigma^d(a)=\sigma^{un+vr}(a)=a\)。反过来 \(d\mid r\)，被 \(\sigma^d\) 固定的元素必被 \(\sigma^r\) 固定。两者不动集相同，而前者正为 \(x^{q^d}-x\) 的根集，即所述子域。
\end{solution}

\begin{exercise}\label{b1:ex:17-generators}
取 \(\mathbb F_{16}^\times\) 的生成元 \(g\)。求本原元的个数，并说明 \(g^3\) 是否为乘法本原元、是否生成扩张 \(\mathbb F_{16}/\mathbb F_2\)。
\end{exercise}
\begin{solution}
乘法群阶为十五，本原元个数为 \(\varphi(15)=8\)。\(g^3\) 的阶为五，所以不是乘法本原元。但真子域只有 \(\mathbb F_2\) 和 \(\mathbb F_4\)，其乘法群阶分别为一和三，都不能含有五阶元。因此 \(g^3\) 不在任何真子域中，仍生成整个域扩张。
\end{solution}

\begin{exercise}\label{b1:ex:17-irreducible-counts}
求 \(\mathbb F_3\) 上首一二、三、四次不可约多项式的个数。再求 \(\mathbb F_{q^6}\) 中对 \(\mathbb F_q\) 恰为六次的元素个数。
\end{exercise}
\begin{solution}
由\eqref{b1:eq:finite-field-irreducible-count}，依次为 \((9-3)/2=3\)、\((27-3)/3=8\)、\((81-9)/4=18\)。六次域中次数小于六的元素恰在 \(\mathbb F_{q^2}\cup\mathbb F_{q^3}\) 中，两域交为 \(\mathbb F_q\)。故六次元素有 \(q^6-q^2-q^3+q\) 个，亦等于 \(6N_q(6)\)。后一等式体现每个六次不可约多项式有六个互异根。
\end{solution}

\begin{exercise}\label{b1:ex:17-explicit-sixteen}
证明 \(f=t^4+t+1\) 在 \(\mathbb F_2\) 上不可约。在 \(F=\mathbb F_2[t]/(f)\) 中令 \(\alpha=\bar t\)，计算 \(\alpha^{-1}\) 和 \((\alpha^2+1)^2\)。再令 \(\beta=\alpha^2+\alpha\)、\(E=\mathbb F_2(\beta)\)。证明 \(E\) 是 \(F\) 中的四元子域，比较平方映射与四次幂映射在 \(E\) 上的作用，并求 \(\alpha\) 在 \(E\) 上的最小多项式。
\end{exercise}
\begin{solution}
零、一都不是根，所以若可约，只能是两个不可约二次因子的乘积。\(\mathbb F_2\) 上唯一首一不可约二次多项式为 \(t^2+t+1\)：一个无根的首一二次式常数项必须为一，再检查一次项系数即可。设其根为 \(u\)，有 \(u^4=u\)，从而 \(f(u)=1\ne0\)，排除了该因子。因此 \(f\) 不可约。在商域中 \(\alpha^4=\alpha+1\)，所以 \(\alpha(\alpha^3+1)=1\)，逆为 \(\alpha^3+1\)；又 \((\alpha^2+1)^2=\alpha^4+1=\alpha\)。

由 \(\alpha^4=\alpha+1\) 得 \(\beta=\alpha^5\)，并有 \(\beta^2=\alpha^4+\alpha^2=\beta+1\)。因此 \(\beta^2+\beta+1=0\)。多项式 \(X^2+X+1\) 在 \(\mathbb F_2\) 中无根，所以 \(\beta\notin\mathbb F_2\)，\(E\) 恰有四个元素。平方映射把 \(\beta\) 送到 \(\beta+1\)，并不逐点固定 \(E\)；四次幂映射则固定 \(E\) 的每个元素，并将 \(\alpha\) 送到 \(\alpha^4=\alpha+1\)。

由于 \([F:E]=2\)，且 \(\alpha\notin E\)，\(\alpha\) 在 \(E\) 上的最小多项式为
\[
m_{\alpha,E}(X)=(X-\alpha)(X-\alpha^4)=X^2+X+\beta.
\]
\end{solution}

\begin{exercise}\label{b1:ex:17-minimal-shift}
在 \(\mathbb F_8=\mathbb F_2(\alpha)\)、\(\alpha^3+\alpha+1=0\) 中，求 \(\beta=\alpha+1\) 的最小多项式。
\end{exercise}
\begin{solution}
\(\alpha^3=\alpha+1=\beta\)，而 \(\alpha\ne1\) 的乘法阶为七。因此 \(\beta^2=\alpha^6=\alpha^2+1\)，\(\beta^3=\alpha^9=\alpha^2\)，得到 \(\beta^3+\beta^2+1=0\)。多项式 \(x^3+x^2+1\) 在零、一处均不为零，故不可约，正是最小多项式。它也由 \(\beta,\beta^2,\beta^4\) 的 Frobenius 轨道给出。
\end{solution}

\begin{exercise}\label{b1:ex:17-irreducibility-test}
设首一 \(f\in\mathbb F_q[x]\) 的次数为 \(n>1\)。证明它不可约当且仅当 \(f\mid x^{q^n}-x\)，并且对每个整除 \(n\) 的素数 \(\ell\)，有 \(\gcd(f,x^{q^{n/\ell}}-x)=1\)。
\end{exercise}
\begin{solution}
若不可约，\cref{b1:prop:finite-field-factorization}立即给出整除关系，且次数 \(n\) 不整除较小的 \(n/\ell\)，所以各最大公因式为一。反之，第一个条件 \(f\mid x^{q^n}-x\) 使 \(f\) 无重因子，并使其每个不可约因子次数 \(d\) 整除 \(n\)。若 \(d<n\)，取 \(n/d\) 的一个素因子 \(\ell\)，则 \(d\mid n/\ell\)，该因子将整除相应最大公因式，矛盾。因此每个因子次数都等于 \(n\)，由 \(\deg f=n\) 知只有一个因子。
\end{solution}

\begin{exercise}\label{b1:ex:17-inseparable-factor}
分解 \(x^6+x^2+1\in\mathbb F_2[x]\)，并确定其在分裂域中各根的重数。
\end{exercise}
\begin{solution}
形式导数为零，且特征二下 \(x^6+x^2+1=(x^3+x+1)^2\)。内层三次式在零、一处都不为零，故不可约。其导数为 \(x^2+1=(x+1)^2\)，而 \(x^3+x+1\) 在 \(x=1\) 处取值为一，所以与导数互素，有三个不同的根。原多项式在分裂域中这三个根的重数各为二。
\end{solution}

\begin{exercise}\label{b1:ex:17-orbit-sum-product}
在 \(\mathbb F_9=\mathbb F_3(u)\)、\(u^2=2\) 中，求 \(1+u\) 的 Frobenius 轨道、轨道元素的和与积，以及它的最小多项式。
\end{exercise}
\begin{solution}
三次幂映射把 \(u\) 送到 \(u^3=2u=-u\)，故轨道为 \(1+u,1-u\)。和为二，积为 \(1-u^2=2\)。最小多项式因而为 \(x^2-2x+2=x^2+x+2\)。\(u\notin\mathbb F_3\)，所以这两个根确实不同，次数为二。
\end{solution}

\begin{exercise}\label{b1:ex:17-factor-nine}
将 \(x^9-x\) 在 \(\mathbb F_3[x]\) 中分解为首一不可约因子的乘积。
\end{exercise}
\begin{solution}
一次因子为 \(x,x-1,x+1\)。首一不可约二次式恰有三个，分别为 \(x^2+1,x^2+x+2,x^2+2x+2\)，在零、一、负一处逐一代入都非零，故都不可约。\cref{b1:prop:finite-field-factorization}说明 \(x^9-x\) 恰为所有次数整除二的首一不可约式各一次的乘积。因此所求分解为上述六个因子的乘积，其总次数为九。
\end{solution}

\begin{exercise}\label{b1:ex:17-evaluation-distance}
在 \(\mathbb F_q\) 中选 \(n\) 个互异点，\(1\leq k\leq n\leq q\)。以次数小于 \(k\) 的多项式在这些点的取值组成求值码。证明不同编码向量之间最少恰有 \(n-k+1\) 个坐标不同，并说明至多能唯一纠正多少处错误。
\end{exercise}
\begin{solution}
两个不同多项式的差次数至多为 \(k-1\)，所以至多在 \(k-1\) 个所选点为零，给出距离至少 \(n-k+1\)。取其中 \(k-1\) 个点，令多项式为对应一次因子的乘积；当 \(k=1\) 时取常数一。它与零多项式的取值恰在 \(n-k+1\) 处不同，所以下界能达到。

记 \(d(v,w)\) 为向量 \(v,w\) 的不同坐标个数，最小距离为 \(\delta=n-k+1\)。若错误数 \(e\) 满足 \(2e<\delta\)，两个候选编码向量不可能同时与收到的向量相差至多 \(e\) 处，否则两者距离至多 \(2e\)。因此可以唯一纠正至多 \(\lfloor(n-k)/2\rfloor\) 处错误。

为证明不能统一保证再多纠正一处，取相距恰为 \(\delta\) 的编码向量 \(c,c'\)。在它们不同的 \(\delta\) 个坐标中，选出 \(\lfloor\delta/2\rfloor\) 个，让接收向量 \(y\) 取 \(c'\) 的值；其余坐标取 \(c\) 的值。于是
\[
d(y,c)=\lfloor\delta/2\rfloor,\qquad
d(y,c')=\lceil\delta/2\rceil
=\lfloor(n-k)/2\rfloor+1.
\]
若允许后一数目的错误，同一个 \(y\) 就可能由 \(c\) 或 \(c'\) 发送，无法保证唯一纠错。例如在 \(\mathbb F_5\) 上取求值点 \(0,1,2,3,4\) 及 \(k=2\)，零多项式与 \(x\) 给出 \(c=(0,0,0,0,0)\)、\(c'=(0,1,2,3,4)\)。接收向量 \(y=(0,1,2,0,0)\) 与两者均相差两处。
\end{solution}

% !TeX root = ../../Book1.tex
% 纲要标题：## 第18章　嵌入的独立性、迹与范数
% 纲要位置：计划纲要.md，第 614 行。

\chapter{嵌入的独立性、迹与范数}
\label{b1:ch:18}
\BookChapterNumber{b1:ch:18}

\begin{chapterabstract}
一个域元素既能作为扩张域中的乘法算子来观察，也能通过不同嵌入得到共轭值。本章证明嵌入的线性独立性，建立两种观点下的迹、范数与判别式公式，并补充后续计算所需的整性工具。
\end{chapterabstract}

\begin{learninggoals}
\begin{itemize}
\item 掌握最短关系法，区别乘法特征的线性独立性与表示特征标的正交性。
\item 从乘法矩阵与最小多项式计算迹、范数，正确处理不可分重数。
\item 证明迹与范数的传递性，用迹配对和判别式检测有限扩张的可分性。
\item 计算有限域上的迹与范数，说明整性如何保证数域中的相应值为整数。
\end{itemize}
\end{learninggoals}

令 \(L=\mathbb Q(\sqrt2)\)。以 \(1,\sqrt2\) 为基，“乘以 \(\sqrt2\)”这一 \(\mathbb Q\)-线性映射的矩阵为
\[
A=\begin{pmatrix}0&2\\1&0\end{pmatrix},\qquad
\det(tI_2-A)=t^2-2.
\]
迹为 \(0\)，行列式为 \(-2\)。这个特征多项式也是 \(\sqrt2\) 在 \(\mathbb Q\) 上的最小多项式，所以乘法算子的两个特征值正好是两个域嵌入像 \(\sqrt2,-\sqrt2\)，它们的和与积给出矩阵的迹与行列式。这里比较的是线性算子的特征值和域元素的共轭像；本章将解释两种计算为何相合。

\ProseParagraphBreak
不过，乘法特征多项式还依赖于整个环境扩张。当把同一个 \(\sqrt2\) 放进 \(\mathbb Q(\sqrt2,\sqrt3)\) 中时，以 \(1,\sqrt2,\sqrt3,\sqrt6\) 为基，乘法矩阵是两个 \(A\) 组成的块对角矩阵，所以特征多项式变为 \((t^2-2)^2\)。最小多项式仍是 \(t^2-2\)；扩大乘法算子作用的空间，使特征值重复出现，并没有改变元素本身满足的最小关系。

\ProseParagraphBreak
这个例子还有一项整数性质：\(\sqrt2\) 满足首一整系数方程 \(X^2-2=0\)，\(1+\sqrt2\) 满足 \(X^2-2X-1=0\)。相应的迹与范数都是整数。哪些域元素具有这样的首一方程，两个这样的元素相加或相乘后又怎样？本章末尾的整性讨论将把这项观察变成可用于后续模素数计算的工具。

\ProseParagraphBreak

本章使用域扩张的次数、嵌入计数、可分次数的塔式公式，以及有限维线性代数中的矩阵迹、行列式与换基。下一章的有限 Galois 对应使用\secref{b1:sec:1801}的嵌入独立性；\secref{b1:sec:1802}至\secref{b1:sec:1803}的迹、范数与判别式用于其他扩张计算，\secref{b1:sec:1804}的整性工具用于后面的模素数约化。

\ProseParagraphBreak
迹、范数和乘法算子的关系，可结合 Roman 的《Field Theory》\cite[Section 8.1]{Roman2006FieldTheory}阅读；整性可继续查阅 Atiyah 与 Macdonald 的《Introduction to Commutative Algebra》\cite[Chapter 5]{AtiyahMacdonald1969CommutativeAlgebra}。本章还区分乘法特征与表示的特征标；要学习后一种概念，可读 Serre 的《Linear Representations of Finite Groups》\cite[Part I, Chapters 1--2, pp. 3--24]{Serre1977FiniteRepresentations}。

% !TeX root = ../../Book1.tex
% 纲要标题：### 18.1 Artin 独立性
% 纲要位置：计划纲要.md，第 616 行。

\section{Artin 独立性}
\label{b1:sec:1801}

% 纲要内容（仅作写作计划，尚非教材正文）：
%
% * 乘法特征的线性无关性
% * 域嵌入的线性无关性
% * 从函数独立性到可逆求值矩阵
% * 与有限群表示特征标正交性的区别
% * 先以最短非平凡线性关系证明任意群上不同乘法特征的有限线性独立性，再得到不同域嵌入的线性独立性和可逆求值矩阵；本节不依赖 Galois 对应，末尾比较函数线性无关与复特征标正交性
%

两个嵌入即使有相同的定义域和值域，也可能作为函数完全不同。我们需要的结论比“两个映射不相等”强：不同嵌入之间不存在系数恒定的非平凡线性关系。稍后可分扩张的迹公式会把迹写成全部嵌入的和；嵌入的线性无关性将保证这个和不是零函数，并由此证明迹配对非退化，见\cref{b1:thm:trace-pairing-separable}。独立性的证明只依赖乘法，因此先在任意群上建立它。

\subsection{乘法特征的线性无关性}
\label{b1:subsec:1801:01}

设 \(G\) 为群，\(\Omega\) 为域。保持乘法的映射 \(\chi:G\rightarrow\Omega^\times\) 称为一个 \(\Omega\) 值的\term[chengfa-tezheng]{乘法特征}[multiplicative character]。全体函数 \(G\rightarrow\Omega\) 可以逐点相加和乘以标量：对函数 \(f,g\) 和 \(c\in\Omega\)，令
\[
(f+g)(x)=f(x)+g(x),\qquad (cf)(x)=c\cdot f(x).
\]
这些运算使它们组成 \(\Omega\) 向量空间。乘法特征只是这个空间中的一类元素，其和一般不再是乘法特征，但仍能在这个函数空间中谈论它们的线性无关性。关系 \(\sum_i c_i\chi_i=0\) 要求同一组常数 \(c_i\) 对每个 \(g\in G\) 都使函数值之和为零。

\ProseParagraphBreak
以下结果称为\term[Artin-dulixing]{Artin 独立性定理}[Artin independence theorem]。
\begin{theorem}[Artin]\label{b1:thm:artin-independence}
设 \(G\) 为群，\(\Omega\) 为域。从 \(G\) 到 \(\Omega^\times\) 的任意一族互异群同态，在全体函数 \(G\rightarrow\Omega\) 组成的 \(\Omega\) 向量空间中线性无关。对无限族，这表示每个有限子族都线性无关。
\end{theorem}
\begin{proof}
若有非平凡关系，从全部这种关系中取非零项数最少者，重新编号写为
\[
c_1\cdot\chi_1(g)+\cdots+c_r\cdot\chi_r(g)=0\qquad(g\in G),
\]
其中每个 \(c_i\ne0\)。因为 \(\chi_1(g)\ne0\)，不可能 \(r=1\)。我们要利用乘法结构消去第 \(r\) 项，同时让至少另一项的系数非零；否则得到零关系，并不会违背最短性。\(\chi_1\ne\chi_r\) 作为函数的含义，就是存在 \(h\in G\) 使 \(\chi_1(h)\ne\chi_r(h)\)。将关系中的 \(g\) 换成 \(hg\)，利用 \(\chi_i(hg)=\chi_i(h)\chi_i(g)\)，再减去原关系的 \(\chi_r(h)\) 倍，得到
\[
\sum_{i=1}^{r-1}c_i\bigl(\chi_i(h)-\chi_r(h)\bigr)\cdot\chi_i(g)=0
\qquad(g\in G).
\]
第 \(r\) 项被消去，而第一项系数 \(c_1(\chi_1(h)-\chi_r(h))\ne0\)，所以这是一个非平凡关系，其非零项数至多为 \(r-1\)，与最小性矛盾。全部操作只涉及有限个项，所以同样排除了无限族中的有限非平凡关系。
\end{proof}
群 \(G\) 不必有限，也不必交换；域的特征没有限制。证明中没有对群元素求平均，也没有除以群阶。

\subsection{域嵌入的线性无关性}
\label{b1:subsec:1801:02}

\begin{corollary}\label{b1:cor:embedding-independence}
设 \(L,\Omega\) 为域。若 \(\sigma_1,\ldots,\sigma_r:L\rightarrow\Omega\) 为互异的域嵌入，则它们在全体函数 \(L\rightarrow\Omega\) 组成的 \(\Omega\) 向量空间中线性无关。
\end{corollary}
\begin{proof}
域嵌入单射，因此将非零元素送到非零元素。故限制 \(\sigma_i|_{L^\times}:L^\times\rightarrow\Omega^\times\) 确实是乘法特征。不同域嵌入必在某个非零元素上不同，因为它们都把零映为零。因此这些限制互异，\cref{b1:thm:artin-independence}排除了其上的线性关系，也就排除了原映射的线性关系。
\end{proof}
例如复数上的恒等映射与共轭映射线性无关：若 \(az+b\bar z=0\) 对所有 \(z\in\mathbb C\) 成立，分别代入 \(z=1,i\)，就有 \(a+b=0,a-b=0\)，从而 \(a=b=0\)。这里可以直接找到两个检验点。一般证明先用最短关系法得到函数独立性，再由下面的线性代数引理选出有限个检验点，不要求事先猜到它们。

\subsection{从函数独立性到可逆求值矩阵}
\label{b1:subsec:1801:04}

最短关系法先把一个未知的关系压缩到不能再短，再借助结构运算消去一项。关键在于消去某项的同时，必须明确指出另一个系数没有消失；否则新关系可能只是恒等的零关系。下面把函数的独立性转成后文可以使用的矩阵。

\begin{lemma}\label{b1:lem:independent-functions-evaluation}
设 \(X\) 为集合，\(\Omega\) 为域，函数 \(f_1,\ldots,f_r:X\rightarrow\Omega\) 在全体函数 \(X\rightarrow\Omega\) 组成的 \(\Omega\) 向量空间中线性无关。则存在 \(x_1,\ldots,x_r\in X\)，使矩阵 \(\bigl(f_i(x_j)\bigr)_{i,j}\) 可逆。
\end{lemma}
\begin{proof}
考虑列向量 \(v(x)=\bigl(f_1(x),\ldots,f_r(x)\bigr)^{\mathsf T}\) 在 \(\Omega^r\) 中的张成空间 \(W\)。若 \(k=\dim W<r\)，取 \(W\) 的基 \(w_1,\ldots,w_k\)。要求行向量 \(c=(c_1,\ldots,c_r)\) 满足 \(cw_j=0\)，只给出 \(k\) 个齐次方程，却有 \(r\) 个未知数，所以存在非零解。这个 \(c\) 也在 \(W\) 的每个向量上取零值，特别有 \(cv(x)=0\)。于是 \(\sum_i c_i\cdot f_i(x)=0\) 对所有 \(x\) 成立，违背独立性。因此这些列向量张成 \(\Omega^r\)，可以从中选出 \(r\) 个组成一组基，所对应的求值矩阵可逆。
\end{proof}
函数独立性排除的是在整个 \(X\) 上成立的关系；这个引理说明，即使 \(X\) 无限，也能从中挑出 \(r\) 个点，让求值矩阵已经足以检验这些系数。所选的点由函数决定，不能任意预先指定。这里的系数属于值域 \(\Omega\)，不必属于嵌入所固定的较小基域。这一点将在不动域定理的次数估计中起作用。

\ProseParagraphBreak
本节的乘法特征是保持乘法的函数，Artin 定理讨论它们的函数线性无关性。第五卷中的表示特征标则由表示矩阵的迹给出，通常不保持乘法；那里的复特征标正交性还需要有限群上的内积。函数线性无关与内积正交是不同的性质。这里不需要内积，也不需要以群阶作分母：即使 \(G\) 有限而 \(\Omega\) 的特征整除 \(|G|\)，独立性定理仍然成立。

\ProseParagraphBreak
本节的嵌入独立性与求值矩阵引理将在\secref{b1:sec:1901}和\secref{b1:sec:1902}用于有限 Galois 对应。\subsecref{b1:subsec:1903:04}的模素数循环型判据则使用\cref{b1:lem:integral-closure-elementary}。

% !TeX root = ../../Book1.tex
% 纲要标题：### 18.2 迹与范数
% 纲要位置：计划纲要.md，第 623 行。

\section{迹与范数}
\label{b1:sec:1802}

% 纲要内容（仅作写作计划，尚非教材正文）：
%
% * 乘法线性算子的迹与行列式
% * 有限扩张中的域迹与域范数
% * 共轭元公式与不可分重数
%

有限扩张 \(L/K\) 中的元素可以通过“乘以这个元素”的线性变换来观察。这个办法不需要扩张可分，也不需要预先列出共轭根。先用矩阵定义迹与范数，再说明共轭公式成立时根应怎样计数。

\subsection{乘法线性算子的迹与行列式}
\label{b1:subsec:1802:01}

设 \([L:K]=n\)，\(a\in L\)。乘法映射
\[
m_a:L\longrightarrow L,\qquad x\longmapsto ax
\]
是 \(K\) 线性的。选定 \(K\) 基以后，它由一个 \(n\) 阶矩阵表示；换基使矩阵相似，所以矩阵迹、行列式和特征多项式与基的选择无关。矩阵迹在相似变换下不变可由 \(\operatorname{tr}(AB)=\operatorname{tr}(BA)\) 得到，后一个等式只须交换有限求和的指标；行列式的不变性来自乘法公式。

\begin{definition}\label{b1:def:field-trace-norm}
设 \(L/K\) 为有限域扩张，\(a\in L\)。令 \(m_a:L\rightarrow L\) 为 \(K\) 线性映射 \(x\mapsto ax\)。把
\[
\operatorname{Tr}_{L/K}(a)=\operatorname{tr}(m_a),\qquad
\operatorname{N}_{L/K}(a)=\det(m_a)
\]
分别称为元素 \(a\) 关于扩张 \(L/K\) 的\term[yu-ji]{域迹}[field trace]与\term[yu-fanshu]{域范数}[field norm]。映射 \(m_a\) 的特征多项式称为 \(a\) 关于 \(L/K\) 的乘法特征多项式，记为 \(\chi_{a,L/K}(t)=\det(tI-m_a)\in K[t]\)。
\end{definition}
\symidx{\(\operatorname{Tr}_{L/K}\)：有限扩张的域迹}
\symidx{\(\operatorname{N}_{L/K}\)：有限扩张的域范数}
\symidx{\(\chi_{a,L/K}(t)\)：乘以 \(a\) 的特征多项式}
这些量的定义都包含整个扩张 \(L/K\)。它们不是只依赖元素最小多项式的绝对数值。

\subsection{有限扩张中的域迹与域范数}
\label{b1:subsec:1802:02}

\begin{proposition}\label{b1:prop:trace-norm-basic}
设 \(L/K\) 为 \(n\) 次有限域扩张。域迹 \(\operatorname{Tr}_{L/K}:L\rightarrow K\) 是 \(K\) 线性映射。对任意 \(a,b\in L\)，域范数满足
\[
\operatorname{N}_{L/K}(ab)=\operatorname{N}_{L/K}(a)\operatorname{N}_{L/K}(b),
\qquad
\operatorname{N}_{L/K}(a)=0\Longleftrightarrow a=0.
\]
若 \(c\in K\)，则 \(\operatorname{Tr}_{L/K}(c)=nc\)、\(\operatorname{N}_{L/K}(c)=c^n\)。
\end{proposition}
\begin{proof}
等式 \(m_{a+b}=m_a+m_b\)、\(m_{ca}=c m_a\) 和迹的线性给出第一项。又有 \(m_{ab}=m_am_b\)，行列式的乘法公式给出范数的乘法性。若 \(a\ne0\)，映射 \(m_a\) 的逆为 \(m_{a^{-1}}\)，所以其行列式非零；若 \(a=0\)，零矩阵的行列式为零。基域元素对应标量矩阵 \(cI_n\)，给出最后两式。
\end{proof}
设 \(a\) 在 \(K\) 上的首一最小多项式为
\[
f(t)=t^d+c_{d-1}t^{d-1}+\cdots+c_0,
\qquad r=[L:K(a)].
\]
令 \(E=K(a)\)。先把乘法算子限制为 \(m_a|_E:E\rightarrow E\)。在 \(E/K\) 的基 \(1,a,\ldots,a^{d-1}\) 下，这个限制算子将前 \(d-1\) 个基向量依次移向后一个，将末项送到 \(-c_0-\cdots-c_{d-1}a^{d-1}\)，矩阵因而是 \(f\) 的友矩阵 \(C_f\)。沿 \(tI-C_f\) 的末列展开并整理，得到 \(\det(tI-C_f)=t^d+c_{d-1}t^{d-1}+\cdots+c_0=f(t)\)。

\ProseParagraphBreak

再取 \(L/E\) 的一组基 \(v_1,\ldots,v_r\)。由塔式公式的基构造，
\[
v_1,av_1,\ldots,a^{d-1}v_1;\quad\ldots;\quad
v_r,av_r,\ldots,a^{d-1}v_r
\]
是 \(L/K\) 的一组基。每个子空间 \(Ev_j\) 都被 \(m_a:L\rightarrow L\) 保持，并且在所列基下具有同一个矩阵 \(C_f\)。因此，整个乘法矩阵为 \(\operatorname{diag}(C_f,\ldots,C_f)\)，共有 \(r\) 个块，从而
\begin{equation}\label{b1:eq:trace-norm-minimal-polynomial}
\chi_{a,L/K}(t)=f(t)^r,\qquad
\operatorname{Tr}_{L/K}(a)=-r c_{d-1},\qquad
\operatorname{N}_{L/K}(a)=(-1)^{dr}c_0^r.
\end{equation}
例如取非平方的 \(D\in\mathbb Q\)，令 \(L=\mathbb Q(\sqrt D)\)，并取 \(a,b\in\mathbb Q\)。乘以 \(a+b\sqrt D\) 的矩阵为
\[
\begin{pmatrix}a&bD\\ b&a\end{pmatrix},
\]
所以迹为 \(2a\)，范数为 \(a^2-Db^2\)。当 \(D<0\) 时，这个范数与复共轭有关，但一般域范数并不是分析学意义的非负长度。

\subsection{共轭元公式与不可分重数}
\label{b1:subsec:1802:03}

\begin{theorem}\label{b1:thm:trace-norm-embeddings}
设 \(L/K\) 为有限域扩张，\(\Omega\) 为包含 \(L\) 的 \(K\) 的代数闭包，\(\sigma_1,\ldots,\sigma_s\) 是全部不同的 \(K\) 嵌入 \(L\rightarrow\Omega\)。记 \(e=[L:K]_{\mathrm i}=[L:K]/s\) 为不可分次数。则对每个 \(a\in L\) 都有
\begin{align*}
\chi_{a,L/K}(t)&=\prod_{i=1}^s\bigl(t-\sigma_i(a)\bigr)^e,\\
\operatorname{Tr}_{L/K}(a)&=e\sum_{i=1}^s\sigma_i(a),&
\operatorname{N}_{L/K}(a)&=\prod_{i=1}^s\sigma_i(a)^e.
\end{align*}
特别地，扩张可分时 \(e=1\)，迹是全部嵌入像的和，范数是其积。
\end{theorem}
\begin{proof}
令 \(E=K(a)\)、\(f=m_{a,K}\)、\(d=[E:K]\)。最小多项式 \(f\) 有 \(d_s=[E:K]_{\mathrm s}\) 个不同根；\cref{b1:prop:inseparable-polynomial-shape}说明正特征下各根重数相同，特征零下各根重数为一，所以统一为 \(d/d_s\)。每个 \(K\) 嵌入 \(E\rightarrow\Omega\) 都能延拓到 \(L\)，且延拓的个数恰为 \(r_s=[L:E]_{\mathrm s}\)，这是\cref{b1:thm:separable-degree-tower}证明中的嵌入计数。因此 \(s=d_s r_s\)，而 \(\sigma_i(a)\) 中每个不同根出现 \(r_s\) 次。令 \(r=[L:E]\)，则
\[
e r_s=\frac{dr}{d_s r_s}r_s=\frac d{d_s}r.
\]
右端恰是 \(f(t)^r\) 中每个不同根的重数。由\eqref{b1:eq:trace-norm-minimal-polynomial}，这就证明特征多项式公式。比较 \(t^{[L:K]-1}\) 的系数与常数项，分别得到迹和范数的公式。
\end{proof}
公式中的重复有两个来源，即不同嵌入在指定元素上取同一值，以及每个嵌入像还要按不可分次数 \(e\) 计重。即使扩张可分，第一种重复也可能出现。取 \(L=\mathbb Q(\sqrt2,\sqrt3)\)、\(K=\mathbb Q\)、\(a=\sqrt2\)。分别为 \(\sqrt2\) 和 \(\sqrt3\) 选择正负像，得到四个不同的 \(K\) 嵌入；它们在 \(a\) 上的取值为 \(\sqrt2,\sqrt2,-\sqrt2,-\sqrt2\)。此时 \(e=1\)，但 \(\chi_{a,L/K}(t)=(t^2-2)^2\)。因此，求迹或范数时仍须保留不同嵌入带来的重复取值。

\ProseParagraphBreak

另一种重复来自不可分次数。设 \(p\) 为素数，取在 \(\mathbb F_p\) 上超越的元素 \(u\)，并令 \(L=\mathbb F_p(u)\)、\(K=\mathbb F_p(u^p)\)。由\eqref{b1:eq:rational-power-degree}知 \([L:K]=p\)；每个 \(x\in L\) 都满足 \(x^p\in K\)，故这是纯不可分扩张，只有一个 \(K\) 嵌入，而 \(e=p\)。于是
\[
\chi_{u,L/K}(t)=(t-u)^p=t^p-u^p,\qquad
\operatorname{Tr}_{L/K}(u)=pu=0,\qquad
\operatorname{N}_{L/K}(u)=u^p.
\]
事实上，这个扩张中每个元素的迹都是零。唯一嵌入的像若不按 \(e\) 计重，就无法得到这些结果。

% !TeX root = ../../Book1.tex
% 纲要标题：### 18.3 传递性与判别式
% 纲要位置：计划纲要.md，第 629 行。

\section{传递性与判别式}
\label{b1:sec:1803}

% 纲要内容（仅作写作计划，尚非教材正文）：
%
% * 迹、范数的塔式公式
% * 迹配对及可分性的判别；证明迹对偶基的存在、唯一性和坐标展开公式
% * 基的判别式及换基公式；区分在代数闭包中为平方与在基域中为平方，以及基判别式和多项式判别式
%

在一座域扩张塔中，先把乘法看作中间域上的线性变换，再把标量限制到基域，所得迹与范数应当与一步计算相同。迹还可以比较两个元素：对乘积取迹便得到一个双线性配对。这个配对是否退化，恰好检测扩张的可分性。

\subsection{迹、范数的塔式公式}
\label{b1:subsec:1803:01}

\begin{theorem}\label{b1:thm:trace-norm-tower}
若 \(K\subseteq E\subseteq L\) 为有限域扩张塔，则对 \(a\in L\) 有
\begin{align*}
\operatorname{Tr}_{L/K}(a)&=\operatorname{Tr}_{E/K}\bigl(\operatorname{Tr}_{L/E}(a)\bigr),\\
\operatorname{N}_{L/K}(a)&=\operatorname{N}_{E/K}\bigl(\operatorname{N}_{L/E}(a)\bigr).
\end{align*}
\end{theorem}
\begin{proof}
取 \(E/K\) 的基 \(u_1,\ldots,u_m\) 与 \(L/E\) 的基 \(v_1,\ldots,v_r\)。由\cref{b1:thm:tower-law}的基构造，按
\[
u_1v_1,\ldots,u_mv_1;\quad\ldots;\quad u_1v_r,\ldots,u_mv_r
\]
排列的全部乘积为 \(L/K\) 的基。对 \(d\in E\)，记 \(C(d)\) 为乘法映射 \(x\mapsto dx\) 在基 \(u_1,\ldots,u_m\) 下的 \(K\) 矩阵。设乘以 \(a\) 的 \(E\) 矩阵为 \(A=(a_{ij})\)，即 \(av_j=\sum_i a_{ij}v_i\)。那么输入的第 \(j\) 个坐标属于 \(E\)，乘以 \(a\) 后对第 \(i\) 个坐标的贡献是乘以 \(a_{ij}\)；把这个 \(E\) 坐标展开为 \(m\) 个 \(K\) 坐标，便得到块 \(C(a_{ij})\)。因此乘以 \(a\) 的 \(K\) 矩阵是 \(A_K=(C(a_{ij}))_{i,j}\)。例如 \(r=2\) 时，这个替换为
\[
\begin{pmatrix}a_{11}&a_{12}\\a_{21}&a_{22}\end{pmatrix}
\quad\longmapsto\quad
\begin{pmatrix}C(a_{11})&C(a_{12})\\C(a_{21})&C(a_{22})\end{pmatrix}.
\]
它的迹只取对角块的迹，故一般情形下总迹为
\[
\sum_j\operatorname{Tr}_{E/K}(a_{jj})
=\operatorname{Tr}_{E/K}\Bigl(\sum_j a_{jj}\Bigr),
\]
即迹的塔式公式。

范数公式所需的是一个对任意 \(A\in M_r(E)\) 都成立的矩阵等式：
\[
\det_K A_K=\operatorname{N}_{E/K}(\det_E A).
\]
这里 \(A_K=(C(a_{ij}))_{i,j}\) 表示同一个 \(E\) 线性映射改看成 \(K\) 线性映射后的矩阵。若 \(A\) 奇异，非零的 \(E\) 核向量展开后仍是非零的 \(K\) 核向量，等式两边皆为零。若 \(A\) 可逆，可用初等行变换化为单位矩阵，其中行交换只用相邻两行的交换即可。加上另一行的 \(c\) 倍，在 \(K\) 矩阵中对应一个对角块均为单位矩阵、只有一个非对角块为 \(C(c)\) 的块初等变换，其行列式为一；而 \(\det_E A\) 也不变。把一行乘以 \(c\ne0\)，对应块行乘以 \(C(c)\)，故 \(K\) 行列式乘以 \(\det C(c)=\operatorname{N}_{E/K}(c)\)；另一边则由范数的乘法性受到同一个因子改变。

交换相邻两行，在 \(K\) 上交换两个各有 \(m\) 行的块。一个块中的每行都要越过另一个块中的每行，共作 \(m^2\) 次普通行交换，符号为 \((-1)^{m^2}\)。由于 \(m^2-m=m(m-1)\) 为偶数，这个符号等于 \((-1)^m\)。另一方面，\cref{b1:prop:trace-norm-basic}对基域元素给出 \(\operatorname{N}_{E/K}(c)=c^m\)，其中 \(m=[E:K]\)，所以 \(\operatorname{N}_{E/K}(-1)=(-1)^m\)。于是两种行列式在每步受到相同的非零因子改变，单位矩阵处又均为一，故所需等式成立。再取 \(A\) 为 \(m_a\) 的 \(E\) 矩阵，便得到范数公式。这个论证不要求可分性。
\end{proof}

\subsection{迹配对及可分性的判别}
\label{b1:subsec:1803:02}

\begin{definition}\label{b1:def:trace-pairing}
设 \(L/K\) 为有限域扩张。双线性映射
\[
L\times L\longrightarrow K,\qquad (x,y)\longmapsto\operatorname{Tr}_{L/K}(xy)
\]
称为 \(L/K\) 的\term[ji-peidui]{迹配对}[trace pairing]。如果对每个非零 \(x\in L\) 都存在 \(y\in L\) 使 \(\operatorname{Tr}_{L/K}(xy)\neq0\)，则称此配对非退化。
\end{definition}
域乘法交换，所以迹配对满足 \(\operatorname{Tr}_{L/K}(xy)=\operatorname{Tr}_{L/K}(yx)\)。因此非退化条件只检查一个变量便已足够；交换两个变量会得到同一个条件。

\begin{theorem}\label{b1:thm:trace-pairing-separable}
设 \(L/K\) 为有限域扩张。扩张 \(L/K\) 可分，当且仅当其迹配对非退化。
\end{theorem}
\begin{proof}
若扩张可分，由\cref{b1:thm:trace-norm-embeddings}，迹函数为全部不同 \(K\) 嵌入的和。\cref{b1:cor:embedding-independence}说明这个和不是零函数，因为其各系数均为非零的一。于是存在 \(z\in L\) 使 \(\operatorname{Tr}_{L/K}(z)\ne0\)。这一个元素已经足以检验配对：对任意 \(x\ne0\)，乘法 \(y\mapsto xy\) 是 \(L\) 的双射，取 \(y=x^{-1}z\)，便有 \(\operatorname{Tr}_{L/K}(xy)\ne0\)。

若扩张不可分，\cref{b1:prop:separable-inseparable-decomposition}说明其特征为某个素数 \(p\)，不可分次数 \(e\) 是大于一的 \(p\) 的幂。由\cref{b1:thm:trace-norm-embeddings}，迹函数是某个和的 \(e\) 倍，而 \(e\cdot1_K=0\)，因而迹函数恒为零。所以整个迹配对为零，必退化。
\end{proof}
有限不可分扩张的迹映射恒为零，而有限可分扩张的迹映射非零，即使它的次数被特征整除。例如，取 \(L=\mathbb F_4=\mathbb F_2(u)\)，其中 \(u^2+u+1=0\)。这个多项式的导数为一，故扩张可分；其次数与特征均为二。由共轭公式，
\[
\operatorname{Tr}_{L/\mathbb F_2}(1)=0,\qquad
\operatorname{Tr}_{L/\mathbb F_2}(u)=u+u^2=1.
\]
所以 \(\operatorname{Tr}_{L/K}(1)=[L:K]\cdot1_K=0\) 只说明单位元的迹为零，不能据此断言所有元素的迹为零。

\begin{corollary}\label{b1:cor:trace-dual-basis}
设 \(L/K\) 为 \(n\) 次有限可分域扩张，\(b_1,\ldots,b_n\) 为 \(L/K\) 的一组基。存在唯一一组元素 \(c_1,\ldots,c_n\in L\)，满足
\[
\operatorname{Tr}_{L/K}(b_i c_j)=\delta_{ij}\qquad(1\le i,j\le n),
\]
其中 \(\delta_{ij}\) 在 \(i=j\) 时为一，否则为零。这组元素也是 \(L/K\) 的基，称为 \(b_1,\ldots,b_n\) 关于迹配对的对偶基。任意 \(x\in L\) 都满足
\[
x=\sum_{j=1}^n\operatorname{Tr}_{L/K}(xc_j)b_j.
\]
\end{corollary}
\begin{proof}
记 \(b=(b_1,\ldots,b_n)\)，并令 \(T_b=(\operatorname{Tr}_{L/K}(b_i b_j))_{i,j}\)。由\cref{b1:thm:trace-pairing-separable}，迹配对非退化。如果 \(T_b\lambda=0\)，则 \(y=\sum_j\lambda_jb_j\) 与每个 \(b_i\) 配对为零，因而与它们的每个 \(K\) 线性组合配对为零。非退化性给出 \(y=0\)，再由 \(b\) 是基得 \(\lambda=0\)。因此 \(T_b\) 可逆；这就是非退化性在有限维坐标中的表达。

令
\[
c_j=\sum_i(T_b^{-1})_{ij}b_i.
\]
等式 \(T_bT_b^{-1}=I_n\) 恰给出 \(\operatorname{Tr}_{L/K}(b_i c_j)=\delta_{ij}\)。任何满足此式的 \(c_j\) 都必须具有 \(T_b^{-1}\) 的第 \(j\) 列坐标，故唯一。若 \(\sum_j\mu_jc_j=0\)，与 \(b_i\) 配对便得 \(\mu_i=0\)，所以这 \(n\) 个元素线性无关，也是 \(L/K\) 的基。最后，若 \(x=\sum_i\lambda_i b_i\)，则 \(\operatorname{Tr}_{L/K}(xc_j)=\lambda_j\)，得到所需的展开式。
\end{proof}

\subsection{基的判别式及换基公式}
\label{b1:subsec:1803:03}

设 \(L/K\) 为 \(n\) 次有限域扩张，\(b=(b_1,\ldots,b_n)\) 是 \(L\) 在 \(K\) 上的一组基。迹配对在这组基下的矩阵为
\[
T_b=\bigl(\operatorname{Tr}_{L/K}(b_i b_j)\bigr)_{i,j},
\]
也称为它的 Gram 矩阵。把
\[
\operatorname{disc}_{L/K}(b)=\det T_b
\]
称为基 \(b\) 关于扩张 \(L/K\) 的\term[ji-panbieshi]{基的判别式}[discriminant of a basis]。
\symidx{\(\operatorname{disc}_{L/K}(b)\)：基 \(b\) 的判别式}
这个行列式用一个标量检测迹配对是否退化。若 \(x,y\in L\) 在基 \(b\) 下的坐标列分别为 \(\xi,\eta\)，则 \(\operatorname{Tr}_{L/K}(xy)=\xi^{\mathsf T}T_b\eta\)。因此，存在非零 \(x\) 与所有 \(y\) 配对为零，当且仅当存在非零 \(\xi\) 满足 \(\xi^{\mathsf T}T_b=0\)，也就是 \(T_b\) 奇异。

\ProseParagraphBreak
若另一组基满足 \(c_j=\sum_i b_i p_{ij}\)，记换基矩阵为 \(P\in\operatorname{GL}_n(K)\)。双线性给出迹矩阵关系 \(T_c=P^{\mathsf T}T_bP\)，所以
\begin{equation}\label{b1:eq:discriminant-change-basis}
\operatorname{disc}_{L/K}(c)=(\det P)^2\operatorname{disc}_{L/K}(b).
\end{equation}
由于 \(\det P\ne0\)，判别式是否为零与基的选择无关。由\cref{b1:thm:trace-pairing-separable}，扩张可分当且仅当一组基的判别式非零，此时每组基的判别式都非零。不同基的非零判别式相差一个非零平方因子，所以它们在商群 \(K^\times/(K^\times)^2\) 中给出同一个类；这里 \((K^\times)^2\) 是 \(K\) 中所有非零平方组成的子群。

\ProseParagraphBreak
在前面的 \(\mathbb F_4/\mathbb F_2\) 例子中，\(u^2=u+1\) 又给出 \(\operatorname{Tr}_{L/\mathbb F_2}(u^2)=1\)。于是基 \(1,u\) 的迹矩阵为
\[
T_{(1,u)}=\begin{pmatrix}0&1\\1&1\end{pmatrix},\qquad \det T_{(1,u)}=1\in\mathbb F_2.
\]
迹配对非退化，但非零元素 \(1\) 与自身配对为零；非退化并不要求每个非零元素与自身配对都非零。

\ProseParagraphBreak
在可分情形，取包含 \(L\) 的 \(K\) 的代数闭包 \(\Omega\)。可分性保证全部不同的 \(K\) 嵌入 \(\sigma_1,\ldots,\sigma_n:L\rightarrow\Omega\) 恰有 \(n=[L:K]\) 个，所以嵌入求值矩阵 \(M=\bigl(\sigma_i(b_j)\bigr)_{i,j}\in M_n(\Omega)\) 是方阵。由共轭公式，
\[
(M^{\mathsf T}M)_{jk}
=\sum_i\sigma_i(b_j)\sigma_i(b_k)
=\sum_i\sigma_i(b_jb_k)
=\operatorname{Tr}_{L/K}(b_jb_k).
\]
因此 \(T_b=M^{\mathsf T}M\)，判别式为 \((\det M)^2\)。不可分时嵌入数小于 \(n\)，求值矩阵不再是 \(n\) 阶方阵，且迹公式还含不可分重数，不能直接沿用这里的等式。此处的平方是在 \(\Omega\) 中取的；\(\det M\) 未必属于 \(K\)，不能据此断言基的判别式是 \(K\) 中的平方。例如，对 \(\mathbb Q(\sqrt2)/\mathbb Q\) 的基 \(1,\sqrt2\)，求值矩阵的行列式为 \(-2\sqrt2\)，其平方为 \(8\)，而 \(8\) 不是有理数的平方。

\ProseParagraphBreak
若 \(L=K(a)\) 可分，取幂基 \(1,a,\ldots,a^{n-1}\)，并记 \(a\) 的全部共轭根为 \(a_1,\ldots,a_n\)。此时求值矩阵为
\[
M=\begin{pmatrix}
1&a_1&\cdots&a_1^{n-1}\\
1&a_2&\cdots&a_2^{n-1}\\
\vdots&\vdots&&\vdots\\
1&a_n&\cdots&a_n^{n-1}
\end{pmatrix}.
\]
这是范德蒙德矩阵。\footnote{范德蒙德（Alexandre-Théophile Vandermonde，1735-1796），法国数学家，研究行列式与方程理论。}其行列式为 \(\prod_{i<j}(a_j-a_i)\)，故
\begin{equation}\label{b1:eq:power-basis-discriminant}
\operatorname{disc}_{L/K}(1,a,\ldots,a^{n-1})
=\prod_{i<j}(a_j-a_i)^2.
\end{equation}
若 \(d\in\mathbb Q\) 不是平方数，则 \(\mathbb Q(\sqrt d)/\mathbb Q\) 的基 \(1,\sqrt d\) 的迹矩阵为 \(\operatorname{diag}(2,2d)\)，判别式为 \(4d\)。基的判别式依赖扩张与所选基，不能直接把任意一组有理基的判别式称为数域整数环的判别式。

\ProseParagraphBreak
根差乘积还给出一个只依赖多项式的量。设 \(K\) 为域，\(f\in K[x]\) 为次数 \(m\ge1\) 的首一多项式，并把它在一个分裂域中的全部根按重数列为 \(a_1,\ldots,a_m\)。把 \(\prod_{i<j}(a_j-a_i)^2\) 称为 \(f\) 的\term[duoxiangshi-panbieshi]{判别式}[discriminant of a polynomial]，记作 \(\operatorname{disc}(f)\)。与 \(\operatorname{disc}_{L/K}(b)\) 不同，这里所需的数据只是多项式 \(f\)，并未选定一个域扩张的基。

\ProseParagraphBreak
交换变量会使 \(\prod_{i<j}(X_j-X_i)\) 至多改变符号，所以其平方是整数系数的对称多项式。由\cref{b1:thm:fundamental-symmetric-polynomials}取系数环为 \(\mathbb Z\)，存在 \(Q\in\mathbb Z[T_1,\ldots,T_m]\)，使它等于 \(Q(e_1,\ldots,e_m)\)，其中 \(e_i\) 是初等对称多项式。若
\[
f(x)=x^m+c_{m-1}x^{m-1}+\cdots+c_0,
\]
根与系数的 Vi\`ete 关系给出 \(e_i(a_1,\ldots,a_m)=(-1)^i c_{m-i}\)。因此
\[
\operatorname{disc}(f)=Q(-c_{m-1},c_{m-2},\ldots,(-1)^m c_0)\in K,
\]
并且这个值不依赖根的排列或所选分裂域。若有两个根重合，根差乘积为零；若所有根互异，每个根差都非零，域中的乘积也非零。故 \(\operatorname{disc}(f)\ne0\) 当且仅当 \(f\) 可分。

\ProseParagraphBreak
当 \(L=K(a)\)、\(b=(1,a,\ldots,a^{n-1})\)、\(f=m_{a,K}\) 时，基的判别式与多项式的判别式相等：可分情形由\eqref{b1:eq:power-basis-discriminant}得到；不可分情形中，迹配对为零，而最小多项式有重根，两种判别式都为零。对于任意其他基，仍须通过换基公式计算，不能把这两个定义直接视为同一个对象。

% !TeX root = ../../Book1.tex
% 纲要标题：### 18.4 迹与范数的算术应用
% 纲要位置：计划纲要.md，第 635 行。

\section{迹与范数的算术应用}
\label{b1:sec:1804}

% 纲要内容（仅作写作计划，尚非教材正文）：
%
% * 有限域中的迹和范数
% * 整元素的迹、范数与整性；区分在子环上整且属于基域与属于该子环
% * 数域整数环与判别式的例子；整基和分歧理论仅说明后续方向
% * 整性的传递性及其限制；补全二次情形的反推，并在正文证明迹、范数为整数仍不整的三次反例
% * 所需整元素的首一方程定义及整性的基本封闭性在本节就地准备；解释有限生成集下的行列式技巧及系数矩阵非唯一性，通过共轭元或有限生成模方法证明迹与范数的整性，不以前置第三卷整扩张理论为条件
%
% ---
%

迹和范数把扩张域中的计算带回基域，但它们的像并非在所有场合都一样。有限域中可以精确求出像与核；数域中则要同时考虑元素是否满足整数系数的首一方程。本节从首一方程出发，建立所需的整性事实。

\subsection{有限域中的迹和范数}
\label{b1:subsec:1804:01}

\begin{proposition}\label{b1:prop:finite-field-trace-norm}
设 \(q\) 为素数幂，\(n\geq1\) 为整数，并令 \(L=\mathbb F_{q^n}\)、\(K=\mathbb F_q\)。则对每个 \(a\in L\) 都有
\[
\operatorname{Tr}_{L/K}(a)=a+a^q+\cdots+a^{q^{n-1}},\qquad
\operatorname{N}_{L/K}(a)=a^{1+q+\cdots+q^{n-1}}.
\]
迹作为加法群同态满射到 \(K\)，其核有 \(q^{n-1}\) 个元素；范数把 \(L^\times\) 满射到 \(K^\times\)，其核有 \((q^n-1)/(q-1)\) 个元素。
\end{proposition}
\begin{proof}
有限域的 Frobenius 映射是有限集合上的单射，故满射；\cref{b1:prop:perfect-frobenius}遂说明 \(K\) 完美，所以有限扩张 \(L/K\) 可分。由\cref{b1:thm:separable-embeddings}，到代数闭包的 \(K\) 嵌入恰有 \(n\) 个；\cref{b1:thm:finite-field-automorphisms}给出的 \(n\) 个互异自同构 \(a\mapsto a^{q^i}\) 因而已经包含全部嵌入。对第十七章得到的这组 Frobenius 像应用\cref{b1:thm:trace-norm-embeddings}，求和、求积便得到两条公式。若 \(a\) 的 Frobenius 轨道长度小于 \(n\)，这 \(n\) 项中会有重复；公式仍按全部嵌入计数。

由\cref{b1:thm:trace-pairing-separable}，迹不是零映射。其像是 \(K\) 的非零 \(K\) 子空间，而 \(\dim_K K=1\)，所以像只能是整个 \(K\)。再由维数公式，迹的核维数为 \(n-1\)，从而有 \(q^{n-1}\) 个元素。

取 \(L^\times\) 的生成元 \(g\)，令 \(N=q^n-1\)、\(h=N/(q-1)\)。范数公式给出 \(\operatorname{N}_{L/K}(g)=g^h\)，由循环群中元素幂的阶公式及 \(h\mid N\)，有
\[
\operatorname{ord}(g^h)
=\frac{N}{\gcd(N,h)}
=\frac Nh
=q-1.
\]
由于范数的值属于 \(K^\times\)，其像包含这个阶为 \(q-1\) 的元素，故范数满射到 \(K^\times\)。第一同构定理给出 \(L^\times/\ker\operatorname{N}_{L/K}\cong K^\times\)，因此
\[
|\ker\operatorname{N}_{L/K}|
=\frac{|L^\times|}{|K^\times|}
=\frac{q^n-1}{q-1}.
\]
\end{proof}
例如在 \(\mathbb F_4=\mathbb F_2(u)\)、\(u^2+u+1=0\) 中，\(\operatorname{Tr}(1)=0\)，但 \(\operatorname{Tr}(u)=u+u^2=1\)；这正展示了迹非零并不要求单位元的迹非零。

\subsection{整元素的迹与范数}
\label{b1:subsec:1804:02}

\begin{definition}\label{b1:def:integral-element-field}
设 \(A\) 是域 \(L\) 的子环。如果 \(a\in L\) 满足某个次数 \(d\geq1\)、系数 \(c_0,\ldots,c_{d-1}\in A\) 的首一方程
\[
a^d+c_{d-1}a^{d-1}+\cdots+c_0=0,
\]
则称 \(a\) 为 \(A\) 上的\term[zheng-yuansu]{整元素}[integral element]，也说它在 \(A\) 上整。如果代数数 \(a\in\overline{\mathbb Q}\) 在 \(\mathbb Z\) 上整，则称 \(a\) 为\term[daishu-zhengshu]{代数整数}[algebraic integer]。
\end{definition}
由于 \(A\) 是域的子环，它是整环，其分式域 \(F=\operatorname{Frac}(A)\) 也包含于 \(L\)。在 \(A\) 上整蕴含在 \(F\) 上代数，但整性要求方程首一，而且系数都在 \(A\) 中，不能仅属于 \(F\)。例如 \(1/2\) 在 \(\mathbb Q\) 上代数，却不在 \(\mathbb Z\) 上整：若有首一 \(d\) 次整数方程，乘以 \(2^d\) 后会得到一加上一个偶数等于零，矛盾。当 \(A=F\) 本身是域时，非零多项式可除以首项系数化为首一，整性才与代数性等价。

\ProseParagraphBreak
设 \(K/\mathbb Q\) 为域扩张。如果 \([K:\mathbb Q]<\infty\)，则称 \(K\) 为\term[shuyu]{数域}[number field]。为将整性用于迹与范数，先给出求和、求积时所需的封闭性。

\ProseParagraphBreak
下面的\cref{b1:lem:integral-closure-elementary}还有一个后续用途：\subsecref{b1:subsec:1903:04}会用根生成环的有限性取得有限剩余域，并将 Frobenius 的作用提升到原分裂域。后面的整数环与判别式例子则说明整性在具体数域中的表现。
\begin{lemma}\label{b1:lem:integral-closure-elementary}
设 \(A\) 是域 \(L\) 的子环，\(a_1,\ldots,a_s\in L\) 都在 \(A\) 上整。则存在有限多个 \(v_1,\ldots,v_m\in A[a_1,\ldots,a_s]\)，使
\[
A[a_1,\ldots,a_s]=Av_1+\cdots+Av_m.
\]
也就是说，其中每个元素都是某个 \(\sum_{i=1}^m c_iv_i\)、\(c_i\in A\)；这称为有限生成 \(A\) 模，不要求系数表示唯一。其中每个元素都在 \(A\) 上整。特别地，\(L\) 中在 \(A\) 上整的元素的和、差与积仍在 \(A\) 上整。
\end{lemma}
\begin{proof}
首一方程使每个 \(a_i\) 的高次幂都能降次，因而先把任意多项式表达压缩成有限多个单项式的线性组合。具体地，设 \(a_i\) 满足次数为 \(d_i\) 的首一方程，反复用这些方程约去高次幂，\(A[a_1,\ldots,a_s]\) 就由 \(a_1^{e_1}\cdots a_s^{e_s}\)、\(0\leq e_i<d_i\)，作为 \(A\) 模生成。记这些生成元为 \(v_1,\ldots,v_m\)，其中取 \(v_1=1\)。这个生成元稍后会让“同时零化所有生成元”的标量本身等于零。

现在给定 \(b\in A[a_1,\ldots,a_s]\)，乘以 \(b\) 仍把这个环送入自身，所以每个 \(bv_i\) 都能写成生成元的线性组合。任选一组系数使 \(bv_i=\sum_j c_{ij}v_j\)、\(c_{ij}\in A\)，令 \(C=(c_{ij})\)。这里的生成元未必是一组基，系数表示和 \(C\) 都可能不唯一；\(C\) 只记录所选的乘法关系，不能直接把它的迹和行列式解释为域迹与域范数。

令 \(v=(v_i)^{\mathsf T}\)，这些关系写成 \((bI-C)v=0\)。由行列式的余子式展开，伴随矩阵满足 \(\operatorname{adj}(bI-C)(bI-C)=\det(bI-C)I\)；乘到上式得到 \(\det(bI-C)v_i=0\) 对每个 \(i\) 成立。特别地，取 \(v_1=1\)，便得 \(\det(bI-C)\cdot1=0\)，即 \(\det(bI-C)=0\)。多项式 \(\det(tI-C)\in A[t]\) 首一，故 \(b\) 在 \(A\) 上整。整个论证只用有限生成集和余子式恒等式，没有假设这个模是自由模。
\end{proof}
这个行列式技巧从有限个乘法关系得到首一方程，其一般模论形式将在第三卷系统讨论。生成元间的关系也会影响所选的系数矩阵，下面的计算说明这种选择不唯一的含义。

\ProseParagraphBreak
例如用 \(v_1=v_2=1\) 生成 \(\mathbb Z\)，考虑乘以二。按照证明中每一行表示 \(2v_i\) 的约定，可以分别选取
\[
C_1=\begin{pmatrix}2&0\\2&0\end{pmatrix},
\qquad C_2=\begin{pmatrix}2&0\\0&2\end{pmatrix}.
\]
两者都满足 \((2I-C_j)(1,1)^{\mathsf T}=0\)，但所得的首一多项式分别为
\[
\det(tI-C_1)=t(t-2),\qquad
\det(tI-C_2)=(t-2)^2.
\]
它们不同，却都在 \(t=2\) 处为零。这里 \(C_1\) 的迹为二，\(C_2\) 的迹为四，而乘以二在一维 \(\mathbb Q\) 空间上的迹是二。这说明任意选择的系数矩阵的迹不能直接当作前文的域迹；前文使用的是基下的乘法矩阵。行列式技巧得到的是某个首一零化多项式，没有在这个生成集上定义一个内在确定的特征多项式。

\begin{proposition}\label{b1:prop:integral-trace-norm}
设 \(L/K\) 为有限扩张，\(A\subseteq K\) 为子环。若 \(a\in L\) 在 \(A\) 上整，则 \(\operatorname{Tr}_{L/K}(a)\) 和 \(\operatorname{N}_{L/K}(a)\) 也是 \(K\) 中在 \(A\) 上整的元素。特别地，数域中的代数整数对 \(\mathbb Q\) 的迹和范数都是整数。
\end{proposition}
\begin{proof}
取包含 \(L\) 的代数闭包 \(\Omega\)。设 \(a\) 满足首一方程 \(a^d+c_{d-1}a^{d-1}+\cdots+c_0=0\)，其中 \(c_i\in A\)。每个 \(K\) 嵌入 \(\sigma:L\rightarrow\Omega\) 都逐点固定 \(A\subseteq K\)，所以对方程应用 \(\sigma\) 得到
\[
\sigma(a)^d+c_{d-1}\sigma(a)^{d-1}+\cdots+c_0=0.
\]
因此每个共轭像也在 \(A\) 上整。由\cref{b1:thm:trace-norm-embeddings}，迹和范数是这些共轭像按不可分重数形成的有限和、积。\cref{b1:lem:integral-closure-elementary}保证它们整，而定义保证它们属于 \(K\)。

若 \(K=\mathbb Q,A=\mathbb Z\)，目前只知迹与范数是在 \(\mathbb Z\) 上整的有理数，还须说明这种有理数只能是整数。写既约分数 \(r/s\)、\(s>0\)，把它满足的首一 \(d\) 次整数方程乘以 \(s^d\)，得到 \(s\mid r^d\)。由 \(\gcd(r,s)=1\)，只能 \(s=1\)。所以迹与范数为整数。
\end{proof}
一般情形下，命题给出的结论是在 \(A\) 上整且属于 \(K\)，并不保证属于 \(A\)。数域情形能得到整数，额外使用了刚才证明的“在 \(\mathbb Z\) 上整的有理数都是整数”；换成任意子环，这一步可能失败。

\subsection{数域整数环与判别式的例子}
\label{b1:subsec:1804:03}

对于数域 \(K\)，集合
\[
\mathcal O_K=\{\,a\in K\mid a\;\text{在}\;\mathbb Z\;\text{上整}\,\}
\]
称为它的\term[zhengshu-huan]{整数环}[ring of integers]。这里仅取域 \(K\) 内的代数整数。
\symidx{\(\mathcal O_K\)：数域 \(K\) 的整数环}
\cref{b1:lem:integral-closure-elementary}说明它对加、减、乘封闭，并含一，所以确实是环。成为环不要求对非零元素取逆封闭；例如二在 \(\mathcal O_K\) 中，\(1/2\) 却不在其中，故 \(\mathcal O_K\) 不是域。

\ProseParagraphBreak
在 \(K=\mathbb Q(\sqrt5)\) 中，\(\sqrt5\) 满足 \(t^2-5=0\)，而 \(\omega=(1+\sqrt5)/2\) 满足
\[
t^2-t-1=0,
\]
所以两者都是代数整数。比较 \(1,\sqrt5\) 的有理系数可知 \(\omega\notin\mathbb Z[\sqrt5]\)，因而 \(\mathbb Z[\sqrt5]\) 没有包含 \(K\) 内全部代数整数。同样地，对于无平方因子整数 \(d\equiv1\pmod4\)，元素 \((1+\sqrt d)/2\) 满足首一整数方程 \(t^2-t+(1-d)/4=0\)。

\ProseParagraphBreak
这个例子也能说明\cref{b1:prop:integral-trace-norm}中的结论不能改成“迹与范数属于 \(A\)”。取 \(A=\mathbb Z[\sqrt5]\)、\(K=L=\mathbb Q(\sqrt5)\)，并令 \(a=\omega\)。由 \(a^2-a-1=0\)，元素 \(a\) 在 \(A\) 上整，但 \(a\notin A\)。在恒等扩张 \(K/K\) 中，
\[
\operatorname{Tr}_{K/K}(a)=\operatorname{N}_{K/K}(a)=a\notin A.
\]

\ProseParagraphBreak
迹配对把数域的线性结构与整数环的算术联系起来。例如在 \(\mathbb Q(\sqrt5)/\mathbb Q\) 中，\(\omega^2=\omega+1\) 给出 \(\operatorname{Tr}(\omega)=1\)、\(\operatorname{Tr}(\omega^2)=3\)。两组基的迹矩阵分别为
\[
T_{(1,\omega)}=\begin{pmatrix}2&1\\1&3\end{pmatrix},\qquad
T_{(1,\sqrt5)}=\begin{pmatrix}2&0\\0&10\end{pmatrix},
\]
故判别式分别为五和二十。由 \(\sqrt5=2\omega-1\)，从前一组基到后一组基的换基矩阵行列式为二，因此 \(20=2^2\cdot5\)，符合\eqref{b1:eq:discriminant-change-basis}。这里首先算的是两组有理基的判别式。若一组基还能让 \(\mathcal O_K\) 中每个元素唯一写成这些基元素的整数线性组合，它就是整数环的整基，其判别式才称为数域判别式。关于整基的存在、素理想分解，以及数域判别式和分歧之间的关系，本节尚未建立相应的整数环结构定理。特别不能把任意一组有理基的判别式中的每个素因子都直接宣布为数域中的分歧素数。

\subsection{整性的传递性及其限制}
\label{b1:subsec:1804:04}

代数性与整性都可以通过方程传递，但整性的证明还必须保持首一和系数所在的环。一条方程只含有限多个系数；即使中间环很大，也只需把这些系数生成的子环纳入论证。

\begin{proposition}\label{b1:prop:integrality-transitive-elementary}
设 \(\Omega\) 为域，\(A\subseteq B\subseteq\Omega\) 为子环，并且 \(B\) 中每个元素都在 \(A\) 上整。如果 \(c\in\Omega\) 在 \(B\) 上整，则 \(c\) 在 \(A\) 上整。
\end{proposition}
\begin{proof}
\(B\) 本身未必由有限多个元素生成，但 \(c\) 满足的一条首一方程只含有限多个系数，记为 \(b_1,\ldots,b_s\)。这与代数性传递时截取有限多个系数的理由相同。由\cref{b1:lem:integral-closure-elementary}，\(B_0=A[b_1,\ldots,b_s]\) 有有限个 \(A\) 模生成元 \(u_1,\ldots,u_m\)，并可取 \(u_1=1\)。若该方程次数为 \(d\)，降次可知 \(B_0[c]\) 由 \(1,c,\ldots,c^{d-1}\) 作为 \(B_0\) 模生成；再把每个系数写成 \(u_i\) 的 \(A\) 线性组合，便得到
\[
B_0[c]=\sum_{i=1}^m\sum_{j=0}^{d-1}A u_i c^j.
\]
这给出有限个 \(A\) 模生成元 \(u_i c^j\)，其中 \(u_1c^0=1\)。乘以 \(c\) 保持环 \(B_0[c]\)，所以对这组生成元应用\cref{b1:lem:integral-closure-elementary}证明中的伴随矩阵论证，便得到 \(c\) 的首一 \(A\) 方程。
\end{proof}
整元素的逆未必整，例如二是代数整数而 \(1/2\) 不是。迹、范数均为整数也不能在任意次数下反推元素整，因为它们只给出乘法特征多项式的 \(t^{n-1}\) 系数和常数项。对于二次数域 \(K/\mathbb Q\)，这两个系数已经是全部非首项系数：若 \(a\in K\) 的迹、范数均为整数，则
\[
t^2-\operatorname{Tr}_{K/\mathbb Q}(a)t+\operatorname{N}_{K/\mathbb Q}(a)
\]
是首一整数多项式。由\eqref{b1:eq:trace-norm-minimal-polynomial}，它是 \(a\) 的最小多项式的正整数次幂，因此在 \(a\) 处为零，\(a\) 是代数整数。次数至少为三时，迹和范数没有控制其余中间系数，这个反推可能失败。

\ProseParagraphBreak
例如在 \(\overline{\mathbb Q}\) 中取 \(a\) 为 \(f(t)=t^3-\dfrac12t-1\) 的一个根，并令 \(K=\mathbb Q(a)\)。整数多项式 \(2t^3-t-2\) 的有理根只能是 \(\pm1,\pm2,\pm\dfrac12\)，逐个代入都不是根。三次多项式在 \(\mathbb Q\) 上可约必有有理根，所以 \(f\) 不可约，是 \(a\) 的最小多项式，且 \([K:\mathbb Q]=3\)。它也是 \(a\) 在这个三次扩张中的乘法特征多项式，故
\[
\operatorname{Tr}_{K/\mathbb Q}(a)=0,
\qquad
\operatorname{N}_{K/\mathbb Q}(a)=1.
\]
但 \(a\) 不是代数整数。否则 \(f\) 的三个根都是 \(a\) 的嵌入像，也都是代数整数；\cref{b1:lem:integral-closure-elementary}说明由它们的和、积形成的初等对称多项式值仍是代数整数。因此 \(f\) 的所有系数都在 \(\mathbb Z\) 上整，而它们属于 \(\mathbb Q\)，由\cref{b1:prop:integral-trace-norm}证明中的既约分数论证，应全为整数。这与系数 \(-\dfrac12\) 矛盾。


\begin{chaptersummary}
不同域嵌入之间不存在常系数的非平凡线性关系，由此可以证明有限可分扩张的迹不是零映射。迹与范数定义为乘法算子的迹与行列式；共轭公式须同时计入嵌入的重复取值与不可分重数。塔式公式适用于所有有限扩张，而迹配对非退化恰好刻画可分性。判别式是迹矩阵的行列式，换基时乘以换基行列式的平方。

整元素对加法和乘法封闭，整元素的迹与范数也保持整性；数域中的代数整数对 \(\mathbb Q\) 的迹与范数是整数。下一章先用嵌入独立性计算不动域的次数，再建立有限 Galois 对应。本章的整性结论还会用于随后按素数约化的计算。
\end{chaptersummary}

\BookExercises

\cref{b1:ex:18-characters-integers}直接运用嵌入独立性；\cref{b1:ex:18-evaluation-matrix}先计算求值矩阵，再在\secref{b1:sec:1803}建立的判别式公式下解释其行列式。迹与范数的计算从\cref{b1:ex:18-cubic-trace-norm}开始，末几题则使用\secref{b1:sec:1804}的整性结论。

\begin{exercise}\label{b1:ex:18-characters-integers}
设 \(a_1,\ldots,a_r\) 为域 \(\Omega\) 中不同的非零元素。对整数 \(n\)，令 \(\chi_i(n)=a_i^n\)。证明这些函数线性无关，并说明只检验 \(n=0,\ldots,r-1\) 为何已经足够。
\end{exercise}
\begin{solution}
它们是加法群 \(\mathbb Z\) 到 \(\Omega^\times\) 的不同乘法特征，可直接用\cref{b1:thm:artin-independence}。若只取指定的 \(r\) 个整数，求值矩阵为 Vandermonde 矩阵，其行列式为 \(\prod_{i<j}(a_j-a_i)\ne0\)。因此在这些点成立的 \(r\) 个齐次方程已经迫使全部系数为零。能预先指定这组整数，是幂函数形成 Vandermonde 矩阵的特殊性质；一般独立函数只保证存在某些使求值矩阵可逆的点。
\end{solution}

\begin{exercise}\label{b1:ex:18-evaluation-matrix}
设 \(L=\mathbb Q(\sqrt2,\sqrt3)\)，基 \(b=(1,\sqrt2,\sqrt3,\sqrt6)\)。

1. 将独立改变两个平方根符号的四个 \(\mathbb Q\) 嵌入作为行，写出求值矩阵 \(M=(\sigma_i(b_j))\)，并求 \((\det M)^2\)。

2. 利用\secref{b1:sec:1803}的判别式公式求基 \(b\) 的判别式，再直接计算迹矩阵 \(T_b=(\operatorname{Tr}_{L/\mathbb Q}(b_i b_j))\) 的行列式加以复核。
\end{exercise}
\begin{solution}
1. \cref{b1:prop:biquadratic-example}给出扩张次数为四及题中所列的基。取行的符号依次为 \((+,+)\)、\((-,+)\)、\((+,-)\)、\((-,-)\)，矩阵可写成
\[
\begin{pmatrix}1&1&1&1\\1&-1&1&-1\\1&1&-1&-1\\1&-1&-1&1\end{pmatrix}
\operatorname{diag}(1,\sqrt2,\sqrt3,\sqrt6).
\]
把左矩阵记为 \(H\)。有 \(HH^{\mathsf T}=4I_4\)，取行列式得到 \((\det H)^2=4^4=256\)；右矩阵的行列式为六。总行列式平方为 \(256\cdot36=9216\)。题目与后面的判别式都只需要这个平方，所以无须另外决定 \(\det H\) 的正负号。

2. 扩张可分，\subsecref{b1:subsec:1803:03}给出 \(T_b=M^{\mathsf T}M\)，所以基的判别式为 \(9216\)。直接计算得 \(T_b=\operatorname{diag}(4,8,12,24)\)，其行列式同为 \(4\cdot8\cdot12\cdot24=9216\)。
\end{solution}

\begin{exercise}\label{b1:ex:18-cubic-trace-norm}
令 \(a=\sqrt[3]2\)，\(L=\mathbb Q(a)\)。求 \(a\)、\(a^2\)、\(1+a\) 对 \(\mathbb Q\) 的迹与范数。
\end{exercise}
\begin{solution}
\(t^3-2\) 由素数二的 Eisenstein 判据不可约，所以 \([L:\mathbb Q]=3\)。由\eqref{b1:eq:trace-norm-minimal-polynomial}，\(a\) 的迹为零、范数为二。元素 \(a^2\) 仍生成 \(L\)，因为 \(a=(a^2)^2/2\)，故其最小多项式也为三次。它满足首一三次式 \(t^3-4\)，最小多项式整除这个式子，二者同次数且首一，因而相等，给出迹零、范数四。又 \(a=(1+a)-1\)，所以 \(1+a\) 也生成 \(L\)，其最小多项式仍为三次。同样，首一三次零化式 \((t-1)^3-2=t^3-3t^2+3t-3\) 就是它的最小多项式，故迹三、范数三。
\end{solution}

\begin{exercise}\label{b1:ex:18-ambient-extension}
比较 \(\sqrt2\) 在 \(\mathbb Q(\sqrt2)/\mathbb Q\) 与 \(\mathbb Q(\sqrt2,\sqrt3)/\mathbb Q\) 中的迹、范数和乘法特征多项式。
\end{exercise}
\begin{solution}
前一个扩张的最小多项式与特征多项式均为 \(t^2-2\)，迹为零，范数为负二。后一个扩张对 \(\mathbb Q(\sqrt2)\) 的次数为二，因此\eqref{b1:eq:trace-norm-minimal-polynomial}给出特征多项式 \((t^2-2)^2\)，迹仍为零，范数变为四。元素本身和最小多项式都未变，但所处的乘法空间维数不同。
\end{solution}

\begin{exercise}\label{b1:ex:18-not-ring-maps}
在 \(\mathbb F_4=\mathbb F_2(u)\)、\(u^2+u+1=0\) 中，计算全部元素对 \(\mathbb F_2\) 的迹与范数，并用 \(u,u+1\) 检验迹是否保持乘法、范数是否保持加法。再对纯不可分扩张 \(L=\mathbb F_2(v)\)、\(K=\mathbb F_2(v^2)\)（\(v\) 为超越元），判断迹是否保持乘法、范数是否保持加法，以及它们是否固定 \(K\)。
\end{exercise}
\begin{solution}
在四元域中，\(\operatorname{Tr}(x)=x+x^2\)、\(\operatorname{N}(x)=x^3\)。因此零和一的迹都为零，\(u,u+1\) 的迹都为一；零的范数为零，其余三个元素的范数都为一。因为 \(u(u+1)=1\)，有
\[
\operatorname{Tr}(u(u+1))=0\ne1=\operatorname{Tr}(u)\operatorname{Tr}(u+1).
\]
而 \(u+(u+1)=1\)，故
\[
\operatorname{N}(u+(u+1))=1\ne0=\operatorname{N}(u)+\operatorname{N}(u+1).
\]
这里迹仍保持加法和 \(\mathbb F_2\) 标量乘法，范数仍保持乘法；它们缺少的是另一个运算的相容性。

第二个扩张的次数为二，且为纯不可分扩张，到代数闭包的唯一 \(K\) 嵌入为包含映射。因此共轭公式给出 \(\operatorname{Tr}_{L/K}(x)=2x=0\)、\(\operatorname{N}_{L/K}(x)=x^2\)。零迹确实保持乘法，但不保持单位，也不固定基域。范数在特征二下还保持加法和单位，因而是域同态；然而 \(\operatorname{N}_{L/K}(v^2)=v^4\ne v^2\)，所以它也不逐点固定 \(K\)。一般不相容的运算，在特殊扩张中可能相容，须保留结论的适用范围。
\end{solution}

\begin{exercise}\label{b1:ex:18-tower-calculation}
在 \(L=\mathbb Q(\sqrt2,\sqrt3)\) 中，沿中间域 \(E=\mathbb Q(\sqrt2)\) 计算 \(a=1+\sqrt2+\sqrt3\) 对 \(\mathbb Q\) 的迹与范数。
\end{exercise}
\begin{solution}
在 \(L/E\) 中共轭把 \(\sqrt3\) 变号，所以
\[
\operatorname{Tr}_{L/E}(a)=2+2\sqrt2,\qquad
\operatorname{N}_{L/E}(a)=(1+\sqrt2)^2-3=2\sqrt2.
\]
再对 \(E/\mathbb Q\) 计算，前者迹为四，后者范数为 \((2\sqrt2)(-2\sqrt2)=-8\)。由\cref{b1:thm:trace-norm-tower}，这就是 \(a\) 在整个四次扩张中的迹与范数。负值并不矛盾：域范数是乘法算子的行列式，不是分析学中衡量长度的范数。
\end{solution}

\begin{exercise}\label{b1:ex:18-purely-inseparable}
设 \(u\) 为超越元，\(L=\mathbb F_p(u)\)、\(K=\mathbb F_p(u^p)\)。证明对每个 \(a\in L\)，有 \(\operatorname{Tr}_{L/K}(a)=0\)、\(\operatorname{N}_{L/K}(a)=a^p\)。
\end{exercise}
\begin{solution}
每个有理函数的 \(p\) 次幂属于 \(K\)，扩张为 \(p\) 次纯不可分扩张。到代数闭包的唯一 \(K\) 嵌入为包含映射，因为 \(u\) 的像必须满足 \(t^p-u^p=(t-u)^p=0\)。\cref{b1:thm:trace-norm-embeddings}中的 \(s=1,e=p\)，故迹为 \(pa=0\)，范数为 \(a^p\)。这也适用于 \(a=0\)。
\end{solution}

\begin{exercise}\label{b1:ex:18-trace-eight}
在 \(\mathbb F_8=\mathbb F_2(\alpha)\)、\(\alpha^3+\alpha+1=0\) 中，列出对 \(\mathbb F_2\) 的迹为零的全部元素，并求每个非零元素的范数。
\end{exercise}
\begin{solution}
\(\operatorname{Tr}(1)=1\)，而 \(\operatorname{Tr}(\alpha)=\alpha+\alpha^2+\alpha^4=0\)。特征为二，且 \(x^8=x\)，所以对任意 \(x\in\mathbb F_8\)，有
\[
\operatorname{Tr}(x)^2=x^2+x^4+x^8
=x^2+x^4+x=\operatorname{Tr}(x^2).
\]
故 \(\operatorname{Tr}(\alpha^2)=0\)。由迹的 \(\mathbb F_2\) 线性，\(a+b\alpha+c\alpha^2\) 的迹为 \(a\)，核正是 \(0,\alpha,\alpha^2,\alpha+\alpha^2\)。非零元的范数为其七次幂，均等于一。
\end{solution}

\begin{exercise}\label{b1:ex:18-trace-dual-basis}
在 \(\mathbb F_8=\mathbb F_2(\alpha)\)、\(\alpha^3+\alpha+1=0\) 中，求基 \(b=(1,\alpha,\alpha^2)\) 关于迹配对的对偶基，并用它恢复任意元素 \(x=A+B\alpha+C\alpha^2\) 的三个坐标。再求基 \(b'=(1,\alpha^2,\alpha^4)\) 的对偶基，解释这与平方自同构有什么关系。
\end{exercise}
\begin{solution}
由前题的计算，\(\operatorname{Tr}(1)=1\)、\(\operatorname{Tr}(\alpha)=\operatorname{Tr}(\alpha^2)=0\)。又 \(\alpha^3=\alpha+1\)、\(\alpha^4=\alpha^2+\alpha\)，故 \(\operatorname{Tr}(\alpha^3)=1\)、\(\operatorname{Tr}(\alpha^4)=0\)。因此迹矩阵为
\[
T_b=\begin{pmatrix}1&0&0\\0&0&1\\0&1&0\end{pmatrix},\qquad T_b^{-1}=T_b.
\]
\cref{b1:cor:trace-dual-basis}给出对偶基 \(c=(1,\alpha^2,\alpha)\)，以及
\[
A=\operatorname{Tr}(x),\qquad
B=\operatorname{Tr}(x\alpha^2),\qquad
C=\operatorname{Tr}(x\alpha).
\]
可逆的迹矩阵使配对能够逐个读取坐标，这就是非退化性的具体作用。

平方自同构固定 \(\mathbb F_2\)，把 \(b\) 送到 \(b'\)，所以 \(b'\) 也是基。由前题 \(\operatorname{Tr}(z^2)=\operatorname{Tr}(z)^2\)，有
\[
\operatorname{Tr}(b_i^2c_j^2)=\operatorname{Tr}(b_ic_j)^2=\delta_{ij}.
\]
因此 \(b'\) 的唯一对偶基为 \(c'=(1,\alpha^4,\alpha^2)\)。这里对基施行自同构，也同时将其迹对偶基送到新基的对偶基。
\end{solution}

\begin{exercise}\label{b1:ex:18-cubic-discriminant}
对 \(a=\sqrt[3]2\)，直接用迹矩阵计算基 \(1,a,a^2\) 的判别式。
\end{exercise}
\begin{solution}
由 \(a^3=2\) 以及 \(\operatorname{Tr}(a)=\operatorname{Tr}(a^2)=0\)，有 \(\operatorname{Tr}(1)=3\)、\(\operatorname{Tr}(a^3)=6\)、\(\operatorname{Tr}(a^4)=2\operatorname{Tr}(a)=0\)。所以迹矩阵为
\[
\begin{pmatrix}3&0&0\\0&0&6\\0&6&0\end{pmatrix},
\]
行列式为 \(-108\)。非零性符合特征零下的可分性；本章的判别只问行列式是否为零，负号不表示配对退化，也不能仅凭这个行列式的负号就断言配对负定。
\end{solution}

\begin{exercise}\label{b1:ex:18-integral-unit}
设 \(K\) 为数域，\(a\in K\) 为非零代数整数。证明 \(a^{-1}\) 也是代数整数，当且仅当 \(\operatorname{N}_{K/\mathbb Q}(a)=\pm1\)。
\end{exercise}
\begin{solution}
若二者均为代数整数，由\cref{b1:prop:integral-trace-norm}，它们的范数是互为倒数的整数，只能各为正一或负一。反向先说明 \(a\) 的乘法特征多项式系数是整数：数域扩张可分，\cref{b1:thm:trace-norm-embeddings}说明其根是全部嵌入像，按嵌入逐个计数。不同嵌入可能在 \(a\) 上取同一值，但这时该值在乘法特征多项式中也重复出现。每个像都在 \(\mathbb Z\) 上整，系数是这些像按重数计的初等对称和，依\cref{b1:lem:integral-closure-elementary}仍整，而本来属于 \(\mathbb Q\)，所以都是整数。若范数为 \(\pm1\)，这个首一多项式的常数项为 \(\pm1\)。把其在 \(a\) 处的等式除以 \(a^{[K:\mathbb Q]}\)，再乘以该常数项的符号，便得到 \(a^{-1}\) 满足的首一整数方程，故其整。
\end{solution}

\begin{exercise}\label{b1:ex:18-quadratic-integrality}
设 \(K/\mathbb Q\) 为二次扩张，\(a\in K\)。证明 \(a\) 为代数整数当且仅当它的迹、范数均为整数。以 \(a=(1+\sqrt{13})/2\) 为例求出二者和基 \(1,a\) 的判别式。
\end{exercise}
\begin{solution}
必要性由\cref{b1:prop:integral-trace-norm}。反向，由\eqref{b1:eq:trace-norm-minimal-polynomial}，\(a\) 满足乘法特征多项式 \(t^2-\operatorname{Tr}(a)t+\operatorname{N}(a)\)。若 \(a\notin\mathbb Q\)，它生成这个二次扩张，特征多项式就是二次最小多项式；若 \(a\in\mathbb Q\)，最小多项式为 \(t-a\)，特征多项式则是 \((t-a)^2\)。两种情形都给出一条首一方程，若迹与范数是整数，其系数便全在 \(\mathbb Z\) 中，因此 \(a\) 整。所给元素满足 \(t^2-t-3=0\)，迹为一，范数为负三。两个共轭之差为 \(\sqrt{13}\)，故基判别式为十三。
\end{solution}

\begin{exercise}\label{b1:ex:18-integrality-counterexample}
设 \(a\) 为 \(t^3-\tfrac12t-1\) 的一个根，\(L=\mathbb Q(a)\)。分别计算 \(2a\) 和 \(a^4\) 对 \(\mathbb Q\) 的迹、范数并判断它们是否为代数整数。为 \(2a\) 写出首一整数方程，并说明为什么整数倍的整性不能反推原元素的整性。
\end{exercise}
\begin{solution}
\subsecref{b1:subsec:1804:04}的三次例子已证明 \([L:\mathbb Q]=3\)，\(\operatorname{Tr}(a)=0\)、\(\operatorname{N}(a)=1\)，而 \(a\) 不是代数整数。令 \(b=2a\)，由 \(a^3=\tfrac12a+1\) 得
\[
b^3-2b-8=0.
\]
所以 \(b\) 是代数整数。迹的线性与范数的乘法性给出 \(\operatorname{Tr}(2a)=0\)、\(\operatorname{N}(2a)=2^3\operatorname{N}(a)=8\)。这已经表明整数倍整，并不保证原元素整。

由根与系数的关系，\(\operatorname{Tr}(a^2)=\operatorname{Tr}(a)^2-2(-1/2)=1\)。也可在幂基 \(1,a,a^2\) 下直接计算乘以 \(a^2\) 的矩阵迹，得到同样的值。再由 \(a^4=\tfrac12a^2+a\)，有
\[
\operatorname{Tr}(a^4)=\tfrac12,\qquad \operatorname{N}(a^4)=\operatorname{N}(a)^4=1.
\]
代数整数对 \(\mathbb Q\) 的迹必须是整数，所以 \(a^4\) 不是代数整数。若 \(a\) 整，其四次幂也整；因此这一计算还提供了排除 \(a\) 整性的另一种办法。迹和范数只给出乘法特征多项式的两个系数，检查高次幂则可能显露原先未见的分母。
\end{solution}

\begin{exercise}\label{b1:ex:18-integrality-tower}
设 \(b=\sqrt2\)，\(a^2=1+b\)。不用数域整数环的结构定理，证明 \(a\) 和 \(a^{-1}\) 都是代数整数，并写出各自的一条首一整数方程。
\end{exercise}
\begin{solution}
\(b\) 整，而 \(a\) 满足 \(t^2-(1+b)=0\)，由\cref{b1:prop:integrality-transitive-elementary}知 \(a\) 整。具体消去 \(b\) 得 \((a^2-1)^2=2\)，即 \(a^4-2a^2-1=0\)。常数项为负一，\(a\ne0\)。除以 \(a^4\) 后可写为 \((a^{-1})^4+2(a^{-1})^2-1=0\)，所以逆元也整。这里没有用到任何整基或唯一分解结论。
\end{solution}

% !TeX root = ../../Book1.tex
% 纲要标题：## 第19章　有限 Galois 理论
% 纲要位置：计划纲要.md，第 644 行。

\chapter{有限 Galois 理论}
\label{b1:ch:19}
\BookChapterNumber{b1:ch:19}

\begin{chapterabstract}
有限 Galois 理论用域自同构群研究扩张中的中间域。本章以正规性和可分性确定适用范围，证明 Artin 不动域定理与有限 Galois 对应，并在具体多项式的分裂域中计算这一对应。
\end{chapterabstract}

\begin{learninggoals}
\begin{itemize}
\item 证明 Artin 不动域定理，说明嵌入独立性和线性关系下降分别承担哪个次数估计。
\item 使用 Galois 对应计算中间域次数，辨认正规子扩张并构造限制同态。
\item 从分裂域、根的置换和判别式识别低次方程的 Galois 群。
\item 正确计算合成域与交域，区别实际子域、共轭子域和抽象同构类型，并用正规化子计算中间域的自同构。
\end{itemize}
\end{learninggoals}

在 \(L=\mathbb Q(\sqrt2,\sqrt3)\) 中，可以分别改变两个平方根的符号而保持有理数不动。这给出四个域自同构。其中固定 \(\sqrt2\) 的两个自同构组成一个子群，其共同不动域为 \(\mathbb Q(\sqrt2)\)；固定 \(\sqrt6\) 的两个自同构则给出共同不动域 \(\mathbb Q(\sqrt6)\)。于是同一扩张中的子域，可以通过“哪些对称使其中每个元素不变”来辨认。子域包含的元素越多，要逐点固定它就须满足越多约束，相应的自同构群便越小。一般有限扩张是否也有足够多的自同构，使所有中间域都能如此找回？正规性、可分性与不动域定理将给出这一对应成立的条件。

\ProseParagraphBreak

不动域定理与 Galois 对应使用\cref{b1:thm:embedding-extension}的嵌入延拓、\cref{b1:thm:normal-embedding-criterion}的正规性判据、\cref{b1:thm:separable-embeddings}的嵌入计数和\cref{b1:cor:embedding-independence}的独立性。具体多项式的群计算还使用\secref{b1:sec:1803}的判别式，以及有限域的 Frobenius、中国剩余定理和\cref{b1:lem:integral-closure-elementary}来证明模素数循环型判据。群论方面需要子群指数、正规子群、商群与置换作用。

\ProseParagraphBreak
Galois 对应及多项式的 Galois 群计算，可参阅 Roman 的《Field Theory》\cite[Sections 6.2--6.5, 7.1, 7.4--7.5]{Roman2006FieldTheory}。Lang 的《Algebra》\cite[Chapter VI, \S\S 1--2]{Lang2002}可供比较域与自同构群的结构关系；阅读具体算例时，宜把群的计算和中间域的实际生成元一并写出。

% !TeX root = ../../Book1.tex
% 纲要标题：### 19.1 Galois 扩张
% 纲要位置：计划纲要.md，第 646 行。

\section{Galois 扩张}
\label{b1:sec:1901}

% 纲要内容（仅作写作计划，尚非教材正文）：
%
% * 正规且可分的有限扩张
% * 自同构群与不动域
% * Artin 不动域定理；展开最少非零坐标、归一化与系数下降，并在映射空间中解释独立性给出的维数估计
% * 调用第18.1节已证的域嵌入线性独立性；先完成有限自同构群 \(G\) 作用下 \([L:L^G]=|G|\) 的证明，再进入 Galois 对应
% * 乘法特征与域嵌入独立性的主要证明统一放在第18.1节；本节专门承担 Artin 不动域定理，说明线性独立性在扩张次数估计中的作用
%

多项式根的置换并不都能保持根之间的代数关系。域自同构挑出的正是那些合法的置换。有限 Galois 理论的第一步，是说明什么时候自同构足够多，以及一组有限自同构能够固定多大的子域。

\subsection{有限正规可分扩张}
\label{b1:subsec:1901:01}

\begin{definition}\label{b1:def:finite-galois-extension}
设 \(L/K\) 为有限域扩张。如果 \(L/K\) 同时可分且正规，即每个元素在 \(K\) 上的最小多项式没有重根，并且每个在 \(L\) 中有根的 \(K\) 上不可约多项式都在 \(L\) 中分裂，则称 \(L/K\) 为\term[Galois-kuozhang]{Galois 扩张}[Galois extension]。
\end{definition}
这两个条件分别处理嵌入的数量与像的位置。可分性使最小多项式的根互不相同，保证到代数闭包的不同 \(K\) 嵌入数达到扩张次数；正规性则保证这些嵌入把整个 \(L\) 映到自身。因此两者结合，才把足够多的共轭根选择变成 \(L\) 的自同构。

\ProseParagraphBreak
可分多项式 \(f\in K[x]\) 的分裂域在 \(K\) 上有限，由\cref{b1:prop:splitting-field-exists}得到；它由 \(f\) 的根生成，这些根的最小多项式都是 \(f\) 的因子，因而均可分。\cref{b1:thm:separable-embeddings}保证整个扩张可分，而\cref{b1:thm:normal-embedding-criterion}保证分裂域正规，故它为有限 Galois 扩张。反过来，设 \(L/K\) 有限、正规且可分，取有限生成元 \(L=K(a_1,\ldots,a_r)\)。每个 \(m_{a_i,K}\) 都可分；将这些首一不可约多项式中相同者只保留一次，所得乘积 \(f\) 仍可分，因为两个不同的首一不可约多项式不能有共同根。正规性又保证它们的全部根都属于 \(L\)。这些根生成的域因此包含于 \(L\)，同时包含各 \(a_i\)，也就包含 \(K(a_1,\ldots,a_r)=L\)。所以 \(L\) 恰为 \(f\) 的分裂域。这证明了有限 Galois 扩张也可以等价地描述为可分多项式的分裂域。

\ProseParagraphBreak
\(\mathbb Q(\sqrt2)/\mathbb Q\) 为 Galois 扩张，而 \(\mathbb Q(\sqrt[3]2)/\mathbb Q\) 可分但不正规，因为后者为实域，未包含 \(x^3-2\) 的两个非实根。正特征下，正规性也不能替代可分性：若 \(u\) 在 \(\mathbb F_p\) 上超越，取 \(L=\mathbb F_p(u)\)、\(K=\mathbb F_p(u^p)\)，则 \([L:K]=p\)，且 \(u\) 的最小多项式在 \(L[X]\) 中分解为
\[
X^p-u^p=(X-u)^p.
\]
它的系数属于 \(K\)，全部根已在 \(L\)，并由这些根生成 \(L\)，所以 \(L\) 是它的分裂域，扩张正规。但这 \(p\) 个根全部重合，\(u\) 不可分，故这个扩张不是 Galois 扩张。

\subsection{自同构群与不动域}
\label{b1:subsec:1901:02}

设 \(L/K\) 为域扩张。保持 \(K\) 中每个元素不变的全部 \(L\) 自同构，在复合下组成群 \(\operatorname{Aut}_K(L)\)。这个群对任意扩张都有定义。如果 \(L/K\) 是 Galois 扩张，则把它记为 \(\operatorname{Gal}(L/K)\)，称为扩张 \(L/K\) 的\term[Galois-qun]{Galois 群}[Galois group]。本书在未满足 Galois 条件时仍使用 \(\operatorname{Aut}_K(L)\) 这一记号；这并不表示该扩张没有自同构群。
\symidx{\(\operatorname{Aut}_K(L)\)：保持基域 \(K\) 的 \(L\) 自同构群}
\symidx{\(\operatorname{Gal}(L/K)\)：Galois 扩张的自同构群}

\ProseParagraphBreak
设 \(H\leq\operatorname{Aut}(L)\) 为域 \(L\) 的自同构群的子群。把子域
\[
L^H=\bigl\{\,a\in L\mid \sigma(a)=a\text{ 对每个 }\sigma\in H\,\bigr\}
\]
称为 \(H\) 在 \(L\) 中的\term[budong-yu]{不动域}[fixed field]。
\symidx{\(L^H\)：子群 \(H\) 的不动域}
自同构保持域的运算，因此共同不动元对加、减、乘以及非零元取逆都封闭，确实组成子域。若 \(H\subseteq\operatorname{Aut}_K(L)\)，则 \(K\subseteq L^H\)。

\begin{proposition}\label{b1:prop:galois-automorphism-count}
若 \(L/K\) 为有限 Galois 扩张，则 \(|\operatorname{Gal}(L/K)|=[L:K]\)。
\end{proposition}
\begin{proof}
由\cref{b1:thm:separable-embeddings}，可分性保证 \(L\) 到代数闭包的 \(K\) 嵌入恰有 \([L:K]\) 个。由\cref{b1:thm:normal-embedding-criterion}，正规性保证每个这样的嵌入把 \(L\) 映到自身；有限维单射的像与 \(L\) 维数相同，所以它是自同构。这些嵌入因而恰为 Galois 群的全部元素。
\end{proof}

\subsection{Artin 不动域定理}
\label{b1:subsec:1901:03}

前面从有限扩张 \(L/K\) 出发研究它的自同构；以下\term[Artin-budongyu-dingli]{Artin 不动域定理}[Artin fixed field theorem]反过来，先给域 \(L\) 的有限自同构群 \(G\)，再取 \(K=L^G\)。这里没有预先假设 \(L/K\) 有限：群 \(G\) 的有限性将迫使这个不动域扩张成为有限 Galois 扩张，其次数恰为 \(|G|\)。
\begin{theorem}[Artin]\label{b1:thm:artin-fixed-field}
设 \(G\) 是域 \(L\) 的一个有限自同构群，\(K=L^G\)。则 \(L/K\) 是有限 Galois 扩张，并且
\[
[L:K]=|G|,\qquad \operatorname{Gal}(L/K)=G.
\]
\end{theorem}
\begin{proof}
记 \(G=\{\sigma_1,\ldots,\sigma_n\}\)。次数等式包含两个方向：我们要用线性方程控制 \([L:K]\leq n\)，再用嵌入独立性得到 \(n\leq[L:K]\)。上界的含义是，任意 \(n+1\) 个元素 \(x_1,\ldots,x_{n+1}\in L\) 都在 \(K\) 上线性相关。考察以 \(L\) 为系数域、以 \(c_1,\ldots,c_{n+1}\in L\) 为未知量的齐次方程组：
\[
\sum_{j=1}^{n+1}\sigma_i(x_j)c_j=0\qquad(1\leq i\leq n).
\]
因为只有 \(n\) 个方程，却有 \(n+1\) 个未知量，存在非零的 \(L\) 解。但这些系数暂时未必在 \(K\) 中，还不能用它得到所需的 \(K\) 线性关系。为了使系数下降到不动域，从全部非零解中取非零坐标数最少者，重新编号，使它的非零坐标恰为前 \(r\) 个。再把整个解除以 \(c_r\)，仍得一个解，并可设 \(c_r=1\)。这个归一化使任何 \(\tau\in G\) 都固定该坐标。

将 \(\tau\) 作用到全部方程，得到
\[
\sum_{j=1}^{n+1}(\tau\sigma_i)(x_j)\tau(c_j)=0
\qquad(1\leq i\leq n).
\]
左乘映射 \(\sigma\mapsto\tau\sigma\) 是 \(G\) 的一个排列，逆映射为左乘 \(\tau^{-1}\)；所以上式仅把原来各方程重新排列，向量 \(\bigl(\tau(c_j)\bigr)_j\) 也是原方程组的解。两解之差在第 \(r\) 个坐标为 \(\tau(1)-1=0\)，在 \(r\) 后也均为零。若差非零，它至多有 \(r-1\) 个非零坐标，违背原解的最小选择。因此
\[
\tau(c_j)=c_j\quad(\tau\in G),\qquad c_j\in L^G=K.
\]
取恒等自同构对应的方程，就有 \(\sum_j x_jc_j=0\)，而 \(c_r=1\) 保证这个 \(K\) 系数关系非平凡。任意 \(n+1\) 个元素都如此，故 \([L:K]\leq n\)，特别地，扩张有限。

记 \(d=[L:K]\)，并把所有 \(K\) 线性映射 \(L\rightarrow L\) 组成的空间记为 \(\mathcal V\)。其加法与 \(L\) 标量乘法都是逐点定义的，特别是
\[
(\lambda f)(x)=\lambda f(x)\qquad(\lambda\in L,\ f\in\mathcal V).
\]
乘上 \(\lambda\) 后仍是 \(K\) 线性映射，因为 \(L\) 的乘法交换。若 \(b_1,\ldots,b_d\) 是 \(L/K\) 的一组基，一个 \(K\) 线性映射 \(f\) 完全由 \(f(b_1),\ldots,f(b_d)\in L\) 决定，而且任意指定这 \(d\) 个像都给出唯一的 \(K\) 线性映射。因此求值映射
\[
\mathcal V\longrightarrow L^d,\qquad
f\longmapsto\bigl(f(b_1),\ldots,f(b_d)\bigr)
\]
是 \(L\) 线性同构，故 \(\dim_L\mathcal V=d\)。每个 \(\sigma_i\) 固定 \(K\)，所以属于 \(\mathcal V\)。\cref{b1:cor:embedding-independence}在此说明：若 \(\sum_i\lambda_i\cdot\sigma_i(x)=0\) 对所有 \(x\in L\) 成立，其中 \(\lambda_i\in L\)，则全部 \(\lambda_i=0\)。因此这 \(n\) 个不同自同构是 \(d\) 维 \(L\) 向量空间中的线性无关向量，得到 \(n\leq d\)。结合上界，\(d=n\)。

对任意 \(a\in L\)，把 \(G\) 轨道中的互异元素记作 \(a_1,\ldots,a_s\)，并构造轨道多项式
\[
P_a(t)=\prod_{j=1}^s(t-a_j).
\]
任何 \(\tau\in G\) 都把轨道映到自身且映满，所以只会排列这些根。\(P_a\) 的系数由它们的初等对称式给出，因而被每个 \(\tau\) 固定，属于 \(L^G=K\)。又因为只将互异轨道元素各取一次，\(P_a\) 没有重根，并已在 \(L[t]\) 中分裂。恒等自同构保证 \(a\) 在轨道中，因此 \(P_a(a)=0\)，\(a\) 的最小多项式整除 \(P_a\)，也就可分并在 \(L\) 中分裂。任意 \(a\) 均如此，说明 \(L/K\) 正规且可分。最后由\cref{b1:prop:galois-automorphism-count}，其全部 \(K\) 自同构只有 \(n\) 个；已经包含的 \(G\) 也有 \(n\) 个元素，所以它就是全群。
\end{proof}

若给定的是一个抽象有限群 \(\Gamma\) 对 \(L\) 的作用，作用同态为 \(\rho:\Gamma\rightarrow\operatorname{Aut}(L)\)，则共同不动元只取决于像群 \(\rho(\Gamma)\)。上面的定理应用于这个实际自同构群，给出
\[
[L:L^{\rho(\Gamma)}]=|\rho(\Gamma)|
=\frac{|\Gamma|}{|\ker\rho|}.
\]
最后一个等号是群同态的第一同构定理。只有作用忠实、\(\ker\rho=1\) 时，才能直接以抽象群 \(\Gamma\) 的阶作为扩张次数；这正是前面群作用通过核的商群给出忠实作用的原则。

\subsection{不动域扩张次数的计算}
\label{b1:subsec:1901:04}

对有限 Galois 扩张 \(L/K\)，令 \(G=\operatorname{Gal}(L/K)\)。\cref{b1:thm:artin-fixed-field}与\cref{b1:prop:galois-automorphism-count}分别给出
\[
[L:L^G]=|G|=[L:K].
\]
因 \(K\subseteq L^G\)，塔式公式迫使 \([L^G:K]=1\)，所以 \(L^G=K\)。这说明基域恰能从全部自同构中恢复。

\ProseParagraphBreak
取特征不为二的域 \(k\) 及在 \(k\) 上超越的元素 \(t\)。虽然 \(k(t)/k\) 不是代数扩张，仍可对 \(k(t)\) 的有限自同构群使用不动域定理。令 \(s\) 为逐点固定 \(k\) 且满足 \(s(t)=-t\) 的 \(k(t)\) 自同构，它生成二阶群 \(G=\langle s\rangle\)。群 \(G\) 的不动域包含 \(k(t^2)\)。元素 \(t\) 满足 \(x^2-t^2=0\)，但不属于 \(k(t^2)\)：若 \(t=f(t^2)/g(t^2)\)，其中 \(f,g\in k[x]\)、\(g\ne0\)，则 \(t\cdot g(t^2)=f(t^2)\) 的两侧分别只含奇次幂和偶次幂，且均不为零，矛盾。因此 \([k(t):k(t^2)]=2\)；与不动域定理给出的 \([k(t):k(t)^G]=2\) 比较，得到 \(k(t)^G=k(t^2)\)。

\subsection{自同构数判据与次数估计}
\label{b1:subsec:1901:05}

\begin{corollary}\label{b1:cor:galois-automorphism-criterion}
设 \(L/K\) 为有限域扩张。则 \(L/K\) 为 Galois 扩张，当且仅当 \(|\operatorname{Aut}_K(L)|=[L:K]\)。
\end{corollary}
\begin{proof}
若 \(L/K\) 为 Galois 扩张，结论由\cref{b1:prop:galois-automorphism-count}给出。反过来，令 \(H=\operatorname{Aut}_K(L)\)。假设 \(|H|=[L:K]\)，则\cref{b1:thm:artin-fixed-field}给出 \([L:L^H]=|H|=[L:K]\)。由于 \(K\subseteq L^H\)，塔式公式说明 \(L^H=K\)。该定理还保证 \(L/L^H\) 正规且可分，故 \(L/K\) 为 Galois 扩张。
\end{proof}

\cref{b1:thm:artin-fixed-field}中的两个次数估计有不同来源。线性方程加最短支集论证，把 \(L\) 系数的关系下降为不动域系数的关系，给出上界；嵌入独立性则保证自同构不可能多于维数，给出下界。轨道多项式至多有 \(|G|\) 个根，能说明每个单独的 \(K(a)/K\) 次数至多为 \(|G|\)。但在尚未证明 \(L\) 由一个元素生成时，这只控制各个单扩张，不能直接控制整个 \(L/K\) 的维数。线性方程法处理的是任意 \(|G|+1\) 个元素之间的关系，因而直接控制整个空间，不需要先知道扩张单生成。

% !TeX root = ../../Book1.tex
% 纲要标题：### 19.2 Galois 对应
% 纲要位置：计划纲要.md，第 654 行。

\section{Galois 对应}
\label{b1:sec:1902}

% 纲要内容（仅作写作计划，尚非教材正文）：
%
% * 中间域与子群的反向对应
% * 次数、指数与群阶
% * 正规子群、正规子扩张与商群
% * 紧接对应定理，以 \(X^3-2\) 的分裂域枚举全部子群、不动域、次数、正规性和包含关系
%

现在固定有限 Galois 扩张 \(L/K\)，记 \(G=\operatorname{Gal}(L/K)\)。给定中间域 \(E\)，考察逐点固定 \(E\) 的自同构所组成的子群 \(\operatorname{Gal}(L/E)\)；给定子群，则取其共同不动元。两个构造反转包含关系；要证明它们互逆，还须使用不动域定理中的精确次数。

\subsection{中间域与子群的反向对应}
\label{b1:subsec:1902:01}

下述对应称为\term[Galois-duiying]{Galois 对应}[Galois correspondence]。
\begin{theorem}\label{b1:thm:finite-galois-correspondence}
有限 Galois 扩张 \(L/K\) 的中间域与 \(G=\operatorname{Gal}(L/K)\) 的子群，通过
\[
E\longmapsto\operatorname{Gal}(L/E),\qquad H\longmapsto L^H
\]
建立互逆的、反转包含关系的一一对应。对任意中间域 \(E\)，扩张 \(L/E\) 都是 Galois 扩张。
\end{theorem}
\begin{proof}
先说明最后一句。每个元素在 \(E\) 上的最小多项式整除它在 \(K\) 上的最小多项式，故可分；任意 \(E\) 嵌入也是 \(K\) 嵌入，所以由\cref{b1:thm:normal-embedding-criterion}，它保持 \(L\)，故 \(L/E\) 正规。于是 \(H=\operatorname{Gal}(L/E)\) 满足 \(|H|=[L:E]\)。不动域定理给出 \([L:L^H]=|H|\)，结合 \(E\subseteq L^H\) 以及塔式公式，得到 \(E=L^H\)。

反向给定任意 \(H\leq G\)，\cref{b1:thm:artin-fixed-field}直接给出 \(\operatorname{Gal}(L/L^H)=H\)。这证明两个构造互逆。若 \(E_1\subseteq E_2\)，固定较大域的条件更强，所以 \(\operatorname{Gal}(L/E_2)\subseteq\operatorname{Gal}(L/E_1)\)；子群越大，不动域越小。故对应反转包含。
\end{proof}
不能把最后一句反过来读成每个 \(E/K\) 都正规。例如三次多项式的非正规根域可以作为其 Galois 分裂域的中间域。

\subsection{次数、指数与群阶}
\label{b1:subsec:1902:02}

对对应的 \(E=L^H\)，不动域定理和塔式公式分别给出
\begin{equation}\label{b1:eq:galois-degree-index}
[L:E]=|H|,\qquad [E:K]=[G:H].
\end{equation}
若 \(E_1\subseteq E_2\) 对应 \(H_1\supseteq H_2\)，则 \([E_2:E_1]=[H_1:H_2]\)。这一等式计算的是群的指数，不要求相应子群正规。

\ProseParagraphBreak
例如 \(L=\mathbb F_{q^6}\)、\(K=\mathbb F_q\) 时，\cref{b1:thm:finite-field-automorphisms}给出 \(G=\langle\sigma\rangle\cong C_6\)。指数为 \(d\) 的子群 \(\langle\sigma^d\rangle\) 对应 \(\mathbb F_{q^d}\)，其中 \(d\mid6\)。这与\cref{b1:thm:finite-field-subfields}的直接计算一致；有限域的子域分类并不依赖本定理，因而这里是对两种方法的比较。

\subsection{正规子群、正规子扩张与商群}
\label{b1:subsec:1902:03}

\begin{theorem}\label{b1:thm:galois-quotient}
设 \(L/K\) 为有限 Galois 扩张，\(G=\operatorname{Gal}(L/K)\)，\(H\leq G\)，并令 \(E=L^H\)。则 \(E/K\) 为 Galois 扩张当且仅当 \(H\) 是 \(G\) 的正规子群。此时限制映射 \(G\rightarrow\operatorname{Gal}(E/K)\) 是满同态，且
\[
\operatorname{Gal}(E/K)\cong G/H.
\]
\end{theorem}
\begin{proof}
对任意 \(\sigma\in G\)，直接检验得到
\[
\sigma(E)=L^{\sigma H\sigma^{-1}}.
\]
事实上，若 \(x\in E\)，则 \(\sigma h\sigma^{-1}\) 固定 \(\sigma(x)\)；反向对 \(\sigma^{-1}\) 作同样检验。若 \(H\) 正规，则每个 \(\sigma\) 都保持 \(E\)。任意 \(E\) 的 \(K\) 嵌入能由\cref{b1:thm:embedding-extension}延拓到 \(L\) 的 \(K\) 嵌入，而 \(L/K\) 正规使此延拓属于 \(G\)。它保持 \(E\)，所以嵌入判据说明 \(E/K\) 正规；\(E/K\) 的可分性来自 \(L/K\)，于是为 Galois 扩张。

若 \(E/K\) 正规，每个 \(\sigma\in G\) 都保持 \(E\)，从而 \(L^{\sigma H\sigma^{-1}}=E=L^H\)。Galois 对应的单射性给出 \(\sigma H\sigma^{-1}=H\)，所以 \(H\) 正规。此时限制映射 \(G\rightarrow\operatorname{Gal}(E/K)\) 的核是恰好固定 \(E\) 的 \(H\)；任意目标自同构均能延拓到 \(L\)，故映射满射。\cref{b1:thm:first-isomorphism}遂把商群 \(G/H\) 同构到这个实际像。
\end{proof}
正规性确保限制后仍是 \(E\) 到自身的映射；若 \(E/K\) 不正规，某些 \(\sigma\) 会把它送到另一个共轭子域，不能对整个 \(G\) 定义到 \(\operatorname{Aut}_K(E)\) 的限制同态。

\ProseParagraphBreak
对上述 \(E=L^H\)，逐点固定 \(E\) 与将 \(E\) 保持为集合是不同的条件。记 \(N_G(H)=\{\,\sigma\in G\mid\sigma H\sigma^{-1}=H\,\}\) 为 \(H\) 在 \(G\) 中的正规化子。Galois 对应及证明中的共轭公式分别给出
\[
\{\,\sigma\in G\mid\sigma|_E=\operatorname{id}_E\,\}
=\operatorname{Gal}(L/E)=H,
\qquad
\{\,\sigma\in G\mid\sigma(E)=E\,\}=N_G(H).
\]
恰是 \(N_G(H)\) 中的自同构才能限制成 \(E\) 的 \(K\) 自同构；\cref{b1:ex:19-normalizer-automorphisms}将用它计算 \(\operatorname{Aut}_K(E)\)。

\subsection{中间域的枚举}
\label{b1:subsec:1904:01}

先以 \(X^3-2\) 的分裂域把对应定理用在一个非交换 Galois 群上：求出全部子群及其不动域，并核对次数、正规性和包含关系。

\begin{proposition}\label{b1:prop:cubic-splitting-field-correspondence}
令 \(a=\sqrt[3]2\in\mathbb R\)，\(\zeta\) 为本原三次单位根，并置 \(L=\mathbb Q(a,\zeta)\)。则 \([L:\mathbb Q]=6\)，且 \(G=\operatorname{Gal}(L/\mathbb Q)\cong S_3\)。除 \(\mathbb Q\) 与 \(L\) 外，全部中间域为唯一的二次域 \(\mathbb Q(\zeta)\) 和三个互异的三次域 \(\mathbb Q(\zeta^j a)\)、\(j=0,1,2\)。前者在 \(\mathbb Q\) 上正规，后三者均不正规。
\end{proposition}
\begin{proof}
由\cref{b1:thm:eisenstein}，\(X^3-2\) 在 \(\mathbb Q\) 上不可约，故 \([\mathbb Q(a):\mathbb Q]=3\)。域 \(\mathbb Q(a)\) 包含于 \(\mathbb R\)，而 \(\zeta\) 非实且满足 \(X^2+X+1=0\)，所以 \([L:\mathbb Q(a)]=2\)。这是一个特征零多项式的分裂域，因而为六次 Galois 扩张。群 \(G\) 忠实置换三个根 \(a,\zeta a,\zeta^2a\)，由阶数知它给出全部 \(S_3\)。

记循环置换三个根的自同构为 \(s\)，复共轭为 \(t\)，则
\[
s(a)=\zeta a,\quad s(\zeta)=\zeta,\qquad
t(a)=a,\quad t(\zeta)=\zeta^{-1}.
\]
它们满足 \(s^3=t^2=1\)、\(tst=s^{-1}\)，且生成 \(G\)。\(G\) 的全部子群为 \(G\)、\(\langle s\rangle\)、\(\langle s^{2j}t\rangle\)（\(j=0,1,2\)）及 \(1\)：真非平凡子群只能有二阶或三阶，分别由三个换位或一个三循环生成。

自同构 \(s\) 固定 \(\zeta\)，而 \(s^{2j}t\) 固定 \(\zeta^j a\)。由\eqref{b1:eq:galois-degree-index}，这两个子群的不动域在 \(\mathbb Q\) 上分别有次数二与三；\(\mathbb Q(\zeta)\) 与 \(\mathbb Q(\zeta^j a)\) 也分别有这些次数，故相应包含必为相等。对应的三个二阶子群互异，所以三个三次域也互异。子群 \(\langle s\rangle\) 正规，三个二阶子群彼此共轭而不正规，\cref{b1:thm:galois-quotient}遂给出各中间域的正规性。
\end{proof}

例如取 \(E=\mathbb Q(\zeta)\)。全部六个 \(G\) 中的自同构都将 \(E\) 保持为集合，但只有 \(\langle s\rangle\) 的三个元素逐点固定 \(E\)；复共轭 \(t\) 把 \(\zeta\) 送到 \(\zeta^{-1}\)，所以不逐点固定它。

\cref{b1:tab:cubic-galois-correspondence}集中列出上述对应。对每一行都有 \([E:\mathbb Q]=[G:H]\)、\([L:E]=|H|\)；最后一列同时说明 \(H\) 是否为 \(G\) 的正规子群，因为这里的全部扩张均可分。
\begin{table}[htbp]
\centering
\caption[\(X^3-2\) 分裂域的 Galois 对应]{\(X^3-2\) 分裂域的 Galois 对应。其中 \(j=0,1,2\)。}
\label{b1:tab:cubic-galois-correspondence}
\begin{tabular}{@{}llccl@{}}
\toprule
子群 \(H\) & 不动域 \(E=L^H\) & \([E:\mathbb Q]\) & \([L:E]\) & \(E/\mathbb Q\) 正规 \\
\midrule
\(G\cong S_3\) & \(\mathbb Q\) & \(1\) & \(6\) & 是 \\
\(\langle s\rangle\cong C_3\) & \(\mathbb Q(\zeta)\) & \(2\) & \(3\) & 是 \\
\(\langle s^{2j}t\rangle\cong C_2\) & \(\mathbb Q(\zeta^j a)\) & \(3\) & \(2\) & 否 \\
\(1\) & \(L\) & \(6\) & \(1\) & 是 \\
\bottomrule
\end{tabular}
\end{table}

\cref{b1:fig:cubic-intermediate-fields}将表中的三个三次域分别画出；图中较低的域包含于较高的域，边上的数字是相邻两域的扩张次数。
\begin{figure}[htbp]
\centering
\begin{tikzpicture}[x=1cm,y=1cm,every node/.style={font=\small}]
\node (L) at (0,1.8) {\(L\)};
\node (Z) at (-5.5,0) {\(\mathbb Q(\zeta)\)};
\node (A) at (-1.8,0) {\(\mathbb Q(a)\)};
\node (B) at (1.8,0) {\(\mathbb Q(\zeta a)\)};
\node (C) at (5.5,0) {\(\mathbb Q(\zeta^2a)\)};
\node (Q) at (0,-1.8) {\(\mathbb Q\)};
\draw (L) -- node[pos=.55,fill=white,inner sep=1pt] {\scriptsize 3} (Z);
\draw (L) -- node[pos=.55,fill=white,inner sep=1pt] {\scriptsize 2} (A);
\draw (L) -- node[pos=.55,fill=white,inner sep=1pt] {\scriptsize 2} (B);
\draw (L) -- node[pos=.55,fill=white,inner sep=1pt] {\scriptsize 2} (C);
\draw (Z) -- node[pos=.45,fill=white,inner sep=1pt] {\scriptsize 2} (Q);
\draw (A) -- node[pos=.45,fill=white,inner sep=1pt] {\scriptsize 3} (Q);
\draw (B) -- node[pos=.45,fill=white,inner sep=1pt] {\scriptsize 3} (Q);
\draw (C) -- node[pos=.45,fill=white,inner sep=1pt] {\scriptsize 3} (Q);
\end{tikzpicture}
\caption[\(X^3-2\) 分裂域的中间域格]{\(X^3-2\) 分裂域的中间域格。边上的数字为较高的域在较低的域上的次数。}
\label{b1:fig:cubic-intermediate-fields}
\end{figure}

图中的四个非平凡真中间域两两互不包含：二次与三次之间由次数的整除性排除包含，三个不同的三次域之间也不能包含。把每条包含关系反向，就得到 \(G\) 的子群格。例如 \(\mathbb Q\subset\mathbb Q(\zeta)\subset L\) 对应 \(G\supset\langle s\rangle\supset1\)。这个枚举穷尽所有中间域，因为已经列出了全部子群；一般地，有限 Galois 扩张的中间域数必有限。

% !TeX root = ../../Book1.tex
% 纲要标题：### 19.3 Galois 群的计算与识别
% 纲要位置：计划纲要.md，第 660 行。

\section{Galois 群的计算与识别}
\label{b1:sec:1903}

% 纲要内容（仅作写作计划，尚非教材正文）：
%
% * 二次、多二次扩张；三次式 \(X^3-2\) 的完整对应计算见第19.2节
% * 四次方程的典型 Galois 群
% * 从置换作用和分解型识别子群
% * 模素数分解与循环型：在此完整证明约化判据，第21.3节直接引用
%

识别 Galois 群通常需要两类信息：扩张次数确定群阶，群对根的作用将它表示为对称群的子群。以下每个例子都先明确分裂域，再计算自同构；不能只写出一个抽象群名而略去它怎样作用。

\subsection{二次与多二次扩张}
\label{b1:subsec:1903:01}

特征不为二时，若 \(d\in K^\times\) 不是平方，\(K(\sqrt d)/K\) 为二次 Galois 扩张，非平凡自同构把 \(\sqrt d\) 送到 \(-\sqrt d\)。更一般的可分二次多项式 \(x^2+bx+c\) 的两根之和为 \(-b\)，所以添加一根后另一根已在域内；即使特征为二，可分二次扩张也仍为 Galois 扩张。

\ProseParagraphBreak
考虑 \(L=\mathbb Q(\sqrt2,\sqrt3)\)。若 \(\sqrt3=a+b\sqrt2\)，其中 \(a,b\in\mathbb Q\)，平方比较得 \(2ab=0\)、\(a^2+2b^2=3\)。若 \(b=0\)，三为有理数平方；若 \(a=0\)，\(3/2\) 为有理数平方；两者都与素因数指数的奇偶性矛盾。因此 \([L:\mathbb Q]=4\)，基可取 \(1,\sqrt2,\sqrt3,\sqrt6\)。独立改变两个平方根的符号保持全部乘法关系，给出四个自同构，所以
\[
\operatorname{Gal}(L/\mathbb Q)\cong C_2\times C_2.
\]
三个二元子群的不动域分别为 \(\mathbb Q(\sqrt2)\)、\(\mathbb Q(\sqrt3)\)、\(\mathbb Q(\sqrt6)\)。例如同时改变两者符号的自同构固定 \(\sqrt6\)，而不动域次数为二，故恰为最后一个域。

\ProseParagraphBreak
三次式 \(X^3-2\) 的六次分裂域、\(S_3\) 作用和三个非正规三次根域，已经在\cref{b1:prop:cubic-splitting-field-correspondence}中逐一算出。该例说明，求得 Galois 群以后，还须确定各自同构对根的具体作用，才能识别它固定的域。后文将用判别式区分可分不可约三次式的 \(S_3\) 与 \(C_3\) 情形。

\subsection{四次方程的 Galois 群}
\label{b1:subsec:1903:02}

四次不可约多项式的分裂域次数未必为四。三个典型计算分别给出不同的群。

\ProseParagraphBreak
先取 \(f=x^4-10x^2+1\)。它的根为 \(\pm\sqrt2\pm\sqrt3\)。令 \(a=\sqrt2+\sqrt3\)，则 \(a^{-1}=\sqrt3-\sqrt2\)，所以
\[
\sqrt3=\frac{a+a^{-1}}2,\qquad \sqrt2=\frac{a-a^{-1}}2.
\]
由上一小节的次数计算，\([\mathbb Q(a):\mathbb Q]=4\)。因此 \(f\) 是 \(a\) 的最小多项式，且其分裂域就是 \(\mathbb Q(\sqrt2,\sqrt3)\)，Galois 群为四元群 \(C_2\times C_2\)。

\ProseParagraphBreak
再取 \(f=x^4+ x^3+x^2+x+1\)，其根为非一的五次单位根。变元平移后的 \(f(x+1)\) 由素数五满足\cref{b1:thm:eisenstein}，故 \(f\) 不可约。若 \(\zeta\) 是一个根，所有根 \(\zeta,\zeta^2,\zeta^3,\zeta^4\) 都在 \(\mathbb Q(\zeta)\) 中，扩张次数为四。自同构由 \(\zeta\mapsto\zeta^j\) 给出，复合对应指数模五相乘。模五的二依次给出 \(2,4,3,1\)，所以这个 Galois 群为 \(C_4\)。

\ProseParagraphBreak
最后取 \(f=x^4-2\)，令 \(a=\sqrt[4]2>0\)。分裂域为 \(L=\mathbb Q(a,i)\)。Eisenstein 判据给出实子域 \(\mathbb Q(a)\) 的次数为四；加入 \(i\) 后次数变为八。设 \(r\) 固定 \(i\) 而把 \(a\) 送到 \(ia\)，设 \(s\) 为复共轭。前者确实存在，因为 \([L:\mathbb Q(i)]=4\)，\(a\) 在 \(\mathbb Q(i)\) 上的最小多项式仍是 \(x^4-2\)，可以把 \(a\) 映到另一个根 \(ia\)。直接作用于生成元可验
\[
r^4=s^2=1,\qquad srs=r^{-1}.
\]
四个 \(r^j\) 固定 \(i\)，四个 \(sr^j\) 把 \(i\) 变号，彼此不同，已经列出全部八个自同构。因此群为本书记号下的 \(D_8\)。

\subsection{置换作用、分解型与子群识别}
\label{b1:subsec:1903:03}

若首一可分多项式 \(f\in K[x]\) 的分裂域为 \(L\)，其 Galois 群在各根上的作用忠实，因为固定全部根就固定它们生成的域。一个不可约因子的根构成一个轨道：任意两根由简单扩张同构对应，该同构能由\cref{b1:thm:embedding-extension}延拓到 \(L\)，正规性保证延拓为自同构。因此在 \(K\) 上各不可约因子的次数，正是全部根分成的轨道大小。特别地，不可约多项式给出传递作用。

\ProseParagraphBreak
还可利用判别式辨认置换的奇偶性。
\begin{proposition}[判别式判据]\label{b1:prop:discriminant-alternating-group}
设 \(K\) 是特征不为二的域，\(f\in K[X]\) 是次数为 \(n\geq1\) 的首一可分多项式，\(L\) 是 \(f\) 在 \(K\) 上的分裂域。令 \(G=\operatorname{Gal}(L/K)\)，并通过它在 \(f\) 的全部根上的作用将其视为 \(S_n\) 的子群。则
\begin{equation}\label{b1:eq:discriminant-alternating-group}
G\subseteq A_n\quad\Longleftrightarrow\quad
\operatorname{disc}(f)\text{ 是 }K\text{ 中的平方}.
\end{equation}
\end{proposition}
\begin{proof}
把 \(f\) 的互异根记为 \(a_1,\ldots,a_n\)，并置 \(\delta=\prod_{i<j}(a_j-a_i)\)。交换两个根使 \(\delta\) 变号，因此任意 \(\sigma\in G\) 满足
\[
\sigma(\delta)=\operatorname{sgn}(\sigma)\delta.
\]
\(\delta^2=\operatorname{disc}(f)\) 属于 \(K\)，且它是 \(K\) 中的平方当且仅当 \(\delta\in K\)：若 \(\delta^2=c^2\)，则 \(\delta=\pm c\)；反向显然。由于 \(\delta\ne0\) 且特征不为二，上式又说明 \(G\subseteq A_n\) 当且仅当 \(G\) 的每个元素都固定 \(\delta\)。由 \(L^G=K\) 即得结论。
\end{proof}

特征二的限制确有必要。例如 \(f=X^2+X+1\in\mathbb F_2[X]\) 不可约且可分，其分裂域为 \(\mathbb F_4\)；非平凡自同构交换两个根，故 \(G=S_2\not\subseteq A_2\)。但 \(\operatorname{disc}(f)=1\) 是 \(\mathbb F_2\) 中的平方。此时 \(1=-1\)，从 \(\sigma(\delta)=\operatorname{sgn}(\sigma)\delta\) 不能辨出置换的奇偶性。

若 \(\operatorname{char}K\ne2\)，一个 \(K\) 上的可分不可约三次多项式的 Galois 群是 \(S_3\) 的传递子群，其阶被三整除且整除六，所以只有 \(A_3\cong C_3\) 与 \(S_3\) 两种可能；\cref{b1:prop:discriminant-alternating-group}用判别式是否为平方将它们区分。

\ProseParagraphBreak
这里的轨道大小取自原基域上的不可约分解。若改在有限域中分解整数多项式，则可以借助 Frobenius 自同构得到原 Galois 群中单个元素的循环型。下面建立两者的联系。

\subsection{模素数分解与循环型}
\label{b1:subsec:1903:04}

对一个整数多项式，模素数后的因子次数往往容易计算。要把这些数据用于特征零中的 Galois 群，须使有限域中置换根的 Frobenius 自同构来自原分裂域的自同构。以下证明用全部根生成的整环及其有限剩余域实现这一点；用到的工具是\cref{b1:lem:integral-closure-elementary}的整性封闭性、\cref{b1:thm:crt}的中国剩余定理，以及已经建立的有限域和有限 Galois 理论。

\begin{lemma}\label{b1:lem:reduction-cycle-type}
设 \(f\in\mathbb Z[X]\) 为首一多项式，并且 \(f\) 在 \(\mathbb Q\) 上无重根。令 \(L\) 为 \(f\) 在 \(\mathbb Q\) 上的分裂域，\(G=\operatorname{Gal}(L/\mathbb Q)\)。设素数 \(p\) 使 \(f\) 的模 \(p\) 约化 \(\bar f\in\mathbb F_p[X]\) 也无重根。若 \(\bar f\) 的不同首一不可约因子的次数为 \(d_1,\ldots,d_r\)，则 \(G\) 在 \(f\) 的全部根上的置换作用中包含一个循环长度为 \(d_1,\ldots,d_r\) 的元素。
\end{lemma}
\begin{proof}
设分裂域为 \(L\)，根为 \(\alpha_1,\ldots,\alpha_n\)，\(G=\operatorname{Gal}(L/\mathbb Q)\)，并取环 \(A=\mathbb Z[\alpha_1,\ldots,\alpha_n]\)。所有根都满足首一整数方程，\cref{b1:lem:integral-closure-elementary}说明 \(A\) 由有限多个元素作为 \(\mathbb Z\) 模生成，且每个元素都在 \(\mathbb Z\) 上整。\(G\) 置换这些根，所以保持 \(A\)。由 \(L^G=\mathbb Q\)，\(A^G\) 中的元素为有理整元素，只能是整数：一个既约分数若满足首一整数方程，清除分母后就要求分母整除分子的某次幂，故分母为一。因此 \(A^G=\mathbb Z\)。

理想 \(pA\) 为真理想，否则 \(1/p\in A\)，与刚才的有理整元素判别矛盾。由\cref{b1:prop:maximal-ideal-exists}取极大理想 \(\mathfrak m\supseteq pA\)。商域 \(k=A/\mathfrak m\) 有限，因为 \(A/pA\) 由有限多个元素作为 \(\mathbb F_p\) 向量空间生成，而 \(k\) 是它的商。令
\[
D=\{\,\sigma\in G\mid\sigma(\mathfrak m)=\mathfrak m\,\},
\]
这是 \(\mathfrak m\) 在 \(G\) 作用下的稳定子群，故为子群。记商映射为 \(q:A\longrightarrow k=A/\mathfrak m\)。对 \(\sigma\in D\) 定义
\[
\rho(\sigma)\bigl(q(a)\bigr)=q\bigl(\sigma(a)\bigr)\qquad(a\in A).
\]
由于 \(\sigma(\mathfrak m)=\mathfrak m\)，此式与代表元无关，且逆映射由 \(\sigma^{-1}\) 给出；\(\sigma\) 固定整数，所以 \(\rho(\sigma)\) 固定 \(\mathbb F_p\)。于是得到同态 \(\rho:D\longrightarrow\operatorname{Gal}(k/\mathbb F_p)\)，记其像为 \(H\)。现在证明 \(H\) 是全群。

任取 \(x\in k^H\)。\(G\) 作用下 \(\mathfrak m\) 的轨道由有限个不同的极大理想组成，它们两两互素。\cref{b1:thm:crt}给出 \(a\in A\)，在 \(\mathfrak m\) 处的余数为 \(x\)，在其余轨道理想处的余数为零。多项式
\[
P(T)=\prod_{\sigma\in G}\bigl(T-\sigma(a)\bigr)
\]
的系数属于 \(A^G=\mathbb Z\)。若 \(\sigma\in D\)，因 \(x\) 被 \(H\) 固定，\(\sigma(a)\) 模 \(\mathfrak m\) 为 \(x\)；若 \(\sigma\notin D\)，则 \(a\) 模 \(\sigma^{-1}(\mathfrak m)\) 为零，故 \(\sigma(a)\) 模 \(\mathfrak m\) 为零。写 \(|D|=p^e u\)、\(p\nmid u\)，约化后得到
\[
\bar P(T)=T^{|G|-|D|}(T-x)^{|D|}
=T^{|G|-|D|}\left(T^{p^e}-x^{p^e}\right)^u.
\]
若直接看 \(T^{|G|-1}\) 的系数，只得到 \(-|D|x\)，而 \(p\) 可能整除 \(|D|\)。写成上式后，降次 \(p^e\) 的项保留下可逆的因子 \(u\)：\(T^{|G|-p^e}\) 的系数为 \(-u x^{p^e}\)。\(\bar P\) 的全部系数在 \(\mathbb F_p\) 中，故 \(x^{p^e}=c\in\mathbb F_p\)。于是 \((x-c)^{p^e}=0\)，得到 \(x=c\)。所以 \(k^H=\mathbb F_p\)。有限域扩张 \(k/\mathbb F_p\) 为 Galois 扩张，\cref{b1:thm:finite-galois-correspondence}迫使 \(H=\operatorname{Gal}(k/\mathbb F_p)\)。

因此有限域 \(k\) 的 Frobenius 自同构 \(\operatorname{Frob}_p:b\mapsto b^p\) 在 \(D\) 中有一个提升 \(\sigma\)。提升满足
\[
q\circ\sigma|_A=\operatorname{Frob}_p\circ q,
\qquad\text{即}\qquad
q\bigl(\sigma(\alpha_i)\bigr)=q(\alpha_i)^p\quad(1\leq i\leq n).
\]
等式 \(f(X)=\prod_i(X-\alpha_i)\) 经 \(q\) 约化后表明，\(q(\alpha_i)\) 是 \(\bar f\) 的全部根；因 \(\bar f\) 无重根，原根集到这些约化根的映射是双射。由\cref{b1:thm:finite-field-subfields}及\cref{b1:thm:finite-field-automorphisms}，次数为 \(d_j\) 的不可约因子的根恰组成长度 \(d_j\) 的 Frobenius 轨道：一个根生成的域有 \(p^{d_j}\) 个元素，而固定该根的最小正 Frobenius 次幂正是 \(d_j\)。上述相容等式使 \(\sigma\) 在原根集上的置换与 Frobenius 在约化根集上的置换对应，故循环长度也相同。
\end{proof}

因此，每次无重根的模素数分解都给出群中的一种循环型。将不同素数提供的信息与传递性、判别式判据结合，往往能确定整个群。\cref{b1:prop:quintic-x5-x-1}将用这一方法计算 \(X^5-X-1\) 的 Galois 群。

\ProseParagraphBreak
以先前的 \(f=X^3-2\) 为例，模 \(5\) 有
\[
\bar f=(X-3)(X^2+3X+4).
\]
二次因子的判别式为 \(3\in\mathbb F_5\)，不是平方，故因子次数为 \(1,2\)，且约化无重根。上述引理给出原 Galois 群中的一个换位；这与前面直接算出的 \(S_3\) 相符。这里算因子次数是在有限域中完成的，换位属于特征零的分裂域自同构群则依赖刚才的提升论证。

% !TeX root = ../../Book1.tex
% 纲要标题：### 19.4 中间域的合成、交与嵌入位置
% 纲要位置：计划纲要.md，第 666 行。

\section{中间域的合成、交与嵌入位置}
\label{b1:sec:1904}

% 纲要内容（仅作写作计划，尚非教材正文）：
%
% * 合成域与交域；解释一侧 Galois 假设、限制映射及反向的群运算，说明非正规扩张的次数公式可能失败
% * 同构类型与嵌入位置的区别；展开同构延拓与共轭子群的对应
% * 在正文证明正规化子商群给出中间域自同构群，并解释嵌入数、共轭子域数与自同构数的关系；相应习题改为四次分裂域的具体计算
%
% ---
%

\cref{b1:prop:cubic-splitting-field-correspondence}已将一个分裂域的中间域全部列出。现在考察两个中间域的合成与交，以及同构的域在同一个环境域中为何仍可能占据不同位置。

\subsection{合成域与交域}
\label{b1:subsec:1904:02}

在一个共同环境域中，两个子域 \(L,E\) 的合成域 \(LE\) 指包含 \(L\cup E\) 的最小子域。若 \(L/K\) 是一个可分多项式的分裂域，换基域到 \(E\) 后，把同一批根加入 \(E\) 得到的就是 \(LE\)。因此，要使 \(LE/E\) 成为 Galois 扩张，只需在 \(L/K\) 这一侧具有正规性与可分性，并不要求 \(E/K\) 也 Galois。

\begin{proposition}\label{b1:prop:galois-compositum}
在 \(K\) 的一个共同代数闭包内，设 \(L/K\) 是有限 Galois 扩张，\(E/K\) 是有限扩张，令 \(LE\) 为由 \(L\cup E\) 生成的子域，并记 \(D=L\cap E\)。则 \(LE/E\) 为 Galois 扩张，限制映射给出同构
\[
\operatorname{Gal}(LE/E)\cong\operatorname{Gal}(L/D).
\]
因此
\[
[LE:E]=[L:D],\qquad
[LE:K]=\frac{[L:K][E:K]}{[D:K]}.
\]
\end{proposition}
\begin{proof}
把 \(L\) 写为某个可分 \(K\) 多项式 \(f\) 的分裂域。把 \(f\) 看作 \(E[x]\) 中的多项式，它仍没有重根；添加其全部根到 \(E\) 后得到的正是 \(LE\)，所以 \(LE/E\) 为有限 Galois 扩张。若 \(\sigma\in\operatorname{Gal}(LE/E)\)，则它固定 \(E\)，因而也固定 \(K\subseteq E\)。其限制 \(\sigma|_L:L\rightarrow LE\) 是一个 \(K\) 嵌入，\(L/K\) 的正规性给出 \(\sigma(L)=L\)，所以才有同态
\[
\operatorname{res}:\operatorname{Gal}(LE/E)\longrightarrow\operatorname{Gal}(L/K),
\qquad \sigma\longmapsto\sigma|_L.
\]
若限制是恒等，\(\sigma\) 就逐点固定 \(L\cup E\)。由于 \(LE\) 由这两个域生成，它也逐点固定整个 \(LE\)，故 \(\sigma=1\)。因此限制映射单射，记像为 \(H\leq\operatorname{Gal}(L/K)\)。

对 \(a\in L\)，因为 \(H\) 是限制映射的像，它被 \(H\) 固定当且仅当它被 \(\operatorname{Gal}(LE/E)\) 的全部元素固定。这个群在 \(LE\) 中的不动域为 \(E\)，所以
\[
L^H=L\cap(LE)^{\operatorname{Gal}(LE/E)}=L\cap E=D.
\]
在 \(L/K\) 的 Galois 对应中，这迫使 \(H=\operatorname{Gal}(L/D)\)，因而限制映射满到所述子群。两边都是 Galois 扩张的群，取阶便有 \([LE:E]=[L:D]\)。再用塔式公式，
\[
[LE:K]=[LE:E][E:K]
=\frac{[L:K]}{[D:K]}[E:K],
\]
得到第二项次数等式。
\end{proof}
若 \(E_1,E_2\) 都是一个有限 Galois 扩张 \(M/K\) 的中间域，对应子群 \(H_1,H_2\)，则
\[
\operatorname{Gal}(M/E_1E_2)=H_1\cap H_2,\qquad
\operatorname{Gal}(M/E_1\cap E_2)=\langle H_1,H_2\rangle.
\]
第一式是因为逐点固定由 \(E_1\cup E_2\) 生成的较大域，等价于同时逐点固定 \(E_1\) 与 \(E_2\)，所以要取两个固定群的交。另一方面，被生成子群 \(\langle H_1,H_2\rangle\) 固定，等价于同时被 \(H_1\) 与 \(H_2\) 固定，因此
\[
M^{\langle H_1,H_2\rangle}=M^{H_1}\cap M^{H_2}=E_1\cap E_2.
\]
再用 Galois 对应的互逆性，得到第二式。合成域是包含两域的最小中间域，交域是包含于两域的最大中间域；反转包含关系后，它们分别对应子群的交与包含两群的最小子群。

\ProseParagraphBreak
取\cref{b1:fig:cubic-intermediate-fields}中的 \(E_1=\mathbb Q(a)\)、\(E_2=\mathbb Q(\zeta a)\)，并沿用\cref{b1:prop:cubic-splitting-field-correspondence}的记号 \(L,G,s,t\)。它们对应 \(H_1=\langle t\rangle\)、\(H_2=\langle s^2t\rangle\)。这两个二阶子群的交为 \(1\)，且两个不同的换位生成 \(G\cong S_3\)，故上述公式给出
\[
E_1E_2=L,\qquad E_1\cap E_2=\mathbb Q.
\]
前一等式也可由 \(a\) 与 \(\zeta a\) 的比得到 \(\zeta\) 而直接看出。因此 \([E_1E_2:\mathbb Q]=6\)，并非 \([E_1:\mathbb Q][E_2:\mathbb Q]=9\)：两个域都不正规时，交域为基域不能保证次数相乘。此处
\[
H_1H_2=\{1,t,s^2t,ts^2t\}=\{1,t,s^2t,s\}
\]
只有四个元素，依 Lagrange 定理不是六阶群 \(G\) 的子群；而 \(\langle H_1,H_2\rangle\) 还包含任意长的有限乘积，恰为 \(G\)。所以交域公式不能把生成子群替换为乘积集合。若其中一群正规，则在取积或取逆时，可以用正规性把另一群的元素移过它的元素，所得仍在乘积集合中，故乘积集合是子群。此时它包含两群，且每个元素都是两群元素的乘积，所以等于生成子群。

\subsection{同构类型与嵌入位置}
\label{b1:subsec:1904:03}

对两个实际子域 \(E_1,E_2\subseteq L\)，要求 \(E_1=E_2\) 是比较其中的元素；要求 \(E_1\cong_K E_2\) 则只要求存在固定 \(K\) 的域同构，允许它们在 \(L\) 中的位置不同。例如，\(\mathbb Q(a)\) 和 \(\mathbb Q(\zeta a)\) 都同构于 \(\mathbb Q[x]/(x^3-2)\)，存在保持 \(\mathbb Q\) 并把 \(a\) 送到 \(\zeta a\) 的同构，但它们是分裂域中的不同子域。

\ProseParagraphBreak
在有限 Galois 扩张 \(L/K\) 中，记 \(G=\operatorname{Gal}(L/K)\)。还有第三个问题：是否存在 \(\sigma\in G\) 使 \(\sigma(E_1)=E_2\)？若有，限制 \(\sigma|_{E_1}:E_1\rightarrow E_2\) 就是 \(K\) 同构。反过来，给定 \(K\) 同构 \(\varphi:E_1\overset{\sim}{\longrightarrow}E_2\)，把 \(E_2\) 包含到一个含 \(L\) 的代数闭包 \(\Omega\) 中，便可将 \(\varphi\) 看作 \(K\) 嵌入 \(E_1\rightarrow\Omega\)。由\cref{b1:thm:embedding-extension}，它能延拓为 \(\widetilde\varphi:L\rightarrow\Omega\)。\(L/K\) 的正规性给出 \(\widetilde\varphi(L)=L\)，所以 \(\widetilde\varphi\in G\)，且 \(\widetilde\varphi(E_1)=E_2\)。因此在这个有限 Galois 环境中，\(K\) 同构与被 \(G\) 中元素互相送到等价，子域相等则是更强的要求。

\ProseParagraphBreak
若 \(H_i=\operatorname{Gal}(L/E_i)\)，由 \(\sigma(E_1)=E_2\) 及\cref{b1:thm:galois-quotient}证明中的共轭公式，
\[
E_2=\sigma(L^{H_1})=L^{\sigma H_1\sigma^{-1}},\qquad
H_2=\sigma H_1\sigma^{-1}.
\]
反过来，共轭子群的共同不动域也由 \(\sigma\) 互相送到。Galois 对应枚举的是实际子域，而按 \(K\) 同构类型分类则须再把共轭子群归为一类。

\begin{proposition}\label{b1:prop:subfield-normalizer-automorphisms}
设 \(L/K\) 为有限 Galois 扩张，\(G=\operatorname{Gal}(L/K)\)，\(H\leq G\)，\(E=L^H\)，并记
\[
N_G(H)=\{\,\sigma\in G\mid\sigma H\sigma^{-1}=H\,\}.
\]
则限制映射给出
\[
\operatorname{Aut}_K(E)\cong N_G(H)/H.
\]
域 \(L\) 中与 \(E\) 在 \(K\) 上同构的不同中间域共有 \([G:N_G(H)]\) 个，而且
\[
[E:K]=[G:N_G(H)]\,|\operatorname{Aut}_K(E)|.
\]
\end{proposition}
\begin{proof}
由\cref{b1:thm:galois-quotient}证明中的公式 \(\sigma(E)=L^{\sigma H\sigma^{-1}}\) 及 Galois 对应的单射性，\(\sigma(E)=E\) 当且仅当 \(\sigma\in N_G(H)\)。因此限制给出群同态
\[
\operatorname{res}:N_G(H)\longrightarrow\operatorname{Aut}_K(E),\qquad
\sigma\longmapsto\sigma|_E,
\]
核为逐点固定 \(E\) 的群 \(\operatorname{Gal}(L/E)=H\)；由正规化子的定义，\(H\) 在 \(N_G(H)\) 中正规。任取 \(\tau\in\operatorname{Aut}_K(E)\)，将它看作到一个含 \(L\) 的代数闭包 \(\Omega\) 的 \(K\) 嵌入，由\cref{b1:thm:embedding-extension}延拓为 \(\widetilde\tau:L\rightarrow\Omega\)。由于 \(L/K\) 正规，\(\widetilde\tau(L)=L\)，故 \(\widetilde\tau\in G\)。它在 \(E\) 上等于 \(\tau\)，所以 \(\widetilde\tau(E)=E\)，即属于 \(N_G(H)\)。限制映射因而满射，第一同构定理给出所述商群同构。

\(G\) 作用在 \(L\) 的中间域集合上，作用为 \(E'\mapsto\sigma(E')\)。由前面的同构延拓论证，与 \(E\) 在 \(K\) 上同构的域恰为这个作用下 \(E\) 的轨道；将 \(E\) 保持为集合的稳定子群恰为 \(N_G(H)\)。\cref{b1:thm:orbit-stabilizer}给出不同域的数目 \([G:N_G(H)]\)。最后，由 Galois 对应与指数乘法公式，
\[
[E:K]=[G:H]=[G:N_G(H)]\,[N_G(H):H]
=[G:N_G(H)]\,|\operatorname{Aut}_K(E)|.
\]
\end{proof}
这个次数等式也可以直接按嵌入来理解。中间扩张 \(E/K\) 可分，所以它到代数闭包的不同 \(K\) 嵌入共有 \([E:K]\) 个。每个嵌入都延拓到 \(L\)，由正规性，其像是 \(L\) 内与 \(E\) 同构的某个子域。固定其中一个像域 \(E'\) 并选一个 \(K\) 同构 \(\varphi:E\rightarrow E'\) 后，全部以 \(E'\) 为像的嵌入恰为 \(\varphi\circ\tau\)，其中 \(\tau\in\operatorname{Aut}_K(E)\)。因此嵌入总数等于共轭副本的数目乘以每个副本内的同构数。自同构要求像仍为原来的 \(E\)，一般嵌入则可以把它送到另一个副本。

\ProseParagraphBreak
在\cref{b1:prop:cubic-splitting-field-correspondence}的三次分裂域中，对 \(E=\mathbb Q(a)\)，对应 \(H=\langle t\rangle\)。正规化 \(H\) 的元素须在共轭下固定它唯一的非单位元 \(t\)；在 \(S_3\) 中，与这个换位交换的元素只有 \(1,t\)，所以 \(N_G(H)=H\)。于是 \(E\) 有三个不同的共轭子域，却只有恒等 \(\mathbb Q\) 自同构，次数公式为 \(3=3\cdot1\)。对 \(E=\mathbb Q(\zeta)\)，对应的三阶子群 \(H=\langle s\rangle\) 正规，故 \(N_G(H)=G\)。这个二次域是同一环境中唯一的同构副本，但有两个 \(\mathbb Q\) 自同构，次数公式为 \(2=1\cdot2\)。两种情形都具有次数那么多的嵌入，只是嵌入是否保持同一子域不同。


\begin{chaptersummary}
有限自同构群的不动域总给出一个次数恰为群阶的 Galois 扩张。由此，中间域与子群反向对应，扩张次数对应群指数，对基域为 Galois 的中间域对应正规子群，限制映射给出相应商群。具体计算时，群的抽象类型还不够：须明确它作用于哪些根、哪些元素被固定。合成域与交域公式需要核对正规可分假设，非正规扩张中的交域为基域并不保证次数相乘。同构的中间域也可能在同一个分裂域中占据不同位置；若 \(E=L^H\)，整体保持 \(E\) 的群为 \(N_G(H)\)，逐点固定 \(E\) 的群为 \(H\)，所以 \(\operatorname{Aut}_K(E)\cong N_G(H)/H\)。下一章将研究 Galois 群为循环群的扩张，考察生成自同构、范数条件与扩张生成元之间的关系。
\end{chaptersummary}

\BookExercises

\begin{exercise}\label{b1:ex:19-characteristic-two}
比较 \(\mathbb F_4/\mathbb F_2\) 与 \(\mathbb F_2(u)/\mathbb F_2(u^2)\)，其中 \(u\) 超越。它们同为二次扩张，分别是否为 Galois 扩张？求各自的基域自同构群。
\end{exercise}
\begin{solution}
前者是 \(x^2+x+1\) 的分裂域，该多项式导数为一，故可分。Galois 群由恒等与平方映射组成，为 \(C_2\)。后者为纯不可分扩张；若嵌入固定 \(u^2\)，则 \(u\) 的像必须满足 \(x^2-u^2=(x-u)^2=0\)，只能仍为 \(u\)。所以自同构群平凡，扩张不可分，不是 Galois 扩张。这个最小多项式虽有重根，却已在扩张域中完全分裂，且其根 \(u\) 生成扩张，所以该扩张仍正规。二次代数扩张的正规性，并不保证它在特征二时也可分。
\end{solution}

\begin{exercise}\label{b1:ex:19-real-fourth-root}
令 \(a=\sqrt[4]2>0\)、\(E=\mathbb Q(a)\)。求 \(\operatorname{Aut}_{\mathbb Q}(E)\) 及其不动域，解释为何这个不动域不是 \(\mathbb Q\)。
\end{exercise}
\begin{solution}
\(x^4-2\) 由素数二的 Eisenstein 判据不可约，是 \(a\) 的最小多项式，故 \([E:\mathbb Q]=4\)。\(E\subseteq\mathbb R\)，所以 \(a\) 的像只能是两个实根 \(a,-a\)。将 \(\mathbb Q[x]/(x^4-2)\) 的生成元分别送到这两个根，所得的像都为 \(E\)，所以确有这两个自同构。因此 \(H=\operatorname{Aut}_{\mathbb Q}(E)\cong C_2\)，其非平凡元 \(s\) 把 \(a\) 变号。

\(\mathbb Q(a^2)=\mathbb Q(\sqrt2)\) 被 \(H\) 固定，且 \([E:\mathbb Q(a^2)]=4/2=2\)。\cref{b1:thm:artin-fixed-field}给出 \([E:E^H]=|H|=2\)，由包含关系和塔式公式可知 \(E^H=\mathbb Q(\sqrt2)\)。该定理保证的是 \(E/\mathbb Q(\sqrt2)\) 为 Galois 扩张；原来的 \(E/\mathbb Q\) 缺少非实共轭根，仍不正规。这也具体说明，在非 Galois 扩张中，\(\mathbb Q\subseteq E^{\operatorname{Aut}_{\mathbb Q}(E)}\) 可以是严格包含。
\end{solution}

\begin{exercise}\label{b1:ex:19-rational-inversion}
设 \(k\) 为任意域，\(t\) 在 \(k\) 上超越。令 \(s\in\operatorname{Aut}_k(k(t))\) 为满足 \(s(t)=t^{-1}\) 的 \(k\) 自同构。求 \(\langle s\rangle\) 的不动域。
\end{exercise}
\begin{solution}
\(s\) 逐点固定 \(k\)，且 \(s^2(t)=t\)，故 \(s^2=\operatorname{id}_{k(t)}\)。又 \(s\ne\operatorname{id}\)，因为 \(t=t^{-1}\) 将使超越元满足 \(t^2=1\)，所以 \(\langle s\rangle\) 的阶为二。令 \(v=t+t^{-1}\)，它被 \(s\) 固定，且 \(t\) 满足 \(x^2-vx+1=0\)，所以 \([k(t):k(v)]\leq2\)。若次数为一，\(s\) 固定 \(v\) 与 \(k\) 就固定 \(t\)，矛盾，故次数为二。与\cref{b1:thm:artin-fixed-field}的次数比较，得到不动域恰为 \(k(t+t^{-1})\)。论证在特征二也成立。
\end{solution}

\begin{exercise}\label{b1:ex:19-faithful-action}
让抽象循环群 \(C_4\) 的生成元通过复共轭作用在 \(\mathbb C\) 上。求不动域，并解释为何不能据此说 \([\mathbb C:\mathbb R]=4\)。
\end{exercise}
\begin{solution}
共同不动元为 \(\mathbb R\)。但作用同态 \(\rho:C_4\rightarrow\operatorname{Aut}(\mathbb C)\) 的核由生成元的平方生成，其像只有恒等和共轭两个元素。正如群作用通过核的商群给出忠实作用，这里的不动元只取决于像群 \(\rho(C_4)\cong C_4/\ker\rho\cong C_2\)。不动域定理计算的是实际互异自同构所组成的群的阶，因而给出 \([\mathbb C:\mathbb R]=2\)。只有当抽象群的作用忠实时，才能直接把它的阶作为这个次数。
\end{solution}

\begin{exercise}\label{b1:ex:19-biquadratic-five}
列出 \(L=\mathbb Q(\sqrt2,\sqrt5)\) 的全部中间域，并写明各自对应的子群。
\end{exercise}
\begin{solution}
若 \(\sqrt5=a+b\sqrt2\)，平方比较给出 \(ab=0\) 及 \(a^2+2b^2=5\)，无有理解，因为五与 \(5/2\) 均非有理数平方。故总次数为四，群由分别变号的 \(s,t\) 生成，满足 \(s^2=t^2=1,st=ts\)。子群为平凡群、全群和 \(\langle s\rangle,\langle t\rangle,\langle st\rangle\)。若 \(s\) 改变 \(\sqrt2\) 的符号、\(t\) 改变 \(\sqrt5\) 的符号，则对应域依次为 \(L,\mathbb Q,\mathbb Q(\sqrt5),\mathbb Q(\sqrt2),\mathbb Q(\sqrt{10})\)。这些子群都是正规子群，所以各中间域对 \(\mathbb Q\) 均为 Galois 扩张。
\end{solution}

\begin{exercise}\label{b1:ex:19-cyclic-twelve}
设有限 Galois 扩张 \(L/K\) 的群为 \(C_{12}=\langle\sigma\rangle\)。求所有可能的中间域次数，判断每个次数的中间域是否唯一，并求四次中间域对 \(K\) 的 Galois 群。
\end{exercise}
\begin{solution}
循环群对每个正因子有唯一相应阶的子群，所以中间域次数为 \(1,2,3,4,6,12\)，每个次数恰有一个。设四次中间域对应 \(H\)，则 \([G:H]=4\)，而 \(|H|=12/4=3\)：对应的是指数为四、阶为三的子群，不是四阶子群。因此 \(H=\langle\sigma^4\rangle\)，且它正规。由\cref{b1:thm:galois-quotient}，四次域的 Galois 群为 \(C_{12}/\langle\sigma^4\rangle\cong C_4\)。
\end{solution}

\begin{exercise}\label{b1:ex:19-mod-primes-s4}
设 \(f=X^4-X-1\in\mathbb Q[X]\)。求 \(f\) 在 \(\mathbb Q\) 上的 Galois 群，并证明所得结论。
\end{exercise}
\begin{solution}
在 \(\mathbb F_2[X]\) 中，\(\bar f=X^4+X+1\) 无根，且除以唯一的首一不可约二次式 \(X^2+X+1\) 余 \(1\)，故 \(\bar f\) 不可约，因而无重根。由于 \(f\) 首一，模 \(2\) 不可约也推出 \(f\) 在 \(\mathbb Q[X]\) 中不可约。循环型判据给出 Galois 群 \(G\leq S_4\) 中一个四轮换。

在 \(\mathbb F_7[X]\) 中，
\[
\bar f=(X-3)(X^3+3X^2+2X-2).
\]
右侧三次因子在 \(0,1,\ldots,6\) 处的值依次为 \(5,4,1,2,6,5,5\)，故在 \(\mathbb F_7\) 中无根，从而不可约；它在 \(3\) 处的值为 \(2\ne0\)，所以约化无重根，判据又给出一个三轮换。由 Lagrange 定理，四轮换和三轮换分别给出 \(4\bigm\vert|G|\) 与 \(3\bigm\vert|G|\)。由于三、四互素，故 \(12\bigm\vert|G|\)；而 \(G\leq S_4\) 又给出 \(|G|\bigm\vert24\)。因此 \(|G|\) 只能为十二或二十四，这并没有断言那两个元素生成的子群恰有十二个元素。

若 \(|G|=12\)，四轮换为奇置换，所以符号同态 \(\operatorname{sgn}|_G:G\rightarrow\{\pm1\}\) 满射，核为 \(G\cap A_4\)。第一同构定理给出 \(|G\cap A_4|=|G|/2=6\)，与\cref{b1:prop:alternating-three-cycles}中 \(A_4\) 没有六阶子群的结论矛盾。因此 \(|G|=24\)，即 \(G=S_4\)。
\end{solution}

\begin{exercise}\label{b1:ex:19-dihedral-quadratics}
在 \(L=\mathbb Q(a,i)\)、\(a^4=2\) 中，设 \(r(a)=ia,r(i)=i\)，\(s\) 为复共轭。证明三个二次中间域为 \(\mathbb Q(i),\mathbb Q(\sqrt2),\mathbb Q(\sqrt{-2})\)，并写出对应子群。
\end{exercise}
\begin{solution}
二次域对应 \(D_8\) 中的指数二子群，即四阶子群。它们恰为 \(\langle r\rangle\)、\(\langle r^2,s\rangle\)、\(\langle r^2,rs\rangle\)：若四阶子群含奇次幂的 \(r\)，就包含 \(\langle r\rangle\)；否则其旋转部分为 \(\{1,r^2\}\)，两个反射的指数奇偶相同，只给出后两种。三群分别固定 \(i,a^2,ia^2\)，这些元素生成的域均为二次域；次数相同使包含成为相等，因此得所列三个域。
\end{solution}

\begin{exercise}\label{b1:ex:19-cyclotomic-five}
令 \(\zeta\) 为本原五次单位根。求 \(\mathbb Q(\zeta)\) 的唯一二次中间域，并写出其一个生成元满足的二次方程。
\end{exercise}
\begin{solution}
正文已证明群为 \(C_4\)，故二次中间域唯一，对应由复共轭生成的二元子群。令 \(u=\zeta+\zeta^{-1}\)，则其为实数且被共轭固定。由 \(1+\zeta+\zeta^2+\zeta^3+\zeta^4=0\)，有 \(\zeta^2+\zeta^{-2}=-1-u\)，故 \(u^2=1-u\)。多项式 \(t^2+t-1\) 无有理根，因而 \(\mathbb Q(u)\) 为二次域，且 \((2u+1)^2=5\)。所求域为 \(\mathbb Q(\sqrt5)\)。
\end{solution}

\begin{exercise}\label{b1:ex:19-reducible-permutation}
求 \((x^2-2)(x^2-3)\) 对 \(\mathbb Q\) 的 Galois 群及其在四个根上的轨道。将它与不可约四次多项式 \(x^4-10x^2+1\) 的情形比较。
\end{exercise}
\begin{solution}
二者分裂域都是 \(\mathbb Q(\sqrt2,\sqrt3)\)，抽象 Galois 群均为 \(C_2\times C_2\)。前一个多项式的根分成两个轨道 \(\{\sqrt2,-\sqrt2\}\)、\(\{\sqrt3,-\sqrt3\}\)，因为有理系数自同构不能把平方为二的元素送到平方为三的元素。后一个多项式的四根为 \(\pm\sqrt2\pm\sqrt3\)，独立变号使其成为一个轨道。这说明相同的抽象群可以在不同根集上给出不同的置换作用。
\end{solution}

\begin{exercise}\label{b1:ex:19-cyclic-cubic}
证明 \(f=x^3-3x-1\) 在 \(\mathbb Q\) 上的 Galois 群为 \(C_3\)。计算判别式时可用全部根 \(a_i\) 上的导数乘积。
\end{exercise}
\begin{solution}
有理根只能为 \(\pm1\)，二者都不是根，故三次式不可约。设三个根为 \(a_1,a_2,a_3\)。由 \(f'(a_i)=\prod_{j\ne i}(a_i-a_j)\)，对三个指标相乘得到 \(\prod_i f'(a_i)=-\operatorname{disc}(f)\)。又 \(f'(x)=3(x-1)(x+1)\)，所以
\[
\prod_i f'(a_i)=27\Bigl(\prod_i(a_i-1)\Bigr)\Bigl(\prod_i(a_i+1)\Bigr)
=27\bigl(-f(1)\bigr)\bigl(-f(-1)\bigr)=-81.
\]
故判别式为 \(81\)，是非零平方。\eqref{b1:eq:discriminant-alternating-group}说明群包含于 \(A_3\)，而不可约性使作用传递，群阶被三整除。因此它就是 \(A_3\cong C_3\)。
\end{solution}

\begin{exercise}\label{b1:ex:19-product-groups}
设 \(L/K,E/K\) 都是有限 Galois 扩张，且 \(L\cap E=K\)。证明 \(LE/K\) 的 Galois 群通过限制同构于 \(\operatorname{Gal}(L/K)\times\operatorname{Gal}(E/K)\)。
\end{exercise}
\begin{solution}
两域都是可分多项式的分裂域。取这两个多项式的全部互异不可约因子的乘积，它仍可分，且分裂域为 \(LE\)，故 \(LE/K\) 为 Galois 扩张。两个限制给出到直积的同态，固定两域的自同构固定其合成域，所以该同态单射。由\cref{b1:prop:galois-compositum}以及交域假设，\([LE:K]=[L:K][E:K]\)。这恰为目标直积群的阶，也等于来源群的阶，所以单射为满射。
\end{solution}

\begin{exercise}\label{b1:ex:19-two-fourth-root-fields}
设 \(a=\sqrt[4]2>0\)，比较 \(E_1=\mathbb Q(a)\) 与 \(E_2=\mathbb Q(ia)\)。证明二者在 \(\mathbb Q\) 上同构而不相等，并求它们的交域、合成域。
\end{exercise}
\begin{solution}
\(a,ia\) 都以不可约多项式 \(x^4-2\) 为最小多项式，所以代入映射给出两个域与 \(\mathbb Q[x]/(x^4-2)\) 的同构。但前者为实域，后者含非实元，故不相等。两者均含 \(a^2=\sqrt2\)，因为 \((ia)^2=-a^2\)。交域对 \(\mathbb Q\) 的次数整除四且至少为二；若为四，则两域相等，矛盾。因此交域恰为 \(\mathbb Q(\sqrt2)\)。合成域含 \((ia)/a=i\)，故为 \(\mathbb Q(a,i)\)。由于 \(\mathbb Q(a)\) 是实域，\(i\notin\mathbb Q(a)\)，而 \(i\) 满足 \(x^2+1=0\)，所以合成域对 \(\mathbb Q(a)\) 的次数为二，总次数为八。

这里直接识别了合成域并独立计算其次数。\(E_1/\mathbb Q\) 与 \(E_2/\mathbb Q\) 都不正规，不能以\cref{b1:prop:galois-compositum}为依据套用次数乘除公式；本题数值恰与 \(4\cdot4/2\) 相同，并不使缺少假设的公式成立。
\end{solution}

\begin{exercise}\label{b1:ex:19-normalizer-automorphisms}
设 \(a=\sqrt[4]2>0\)、\(L=\mathbb Q(a,i)\)，并令 \(r(a)=ia,r(i)=i\)，\(s\) 为复共轭。由正文知 \(G=\operatorname{Gal}(L/\mathbb Q)=\langle r,s\rangle\cong D_8\)。对 \(E=\mathbb Q(a)\)，求逐点固定 \(E\) 的子群 \(H\) 及其正规化子，写出 \(E\) 的全部 \(\mathbb Q\) 自同构和 \(L\) 中与 \(E\) 在 \(\mathbb Q\) 上同构的全部中间域，并核对
\[
[E:\mathbb Q]=[G:N_G(H)]\,|\operatorname{Aut}_{\mathbb Q}(E)|.
\]
再对 \(E'=\mathbb Q(i)\) 完成相同的计算，并比较两个子域的嵌入如何分成不同像域。
\end{exercise}
\begin{solution}
\(s\) 固定 \(a\)，故 \(\langle s\rangle\subseteq H\)。又 \([L:E]=2\)，Galois 对应给出 \(|H|=2\)，所以 \(H=\langle s\rangle\)。由于 \(H\) 只有一个非单位元，将它保持为集合就须固定这个元素。利用 \(srs=r^{-1}\)，有
\[
r^j s r^{-j}=r^{2j}s\qquad(0\leq j<4).
\]
元素 \(r^js\) 的共轭作用在 \(s\) 上也给出 \(r^{2j}s\)，因此当且仅当 \(j\) 为偶数时属于正规化子。于是
\[
N_G(H)=\{1,r^2,s,r^2s\}=\langle r^2,s\rangle.
\]
\cref{b1:prop:subfield-normalizer-automorphisms}给出 \(\operatorname{Aut}_{\mathbb Q}(E)\cong N_G(H)/H\cong C_2\)。两自同构分别把 \(a\) 送到 \(a\) 与 \(-a\)，非平凡者来自 \(r^2\) 的限制。

不同共轭子域共有 \([G:N_G(H)]=8/4=2\) 个，实际为 \(\mathbb Q(a)\) 和 \(\mathbb Q(ia)\)：\(G\) 将 \(a\) 送到 \(\pm a,\pm ia\)，正负号不改变生成的域。这两个域不同，因为前者为实域，后者含非实元。四个 \(\mathbb Q\) 嵌入因而分成两组，每组有两个嵌入，其像为同一个子域，次数公式在这里成为 \(4=2\cdot2\)。

对 \(E'=\mathbb Q(i)\)，\(r\) 固定 \(i\)，而 \([L:E']=4\)，故对应子群为 \(H'=\langle r\rangle\)。它在 \(G\) 中正规，因而 \(N_G(H')=G\)。相应自同构群为 \(G/H'\cong C_2\)，其元素分别将 \(i\) 送到 \(i\) 与 \(-i\)，共轭子域只有 \(E'\) 自身。此时
\[
[E':\mathbb Q]=2=1\cdot2.
\]
两个嵌入都保持同一子域，正是这个二次扩张正规性的表现。
\end{solution}

\begin{exercise}\label{b1:ex:19-translation-fixed-field}
设 \(k\) 为特征 \(p>0\) 的域，\(t\) 在 \(k\) 上超越。对每个 \(c\in\mathbb F_p\)，令 \(\sigma_c\in\operatorname{Aut}_k(k(t))\) 为逐点固定 \(k\)、满足 \(\sigma_c(t)=t+c\) 的自同构。求这些自同构所成群的不动域。
\end{exercise}
\begin{solution}
这 \(p\) 个自同构互异，且 \(\sigma_c\sigma_d=\sigma_{c+d}\)。元素 \(v=t^p-t\) 被全部自同构固定，因为 \(c^p=c\)。又 \(t\) 满足 \(x^p-x-v=0\)，所以 \([k(t):k(v)]\leq p\)。由\cref{b1:thm:artin-fixed-field}，对真正不动域的扩张次数恰为 \(p\)；由于 \(k(v)\) 包含于该不动域，塔式公式迫使两者相等。故不动域为 \(k(t^p-t)\)。
\end{solution}

% !TeX root = ../../Book1.tex
% 纲要标题：## 第20章　循环扩张、Hilbert 第 90 定理与 Kummer 理论
% 纲要位置：计划纲要.md，第 674 行。

\chapter[循环扩张、Hilbert 第 90 定理与 Kummer 理论]{循环扩张、\mbox{Hilbert 第 90 定理}与 Kummer 理论}
\label{b1:ch:20}
\BookChapterNumber{b1:ch:20}

\begin{chapterabstract}
循环扩张中的迹与范数满足特殊的相消关系。Hilbert 第 \(90\) 定理刻画其核，本章据此研究单位根条件下的根式生成与 Kummer 对应，并以 Artin--Schreier 方程处理正特征的循环扩张。
\end{chapterabstract}

\begin{learninggoals}
\begin{itemize}
\item 说明循环扩张的迹与范数公式，并区分乘法群与加法群上的两个核。
\item 掌握 Hilbert 第 90 定理的乘法形式、加法形式及各自的证明，指出可分性和嵌入独立性各在何处使用。
\item 核对 Kummer 理论的特征、单位根及有限性假设，利用配对计算扩张次数。
\item 用乘法或加法商空间中的一维子空间分类素数次循环扩张。
\end{itemize}
\end{learninggoals}

复共轭给出二次扩张 \(\mathbb C/\mathbb R\) 的一个自同构。取 \(b=2+i\)，则 \(b/\overline b=(3+4i)/5\)，这个数与其共轭的乘积为 \(1\)。若先给定一个范数为 \(1\) 的元素，是否总能反过来找到这样的 \(b\)？在循环扩张中，这个问题由 Hilbert 第 \(90\) 定理回答；它同时说明范数条件为何会出现在由幂根生成的扩张中。把乘积换成求和，又会得到一条平行而不相同的迹的结论。

\ProseParagraphBreak

本章使用\cref{b1:thm:finite-galois-correspondence}的固定域对应与\cref{b1:thm:trace-norm-embeddings}的迹、范数公式。乘法形式的 Hilbert 第 90 定理还调用\cref{b1:cor:embedding-independence}；Kummer 对应需要第一编的有限交换群结构和有限维线性代数。下一章的特征零根式可解性判据使用乘法形式、单位根与循环扩张的结果；加法形式则用于本章的 Artin--Schreier 理论。

\ProseParagraphBreak
循环扩张、二项式与 Kummer 理论，可参阅 Roman 的《Field Theory》\cite[Chapter 12; Sections 14.1--14.2, 15.1--15.3]{Roman2006FieldTheory}。若希望继续学习\subsecref{b1:subsec:2002:03}的上同调解释，可在补充群上同调之后阅读 Gille 与 Szamuely 的《Central Simple Algebras and Galois Cohomology》\cite[Example 2.3.4; Section 4.3]{GilleSzamuely2017CentralSimpleAlgebras}，其中还说明 Kummer 理论和 Artin--Schreier 理论与上同调正合列的关系。

% !TeX root = ../../Book1.tex
% 纲要标题：### 20.1 循环扩张
% 纲要位置：计划纲要.md，第 676 行。

\section{循环扩张}
\label{b1:sec:2001}

循环扩张的全部自同构都由一个自同构反复复合得到。对一个 \(n\) 次循环扩张，迹与范数分别按全部 \(n\) 个自同构求和、求积；即使不同自同构在给定元素上的取值相同，也要按出现次数计入。下面的相消计算给出迹核与范数核中的元素，下一节将证明它们恰好构成整个核。

% 纲要内容（仅作写作计划，尚非教材正文）：
%
% * 循环 Galois 群与生成自同构
% * 范数为一的元素
% * 加法与乘法形式的区别
%

\subsection{循环 Galois 群与生成自同构}
\label{b1:subsec:2001:01}

\begin{definition}\label{b1:def:cyclic-extension}
设 \(L/K\) 为有限 Galois 扩张。如果保持 \(K\) 中每个元素不变的全部 \(L\) 自同构组成的群 \(\operatorname{Gal}(L/K)\) 是循环群，则称 \(L/K\) 为\term[xunhuan-kuozhang]{循环扩张}[cyclic extension]。记 \([L:K]=n\)，并选取 \(\operatorname{Gal}(L/K)\) 的一个生成元 \(\sigma\)，于是
\[
G=\operatorname{Gal}(L/K)=\{1,\sigma,\ldots,\sigma^{n-1}\},
\qquad \sigma^n=1.
\]
\end{definition}

这里必须同时要求正规与可分。一个扩张的自同构群偶然是循环群，并不能保证它是循环扩张。例如 \(\mathbb Q(\sqrt[3]{2})/\mathbb Q\) 的自同构群只有恒等映射，但它不是正规扩张。相反，在特征不为 \(2\) 的域上，非平凡二次扩张 \(K(\sqrt d)/K\) 的两个自同构为恒等映射及 \(\sqrt d\mapsto-\sqrt d\)，故为循环扩张。

\ProseParagraphBreak
生成元通常不是唯一的。\(\sigma^r\) 生成同一个 \(n\) 阶群当且仅当 \(\gcd(r,n)=1\)。由\cref{b1:thm:finite-galois-correspondence}，每个 \(n\) 的正因子 \(d\) 对应唯一的 \(d\) 次中间扩张；具体地，\(L^{\langle\sigma^d\rangle}\) 在 \(K\) 上的次数为 \(d\)。这个次数公式使用的是子群指数，而不是子群阶。

\subsection{范数为一的元素}
\label{b1:subsec:2001:02}

设 \(L/K\) 为 \(n\) 次循环扩张，\(\operatorname{Gal}(L/K)=\langle\sigma\rangle\)。对每个 \(x\in L\)，\cref{b1:thm:trace-norm-embeddings}给出
\begin{equation}\label{b1:eq:cyclic-trace-norm}
\operatorname{Tr}_{L/K}(x)=\sum_{i=0}^{n-1}\sigma^i(x),
\qquad
N_{L/K}(x)=\prod_{i=0}^{n-1}\sigma^i(x).
\end{equation}
迹是加法群同态 \(\operatorname{Tr}_{L/K}:L\to K\)；将范数限制于非零元素，则得到群同态 \(N_{L/K}:L^\times\to K^\times\)。若 \(b\in L^\times\)，则
\[
N_{L/K}\Bigl(\frac{b}{\sigma(b)}\Bigr)
=\prod_{i=0}^{n-1}\frac{\sigma^i(b)}{\sigma^{i+1}(b)}=1.
\]
分子、分母逐项相消，末项使用 \(\sigma^n(b)=b\)。因此每个这样的比值都在范数的核中。

\ProseParagraphBreak
以 \(\mathbb C/\mathbb R\) 为例，生成自同构是复共轭，范数为 \(z\bar z\)。非零数 \(b/\bar b\) 总在单位圆上；反过来，单位圆上的每个数能否如此表示，就是乘法形式要回答的问题。这个问题只要求找出一个 \(b\)，不要求唯一：用任意非零实数乘 \(b\)，比值并不改变。

\subsection{加法形式与乘法形式}
\label{b1:subsec:2001:03}

对加法作同样计算，可得
\[
\operatorname{Tr}_{L/K}\bigl(b-\sigma(b)\bigr)
=\sum_{i=0}^{n-1}\bigl(\sigma^i(b)-\sigma^{i+1}(b)\bigr)=0.
\]
两种计算分别给出
\[
\bigl\{\,b/\sigma(b)\mid b\in L^\times\,\bigr\}\subseteq\ker N_{L/K},
\qquad
\{\,b-\sigma(b)\mid b\in L\,\}\subseteq\ker\operatorname{Tr}_{L/K}.
\]
这里的比值映射 \(b\mapsto b/\sigma(b)\) 从 \(L^\times\) 到 \(L^\times\)；必须排除 \(b=0\)，因为分母不能为零。范数映射的目标则是 \(K^\times\)。加法差映射 \(b\mapsto b-\sigma(b)\) 从 \(L\) 到 \(L\)。在特征整除 \(n\) 时，加法形式不能靠除以 \(n\) 来证明；下一节的论证依靠扩张的可分性，对这些特征仍然成立。

% !TeX root = ../../Book1.tex
% 纲要标题：### 20.2 Hilbert 第 90 定理
% 纲要位置：计划纲要.md，第 682 行。

\section{Hilbert 第 90 定理}
\label{b1:sec:2002}

上一节的相消计算给出了范数为一、迹为零的充分条件。\term[Hilbert-di90-dingli]{Hilbert 第 90 定理}[Hilbert's Theorem 90]说明这些条件也是必要的。以下固定有限循环扩张 \(L/K\)，并记 \(\operatorname{Gal}(L/K)=\langle\sigma\rangle\)、\([L:K]=n\)。

% 纲要内容（仅作写作计划，尚非教材正文）：
%
% * 循环扩张的乘法形式及证明
% * 加法形式
% * 第四卷第一群上同调的重新解释
%

\subsection{Hilbert 第 90 定理的乘法形式}
\label{b1:subsec:2002:01}

\begin{theorem}[Hilbert 第 90 定理，乘法形式]\label{b1:thm:hilbert90-multiplicative}
设 \(L/K\) 为 \(n\) 次循环扩张，并选取生成元 \(\sigma\in\operatorname{Gal}(L/K)\)。对 \(a\in L^\times\)，有 \(N_{L/K}(a)=1\) 当且仅当存在 \(b\in L^\times\) 使
\[
a=\frac{b}{\sigma(b)}.
\]
如果 \(b\) 是一个解，则全部解恰为 \(cb\)，其中 \(c\in K^\times\)。
\end{theorem}

为求解 \(a\sigma(b)=b\)，可尝试把 \(b\) 写成各 \(\sigma^i(c)\) 的加权和，并把首项权重取为 \(A_0=1\)。若使等式对每个 \(c\in L\) 成立，逐项配对得到递推 \(A_{i+1}=a\sigma(A_i)\)；递推至第 \(n\) 项时要求 \(A_n=A_0\)，恰是 \(N_{L/K}(a)=1\)。

\begin{proof}
由\eqref{b1:eq:cyclic-trace-norm}的相消计算，右边蕴含左边。现在设 \(N_{L/K}(a)=1\)，令
\[
A_0=1,\qquad A_i=a\cdot\sigma(a)\cdots\sigma^{i-1}(a)\quad(1\leq i\leq n).
\]
于是 \(A_n=1\)，且 \(a\cdot\sigma(A_i)=A_{i+1}\)。考察 \(K\) 线性映射
\[
T:L\longrightarrow L,\qquad
T(c)=\sum_{i=0}^{n-1}A_i\cdot\sigma^i(c).
\]
\cref{b1:cor:embedding-independence}说明不同嵌入 \(1,\sigma,\ldots,\sigma^{n-1}\) 作为 \(L\) 值函数线性无关。由于 \(A_0=1\)，这个映射不恒为零。选取 \(c\) 使 \(b=T(c)\neq0\)，便有
\[
a\cdot\sigma(b)=\sum_{i=0}^{n-1}A_{i+1}\cdot\sigma^{i+1}(c)
=\sum_{i=1}^{n-1}A_i\cdot\sigma^i(c)+c=b.
\]
因此 \(a=b/\sigma(b)\)。若 \(b'\) 也是解，则 \(\sigma(b'/b)=b'/b\)，故 \(b'/b\in L^{\langle\sigma\rangle}=K^\times\)；反之，乘以 \(K^\times\) 中的元素仍为解。
\end{proof}

对 \(\mathbb C/\mathbb R\)，若 \(a\bar a=1\) 且 \(a\neq-1\)，可直接取 \(b=1+a\)，因为 \(a\bar b=a+1=b\)。若 \(a=-1\)，这一选择成为零，应改取 \(b=i\)。证明中要求 \(T(c)\neq0\) 正是为了处理这种退化，不能随意把辅助元素固定为 \(1\)。

\subsection{Hilbert 第 90 定理的加法形式}
\label{b1:subsec:2002:02}

\begin{theorem}[Hilbert 第 90 定理，加法形式]\label{b1:thm:hilbert90-additive}
设 \(L/K\) 为 \(n\) 次循环扩张，并选取生成元 \(\sigma\in\operatorname{Gal}(L/K)\)。对 \(a\in L\)，有 \(\operatorname{Tr}_{L/K}(a)=0\) 当且仅当存在 \(b\in L\) 使 \(a=b-\sigma(b)\)。若 \(b\) 是一个解，则全部解组成陪集 \(b+K\)。
\end{theorem}
\begin{proof}
令 \(D:L\rightarrow L\) 为 \(K\) 线性映射 \(D(b)=b-\sigma(b)\)。其核为固定域 \(K\)，所以像的 \(K\) 维数为 \(n-1\)。另一方面，迹映射是不同嵌入之和；由\cref{b1:cor:embedding-independence}，这不是零映射。它的目标空间 \(K\) 是一维的，因此满射，核的维数也为 \(n-1\)。上一节已算出 \(\operatorname{Tr}_{L/K}\circ D=0\)，故 \(\operatorname{im}D\subseteq\ker\operatorname{Tr}_{L/K}\)。维数相等使这个包含成为等号。两个原像之差属于 \(\ker D=K\)，也就给出最后的断言。
\end{proof}

即使 \(\operatorname{char}K=p\) 且 \(p\mid n\)，迹仍然满射。此时 \(\operatorname{Tr}_{L/K}(1)=n\cdot1_K=0\) 只说明 \(1\) 落入迹的核，不能推出迹映射为零。这个区别在构造 Artin--Schreier 生成元时十分关键。

\subsection{第一群上同调的解释}
\label{b1:subsec:2002:03}

第四卷讨论群上同调时，将用\term[yici-yuquan]{一次余圈}[1-cocycle]与\term[yici-yubian]{一次余边}[1-coboundary]的语言解释这两个定理。这里仅说明有限循环 Galois 扩张中的相应方程，不要求预先学习上同调。

\ProseParagraphBreak
设 \(L/K\) 为有限 Galois 扩张，\(G=\operatorname{Gal}(L/K)\)。如果一族 \((c_g)_{g\in G}\) 中每个 \(c_g\in L^\times\)，并且满足
\[
c_{gh}=c_g\cdot g(c_h)\qquad(g,h\in G)
\]
则称它为乘法一次余圈。如果存在 \(b\in L^\times\)，使 \(c_g=b/g(b)\) 对每个 \(g\in G\) 都成立，则称这族元素为一次余边。当 \(G=\langle\sigma\rangle\)、\(|G|=n\) 时，余圈由 \(a=c_\sigma\) 决定：
\[
c_{\sigma^j}=a\cdot\sigma(a)\cdots\sigma^{j-1}(a)
\qquad(1\le j\le n).
\]
由于 \(\sigma^n=1\)，递推到 \(j=n\) 时必须有 \(N_{L/K}(a)=c_{\sigma^n}=1\)；反过来，这一条件也保证递推式与群关系相容。\cref{b1:thm:hilbert90-multiplicative}因而说明，在上述循环 Galois 群情形，每个乘法一次余圈都是一次余边。加法情形将上式改为 \(c_{gh}=c_g+g(c_h)\)，余边改为 \(b-g(b)\)；与 \(\sigma^n=1\) 相容的条件正是 \(\operatorname{Tr}_{L/K}(c_\sigma)=0\)。\cref{b1:thm:hilbert90-additive}同样说明，在此循环群情形，每个加法一次余圈都是一次余边。第四卷第十一章第 11.2 节将把一次余圈模一次余边所成的商组织成群，并重新解释这里的结论；本章的证明只使用域论。

% !TeX root = ../../Book1.tex
% 纲要标题：### 20.3 Kummer 理论
% 纲要位置：计划纲要.md，第 720 行。

\section{Kummer 理论}
\label{b1:sec:2003}

方程 \(X^n-a=0\) 看起来只涉及开方，但它的全部根还相差单位根。只有先确定基域中是否包含这些单位根，根式扩张与循环群之间的关系才足够直接。\term[Kummer-lilun]{Kummer 理论}[Kummer theory]研究这类根式与交换 Galois 群的联系。本节固定代数闭包 \(\overline K\)，所有扩张都取在其中；除另有说明外，\(n\geq1\)、\(\operatorname{char}K\nmid n\)，并假设 \(K\) 含有全部 \(n\) 次单位根。

% 纲要内容（仅作写作计划，尚非教材正文）：
%
% * 特征不整除 n 且基域含所需单位根的假设
% * 先以 Hilbert 第 90 定理证明循环扩张的根式生成，作为第21.2节的直接准备
% * 单根式扩张次数等于被开方数的商类阶；根式参数选择与 Kummer 配对的良定义
% * 有限 Kummer 对应的双侧单射、阶比较与同时特征基重建；素数次扩张由参数直线分类
% * 在本节定义 \(\mu_n\) 及本原单位根；循环扩张的根式生成不依赖一般 Kummer 对应，也不调用后续分圆多项式不可约性
%

\subsection{特征与单位根假设}
\label{b1:subsec:2003:01}

\begin{definition}\label{b1:def:roots-of-unity}
设 \(K\) 为域，\(\overline K\) 为 \(K\) 的代数闭包，\(n\geq1\) 为整数。如果 \(\zeta\in\overline K^\times\) 满足 \(\zeta^n=1\)，则称 \(\zeta\) 为 \(n\) 次\term[danwei-gen]{单位根}[root of unity]。全体 \(n\) 次单位根组成乘法群
\[
\mu_n\coloneqq\{\,\zeta\in\overline K^\times\mid\zeta^n=1\,\}.
\]
如果 \(\zeta\in\mu_n\) 的乘法阶恰为 \(n\)，则称 \(\zeta\) 为\term[benyuan-danwei-gen]{本原单位根}[primitive root of unity]。
\end{definition}
\symidx{\(\mu_n\)：全体 \(n\) 次单位根组成的群}

在 \(\operatorname{char}K\nmid n\) 的假设下，\(X^n-1\) 与导数 \(nX^{n-1}\) 没有公共根，所以 \(\mu_n\) 有 \(n\) 个元素。\cref{b1:lem:finite-multiplicative-cyclic}说明域的乘法群的有限子群是循环群；\(\mu_n\) 因而是阶为 \(n\) 的循环群，它的任一生成元 \(\zeta\) 都有阶 \(n\)，即为本原 \(n\) 次单位根，且 \(\mu_n=\langle\zeta\rangle\)。假设 \(\mu_n\subseteq K\) 就等价于 \(K\) 含有这样一个 \(\zeta\)。

\ProseParagraphBreak
对 \(a\in K^\times\)，特征条件使 \((X^n-a)'=nX^{n-1}\) 与 \(X^n-a\) 没有公共根，保证多项式可分；若去掉它，在特征 \(p\) 中就有 \(X^p-a=(X-\alpha)^p\)，其中 \(\alpha^p=a\)，只有一个重根，并且 \(X^p-1=(X-1)^p\) 表明不存在阶为 \(p\) 的单位根。单位根条件则保证：已有一个根 \(\alpha\) 后，其余根 \(\zeta\alpha\)（\(\zeta\in\mu_n\)）全部位于 \(K(\alpha)\) 中。若去掉它，\(\mathbb Q(\sqrt[3]{2})/\mathbb Q\) 就是可分而不正规的根式扩张。这两个假设分别排除了重根与遗漏其他共轭根的障碍。

\subsection{循环扩张与根式生成}
\label{b1:subsec:2003:04}

\begin{proposition}\label{b1:prop:single-radical-cyclic}
设 \(K\) 为域，\(n\geq1\) 为整数，\(\operatorname{char}K\nmid n\)，且 \(K\) 包含全部 \(n\) 次单位根。固定 \(K\) 的代数闭包 \(\overline K\)，取 \(a\in K^\times\) 及满足 \(\alpha^n=a\) 的 \(\alpha\in\overline K\)。则 \(K(\alpha)/K\) 是循环 Galois 扩张，其次数整除 \(n\)。映射
\[
\operatorname{Gal}\bigl(K(\alpha)/K\bigr)\longrightarrow\mu_n,
\qquad\sigma\longmapsto\frac{\sigma(\alpha)}{\alpha}
\]
是单同态。
\end{proposition}
\begin{proof}
取 \(\mu_n\) 的生成元 \(\zeta\)。\(X^n-a\) 的根恰为 \(\zeta^j\alpha\)，全在 \(K(\alpha)\) 中；因 \(a\neq0\) 且 \(n\neq0\) 在 \(K\) 中成立，它没有重根。所以 \(K(\alpha)/K\) 是一个可分多项式的分裂域，为有限 Galois 扩张。对其任一自同构 \(\sigma\)，
\[
\sigma(\alpha)^n=\sigma(\alpha^n)=\sigma(a)=a=\alpha^n,
\]
故 \((\sigma(\alpha)/\alpha)^n=1\)，比值属于 \(\mu_n\)。又因 \(\mu_n\subseteq K\)，每个自同构都固定这些比值，所以
\[
\frac{\sigma\tau(\alpha)}{\alpha}
=\sigma\Bigl(\frac{\tau(\alpha)}{\alpha}\Bigr)\frac{\sigma(\alpha)}{\alpha}
=\frac{\tau(\alpha)}{\alpha}\frac{\sigma(\alpha)}{\alpha}.
\]
目标群交换，故上式就是普通群同态的乘法公式。比值为 \(1\) 时，自同构固定生成元 \(\alpha\)，因而为恒等映射，这证明了单射性。Galois 群由此同构于 \(\mu_n\) 的一个子群，所以为循环群，群阶整除 \(|\mu_n|=n\)。由于扩张已经证明为 Galois，\([K(\alpha):K]=|\operatorname{Gal}(K(\alpha)/K)|\)，扩张次数也整除 \(n\)。
\end{proof}

这个同态计算再次用到了单位根条件。在一个更大的分裂域中，若自同构不固定单位根，则 \(\sigma(\tau(\alpha)/\alpha)\) 不能直接换成 \(\tau(\alpha)/\alpha\)，仍须保留自同构对该比值的作用；那正是上一节一次余圈的乘法公式。

\ProseParagraphBreak
根式扩张的次数还可以从被开方数在乘法商群中的阶直接读出。

\begin{corollary}\label{b1:cor:kummer-radical-degree}
设 \(K\) 为域，\(n\geq1\) 为整数，\(\operatorname{char}K\nmid n\)，且 \(K\) 包含全部 \(n\) 次单位根。固定 \(K\) 的代数闭包，取 \(a\in K^\times\)，并在该闭包中取满足 \(\alpha^n=a\) 的根 \(\alpha\)。记 \(K^{\times n}=\{\,c^n\mid c\in K^\times\,\}\)。则
\[
[K(\alpha):K]=\operatorname{ord}_{K^\times/K^{\times n}}([a]).
\]
\end{corollary}
\begin{proof}
记 \(L=K(\alpha)\)、\(d=[L:K]\)，并令 \(e\) 为 \([a]\) 的阶。因为 \([a]^n=[1]\)，\(e\) 是整除 \(n\) 的正整数。由\cref{b1:prop:single-radical-cyclic}，\(L/K\) 循环 Galois，其 Galois 群通过比值映射嵌入 \(\mu_n\)，像群的阶为 \(d\)。因此每个比值 \(\sigma(\alpha)/\alpha\) 的 \(d\) 次幂为一，故 \(\sigma(\alpha^d)=\alpha^d\) 对所有自同构成立。不动域为 \(K\)，所以 \(\alpha^d\in K^\times\)，进而 \(a^d=(\alpha^d)^n\in K^{\times n}\)，即 \(e\mid d\)。

反过来，\(a^e=c^n\) 对某个 \(c\in K^\times\) 成立，故 \((\alpha^e/c)^n=1\)。由 \(\mu_n\subseteq K\) 得 \(\alpha^e/c\in K\)，所以 \(\alpha^e\in K\)。于是 \(\alpha\) 满足 \(K\) 上的 \(e\) 次多项式 \(X^e-\alpha^e\)，其最小多项式次数 \(d\) 不超过 \(e\)。结合 \(e\mid d\) 与 \(d\le e\)，得到 \(d=e\)。
\end{proof}

反过来，循环扩张也可以用一个根式生成。

\begin{theorem}\label{b1:thm:kummer-cyclic}
设 \(L/K\) 是次数为 \(d\) 的循环扩张。如果 \(\operatorname{char}K\nmid d\)，并且 \(K\) 包含全部 \(d\) 次单位根 \(\mu_d\)，则存在 \(\alpha\in L^\times\) 使
\[
L=K(\alpha),\qquad\alpha^d\in K.
\]
\end{theorem}

为寻找这样的 \(\alpha\)，选 Galois 群的生成元 \(\sigma\) 和本原 \(d\) 次单位根 \(\zeta\)，要求 \(\sigma(\alpha)=\zeta\alpha\)，就是在线性代数意义下寻找 \(\sigma\) 的一个特征向量，特征值为 \(\zeta\)。这使 \(\sigma(\alpha^d)=\zeta^d\alpha^d=\alpha^d\)，又把构造化为 \(\alpha/\sigma(\alpha)=\zeta^{-1}\)，恰可调用乘法形式的 Hilbert 第 90 定理。

\begin{proof}
选取 \(\operatorname{Gal}(L/K)\) 的生成元 \(\sigma\)，以及本原 \(d\) 次单位根 \(\zeta\in K\)。因为 \(\zeta\) 属于基域，全部自同构都固定它，所以
\[
N_{L/K}(\zeta^{-1})=\prod_{j=0}^{d-1}\sigma^j(\zeta^{-1})
=\zeta^{-d}=1.
\]
\cref{b1:thm:hilbert90-multiplicative}给出非零 \(\alpha\) 满足 \(\alpha/\sigma(\alpha)=\zeta^{-1}\)，即 \(\sigma(\alpha)=\zeta\alpha\)。所以 \(\alpha^d\) 被 \(\sigma\) 固定，属于 \(K\)。又有 \(\sigma^j(\alpha)=\zeta^j\alpha\) 对 \(0\le j<d\) 成立。\(\sigma^j\) 固定最小多项式 \(m_{\alpha,K}\) 的系数，故这些像全部是它的根；因为 \(\zeta\) 的阶为 \(d\) 且 \(\alpha\ne0\)，它们两两不同。因此 \(\deg m_{\alpha,K}\ge d\)。另一方面，\(K(\alpha)\subseteq L\) 给出 \(\deg m_{\alpha,K}=[K(\alpha):K]\le[L:K]=d\)，故次数相等，\(K(\alpha)=L\)。
\end{proof}

\cref{b1:thm:kummer-cyclic}将用于\secref{b1:sec:2102}的特征零根式可解性判据。多个根式之间的关系还可由有限交换 Galois 群统一描述，下面据此建立 Kummer 配对。

\subsection{根式扩张与 Kummer 配对}
\label{b1:subsec:2003:02}

一个根式的自同构由比值 \(\sigma(\alpha)/\alpha\) 决定。加入多个根式时，我们希望同时记录这些比值，并判断哪些根式增加了扩张次数。若 \(a'=ab^n\)、\(b\in K^\times\)，而 \(\alpha^n=a\)，则 \((b\alpha)^n=a'\)，且 \(K(b\alpha)=K(\alpha)\)。因此应先把被开方数放入商群 \(K^\times/K^{\times n}\)，再用同一自同构在各个根上的比值来比较它们。

\ProseParagraphBreak
为同时处理若干根式，记 \(K^{\times n}=\{\,b^n\mid b\in K^\times\,\}\)。商群 \(K^\times/K^{\times n}\) 中每个元素的阶都整除 \(n\)。取其有限子群 \(B\)，选择生成元的代表 \(a_1,\ldots,a_r\)，并令
\[
L_B=K(\alpha_1,\ldots,\alpha_r),\qquad\alpha_i^n=a_i.
\]
换一个根只会乘以 \(\mu_n\subseteq K\) 中的元素；换一个代表只会再乘以 \(K^\times\) 中的元素。若 \([a]\in B\)，可写 \([a]=[a_1]^{e_1}\cdots[a_r]^{e_r}\)，其中 \(e_i\in\mathbb Z\)。这意味着对某个 \(c\in K^\times\) 有
\[
a=c^n a_1^{e_1}\cdots a_r^{e_r},\qquad
\bigl(c\alpha_1^{e_1}\cdots\alpha_r^{e_r}\bigr)^n=a.
\]
所有 \(\alpha_i\) 均非零，所以负指数也有意义；\(a\) 的任一 \(n\) 次根只与所列的根相差 \(K\) 中的单位根，故都属于 \(L_B\)。换一组生成元时，新生成元代表的根因而全部属于原来的 \(L_B\)，反过来原生成元的根也属于新域，所以两域相等。这说明 \(L_B\) 不依赖代表、根或生成列表的选择。它也是多项式 \(X^n-a_1,\ldots,X^n-a_r\) 的共同分裂域；这些多项式都可分，故 \(L_B/K\) 是有限 Galois 扩张。记 \(G_B=\operatorname{Gal}(L_B/K)\)。

\begin{definition}\label{b1:def:kummer-pairing}
设 \(K\) 为域，\(n\geq1\) 且 \(\operatorname{char}K\nmid n\)，在 \(K\) 的固定代数闭包 \(\overline K\) 中记 \(\mu_n=\{\zeta\in\overline K^\times\mid\zeta^n=1\}\)，并假设 \(\mu_n\subseteq K\)。取乘法商群 \(K^\times/K^{\times n}\) 的有限子群 \(B\)，其中 \(K^{\times n}=\{b^n\mid b\in K^\times\}\)；令 \(L_B\) 为在 \(\overline K\) 中给 \(B\) 的生成元代表添入 \(n\) 次根所得的域，并记 \(G_B=\operatorname{Gal}(L_B/K)\)。对 \(\sigma\in G_B\) 与 \(\xi=[a]\in B\)，取 \(a\in K^\times\) 及满足 \(\alpha^n=a\) 的 \(\alpha\in L_B\)，定义
\[
\langle\sigma,\xi\rangle\coloneqq\frac{\sigma(\alpha)}{\alpha},
\qquad \alpha^n=a.
\]
该值与代表 \(a\) 及根 \(\alpha\) 的选择无关。由此得到的映射 \(G_B\times B\rightarrow\mu_n\) 称为\term[Kummer-peidui]{Kummer 配对}[Kummer pairing]。
\end{definition}

配对的值与根和代表的选择无关。若把 \(\alpha\) 换成 \(\alpha'=\zeta\alpha\)，其中 \(\zeta\in\mu_n\subseteq K\)，则
\[
\frac{\sigma(\alpha')}{\alpha'}=
\frac{\zeta\sigma(\alpha)}{\zeta\alpha}=
\frac{\sigma(\alpha)}{\alpha}.
\]
若把代表换成 \(a'=ab^n\)，其中 \(b\in K^\times\)，可取根 \(b\alpha\)，同样有 \(\sigma(b\alpha)/(b\alpha)=\sigma(\alpha)/\alpha\)。该代表的其他根再由前一种计算处理，故两个选择都不影响配对值。

\ProseParagraphBreak
根的乘积是乘积的根，所以配对在第二个变量上保持乘法；第一个变量上的计算仍须使用 \(\sigma\) 固定 \(\mu_n\) 这一事实，正如\cref{b1:prop:single-radical-cyclic}证明中的比值公式。于是对 \(\sigma,\tau\in G_B\)、\(\xi,\eta\in B\)，有
\[
\begin{aligned}
\langle\sigma\tau,\xi\rangle&=
\langle\sigma,\xi\rangle\langle\tau,\xi\rangle,\\
\langle\sigma,\xi\eta\rangle&=
\langle\sigma,\xi\rangle\langle\sigma,\eta\rangle.
\end{aligned}
\]
因此固定任意一个变量，就得到另一个变量上的群同态。

\subsection{有限 Kummer 对应}
\label{b1:subsec:2003:03}

在 \(\operatorname{char}K\nmid n\) 的假设下，\(\mu_n\) 是阶为 \(n\) 的循环群。设 \(A\) 为有限交换群。对群同态 \(\chi,\psi:A\rightarrow\mu_n\)，\(\operatorname{Hom}(A,\mu_n)\) 的运算是逐点相乘：\((\chi\psi)(x)=\chi(x)\psi(x)\)，单位元是常值 \(1\) 的同态。\(A\) 的指数指使每个元素的该次幂等于 \(1\) 的最小正整数。若其指数整除 \(n\)，则
\begin{equation}\label{b1:eq:finite-character-count}
\bigl|\operatorname{Hom}(A,\mu_n)\bigr|=|A|.
\end{equation}
这里应用\cref{b1:thm:fga-structure}，把有限交换群 \(A\) 写成若干循环群的直积；因 \(A\) 的指数整除 \(n\)，每个循环因子的阶 \(d\) 也整除 \(n\)。从 \(C_d\) 到 \(\mu_n\) 的同态由生成元的像决定，这个像必须满足 \(\zeta^d=1\)。在阶为 \(n\) 的循环群中，满足这个条件的元素有 \(\gcd(d,n)\) 个；只有 \(d\mid n\) 时才恰有 \(d\) 个。当前每个循环因子均满足此条件，各因子的选择独立，相乘得到所需等式。这说明指数条件不能省去；单凭 \(A\) 有限交换，并不能保证有 \(|A|\) 个这样的同态。

\begin{theorem}[Kummer 对应]\label{b1:thm:kummer-correspondence}
设 \(K\) 为域，\(n\geq1\) 为整数，\(\operatorname{char}K\nmid n\)，并且 \(K\) 包含全部 \(n\) 次单位根 \(\mu_n\)。固定 \(K\) 的代数闭包 \(\overline K\)，记 \(K^{\times n}=\{c^n\mid c\in K^\times\}\)。以下两类对象一一对应：
\begin{enumerate}
\item \(K^\times/K^{\times n}\) 的有限子群 \(B\)；
\item \(\overline K\) 中的有限 Galois 扩张 \(L/K\)，其 Galois 群为指数整除 \(n\) 的交换群。
\end{enumerate}
对第一类中的有限子群 \(B\)，令 \(L_B\subseteq\overline K\) 为给 \(B\) 中各元素的代表添入 \(n\) 次根所生成的域。正向对应为 \(B\mapsto L_B\)，反向对应为
\[
B_L=\left(K^\times\cap L^{\times n}\right)/K^{\times n}.
\]
这里 \(L^{\times n}=\{\,\beta^n\mid\beta\in L^\times\,\}\)，所以分子收集 \(K\) 中那些在 \(L\) 内已经具有 \(n\) 次根的非零元素，再模掉它们在 \(K\) 中本来就是 \(n\) 次幂的平凡类。
Kummer 配对给出同构
\[
\operatorname{Gal}(L_B/K)\cong\operatorname{Hom}(B,\mu_n),
\qquad B\cong\operatorname{Hom}\bigl(\operatorname{Gal}(L_B/K),\mu_n\bigr),
\]
特别地，\([L_B:K]=|B|\)。子群包含对应于域包含。
\end{theorem}
\begin{proof}
从有限 \(B\) 出发，配对给出群同态
\[
G_B\longrightarrow\operatorname{Hom}(B,\mu_n),\qquad
\sigma\longmapsto\bigl(\xi\mapsto\langle\sigma,\xi\rangle\bigr).
\]
若 \(\sigma\) 在其核中，则全部选定根的比值 \(\sigma(\alpha_i)/\alpha_i\) 均为一。因此 \(\sigma\) 固定生成域的每个根 \(\alpha_i\)，也固定由它们和 \(K\) 生成的整个 \(L_B\)，故为恒等。这个映射因而单射。目标是逐点相乘的交换群，每个同态的 \(n\) 次幂都是常值一的同态，所以这个单射还说明 \(G_B\) 交换且指数整除 \(n\)。

在另一侧，配对给出 \(B\rightarrow\operatorname{Hom}(G_B,\mu_n)\)。若 \([a]\) 对每个 \(\sigma\in G_B\) 的配对值都是一，则其根 \(\alpha\) 被整个群固定。不动域为 \(K\)，所以 \(\alpha\in K\)，\(a=\alpha^n\) 在 \(K\) 中本来就是 \(n\) 次幂，即 \([a]=[1]\)。这也得到一个单射：任何非平凡被开方数类都能被某个自同构通过配对区别出来。

有限交换群 \(B\) 的指数整除 \(n\)，所以\eqref{b1:eq:finite-character-count}给出 \(|\operatorname{Hom}(B,\mu_n)|=|B|\)，第一个单射推出 \(|G_B|\le|B|\)。刚才已从这个单射得到 \(G_B\) 的交换性与指数条件，因而同一个计数公式也适用于 \(G_B\)，得到 \(|\operatorname{Hom}(G_B,\mu_n)|=|G_B|\)。第二个单射于是推出 \(|B|\le|G_B|\)。两边的有限阶相等，两个单射都为满射，故都是同构。又因 \(L_B/K\) 为 Galois，\([L_B:K]=|G_B|=|B|\)。

现在从第二类扩张 \(L/K\) 出发，记群为 \(G\)。要将它重建为一个根式扩张，只需找到一组 \(K\) 基 \(\beta_i\)，使每个 \(\beta_i^n\) 都在 \(K\) 中。把各 \(\sigma\in G\) 视为 \(K\) 向量空间 \(L\) 上的线性算子，寻找它们的共同特征向量，就能使这些 \(n\) 次幂被整个群固定。由于群的指数整除 \(n\)，各算子满足 \(\sigma^n=1\)；\(X^n-1\) 在 \(K\) 中分裂且无重根，可以用插值多项式把向量分解到其特征空间中。对每个 \(\zeta\in\mu_n\)，令
\[
P_\zeta(X)=\prod_{\substack{\eta\in\mu_n\\\eta\neq\zeta}}
\frac{X-\eta}{\zeta-\eta}.
\]
这就是\cref{b1:thm:lagrange-interpolation}在互异插值点 \(\mu_n\) 上的基多项式，满足 \(P_\zeta(\eta)=1\) 当 \(\eta=\zeta\)，否则为零。因而 \(\sum_\zeta P_\zeta(X)-1\) 是次数小于 \(n\) 却有 \(n\) 个不同根的多项式，只能为零。又有
\[
(X-\zeta)P_\zeta(X)=
\frac{X^n-1}{\prod_{\eta\in\mu_n,\,\eta\ne\zeta}(\zeta-\eta)},
\]
因为分子恰是全部 \(n\) 个一次因子的乘积，分母是非零的基域常数。代入算子 \(\sigma\)，利用 \(\sigma^n=1\)，得到 \((\sigma-\zeta I)P_\zeta(\sigma)=0\)。所以
\[
v=\sum_{\zeta\in\mu_n}P_\zeta(\sigma)v
\]
的每一项都属于 \(\sigma\) 的 \(\zeta\) 特征空间。\(P_\zeta(\sigma)\) 在该空间上为恒等，在其他特征空间上为零；若这些空间中向量之和为零，分别作用各 \(P_\zeta(\sigma)\) 就得到每项都为零。因此这个分解是直和，\(\sigma\) 可对角化。

要得到整个 \(G\) 的同时特征基，还需要群交换。若 \(\sigma(v)=\zeta v\)，则对任意 \(\tau\in G\)，
\[
\sigma(\tau(v))=\tau(\sigma(v))=\tau(\zeta v)=\zeta\tau(v),
\]
因为两算子交换且 \(\zeta\in K\)。因此一个算子的每个特征空间都被其余算子保持，可以在其中对下一算子再作上述直和分解。群 \(G\) 有限，依次处理其全部元素后，便得到 \(L\) 的公共特征空间直和分解；在每个非零公共特征空间中选一组基，合在一起得到同时特征向量基 \(\beta_1,\ldots,\beta_m\)。这里指数条件用于单个算子的分解，交换性用于使这些分解相容。

写 \(\sigma(\beta_i)=\lambda_{\sigma,i}\beta_i\)，其中 \(\lambda_{\sigma,i}\in K\)。因为 \(\beta_i\ne0\)，
\[
\beta_i=\sigma^n(\beta_i)=\lambda_{\sigma,i}^n\beta_i
\]
给出 \(\lambda_{\sigma,i}^n=1\)，即 \(\lambda_{\sigma,i}=\sigma(\beta_i)/\beta_i\in\mu_n\)。于是 \(\sigma(\beta_i^n)=\beta_i^n\) 对所有 \(\sigma\in G\) 成立，故 \(\beta_i^n\in K^\times\)。这些 \(n\) 次幂的类生成一个有限子群 \(B_0\)：各生成元的幂指数都可约为 \(0,\ldots,n-1\)，所以其元素数至多为 \(n^m\)。构造 \(L_{B_0}\) 时可取 \(\beta_i\) 作为这些类的根，因此 \(L_{B_0}\subseteq L\)。另一方面，这个子域含有 \(K\) 和全部 \(\beta_i\)，故含有它们的所有 \(K\) 线性组合，即整个 \(L\)。两边包含给出 \(L_{B_0}=L\)。

\(B_0\) 是利用一次基的选择构造出来的参数群；还须确认它已经包含 \(L\) 中所有被开方数类，即等于由扩张本身定义的 \(B_L\)。对 \(B_L\) 中任意元素，在 \(L\) 中选取其代表的 \(n\) 次根。以同一比值公式得到群同态 \(B_L\rightarrow\operatorname{Hom}(G,\mu_n)\)；若某个类与所有 \(\sigma\in G\) 的配对值均为一，其根便属于不动域 \(K\)，所以该映射单射。目标是有限群，因而 \(B_L\) 有限。由构造可知 \(B_0\subseteq B_L\)；而 \(B_L\) 中各类在 \(L\) 中有根，另选的根只差 \(K\) 中的单位根，因此 \(L_{B_L}\subseteq L\)。于是
\[
L=L_{B_0}\subseteq L_{B_L}\subseteq L.
\]
故 \(L_{B_L}=L\)。对两个有限子群分别应用已证的次数公式，得到 \(|B_L|=[L:K]=|B_0|\)。因为 \(B_0\subseteq B_L\)，有限群的阶相等迫使 \(B_L=B_0\)，所以选取的这组基虽然不唯一，所得到的参数群却是扩张内在确定的全部参数群。

反过来，从任意有限 \(B\) 出发，由各代表在 \(L_B\) 中具有根可知 \(B\subseteq B_{L_B}\)。刚才对第二类扩张的证明应用于 \(L_B/K\)，给出 \(B_{L_B}\) 有限且 \(L_{B_{L_B}}=L_B\)。因此 \(|B_{L_B}|=[L_B:K]=|B|\)，由包含与阶相等得 \(B_{L_B}=B\)。这证明了两个构造互逆。若 \(B\subseteq B'\)，则其代表的根也包含于 \(L_{B'}\)，从而 \(L_B\subseteq L_{B'}\)；反之，若 \(L_B\subseteq L_{B'}\)，一个在前者中已有 \(n\) 次根的基域元素在后者中也有根，故 \(B=B_{L_B}\subseteq B_{L_{B'}}=B'\)。
\end{proof}

这里参数子群 \(B\) 增大，所要添入的根也增多，所以 Kummer 对应中的子群包含与域包含同向。上一章的 Galois 对应则反向：自同构子群越大，共同不动元需要满足的条件越多，不动域越小。两种对应中的子群分别记录被开方数类与自同构，含义不同。

\ProseParagraphBreak
定理要求扩张有限，且 Galois 群的指数整除给定的 \(n\)。与之对应的 \(B\) 是有限子群，但不要求它在整个 \(K^\times/K^{\times n}\) 中具有有限指数。例如取 \(K=\mathbb Q\)、\(n=2\)，则 \(\mathbb Q^\times/\mathbb Q^{\times2}\) 是无限群，但 \(B=\{[1],[2],[3],[6]\}\) 有四个元素，给出 \(L_B=\mathbb Q(\sqrt2,\sqrt3)\)。独立性可由 \(2,3,6\) 都不是有理数的平方核验。

\ProseParagraphBreak
令 \(s\) 只改变 \(\sqrt2\) 的符号，\(t\) 只改变 \(\sqrt3\) 的符号。此时 \(G_B=\{1,s,t,st\}\)、\(\mu_2=\{1,-1\}\)，配对的全部取值为
\[
\begin{array}{c|rrrr}
\langle\sigma,\xi\rangle &[1]&[2]&[3]&[6]\\ \hline
1  & 1& 1& 1& 1\\
s  & 1&-1& 1&-1\\
t  & 1& 1&-1&-1\\
st & 1&-1&-1& 1
\end{array}
\]
每一行是 \(B\rightarrow\mu_2\) 的同态，每一列是 \(G_B\rightarrow\mu_2\) 的同态；各行两两不同，各列也两两不同，具体显示了配对的两侧非退化性。特别地，\([6]\) 的核为 \(\{1,st\}\)，所以 \(\mathbb Q(\sqrt6)=L_B^{\langle st\rangle}\)。这里 \(\{[1],[6]\}\subseteq B\) 给出子域包含 \(\mathbb Q(\sqrt6)\subseteq L_B\)，而子域与其固定子群之间的包含方向相反。

\ProseParagraphBreak
设 \(K\) 为域，\(\ell\) 为素数。乘法商群 \(V=K^\times/K^{\times\ell}\) 的每个元素的阶都整除 \(\ell\)，因而可以把 \(V\) 视为 \(\mathbb F_\ell\) 向量空间；\(V\) 也可能是平凡群。若用加法记号，其运算为
\[
[a]+[b]=[ab],\qquad \bar r\,[a]=[a^r],\qquad 0=[1]
\quad(a,b\in K^\times,\ r\in\mathbb Z).
\]
其中 \(\bar r\) 是整数 \(r\) 在 \(\mathbb F_\ell\) 中的像。因为 \([a]^\ell=[1]\)，若 \(r'\equiv r\pmod\ell\)，则 \([a^{r'}]=[a^r]\)，所以标量乘法只依赖 \(\bar r\)，不依赖其整数代表。分配律与标量乘法的结合律来自交换群的幂运算规则，例如 \((ab)^r=a^rb^r\)、\(a^{r+s}=a^ra^s\)。这里的非零向量是 \([a]\ne[1]\)，即 \(a\in K^\times\setminus K^{\times\ell}\)；只写 \(a\ne0\) 并不足以保证这一点。


\begin{corollary}\label{b1:cor:kummer-prime-classification}
设 \(K\) 为域，\(\ell\) 为不同于 \(\operatorname{char}K\) 的素数，且 \(K\) 包含全部 \(\ell\) 次单位根 \(\mu_\ell\)。在 \(K\) 的一个固定代数闭包中，\(K\) 的 \(\ell\) 次循环扩张由 \(K^\times/K^{\times\ell}\) 的一维 \(\mathbb F_\ell\) 子空间分类。
\end{corollary}
\begin{proof}
素数情形应用\cref{b1:thm:kummer-correspondence}：次数为 \(\ell\) 的域对应阶为 \(\ell\) 的子群。上述乘法商群的指数整除 \(\ell\)，它作为 \(\mathbb F_\ell\) 向量空间时，阶为 \(\ell\) 的子群正是一维子空间。
\end{proof}

因此，若 \(a\in K^\times\setminus K^{\times\ell}\)，同一直线上的 \(\ell-1\) 个非零向量 \([a]^r\)（\(1\le r<\ell\)）都生成同一个参数子群 \(\langle[a]\rangle\)，给出同一个非平凡 \(\ell\) 次扩张。具体地，若 \(\alpha^\ell=a\)、\(b\in K^\times\)，则 \(b\alpha^r\) 的 \(\ell\) 次幂为 \(a^r b^\ell\)。由于 \(r\) 与 \(\ell\) 互素，取整数 \(u,v\) 使 \(ru+\ell v=1\)，就有
\[
\alpha=(\alpha^r)^u a^v\in K(b\alpha^r).
\]
反向包含显然，所以 \(K(b\alpha^r)=K(\alpha)\)。以每个非零类分别标记扩张，会将这一条直线重复计算 \(\ell-1\) 次；一维子空间恰好把这些生成元变化归为同一个参数。这个区别与下一节的加法分类相对应。

% !TeX root = ../../Book1.tex
% 纲要标题：### 20.4 Artin–Schreier 理论
% 纲要位置：计划纲要.md，第 728 行。

\section{Artin–Schreier 理论}
\label{b1:sec:2004}

在特征 \(p\) 中，\(p\) 次开方通常产生纯不可分扩张，不能代替 \(p\) 次循环 Galois 扩张。一个导数恒为 \(-1\) 的多项式提供了正确的可分模型。以下 \(K\) 始终为特征 \(p>0\) 的域，扩张取在一个固定代数闭包中。

% 纲要内容（仅作写作计划，尚非教材正文）：
%
% * 特征 p 中的方程 X^p−X−a；素域线性、根的平移与不可约二分
% * p 次循环扩张的构造与参数直线分类；用有理函数次数解释无穷远极点
% * 与 Kummer 理论的平行及差别；区分纯不可分开方与可分循环步骤
%
% ---
%

\subsection{Artin–Schreier 方程}
\label{b1:subsec:2004:01}

\begin{definition}\label{b1:def:artin-schreier-map}
设 \(K\) 为特征 \(p>0\) 的域。映射
\[
\wp:K\longrightarrow K,\qquad b\longmapsto b^p-b
\]
称为\term[Artin-Schreier-yingshe]{Artin--Schreier 映射}[Artin--Schreier map]。给定 \(a\in K\)，把方程 \(X^p-X-a=0\) 称为相应的\term[Artin-Schreier-fangcheng]{Artin--Schreier 方程}[Artin--Schreier equation]。
\end{definition}
\symidx{\(\wp(b)=b^p-b\)：Artin--Schreier 映射}

在特征 \(p\) 中，二项式公式给出 \((u+v)^p=u^p+v^p\)。又因 \(c^p=c\) 对每个 \(c\in\mathbb F_p\) 成立，有
\[
\wp(u+v)=\wp(u)+\wp(v),\qquad
\wp(cu)=c^pu^p-cu=c\wp(u).
\]
这证明了 \(\wp\) 的 \(\mathbb F_p\) 线性性。

\ProseParagraphBreak
素域 \(\mathbb F_p\) 的 \(p\) 个元素都是 \(X^p-X\) 的根；这个多项式的次数为 \(p\)，在任何扩域中都没有其他根。因此 \(\wp(b)=0\) 当且仅当 \(b\in\mathbb F_p\)，即 \(\ker\wp=\mathbb F_p\)。若 \(K\neq\mathbb F_p\)，取 \(d\in K\setminus\mathbb F_p\)，则 \(\wp(d\cdot1)=d^p-d\neq0=d\wp(1)\)，所以它并不是 \(K\) 线性映射。像 \(\wp(K)\) 是 \(K\) 的加法群中的 \(\mathbb F_p\) 线性子空间，所以 \(K/\wp(K)\) 是 \(\mathbb F_p\) 向量空间。后文的一维子空间正是按这个线性结构理解的；这里没有对乘法取商，也不把它看成商域。

\begin{proposition}\label{b1:prop:artin-schreier-polynomial}
设 \(K\) 为特征 \(p>0\) 的域，\(a\in K\)，记 \(\wp(K)=\{b^p-b\mid b\in K\}\)。则 \(X^p-X-a\) 要么在 \(K\) 上完全分裂，要么在 \(K[X]\) 中不可约。前者当且仅当 \(a\in\wp(K)\)；在后一种情形，该多项式的任一根 \(\alpha\) 都生成 \(p\) 次循环扩张 \(K(\alpha)/K\)，其全部 \(K\) 自同构为 \(\alpha\mapsto\alpha+c\)，其中 \(c\in\mathbb F_p\)。
\end{proposition}
\begin{proof}
若 \(\alpha,\beta\) 都是 \(X^p-X-a\) 的根，则
\[
(\beta-\alpha)^p-(\beta-\alpha)
=(\beta^p-\beta)-(\alpha^p-\alpha)=0.
\]
由 \(X^p-X\) 的根的描述，\(\beta-\alpha\in\mathbb F_p\)。反过来，每个 \(\alpha+c\)、\(c\in\mathbb F_p\)，都满足原方程。因此全部根恰为这 \(p\) 个不同的元素。它们都在 \(L=K(\alpha)\) 中，所以 \(L\) 是该多项式的分裂域；导数为 \(-1\)，又保证它没有重根。由可分多项式的分裂域判据，\(L/K\) 是有限 Galois 扩张，故\cref{b1:prop:galois-automorphism-count}给出 \(|\operatorname{Gal}(L/K)|=[L:K]\)。

对 \(\sigma\in\operatorname{Gal}(L/K)\)，令 \(c_\sigma=\sigma(\alpha)-\alpha\)。因为 \(\sigma(\alpha)\) 也是原方程的根，所以 \(c_\sigma\in\mathbb F_p\)。每个 \(K\) 自同构都固定素域，故
\[
\begin{aligned}
c_{\sigma\tau}
&=\sigma\tau(\alpha)-\alpha\\
&=\sigma\bigl(\tau(\alpha)-\alpha\bigr)+\sigma(\alpha)-\alpha\\
&=c_\tau+c_\sigma.
\end{aligned}
\]
这说明 \(\sigma\mapsto c_\sigma\) 是到 \((\mathbb F_p,+)\) 的群同态。若 \(c_\sigma=0\)，则 \(\sigma\) 固定生成元 \(\alpha\)，因而是恒等映射，所以这个同态是单射。目标群阶为素数，只有平凡子群和整个群，于是 \([L:K]\) 只能为 \(1\) 或 \(p\)。

次数为 \(1\) 时 \(\alpha\in K\)，所有根都在 \(K\) 中，且 \(a=\wp(\alpha)\)。反过来，若 \(a=\wp(b)\)、\(b\in K\)，全部根就是 \(b+c\)、\(c\in\mathbb F_p\)，故多项式完全分裂。次数为 \(p\) 时，\(\alpha\) 的最小多项式次数为 \(p\)，且整除首一 \(p\) 次多项式 \(X^p-X-a\)，因而二者相等，给定多项式不可约。此时上述单同态的像是整个 \((\mathbb F_p,+)\)，所以 Galois 群为循环群，并且每个平移 \(\alpha\mapsto\alpha+c\) 都对应一个 \(K\) 自同构。
\end{proof}

\subsection{素数次循环扩张的分类}
\label{b1:subsec:2004:02}

\term[Artin-Schreier-dingli]{Artin--Schreier 定理}[Artin--Schreier theorem]说明上一小节的方程穷尽了特征 \(p\) 中的循环 \(p\) 次扩张，并给出避免重复的参数空间。

\begin{theorem}[Artin--Schreier]\label{b1:thm:artin-schreier-classification}
设 \(K\) 为特征 \(p>0\) 的域，并固定 \(K\) 的代数闭包 \(\overline K\)，记 \(\wp(b)=b^p-b\)（\(b\in K\)）。每个包含于 \(\overline K\) 的 \(p\) 次循环扩张 \(L/K\) 都形如
\[
L=K(\alpha),\qquad\alpha^p-\alpha=a\in K\setminus\wp(K).
\]
这些位于固定代数闭包中的子扩张与 \(\mathbb F_p\) 向量空间 \(K/\wp(K)\) 的一维子空间一一对应。具体地，对 \(a,a'\in K\setminus\wp(K)\)，分别在 \(\overline K\) 中选取满足 \(\alpha^p-\alpha=a\)、\(\beta^p-\beta=a'\) 的根，则 \(K(\alpha)=K(\beta)\) 当且仅当
\[
a'=ca+\wp(b)\qquad\text{对某个 }c\in\mathbb F_p^\times,\ b\in K.
\]
\end{theorem}

在\cref{b1:thm:kummer-cyclic}的构造中，为得到 \(\sigma(\alpha)=\zeta\alpha\)，把 \(\dfrac{\alpha}{\sigma(\alpha)}=\zeta^{-1}\) 交给乘法形式的 Hilbert 第 90 定理。这里设 \(\sigma\) 是 \(p\) 次循环扩张 \(L/K\) 的生成自同构，改为要求 \(\sigma(\alpha)=\alpha+1\)，也就是 \(\alpha-\sigma(\alpha)=-1\)。这样 \(\alpha^p-\alpha\) 会被 \(\sigma\) 固定，而 \(-1\) 的迹为零，加法形式的 Hilbert 第 90 定理便能给出所需的 \(\alpha\)。

\begin{proof}
设 \(L/K\) 是 \(p\) 次循环扩张，生成自同构为 \(\sigma\)。由于扩张次数和基域特征都为 \(p\)，有 \(\operatorname{Tr}_{L/K}(-1)=-p\cdot1_K=0\)。\cref{b1:thm:hilbert90-additive}不要求扩张次数在 \(K\) 中可逆，其证明没有除以 \(p\)，所以这里仍可应用。它给出 \(\alpha\in L\) 满足 \(\alpha-\sigma(\alpha)=-1\)，即 \(\sigma(\alpha)=\alpha+1\)。于是
\[
\sigma(\alpha^p-\alpha)
=(\alpha+1)^p-(\alpha+1)
=\alpha^p+1-\alpha-1
=\alpha^p-\alpha.
\]
因为 \(\sigma\) 生成整个 Galois 群，这个元素属于固定域 \(K\)，记为 \(a\)。又因 \(\sigma(\alpha)\neq\alpha\)，有 \(\alpha\notin K\)。\cref{b1:prop:artin-schreier-polynomial}因而给出 \(a\notin\wp(K)\) 和 \([K(\alpha):K]=p\)。结合 \(K(\alpha)\subseteq L\) 及 \([L:K]=p\)，得到 \(K(\alpha)=L\)。

若 \(a'=ca+\wp(b)\)，其中 \(c\in\mathbb F_p^\times\)、\(b\in K\)，令 \(\gamma=c\alpha+b\)。由 \(c^p=c\) 得
\[
\gamma^p-\gamma=c(\alpha^p-\alpha)+(b^p-b)=a'.
\]
又因 \(\alpha=c^{-1}(\gamma-b)\)，有 \(K(\gamma)=K(\alpha)\)。两根 \(\beta,\gamma\) 相差 \(\mathbb F_p\) 中的元素，所以 \(K(\beta)=K(\gamma)=K(\alpha)\)。

反过来，设 \(K(\alpha)=K(\beta)=L\)。由\cref{b1:prop:artin-schreier-polynomial}，映射 \(\operatorname{Gal}(L/K)\to(\mathbb F_p,+)\)、\(\tau\mapsto\tau(\alpha)-\alpha\) 是同构。因此可选取 \(\sigma\) 使 \(\sigma(\alpha)=\alpha+1\)；它对应非零元素 \(1\)，故是 Galois 群的生成元。因为 \(a'\in K\)，有
\[
\sigma(\beta)^p-\sigma(\beta)=\sigma(\beta^p-\beta)=a'.
\]
这说明 \(\beta\) 与 \(\sigma(\beta)\) 是同一个 Artin--Schreier 方程的根，于是 \(\sigma(\beta)-\beta=c\in\mathbb F_p\)。若 \(c=0\)，则 \(\beta\) 被生成元 \(\sigma\) 固定，从而 \(\beta\in K\)，与 \(a'\notin\wp(K)\) 矛盾。因此 \(c\neq0\)，且
\[
\sigma(\beta-c\alpha)=\beta+c-c(\alpha+1)=\beta-c\alpha.
\]
故 \(\beta-c\alpha\) 等于某个 \(b\in K\)。代入 \(\beta=c\alpha+b\) 就得到 \(a'=ca+\wp(b)\)。

在商空间 \(K/\wp(K)\) 中，这个参数条件等价于 \([a']=c[a]\)，其中 \(c\in\mathbb F_p^\times\)。因此同一个扩张对应同一个一维子空间中的所有非零向量，而 \(\wp(b)\) 只改变向量的代表。每个非零类 \([a]\) 又由前一命题给出一个 \(p\) 次循环扩张，所以上述等价条件恰好给出所述的一一对应。
\end{proof}

例如在 \(K=\mathbb F_p(t)\) 中，\(X^p-X-t\) 不可约。由\cref{b1:prop:artin-schreier-polynomial}，只需证明 \(t\notin\wp(K)\)。若 \(b=0\)，显然 \(b^p-b\neq t\)；否则把 \(b\) 写成既约分式 \(\dfrac{u(t)}{v(t)}\)，其中 \(u,v\in\mathbb F_p[t]\) 都非零，令 \(m=\deg u-\deg v\)。此时
\[
b^p-b=\frac{u^p-uv^{p-1}}{v^p}.
\]
若 \(m>0\)，则 \(p\deg u>\deg u+(p-1)\deg v\)，所以分子的两项最高次项不会相消，分子次数为 \(p\deg u\)，分母次数为 \(p\deg v\)，次数差为 \(pm\)。约去公因子会使分子、分母次数减少同一个数，不改变这个差；而 \(t\) 的次数差为 \(1\)，不可能等于 \(pm\)。若 \(m\le0\)，分子次数不超过分母次数，或者分子为零，也不可能表示 \(t\)。这就证明了所需结论。通常把非零有理函数 \(\dfrac{u(t)}{v(t)}\) 的 \(\max(\deg u-\deg v,0)\) 称为它在无穷远处的极点阶；上述计算也说明，\(b\) 有正极点阶 \(m\) 时，\(b^p-b\) 的极点阶恰为 \(pm\)。

\subsection{与 Kummer 理论的比较}
\label{b1:subsec:2004:03}

设 \(\ell\) 为不同于 \(p\) 的素数，且 \(K\) 包含全部 \(\ell\) 次单位根。Kummer 理论使用乘法商群 \(K^\times/K^{\times\ell}\)，把它看成 \(\mathbb F_\ell\) 向量空间时，标量作用由取幂给出。方程为 \(X^\ell-a\)，根之间相差 \(\mu_\ell\) 中的乘法因子；选取本原 \(\ell\) 次单位根 \(\zeta\)，生成自同构可写为 \(\sigma(\alpha)=\zeta\alpha\)，其构造来自乘法形式的 Hilbert 第 90 定理。在特征 \(p\) 中，Artin--Schreier 理论使用加法商空间 \(K/\wp(K)\)，它按通常的加法与素域标量作用成为 \(\mathbb F_p\) 向量空间。方程为 \(X^p-X-a\)，根之间相差一个加法常数 \(c\in\mathbb F_p\)，生成自同构可写为 \(\sigma(\alpha)=\alpha+1\)，其构造来自加法形式的 Hilbert 第 90 定理。前者用 \(\dfrac{b}{\sigma(b)}\) 描述范数为一的元素，后者用 \(b-\sigma(b)\) 描述迹为零的元素；两者的非平凡素数次循环扩张都由相应商空间的一维子空间分类。

\ProseParagraphBreak
这也解释了为什么不能把特征零的根式可解性判据原样搬到正特征：特征 \(p\) 中的可分循环 \(p\) 次步骤由 Artin--Schreier 方程描述，\(p\) 次开方却是纯不可分步骤。二者的次数可以相同，自同构行为则完全不同。下一章先在特征零下给出根式可解性的完整等价定理。


\begin{chaptersummary}
Hilbert 第 90 定理将范数为一与共轭比值相联系，将迹为零与共轭差相联系。范数和迹均按全部自同构计算，相同的取值也保留出现次数。乘法证明构造一个非零的加权嵌入和，加法证明比较两个线性空间的维数；两个证明均不需要除以扩张次数。

\ProseParagraphBreak
当基域 \(K\) 的特征不整除 \(n\)，且 \(K\) 含 \(\mu_n\) 时，Kummer 配对把指数整除 \(n\) 的有限交换 Galois 扩张与有限根式数据相配。特征 \(p\) 中则由 \(X^p-X-a\) 代替 \(p\) 次开方。两类素数次扩张的分类参数都是相应商空间中的直线。这个“直线”的条件排除了因改变生成元而重复计算同一个扩张。
\end{chaptersummary}

\BookExercises

\begin{exercise}\label{b1:ex:20-1}
设 \(L=\mathbb F_{q^n}\)、\(K=\mathbb F_q\)。把范数视为 \(N_{L/K}:L^\times\rightarrow K^\times\)，求其核的阶；再证明映射 \(L^\times\rightarrow L^\times\)、\(b\mapsto b/b^q\) 的像恰为此核，并求每个像的原像数。
\end{exercise}
\begin{solution}
\cref{b1:thm:finite-field-automorphisms}给出 Frobenius \(b\mapsto b^q\) 生成 Galois 群，故\eqref{b1:eq:cyclic-trace-norm}中的范数为 \(b^{1+q+\cdots+q^{n-1}}\)。\cref{b1:cor:finite-field-cyclic}说明乘法群为 \(q^n-1\) 阶循环群，所以该幂映射的像阶为 \(q-1\)，核阶为 \((q^n-1)/(q-1)\)。映射 \(b\mapsto b^{1-q}\) 的核为 \(K^\times\)，也有 \(q-1\) 个元素；相消计算使其像包含于范数核，而像的阶已经等于范数核的阶，因而相等。每个非空纤维为核的一个陪集，所以有 \(q-1\) 个元素。
\end{solution}

\begin{exercise}\label{b1:ex:20-2}
设 \(d\in\mathbb Q^\times\) 不是平方。参数化全部有理数解 \((x,y)\) 满足 \(x^2-dy^2=1\)，并解释它们与二次扩张的范数有何关系。
\end{exercise}
\begin{solution}
左边为 \(N_{\mathbb Q(\sqrt d)/\mathbb Q}(x+y\sqrt d)\)。除去解 \((-1,0)\)，可令 \(t=y/(x+1)\)。代入 \((x-1)(x+1)=dy^2\) 并除以 \(x+1\)，得 \(x-1=dt^2(x+1)\)。因为 \(d\) 不是有理平方，\(1-dt^2\neq0\)，故
\[
x=\frac{1+dt^2}{1-dt^2},\qquad y=\frac{2t}{1-dt^2}\quad(t\in\mathbb Q).
\]
反向代入给出恒等式 \(x^2-dy^2=1\)。在扩张域中，这一参数化正是 \((1+t\sqrt d)/(1-t\sqrt d)\) 的有理化结果。
\end{solution}

\begin{exercise}\label{b1:ex:20-3}
令 \(L=\mathbb F_2(\theta)\)，其中 \(\theta^2+\theta+1=0\)。分别求出对 \(\mathbb F_2\) 的迹为零的元素集合、范数为一的元素集合，再求两集合的交，并解 \(b-b^2=1\)。
\end{exercise}
\begin{solution}
域中四个元素为 \(0,1,\theta,\theta+1\)，Frobenius 为 \(b\mapsto b^2\)。迹分别为 \(0,0,1,1\)，所以迹零集合为 \(\{0,1\}\)。每个非零元素的立方为 \(1\)，故范数一集合为 \(\{1,\theta,\theta+1\}\)，两集合的交为 \(\{1\}\)。方程 \(b-b^2=1\) 在特征 \(2\) 中就是 \(b+b^2=1\)，解为 \(\theta,\theta+1\)。这两个解相差素域元素，与\cref{b1:thm:hilbert90-additive}的陪集描述一致。
\end{solution}

\begin{exercise}\label{b1:ex:20-4}
设 \(L/K\) 为 \(n\) 次循环扩张，\(\sigma\) 生成其 Galois 群，\(b\in L^\times\)，并令 \(a=b/\sigma(b)\)。对正整数 \(r\)，证明
\[
\frac b{\sigma^r(b)}=\prod_{j=0}^{r-1}\sigma^j(a).
\]
写出加法版本。若 \(\gcd(r,n)=1\)，说明改用生成元 \(\sigma^r\) 为何不会改变两个定理所描述的核。
\end{exercise}
\begin{solution}
把各 \(\sigma^j(a)=\sigma^j(b)/\sigma^{j+1}(b)\) 相乘，分子分母相消即可。加法版本是 \(b-\sigma^r(b)=\sum_{j=0}^{r-1}\sigma^j\bigl(b-\sigma(b)\bigr)\)。若 \(\gcd(r,n)=1\)，则 \((\sigma^r)^i\)（\(0\leq i<n\)）仍遍历全部自同构，按这 \(n\) 项计算的迹与范数不变。由 Hilbert 第 90 定理的两个版本，两种生成元分别给出的像都等于相同的迹核、范数核。
\end{solution}

\begin{exercise}\label{b1:ex:20-5}
设 \(L=\mathbb Q(\sqrt2,\sqrt3)\)，\(\sigma\) 固定 \(\sqrt3\) 而把 \(\sqrt2\) 变号。构造范数 \(N_{L/\mathbb Q}\) 为 \(1\) 的元素，却不能写成 \(b/\sigma(b)\)。
\end{exercise}
\begin{solution}
取 \(a=(2+\sqrt3)/(2-\sqrt3)\)。它在 \(\mathbb Q(\sqrt3)/\mathbb Q\) 中的范数为 \(1\)，故由\cref{b1:thm:trace-norm-tower}，\(N_{L/\mathbb Q}(a)=1\)。若 \(a=b/\sigma(b)\)，则 \(a\cdot\sigma(a)=1\)。但 \(\sigma(a)=a\)，而 \(a\neq\pm1\)，故 \(a^2\neq1\)，矛盾。范数到 \(\mathbb Q\) 涉及四个自同构，单个二阶 \(\sigma\) 不能生成全部群，因而不满足乘法定理的假设。
\end{solution}

\begin{exercise}\label{b1:ex:20-6}
设 \(\ell\) 为不同于域 \(K\) 特征的素数，且 \(\mu_\ell\subseteq K\)。证明对任意 \(a\in K^\times\)，多项式 \(X^\ell-a\) 要么在 \(K\) 上完全分裂，要么不可约。

再设 \(L/K\) 为 \(d\) 次循环扩张，\(\sigma\) 生成其 Galois 群，\(K\) 含本原 \(d\) 次单位根 \(\zeta\)。给定一组 \(K\) 基 \(e_1,\ldots,e_d\)，令
\[
T_\zeta(x)=\sum_{i=0}^{d-1}\zeta^{-i}\sigma^i(x)\qquad(x\in L).
\]
证明至少有一个 \(T_\zeta(e_j)\) 非零，并证明任取这样的非零值 \(\alpha\)，都有 \(L=K(\alpha)\) 且 \(\alpha^d\in K\)。
\end{exercise}
\begin{solution}
由\cref{b1:prop:single-radical-cyclic}，任一根生成的扩张次数整除 \(\ell\)，故为 \(1\) 或 \(\ell\)。前者中一根属于 \(K\)，其余根由单位根相乘也属于 \(K\)；后者中最小多项式就是整个多项式。

\(T_\zeta\) 是 \(K\) 线性映射。\(\sigma^0,\ldots,\sigma^{d-1}\) 是不同的 \(K\) 嵌入，\cref{b1:cor:embedding-independence}说明它们以非零系数 \(1,\zeta^{-1},\ldots,\zeta^{-(d-1)}\) 所成的和不可能是零映射。因此 \(T_\zeta\) 不会在全部基向量上取零。重新排列求和指标可得 \(\sigma(T_\zeta(x))=\zeta T_\zeta(x)\)。若 \(\alpha=T_\zeta(e_j)\ne0\)，则 \(\alpha^d\) 被 \(\sigma\) 固定，故属于 \(K\)；其共轭元 \(\zeta^i\alpha\)（\(0\leq i<d\)）两两不同，所以 \([K(\alpha):K]\geq d=[L:K]\)，从而 \(L=K(\alpha)\)。
\end{solution}

\begin{exercise}\label{b1:ex:20-7}
求 \(\mathbb Q(\sqrt2,\sqrt3,\sqrt5)/\mathbb Q\) 的次数、Galois 群及全部二次中间域。
\end{exercise}
\begin{solution}
若 \(2^a3^b5^c\) 是有理平方，其中 \(a,b,c\in\{0,1\}\)，把有理数写成互素整数之比并比较各素数指数，得到 \(a=b=c=0\)。因此三个平方类独立，生成的子群阶为 \(8\)。\cref{b1:thm:kummer-correspondence}给出次数 \(8\)、群 \(C_2^3\)。二次中间域对应于平方类空间的一维子空间；在 \(\mathbb F_2\) 上每条直线只有一个非零元素，所以恰为 \(\mathbb Q(\sqrt d)\)，其中 \(d=2,3,5,6,10,15,30\)。
\end{solution}

\begin{exercise}\label{b1:ex:20-8}
设 \(K\) 为域，\(n\geq1\)，\(\operatorname{char}K\nmid n\)，且 \(\mu_n\subseteq K\)。取 \(a\in K^\times\) 及 \(\alpha^n=a\)。证明 \([K(\alpha):K]\) 等于 \([a]\) 在 \(K^\times/K^{\times n}\) 中的阶。再取 \(K=\mathbb Q(i)\)、\(n=4\)，分别以 \(a=2\) 和 \(a=4\) 计算扩张次数。
\end{exercise}
\begin{solution}
记 \(L=K(\alpha)\)、\(d=[L:K]\)，并令 \(e\) 为 \([a]\) 的阶。由\cref{b1:prop:single-radical-cyclic}，\(L/K\) 循环 Galois，且 \(\operatorname{Gal}(L/K)\) 嵌入 \(\mu_n\)。每个 \(\sigma\) 所对应的比值 \(\sigma(\alpha)/\alpha\) 都属于这个阶为 \(d\) 的子群，故 \(\sigma(\alpha^d)=\alpha^d\)。因此 \(\alpha^d\in K^\times\)，于是 \(a^d=(\alpha^d)^n\in K^{\times n}\)，得到 \(e\mid d\)。反过来，\(a^e=c^n\) 对某个 \(c\in K^\times\) 成立，因而 \((\alpha^e/c)^n=1\)。由 \(\mu_n\subseteq K\) 得 \(\alpha^e\in K\)，故 \(d\leq e\)。结合两式，\(d=e\)。

在 \(K=\mathbb Q(i)\) 中，\(4\notin K^{\times4}\)：若 \(c^4=4\)，则 \(c\in\{\pm\sqrt2,\pm i\sqrt2\}\)，这些元素均不在 \(K\) 中。因此 \([4]\) 的阶为 \(2\)，而 \([2]^2=[4]\ne[1]\)、\([2]^4=[16]=[1]\)，所以 \([2]\) 的阶为 \(4\)。相应地，四次方根 \(\alpha^4=2\) 生成四次扩张，\(\beta^4=4\) 生成二次扩张。合数指数下，根式的次数可以是真因子。
\end{solution}

\begin{exercise}\label{b1:ex:20-9}
设 \(K\) 为域，\(n\geq2\)，\(\operatorname{char}K\nmid n\)，且 \(\mu_n\subseteq K\)。在 \(K\) 的一个固定代数闭包中，对 \(K^\times/K^{\times n}\) 的任一有限子群 \(B\)，从每个类 \([a]\in B\) 中任选代表 \(a\in K^\times\) 和一个 \(n\) 次根，记 \(L_B=K(\sqrt[n]{a}\mid [a]\in B)\)。若 \(B_1,B_2\) 是两个这样的子群，证明 \(L_{B_1B_2}=L_{B_1}L_{B_2}\)，以及 \(L_{B_1\cap B_2}=L_{B_1}\cap L_{B_2}\)。
\end{exercise}
\begin{solution}
乘积子群由两个子群的生成元合在一起生成，加入它们的根得到的就是合成域，因此第一式成立。两域的交仍为有限 Galois 扩张，其群是其中任一 Galois 群的商，因而交换且指数整除 \(n\)。对 \(a\in K^\times\)，若其某个 \(n\) 次根分别属于两域，因为所有根只相差 \(K\) 中的单位根，任取一个根就同时属于两域。故交域对应的子群恰为 \(B_1\cap B_2\)，由对应的唯一性得到第二式。
\end{solution}

\begin{exercise}\label{b1:ex:20-10}
设 \(K\) 含本原三次单位根，并取 \(a\in K^\times\setminus K^{\times3}\)。对 \(\alpha^3=a\)，写出 \(K(\alpha)/K\) 的自同构，并证明 \(a^2b^3\)（\(b\in K^\times\)）给出同一个扩张。
\end{exercise}
\begin{solution}
\cref{b1:prop:single-radical-cyclic}及素数次数给出三个自同构 \(\alpha\mapsto\zeta_3^j\alpha\)。元素 \(b\alpha^2\) 的立方为 \(a^2b^3\)。反过来，由 \(\alpha=a^{-1}(\alpha^2)^2\)，得 \(K(b\alpha^2)=K(\alpha)\)。商群中 \([a]\) 和 \([a]^2\) 是同一直线上的两个不同非零元素，因此以单个非零类分类会重复计算这个域。
\end{solution}

\begin{exercise}\label{b1:ex:20-11}
在 \(K=\mathbb F_3(t)\) 上，比较参数 \(t\)、\(2t\)、\(t+t^3-t\) 给出的 Artin--Schreier 扩张，并说明它们非平凡。
\end{exercise}
\begin{solution}
第三个参数就是 \(t^3\)，与 \(t\) 的差为 \(\wp(t)\)；第二个参数是第一个的非零素域倍数。由\cref{b1:thm:artin-schreier-classification}，三者给出同一扩张。具体地，若 \(\alpha^3-\alpha=t\)，则 \(2\alpha\) 对应 \(2t\)，\(\alpha+t\) 对应 \(t^3\)。\(t\) 在无穷远处有一阶极点；若 \(\wp(b)\) 在无穷远处有极点，则其极点阶是 \(3\) 的倍数。故 \(t\notin\wp(K)\)，扩张次数为 \(3\)。
\end{solution}

\begin{exercise}\label{b1:ex:20-12}
设 \(p>2\)、\(K=\mathbb F_p(t)\)。证明 \(t,t^2\) 的类在 \(K/\wp(K)\) 中线性无关，并求由 \(\alpha^p-\alpha=t\)、\(\beta^p-\beta=t^2\) 生成的合成域次数及 Galois 群。列出合成域中全部对 \(K\) 为 \(p\) 次的中间域，并说明为什么没有遗漏。
\end{exercise}
\begin{solution}
任何非零组合 \(ct+dt^2\) 的无穷远极点阶为 \(1\) 或 \(2\)，不被 \(p\) 整除；它不可能等于 \(\wp(b)\)，故两类独立。两个 \(p\) 次循环域因其参数直线不同而不同。把第二个 Artin--Schreier 多项式看作第一个域上的多项式，由\cref{b1:prop:artin-schreier-polynomial}，其根要么已在该域，要么产生 \(p\) 次扩张；前者会使两个原来的 \(p\) 次域相等，已被排除。因此合成域 \(M=K(\alpha,\beta)\) 的次数为 \(p^2\)。它是两个可分多项式的共同分裂域，限制到两域给出到 \(C_p\times C_p\) 的单射，两边阶相同，所以 Galois 群正是该直积。

若 \(c\in\mathbb F_p\)，则 \((\alpha+c\beta)^p-(\alpha+c\beta)=t+ct^2\)，其参数类与 \([t^2]\) 一起代表二维空间 \(\langle[t],[t^2]\rangle_{\mathbb F_p}\) 的全部 \(p+1\) 条直线。因此\cref{b1:thm:artin-schreier-classification}给出 \(p+1\) 个互异的 \(p\) 次中间域
\[
K(\beta),\qquad K(\alpha+c\beta)\quad(c\in\mathbb F_p).
\]
另一方面，\(\operatorname{Gal}(M/K)\cong C_p^2\) 恰有 \(p+1\) 个指数为 \(p\) 的子群；由 Galois 对应，没有其他 \(p\) 次中间域。
\end{solution}

\begin{exercise}\label{b1:ex:20-13}
令 \(K=\mathbb F_{p^r}\)。证明 \(\wp(K)=\ker\operatorname{Tr}_{K/\mathbb F_p}\)，并据此说明 \(K\) 在固定代数闭包中恰有一个 \(p\) 次循环扩张。
\end{exercise}
\begin{solution}
因 \(\ker\wp=\mathbb F_p\)，\(\wp(K)\) 有 \(p^{r-1}\) 个元素。迹的嵌入和公式给出 \(\operatorname{Tr}(b^p)=\operatorname{Tr}(b)\)，因为取 \(p\) 次幂只会循环置换各项，所以 \(\wp(K)\) 包含于迹核。由\cref{b1:cor:embedding-independence}，这些不同嵌入之和不是零映射，而目标空间 \(\mathbb F_p\) 为一维，故迹满射。迹核也有 \(p^{r-1}\) 个元素，两者相等。商 \(K/\wp(K)\) 因而是一维 \(\mathbb F_p\) 空间，只有一条一维子空间。应用\cref{b1:thm:artin-schreier-classification}即可。
\end{solution}

\begin{exercise}\label{b1:ex:20-14}
证明 \(X^p-X-1\) 在 \(\mathbb F_p[X]\) 中不可约，并求其根 \(\alpha\) 到 \(\mathbb F_p\) 的迹和范数。
\end{exercise}
\begin{solution}
\(\wp(\mathbb F_p)=0\)，故\cref{b1:prop:artin-schreier-polynomial}给出不可约性。全部根为 \(\alpha+c\)、\(c\in\mathbb F_p\)，其和为 \(\sum c\)。该和在 \(p>2\) 时为 \(0\)，在 \(p=2\) 时为 \(1\)。范数是根的乘积，由首一多项式常数项公式等于 \((-1)^p(-1)=1\) 于 \(\mathbb F_p\)。迹在奇素数情形为零，不妨碍扩张可分或迹映射满射。
\end{solution}

\begin{exercise}\label{b1:ex:20-15}
设 \(K=\mathbb F_p(t)\)、\(L=K(u)\)，其中 \(u^p=t\)。证明 \([L:K]=p\) 且 \(\operatorname{Aut}_K(L)=1\)，并与参数 \(t\) 的 Artin--Schreier 扩张比较。
\end{exercise}
\begin{solution}
在 \(\mathbb F_p[t][X]\) 中，\(X^p-t\) 对素元 \(t\) 满足\cref{b1:thm:eisenstein}的条件：非首项系数均被 \(t\) 整除，常数项不被 \(t^2\) 整除。因此它在 \(K[X]\) 中不可约，次数为 \(p\)。任一自同构必须把 \(u\) 送到 \(X^p-t=(X-u)^p\) 的根，唯一可能是 \(u\)，所以只能是恒等映射。相比之下，\(X^p-X-t\) 的导数为 \(-1\)，而 \(t\) 的无穷远极点阶为 \(1\)，不能是 \(\wp(b)\) 的极点阶，所以\cref{b1:prop:artin-schreier-polynomial}给出其不可约性；其根生成同次数的可分正规扩张，自同构群为 \(C_p\)。
\end{solution}

% !TeX root = ../../Book1.tex
% 纲要标题：## 第21章　分圆扩张、根式与尺规作图
% 纲要位置：计划纲要.md，第 736 行。

\chapter[分圆扩张、根式与尺规作图]{分圆扩张、\mbox{根式与尺规作图}}
\label{b1:ch:21}
\BookChapterNumber{b1:ch:21}

\begin{chapterabstract}
单位根把正多边形作图转成分圆域的结构问题；根式求解则与多项式 Galois 群的可解性相连。本章建立分圆域及根式可解性理论，比较低次方程与一般五次方程，并以二次扩张塔判断尺规可作性。
\end{chapterabstract}

\begin{learninggoals}
\begin{itemize}
\item 从全部单位根中分离本原单位根，证明并运用分圆多项式的不可约性。
\item 区分根式塔、分裂域和正规闭包，掌握根式可解性判据的两个方向。
\item 解释一般方程的独立系数含义，避免将一般五次结论误用于所有特殊五次方程。
\item 把尺规操作与实二次扩张塔互相转换，判断典型长度、角和正多边形是否可作图。
\end{itemize}
\end{learninggoals}

正五边形与正七边形都可以由单位圆上的根来描述。前者涉及的分圆域在有理数域上次数为 \(4\)，且其 Galois 群为循环群 \(C_4\)，由指数为二的子群可得到两次二次扩张；后者的相应次数为 \(6\)，不可能包含在二次扩张塔的末端。这解释了为什么尺规能作正五边形，却不能作正七边形。这里要先弄清单位根生成了什么域、域的自同构怎样排列这些根，才能把几何作图翻译成域扩张问题。类似地，讨论方程能否用根式求解，也须把允许的求根步骤同相应 Galois 群的结构联系起来。

\ProseParagraphBreak

本章使用\cref{b1:thm:finite-galois-correspondence}、\cref{b1:thm:kummer-cyclic}和可解群的性质；分圆不可约性还需\cref{b1:thm:gauss-descent}。一般五次方程的结论须另算独立参数多项式的 Galois 群，具体多项式 \(X^5-X-1\) 则使用\cref{b1:lem:reduction-cycle-type}和置换群论证。作图判据以平面几何、二次方程和二次扩张塔为基础；正多边形还需分圆域与 Galois 对应。圆周率超越性的证明概要另用复指数函数、路径积分与代数整数范数。

\ProseParagraphBreak
单位根与根式可解性可读 Roman 的《Field Theory》\cite[Sections 11.1--11.2, 13.2--13.7]{Roman2006FieldTheory}。三次方程、实根和对称多项式的中文基础材料可见席南华的《基础代数》第一卷\cite[第7章]{Xi2016BasicAlgebraI}。尺规作图末尾的超越性概要依据 Lindemann 的原文\cite[\S\S 1--4, pp. 214--225]{Lindemann1882Pi}，那里使用的积分估计与本章的代数作图判据分别承担不同的论证。

% !TeX root = ../../Book1.tex
% 纲要标题：### 21.1 分圆多项式
% 纲要位置：计划纲要.md，第 738 行。

\section{分圆多项式}
\label{b1:sec:2101}

上一章在“基域已经含有所需单位根”的条件下讨论开方。现在把基域取为 \(\mathbb Q\)，研究加入单位根本身造成的扩张。需要区分两件事：\(X^n-1\) 收集全部 \(n\) 次单位根，而分圆多项式只收集乘法阶恰为 \(n\) 的那些根。

% 纲要内容（仅作写作计划，尚非教材正文）：
%
% * 单位根与分圆多项式
% * 分圆多项式的不可约性
% * 分圆域及其 Galois 群；一般特征零基域上的单位根扩张只嵌入相应单位群
%

\subsection{单位根与分圆多项式}
\label{b1:subsec:2101:01}

取 \(\zeta_n=e^{2\pi i/n}\)。它是一个本原 \(n\) 次单位根；所有本原 \(n\) 次单位根为 \(\zeta_n^a\)，其中 \(1\leq a\leq n\)、\(\gcd(a,n)=1\)。

\begin{definition}\label{b1:def:cyclotomic-polynomial}
设 \(n\geq1\) 为整数，并在 \(\mathbb C\) 中选取一个本原 \(n\) 次单位根 \(\zeta_n\)。多项式
\[
\Phi_n(X)\coloneqq\prod_{\substack{1\leq a\leq n\\\gcd(a,n)=1}}(X-\zeta_n^a)
\]
不依赖本原单位根 \(\zeta_n\) 的选择，称为第 \(n\) 个\term[fenyuan-duoxiangshi]{分圆多项式}[cyclotomic polynomial]。其次数为 Euler 函数 \(\varphi(n)\)，即模 \(n\) 可逆剩余类的个数。
\end{definition}
\symidx{\(\Phi_n(X)\)：第 \(n\) 个分圆多项式}

若改取另一个本原单位根 \(\zeta_n'=\zeta_n^b\)，则 \(\gcd(b,n)=1\)。在单位群 \((\mathbb Z/n\mathbb Z)^\times\) 中，乘以 \(b\) 是一个置换，所以 \((\zeta_n')^a=\zeta_n^{ba}\) 仍遍历同一组本原根。这就证明了定义中的乘积不依赖 \(\zeta_n\) 的选择。

\ProseParagraphBreak
在特征零中，\(X^n-1\) 与导数 \(nX^{n-1}\) 互素，所以它的 \(n\) 个根彼此不同。每个根的乘法阶又是 \(n\) 的唯一一个正因子，按阶分组便把全部根不重不漏地分到各个 \(\Phi_d\) 中。因此
\begin{equation}\label{b1:eq:cyclotomic-factorization}
X^n-1=\prod_{d\mid n}\Phi_d(X).
\end{equation}
这还给出 \(\Phi_n\in\mathbb Z[X]\)。对 \(n\) 归纳，\(\Phi_1=X-1\)；若真因子对应的多项式均已知为首一整系数多项式，则用它们的乘积去除 \(X^n-1\)。除式首一，每次消去最高项时只须减去除式的一个整系数单项式倍数，无须除以非单位的首项系数，所以商、余式仍在 \(\mathbb Z[X]\) 中。而\eqref{b1:eq:cyclotomic-factorization}说明在 \(\mathbb C[X]\) 中余式为零、商为 \(\Phi_n\)；除法的唯一性便给出结论。

\ProseParagraphBreak
例如 \(\Phi_2=X+1\)、\(\Phi_3=X^2+X+1\)、\(\Phi_4=X^2+1\)、\(\Phi_6=X^2-X+1\)。对素数 \(p\)，有 \(\Phi_p=1+X+\cdots+X^{p-1}\)。不过，整系数性本身不能说明不可约；下一小节需要另外论证所有本原根互为有理数域上的共轭。

\subsection{分圆多项式的不可约性}
\label{b1:subsec:2101:02}

要证明 \(\Phi_n\) 不可约，只须证明一个本原根的最小多项式已经包含全部本原根。我们将证明，它的根取任意不整除 \(n\) 的素数次幂后，仍是同一个最小多项式的根；将互素指数分解成素因子，就能从 \(\zeta_n\) 逐次到达每个 \(\zeta_n^a\)。

\begin{theorem}\label{b1:thm:cyclotomic-irreducible}
设 \(n\geq1\) 为整数，\(\zeta_n\in\mathbb C\) 为本原 \(n\) 次单位根。则第 \(n\) 个分圆多项式 \(\Phi_n(X)\) 在 \(\mathbb Q[X]\) 中不可约。因此 \(\bigl[\mathbb Q(\zeta_n):\mathbb Q\bigr]=\varphi(n)\)。
\end{theorem}
\begin{proof}
设 \(f\) 为 \(\zeta_n\) 在 \(\mathbb Q\) 上的首一最小多项式。它是首一整系数多项式 \(\Phi_n\) 的因子，由\cref{b1:thm:gauss-descent}，可写 \(\Phi_n=fg\)，其中 \(f,g\in\mathbb Z[X]\) 都首一。

设 \(\alpha\) 是 \(f\) 的根，\(q\) 为不整除 \(n\) 的素数。由 \(f\mid\Phi_n\)，\(\alpha\) 的乘法阶为 \(n\)。又因 \(\gcd(q,n)=1\)，\cref{b1:prop:power-order}给出
\[
\operatorname{ord}(\alpha^q)=\frac{n}{\gcd(n,q)}=n,
\]
故 \(\alpha^q\) 仍是 \(\Phi_n\) 的根。若它不是 \(f\) 的根，由 \(\Phi_n=fg\) 就有 \(g(\alpha^q)=0\)。\(f\) 不可约且零化 \(\alpha\)，所以它也是 \(\alpha\) 的首一最小多项式，从而 \(f(X)\mid g(X^q)\) 于 \(\mathbb Q[X]\)。首一除法又说明这个整除关系的商在 \(\mathbb Z[X]\) 中。

把这个整系数整除关系模 \(q\) 约化，并以横线表示所得多项式。特征 \(q\) 的二项式公式以及 \(c^q=c\) 对 \(c\in\mathbb F_q\) 的成立给出
\[
\bar f(X)\mid\bar g(X^q)=\bar g(X)^q.
\]
\(f\) 首一，故约化不会降低次数，\(\bar f\) 仍有正次数，可取它的一个不可约因子 \(h\)。域上的一元多项式环 \(\mathbb F_q[X]\) 是唯一分解整环，\cref{b1:prop:ufd-prime}说明其中的不可约元是素元。因此 \(h\mid\bar g^q\) 迫使 \(h\mid\bar g\)。于是 \(h\) 同时整除 \(\bar f\) 和 \(\bar g\)，给出 \(h^2\mid\bar f\bar g=\overline{\Phi_n}\)。再把\eqref{b1:eq:cyclotomic-factorization}模 \(q\) 约化，得到 \(h^2\mid X^n-1\)。但 \(q\nmid n\)，所以导数 \(nX^{n-1}\) 的系数 \(n\) 非零，而 \(X\) 不整除 \(X^n-1\)；两多项式互素，不可能有重不可约因子，矛盾。这证明了 \(\alpha^q\) 仍为 \(f\) 的根。

任意与 \(n\) 互素的正整数 \(a>1\) 可写成 \(a=q_1\cdots q_s\)，其中每个素数 \(q_j\) 都不整除 \(n\)。从已知的根 \(\zeta_n\) 出发，刚证的结论依次说明
\[
\zeta_n^{q_1},\quad \zeta_n^{q_1q_2},\quad\ldots,\quad\zeta_n^{q_1\cdots q_s}=\zeta_n^a
\]
都是 \(f\) 的根；\(a=1\) 时无需这一步。因此 \(f\) 含有全部 \(\varphi(n)\) 个互异的本原根，而它又整除同为首一的 \(\Phi_n\)，只能有 \(f=\Phi_n\)。次数结论来自\cref{b1:prop:minimal-polynomial-degree}。
\end{proof}

\subsection{分圆域及其 Galois 群}
\label{b1:subsec:2101:03}

\begin{theorem}\label{b1:thm:cyclotomic-galois-group}
设 \(n\geq1\) 为整数，\(\zeta_n\in\mathbb C\) 为本原 \(n\) 次单位根。则\term[fenyuan-yu]{分圆域}[cyclotomic field] \(\mathbb Q(\zeta_n)\) 是 \(X^n-1\) 在 \(\mathbb Q\) 上的分裂域，且
\[
\operatorname{Gal}\bigl(\mathbb Q(\zeta_n)/\mathbb Q\bigr)
\cong(\mathbb Z/n\mathbb Z)^\times,
\qquad \sigma_a(\zeta_n)=\zeta_n^a.
\]
\end{theorem}
\begin{proof}
全部 \(n\) 次单位根都是 \(\zeta_n\) 的幂，所以 \(\mathbb Q(\zeta_n)\) 正是 \(X^n-1\) 的分裂域。它有限且正规，特征零又保证可分，因而为有限 Galois 扩张。自同构保持元素乘法阶，所以每个自同构 \(\sigma\) 都满足 \(\sigma(\zeta_n)=\zeta_n^a\)，其中 \(\gcd(a,n)=1\)。剩余类 \(a\bmod n\) 唯一，而且自同构由生成元 \(\zeta_n\) 的像决定，于是得到单射
\[
\operatorname{Gal}\bigl(\mathbb Q(\zeta_n)/\mathbb Q\bigr)
\longrightarrow(\mathbb Z/n\mathbb Z)^\times.
\]
若 \(\sigma(\zeta_n)=\zeta_n^a\)、\(\tau(\zeta_n)=\zeta_n^b\)，则 \((\sigma\tau)(\zeta_n)=\zeta_n^{ab}\)，故这也是群同态。由\cref{b1:prop:galois-automorphism-count,b1:thm:cyclotomic-irreducible}，源群的阶为
\[
\left|\operatorname{Gal}\bigl(\mathbb Q(\zeta_n)/\mathbb Q\bigr)\right|
=[\mathbb Q(\zeta_n):\mathbb Q]=\varphi(n).
\]
目标单位群的阶按 \(\varphi(n)\) 的定义也为 \(\varphi(n)\)，所以单射必为满射，得到所述同构。
\end{proof}

例如 \(\mathbb Q(\zeta_5)/\mathbb Q\) 的群为 \(C_4\)，而 \(\mathbb Q(\zeta_8)/\mathbb Q\) 的群为 \(C_2\times C_2\)，不能把所有分圆扩张都称为循环扩张。“循环扩张”要求 Galois 群循环；\(\mathbb Q(\zeta_8)\) 虽由一个元素生成，却不满足这个要求。

\ProseParagraphBreak
复共轭 \(c\) 总对应于剩余类 \(-1\)。设 \(n\geq3\)，记 \(t=\zeta_n+\zeta_n^{-1}\)，则 \(c(t)=t\)，所以
\[
\mathbb Q(t)\subseteq\mathbb Q(\zeta_n)^{\langle c\rangle}.
\]
同时 \(\mathbb Q(t)\subseteq\mathbb R\)，而 \(\zeta_n\) 非实，并满足
\[
X^2-(\zeta_n+\zeta_n^{-1})X+1=0;
\]
故 \([\mathbb Q(\zeta_n):\mathbb Q(t)]=2\)。子群 \(\langle c\rangle\) 的阶也为 \(2\)，由\eqref{b1:eq:galois-degree-index}，它的不动域同样满足 \([\mathbb Q(\zeta_n):\mathbb Q(\zeta_n)^{\langle c\rangle}]=2\)。结合上述包含关系和塔式公式，得到这两个中间域相等。

\ProseParagraphBreak
把基域 \(\mathbb Q\) 换成任意特征零域 \(K\) 后，自同构还须固定 \(K\) 中的其他元素，因而不能再保证每个互素幂次都能作为 \(\zeta_n\) 的像。不过，确定生成元像的办法仍给出到单位群的单同态，所以 Galois 群仍然交换。

\begin{proposition}\label{b1:prop:cyclotomic-general-base}
设 \(K\) 为特征零域，\(\Omega\) 为 \(K\) 的代数闭包，\(n\geq1\) 为整数，\(\mu_n\subseteq\Omega\) 为全部 \(n\) 次单位根。选取一个本原 \(n\) 次单位根 \(\zeta_n\in\Omega\)，并令 \(L=K(\mu_n)=K(\zeta_n)\)。则 \(L/K\) 为有限 Galois 扩张，映射
\[
\operatorname{Gal}(L/K)\longrightarrow(\mathbb Z/n\mathbb Z)^\times,
\qquad\sigma\longmapsto a\bmod n\quad
\bigl(\sigma(\zeta_n)=\zeta_n^a\bigr)
\]
是单同态。因此 \(\operatorname{Gal}(L/K)\) 是交换群，且 \([L:K]\) 整除 \(\varphi(n)\)。
\end{proposition}
\begin{proof}
特征零保证 \(X^n-1\) 有 \(n\) 个互异的根，它们组成乘法群 \(\mu_n\)。由\cref{b1:lem:finite-multiplicative-cyclic}，该群循环，存在所述的本原单位根，且全部根都是 \(\zeta_n\) 的幂。因此 \(L=K(\zeta_n)\) 是 \(X^n-1\) 的分裂域，为有限 Galois 扩张。

每个 \(\sigma\in\operatorname{Gal}(L/K)\) 保持 \(\zeta_n\) 的乘法阶，所以 \(\sigma(\zeta_n)=\zeta_n^a\)，其中 \(a\bmod n\) 是一个唯一确定的可逆剩余类。生成元的像决定自同构，故所列映射单射；复合对应指数相乘，故它是同态。其像是交换群 \((\mathbb Z/n\mathbb Z)^\times\) 的子群，因而交换，且阶整除 \(\varphi(n)\)。再由有限 Galois 扩张的群阶等于次数，得到次数的整除结论。
\end{proof}

这个单射在一般基域上不必满射。例如 \(K\) 本来就包含 \(\mu_n\) 时，\(L=K\)，Galois 群平凡；\(n\geq3\) 时，单位群却并不平凡。

% !TeX root = ../../Book1.tex
% 纲要标题：### 21.2 根式可解性
% 纲要位置：计划纲要.md，第 744 行。

\section{根式可解性}
\label{b1:sec:2102}

一个根式公式允许有限次四则运算和开方，也允许在计算途中使用复数。域塔把这样的表达式记录为一串有限扩张；Galois 群则检验这一串扩张是否可能存在。本节的等价定理以特征零为主版本，且把“一个根能够表示”与“多项式的全部根能够表示”明确分开。

% 纲要内容（仅作写作计划，尚非教材正文）：
%
% * 根式塔与可解扩张
% * 根式塔的正规可解包络
% * 根式可解性与可解群：集中陈述特征零判据并连续证明两个方向
% * 主证明完成后讨论特征条件与可分性约定
% * 完整的根式可解性等价定理以特征零为主版本：在加入必要单位根后，调用第20.3节的循环扩张根式描述，沿可解群的素数阶循环商逐层构造；正特征情形单独说明与 Artin–Schreier 扩张的差别。
%

\subsection{根式塔与可解扩张}
\label{b1:subsec:2102:01}

\begin{definition}\label{b1:def:radical-extension}
设 \(K\) 为域，\(r\geq0\) 为整数。给定有限域扩张塔
\[
K=K_0\subseteq K_1\subseteq\cdots\subseteq K_r,
\qquad K_i=K_{i-1}(\alpha_i),\quad\alpha_i^{m_i}\in K_{i-1},
\]
其中每个 \(m_i\geq2\) 都是整数。如果各步都满足所列条件，则称这座塔为\term[genshi-ta]{根式塔}[radical tower]，并称末端扩张 \(K_r/K\) 为\term[genshi-kuozhang]{根式扩张}[radical extension]。给定多项式 \(f\in K[X]\)，如果存在一座从 \(K\) 出发的根式塔，其末端域包含 \(f\) 在一个代数闭包中的全部根，则称 \(f\) \term[genshi-kejie]{根式可解}[solvable by radicals]。
\end{definition}

\(\alpha_i=0\) 或本已属于 \(K_{i-1}\) 的步骤可以删去。开方的指数可以是合数，但能够分拆：若 \(\alpha^{rs}=a\in F\)，其中 \(r,s\geq2\)，令 \(\beta=\alpha^r\)，则 \(\beta^s=a\) 且 \(\alpha^r=\beta\)。于是原来一步 \(F\subseteq F(\alpha)\) 可以改成
\[
F\subseteq F(\beta)\subseteq F(\alpha),
\]
分别取 \(s\) 次方根和 \(r\) 次方根。继续分解指数，就得到一条只取素数次方根、末端域相同的塔。这里并不要求每一步正规，也不要求 \(K_r\) 恰好等于分裂域。例如 \(\mathbb Q(\sqrt[3]{2})\) 是根式扩张，但还没有包含 \(X^3-2\) 的全部根；再加入 \(\zeta_3\) 即可。

\begin{definition}\label{b1:def:solvable-extension}
设 \(K\) 的特征为零，\(E/K\) 为有限扩张。若 \(E/K\) 的正规闭包的 Galois 群是可解群，则称 \(E/K\) 为\term[kejie-kuozhang]{可解扩张}[solvable extension]。
\end{definition}

这里的正规闭包可按\cref{b1:prop:normal-closure}取为有限生成元的最小多项式乘积的分裂域，所以仍是有限正规扩张。特征零又保证它可分，因而确为有限 Galois 扩张。

\ProseParagraphBreak
根式扩张要求这个具体的域 \(E\) 本身能作为一座根式塔的末端；可解扩张只要求它的正规闭包具有可解的 Galois 群。若 \(E/K\) 已经是有限 Galois 扩张，后一个定义中的群就是 \(\operatorname{Gal}(E/K)\)。下面的判据将说明，可解性保证 \(E\) 能包含在某座根式塔的末端域中，而这个末端域可以严格大于 \(E\)。稍后的实三次分圆子域正是 \(E/K\) 可解而 \(E/K\) 本身不是根式扩张的例子。

\subsection{根式塔的正规可解包络}
\label{b1:subsec:2102:02}

根式塔的各层未必正规。单位根保证同一个二项式的根只相差一个已知因子，但补齐各层原方程的根还不够：基域上的共轭变换也会移动方程的系数。因此，要把根式塔放进一个 Galois 扩张，在加入每个根式时还须同时补入共轭方程的根。

\ProseParagraphBreak
例如令 \(\alpha=\sqrt{1+\sqrt2}>0\)。根式塔
\[
\mathbb Q\subseteq\mathbb Q(\sqrt2)\subseteq\mathbb Q(\alpha)
\]
的第二步由 \(X^2-(1+\sqrt2)\) 给出；这里 \(\sqrt2=\alpha^2-1\)，所以前一层确实包含在末端域中。这个二项式的两个根 \(\pm\alpha\) 已经都在末端域中，二次单位根也已经属于 \(\mathbb Q\)。可是将 \(\sqrt2\) 送到 \(-\sqrt2\) 的共轭变换，会把方程变成
\[
X^2-(1-\sqrt2)=0.
\]
它的根是非实数，不属于实域 \(\mathbb Q(\alpha)\)。因此补齐原方程的根仍不能得到 \(\mathbb Q\) 上的正规扩张；还须补入共轭方程的根。一般构造中让前一层的自同构遍历，正是为了包含这些随系数变化而出现的方程。

\begin{lemma}\label{b1:lem:radical-normal-envelope}
特征零域 \(K\) 上的任意有限根式塔，都包含在一个有限 Galois 扩张 \(M/K\) 中，使 \(\operatorname{Gal}(M/K)\) 可解。
\end{lemma}
\begin{proof}
若塔中所有域都等于 \(K\)，取 \(M=K\) 即可。以下略去平凡步骤，记 \(K_i=K_{i-1}(\alpha_i)\)、\(\alpha_i^{m_i}=a_i\neq0\)，并固定一个包含末端域的 \(K\) 的代数闭包 \(\Omega\)。我们将在 \(\Omega\) 中逐层构造 \(M_i\)，使 \(K_i\subseteq M_i\)，而 \(M_i/K\) 始终是群可解的有限 Galois 扩张。

取全部 \(m_i\) 的公倍数 \(m\)，先令 \(M_0=K(\mu_m)\)。由\cref{b1:prop:cyclotomic-general-base}，它是 \(X^m-1\) 的有限 Galois 分裂域，其 Galois 群嵌入交换群 \((\mathbb Z/m\mathbb Z)^\times\)。所以这个群本身交换，特别地可解。

归纳设 \(M_{i-1}/K\) 为有限 Galois 扩张，包含 \(K_{i-1}\)，且群可解。令 \(M_i\) 为 \(M_{i-1}\) 上全部多项式
\[
X^{m_i}-\tau(a_i),\qquad\tau\in\operatorname{Gal}(M_{i-1}/K),
\]
在 \(\Omega\) 中的共同分裂域。这族多项式只有有限多个，每一个又只有有限多个代数根，故 \(M_i/M_{i-1}\) 有限，从而 \(M_i/K\) 有限。取 \(\tau=1\) 时的方程，其根 \(\alpha_i\) 必在共同分裂域中，所以 \(K_i\subseteq M_i\)。让 \(\tau\) 遍历的目的，正是把前例中随 \(\sqrt2\mapsto-\sqrt2\) 出现的新方程一并加入。

任取一个 \(K\) 嵌入 \(\iota:M_i\rightarrow\Omega\)。因 \(M_{i-1}/K\) 正规，它在 \(M_{i-1}\) 上的限制为某个 \(\sigma\in\operatorname{Gal}(M_{i-1}/K)\)。若 \(\beta\) 是 \(X^{m_i}-\tau(a_i)\) 的根，则
\[
\iota(\beta)^{m_i}=\iota\bigl(\tau(a_i)\bigr)=(\sigma\tau)(a_i).
\]
因此 \(\iota\) 将这个方程的根集映到指标为 \(\sigma\tau\) 的方程的根集。因 \(\tau(a_i)\neq0\) 且特征为零，两个根集均有 \(m_i\) 个元素，嵌入又是单射，所以该映射是双射；左乘 \(\sigma\) 同时置换全部指标。于是 \(\iota\) 把这些根和 \(M_{i-1}\) 生成的域 \(M_i\) 映到自身且映满。由\cref{b1:thm:normal-embedding-criterion}，\(M_i/K\) 正规；特征零保证可分，故它是有限 Galois 扩张。

对每个所列多项式选一个根 \(\beta_\tau\)。因为 \(m_i\mid m\)，且 \(M_0\subseteq M_{i-1}\)，有 \(\mu_{m_i}\subseteq\mu_m\subseteq M_{i-1}\)。每个方程的全部根因而为 \(\zeta\beta_\tau\)，其中 \(\zeta\in\mu_{m_i}\)，所以
\[
M_i=M_{i-1}\bigl(\beta_\tau\mid\tau\in\operatorname{Gal}(M_{i-1}/K)\bigr).
\]
对 \(\rho\in\operatorname{Gal}(M_i/M_{i-1})\)，因 \(\rho\) 固定 \(\tau(a_i)\)，比值 \(\rho(\beta_\tau)/\beta_\tau\) 属于 \(\mu_{m_i}\)。由此得到映射
\[
\operatorname{Gal}(M_i/M_{i-1})\longrightarrow
\prod_{\tau}\mu_{m_i},\qquad
\rho\longmapsto\bigl(\rho(\beta_\tau)/\beta_\tau\bigr)_\tau
\]
它与\cref{b1:prop:single-radical-cyclic}中的比值映射相同。具体地，\(\rho\) 固定 \(M_{i-1}\)，故固定全部单位根；对另一个相对自同构 \(\eta\)，有
\[
\frac{\rho\eta(\beta_\tau)}{\beta_\tau}
=\rho\Bigl(\frac{\eta(\beta_\tau)}{\beta_\tau}\Bigr)
\frac{\rho(\beta_\tau)}{\beta_\tau}
=\frac{\eta(\beta_\tau)}{\beta_\tau}
\frac{\rho(\beta_\tau)}{\beta_\tau}.
\]
目标群交换，上式便是逐坐标的同态公式。若全部比值为 \(1\)，\(\rho\) 固定全部生成根，也就固定整个 \(M_i\)，故这个同态单射。单位根的直积是交换群，其子群也交换，所以 \(\operatorname{Gal}(M_i/M_{i-1})\) 实际上是交换群。由于 \(M_{i-1}/K\) 已是 Galois 扩张，\cref{b1:thm:galois-quotient}给出限制映射的正合列
\[
1\longrightarrow\operatorname{Gal}(M_i/M_{i-1})
\longrightarrow\operatorname{Gal}(M_i/K)
\longrightarrow\operatorname{Gal}(M_{i-1}/K)\longrightarrow1.
\]
其核是刚证明交换的相对 Galois 群，商是归纳中已知可解的 \(\operatorname{Gal}(M_{i-1}/K)\)。由\cref{b1:thm:solvable-closure}，\(\operatorname{Gal}(M_i/K)\) 可解。这就保持了构造中所需的全部条件；取 \(M=M_r\)，它包含根式塔的末端域，因而得到结论。
\end{proof}

\subsection{根式可解性与可解群}
\label{b1:subsec:2102:04}

\begin{theorem}[根式可解性判据]\label{b1:thm:radical-solvability}
设 \(K\) 的特征为零，\(f\in K[X]\) 非常数，\(L\) 为其分裂域。则 \(f\) 根式可解当且仅当 \(\operatorname{Gal}(L/K)\) 是可解群。
\end{theorem}

由根式塔推出群可解，所用的是正规可解包络及分裂域群作为商群的性质；由群可解构造根式塔，所用的则是素数阶循环因子和 Kummer 定理。后一方向须先加入单位根，才能把循环扩张转成开方。

\begin{proof}
先设 \(f\) 的全部根包含在某条根式塔中。由\cref{b1:lem:radical-normal-envelope}，该塔包含于有限 Galois 扩张 \(M/K\)，其群可解。\(L\) 由 \(f\) 的全部根在 \(K\) 上生成，故 \(L\subseteq M\)。又因 \(L/K\) 是特征零域上多项式的分裂域，它也为有限 Galois 扩张。对 \(M/K\) 的中间域 \(L\) 应用\cref{b1:thm:galois-quotient}，\(\operatorname{Gal}(M/L)\) 正规于 \(\operatorname{Gal}(M/K)\)，限制同态满射到 \(\operatorname{Gal}(L/K)\)，且
\[
\operatorname{Gal}(L/K)\cong
\operatorname{Gal}(M/K)/\operatorname{Gal}(M/L).
\]
由\cref{b1:thm:solvable-closure}，这个商群可解，得到必要性。

反过来，设 \(G=\operatorname{Gal}(L/K)\) 可解。取 \(m=|G|\)。若 \(m=1\)，由 \([L:K]=|G|\) 已有 \(L=K\)，无须开方。以下设 \(m>1\)，选取本原 \(m\) 次单位根 \(\zeta_m\)，令 \(F=K(\mu_m)=K(\zeta_m)\)、\(E=LF\)。因为 \(\zeta_m^m=1\in K\)，这一个特定的 \(m\) 次方根就给出根式扩张 \(F/K\)。\(E/F\) 是 \(f\) 在 \(F\) 上的分裂域，有限且 Galois。对任意 \(\sigma\in\operatorname{Gal}(E/F)\)，因 \(K\subseteq F\)，\(\sigma|_L\) 是一个 \(K\) 嵌入；\(L/K\) 的正规性又保证 \(\sigma(L)=L\)，所以限制给出同态
\[
\operatorname{Gal}(E/F)\hookrightarrow\operatorname{Gal}(L/K),
\]
如果 \(\sigma\) 的限制为恒等，它既固定 \(L\) 又固定 \(F\)，便固定这两个域生成的 \(E=LF\)，所以核平凡。由这个单同态及\cref{b1:thm:solvable-closure}，\(H=\operatorname{Gal}(E/F)\) 可解，且其阶整除 \(m\)。

由\cref{b1:cor:solvable-composition-factors}，有限可解群 \(H\) 有素数阶循环因子的次正规列。其产生过程是：先由\cref{b1:prop:derived-quotient}取交换因子的导出列，再在每个非平凡有限交换因子中取极大真子群，并对阶较小的子群继续做同样的选择。交换群的子群都正规，极大性使所得商为单群；由\cref{b1:prop:abelian-simple}，每个商都是素数阶循环群。通过商映射取这些子群的原像，就把导出列细化为
\[
H=H_0\trianglerighteq H_1\trianglerighteq\cdots\trianglerighteq H_s=1,
\qquad H_{i-1}/H_i\cong C_{\ell_i}.
\]
\cref{b1:thm:finite-galois-correspondence}给出反向域塔 \(F=E_0\subseteq E_1\subseteq\cdots\subseteq E_s=E\)，其中 \(E_i=E^{H_i}\)。相邻的包含 \(H_i\subseteq H_{i-1}\) 因而对应 \(E_{i-1}=E^{H_{i-1}}\subseteq E^{H_i}=E_i\)。为确定这一步的 Galois 群，把基域换成 \(E_{i-1}\)：\(E/E_{i-1}\) 仍是有限 Galois 扩张，且
\[
\operatorname{Gal}(E/E_{i-1})=H_{i-1}.
\]
在这个群中，\(H_i\trianglelefteq H_{i-1}\)，其不动域为 \(E_i\)，所以\cref{b1:thm:galois-quotient}给出 \(E_i/E_{i-1}\) 为 Galois 扩张，且
\[
\operatorname{Gal}(E_i/E_{i-1})\cong H_{i-1}/H_i\cong C_{\ell_i}.
\]
这里所需的正规性是 \(H_i\trianglelefteq H_{i-1}\)，并不要求 \(H_i\trianglelefteq H\)：每一步 \(E_i/E_{i-1}\) 都是循环 Galois 扩张，而 \(E_i/F\) 可以不正规。素数 \(\ell_i\) 整除 \(|H|\)，\(|H|\) 又整除 \(m\)，故 \(\ell_i\mid m\)。因此 \(F\) 及所有 \(E_{i-1}\) 都含 \(\mu_{\ell_i}\)，特征零又使 \(\ell_i\) 在这些域中非零。对每一步应用\cref{b1:thm:kummer-cyclic}，得到
\[
E_i=E_{i-1}(\gamma_i),\qquad \gamma_i^{\ell_i}\in E_{i-1}.
\]
把开头的 \(K\subseteq F\) 接上这些步骤，便得到包含 \(L\) 全部元素、从而包含 \(f\) 全部根的根式塔。
\end{proof}

\begin{corollary}\label{b1:cor:irreducible-one-radical-root}
设 \(K\) 为特征零域，\(f\in K[X]\) 为不可约多项式。若 \(f\) 的一个根属于某条从 \(K\) 出发的根式塔的末端，则 \(f\) 根式可解。
\end{corollary}
\begin{proof}
由\cref{b1:lem:radical-normal-envelope}，可将该根式塔置于有限 Galois 扩张 \(M/K\) 中，且 \(\operatorname{Gal}(M/K)\) 可解。由于 \(f\) 不可约，它与所给根的最小多项式只相差一个非零常数因子。\(M/K\) 正规，故这个最小多项式的全部共轭根都在 \(M\) 中，即 \(f\) 在 \(M\) 中完全分裂。其分裂域 \(L\) 是 \(M\) 的中间域；由\cref{b1:thm:galois-quotient}，\(\operatorname{Gal}(L/K)\) 是 \(\operatorname{Gal}(M/K)\) 的商群，也可解。再由\cref{b1:thm:radical-solvability}，\(f\) 根式可解。
\end{proof}

反向蕴含从根式可解的定义直接得到。因此，对特征零域上的不可约多项式，一个根能否用根式表示与全部根能否用根式表示是同一个存在性问题。不可约条件不能删去：若 \(f\) 可约，已知的一个根只约束它所属的那个不可约因子，不能据此判断其余因子的根式可解性。

\ProseParagraphBreak
可解扩张不一定自身是根式扩张。取本原七次单位根 \(\zeta_7=e^{2\pi i/7}\)，并令
\[
E=\mathbb Q(\zeta_7+\zeta_7^{-1})\subseteq\mathbb R.
\]
由\cref{b1:thm:cyclotomic-galois-group}，\(\operatorname{Gal}(\mathbb Q(\zeta_7)/\mathbb Q)\cong(\mathbb Z/7\mathbb Z)^\times\cong C_6\)，其中 \(3\bmod7\) 的阶为 \(6\)。复共轭的不动域是 \(E\)，它生成的二阶子群正规于 \(C_6\)，所以\cref{b1:thm:galois-quotient}给出 \(E/\mathbb Q\) 为 Galois 扩张，且 \(\operatorname{Gal}(E/\mathbb Q)\cong C_3\)。因此 \(E/\mathbb Q\) 是可解扩张，但以下论证将说明：开方若始终限制在这个实域 \(E\) 内，就不能生成 \(E\)。

\ProseParagraphBreak
若 \(0\neq\beta\in E\)，且对某个整数 \(m\geq2\) 有 \(\beta^m\in\mathbb Q\)，则对任意 \(\sigma\in\operatorname{Gal}(E/\mathbb Q)\)，比值 \(\sigma(\beta)/\beta\) 是实数域 \(E\) 中的 \(m\) 次单位根，因而属于 \(\{1,-1\}\subseteq\mathbb Q\)。这些比值被每个自同构固定，所以与前面根式扩张中的比值计算一样，\(\sigma\mapsto\sigma(\beta)/\beta\) 是到 \(\{1,-1\}\) 的群同态。它的像的阶必须同时整除 \(3\) 和 \(2\)，故只能为平凡群。这说明全部自同构都固定 \(\beta\)，由 Galois 对应得到 \(\beta\in\mathbb Q\)。

\ProseParagraphBreak
若某条根式塔从 \(\mathbb Q\) 出发而以 \(E\) 为末端，则各层都在 \(E\) 内，第一步非平凡扩张须由这样的 \(\beta\in E\setminus\mathbb Q\) 生成，矛盾。所以 \(E/\mathbb Q\) 本身不是根式扩张。另一方面，\(\zeta_7+\zeta_7^{-1}\) 生成 \(E\)，而 \(E/\mathbb Q\) 正规，故它的最小多项式的全部根都在 \(E\) 内，并生成 \(E\)。这个多项式的分裂域便是 \(E\)，由\cref{b1:thm:radical-solvability}，它根式可解。因此 \(E\) 仍包含于某条从 \(\mathbb Q\) 出发的根式塔的末端，只是那个末端域必须严格大于 \(E\)。

\subsection{特征条件与可分性}
\label{b1:subsec:2102:03}

特征零保证每个不可约多项式可分，并允许所取根式指数作为非零域元素使用。\(f\) 本身可以含重因子；删去重数既不改变分裂域，也不改变根式可解性。因此定理并不另要求 \(f\) 的各因子互异。

\ProseParagraphBreak
在特征 \(p\) 中，若只允许指数与 \(p\) 互素的根式，\cref{b1:lem:radical-normal-envelope}的证明仍能构造可解的正规包络。但循环 \(p\) 次扩张由 \(X^p-X-a\) 产生，不能调用\cref{b1:thm:kummer-cyclic}把这一步变为 \(p\) 次开方。另一方面，\(X^p-a\) 的根给出纯不可分扩张；纯不可分扩张的每个元素在基域上的共轭只有自身，所以其基域自同构群必为平凡群。当扩张非平凡时，它不是有限 Galois 扩张，不能用有限 Galois 对应来分析。正特征下需要把 Kummer 步骤、Artin--Schreier 步骤和纯不可分步骤分别处理。

\ProseParagraphBreak
可解群的次正规列仍能指导正特征下的构造：若某个素数阶因子等于特征，使用\cref{b1:thm:artin-schreier-classification}。得到的是允许 Kummer 和 Artin--Schreier 两类步骤的扩张塔，与只允许开方的根式塔应当区分。

% !TeX root = ../../Book1.tex
% 纲要标题：### 21.3 经典方程
% 纲要位置：计划纲要.md，第 752 行。

\section{经典方程}
\label{b1:sec:2103}

二次、三次和四次方程的根式公式并非互不相关的技巧。它们背后的 Galois 群都嵌在可解的低阶对称群中；五次开始，对称群的可解性失效。不过，“一般方程没有通用根式公式”不表示每个具体五次方程都不能用根式求解。

% 纲要内容（仅作写作计划，尚非教材正文）：
%
% * 三次、四次方程的群论解释；从三次根的循环作用理解配对公式，从四根的三种配对构造辅助三次方程并联系 Ferrari 配方
% * 一般五次方程不可根式求解
% * 一个具体不可解五次方程：以 \(x^5-x-1\) 的模素数分解判定 Galois 群，引用第19.3节已证的循环型判据
% * 可解但非交换的 Galois 群：回引三次分裂域的群计算，解释正规子群链、单位根和根式塔
%

\subsection{三次、四次方程的群论解释}
\label{b1:subsec:2103:01}

设基域特征为零。若次数为 \(d\) 的多项式有 \(s\) 个不同的根，则其 Galois 群忠实作用在这 \(s\) 个根组成的集合上，因而嵌入 \(S_s\)，再将 \(S_s\) 看成 \(S_d\) 的子群即可。这里置换的是不同的根；即使按重数列出 \(d\) 个根，也不能据此增加域自同构。式\eqref{b1:eq:small-symmetric-derived}已经给出
\[
S_3\trianglerighteq A_3\trianglerighteq1,
\qquad S_4\trianglerighteq A_4\trianglerighteq V_4\trianglerighteq1
\]
都有交换因子。\cref{b1:thm:solvable-closure}把可解性传给子群，再由\cref{b1:thm:radical-solvability}得到三次、四次方程都根式可解。

\ProseParagraphBreak
可解性判据说明根式表达存在，却没有直接给出求根的步骤。三次方程的 Cardano 公式可以从根的置换中构造出来。将一般三次方程除以首项系数，再平移变量消去二次项，得到 \(y^3+py+q=0\)。设其三个根为 \(r_1,r_2,r_3\)，按重数列出，于是 \(r_1+r_2+r_3=0\)。先加入本原三次单位根 \(\omega\)。我们希望选取根的线性组合，使三轮换的作用成为乘以 \(\omega\) 或 \(\omega^2\)，从而把循环作用化为特征向量上的标量作用。下面两组系数正是为此选择的：
\[
u=\frac{r_1+\omega^2r_2+\omega r_3}{3},\qquad
v=\frac{r_1+\omega r_2+\omega^2r_3}{3}.
\]
由 \(1+\omega+\omega^2=0\)，反过来有
\[
r_1=u+v,\qquad r_2=\omega u+\omega^2v,\qquad
r_3=\omega^2u+\omega v.
\]
在固定 \(\omega\) 的条件下，循环改标 \((r_1,r_2,r_3)\mapsto(r_2,r_3,r_1)\) 使 \((u,v)\mapsto(\omega u,\omega^2v)\)，而交换 \(r_2,r_3\) 使 \(u,v\) 交换。因此 \(u^3,v^3\) 不受三轮换影响，却仍可能互换。这里的改标首先是对根的形式重排，并没有断言某个具体方程一定存在实现它的域自同构。对无重根方程，加入 \(\omega\) 后的 Galois 群中实际出现的置换才按这些公式作用；这个群可以严格小于 \(S_3\)。就上述形式置换而言，商群 \(S_3/A_3\cong C_2\) 记录 \(u^3,v^3\) 的交换，因而先用至多一次开平方将这两个量区分；随后 \(A_3\cong C_3\) 在 \(u,v\) 上分别乘以 \(\omega,\omega^2\)，对应于开立方恢复它们。

\ProseParagraphBreak
Vieta 公式给出 \(r_1+r_2+r_3=0\) 与 \(r_1r_2+r_1r_3+r_2r_3=p\)，所以
\[
r_1^2+r_2^2+r_3^2
=(r_1+r_2+r_3)^2-2(r_1r_2+r_1r_3+r_2r_3)=-2p.
\]
展开 \(u,v\) 的乘积并使用 \(1+\omega+\omega^2=0\)，便得
\[
9uv=r_1^2+r_2^2+r_3^2-(r_1r_2+r_1r_3+r_2r_3)
=-2p-p=-3p,
\]
所以 \(uv=-p/3\)，又有 \(u^3+v^3=r_1^3-3uvr_1=r_1^3+pr_1=-q\)。故 \(u^3,v^3\) 是二次方程
\[
T^2+qT-\frac{p^3}{27}=0
\]
的两个根。这就解释了 Cardano 公式中的平方根为何先于立方根出现。当 \(p=0\) 时直接开立方即可。当 \(p\neq0\) 时，设
\[
\Delta=\frac{q^2}{4}+\frac{p^3}{27},\qquad
U=-\frac q2+\sqrt\Delta,\qquad V=-\frac q2-\sqrt\Delta.
\]
有 \(U+V=-q\)、\(UV=-p^3/27\neq0\)。任取 \(u^3=U\)，再按乘积条件定义 \(v=-p/(3u)\)，则 \(v^3=V\)。于是
\[
(u+v)^3+p(u+v)+q=u^3+v^3+(3uv+p)(u+v)+q=0.
\]
利用已加入的 \(\omega\)，三个根为 \(\omega^j u+\omega^{-j}v\)、\(j=0,1,2\)，重根情形按重数理解。不能独立选取两个立方根后直接相加；\(uv=-p/3\) 是公式成立的必要配对条件。

\ProseParagraphBreak
实系数方程的计算途中也可能出现复数。例如 \(y^3-3y+1=0\) 的 \(\Delta=1/4-1=-3/4\)，选取复平方根后有 \(U=-1/2+i\sqrt3/2\)。因为 \(|U|=1\)，任一立方根 \(u\) 也满足 \(|u|=1\)，配对条件给出 \(v=1/u=\bar u\)。又因 \(\omega^{-j}=\overline{\omega^j}\)，三个表达式 \(\omega^j u+\omega^{-j}v\) 全为实数，计算却经过了复数中间量。这正说明根式可解允许使用复数中间量，并不要求每次开方都留在实域中。

\ProseParagraphBreak
四次方程也能通过一个辅助三次方程归结为二次方程。经平移后，它成为 \(y^4+py^2+qy+r=0\)。若 \(q=0\)，先将 \(y^2\) 作为未知量解二次方程。若 \(q\neq0\)，希望把四次式写成两个平方之差，从而分解为两个二次式。引入待定参数 \(m\)，有
\[
y^4+py^2+qy+r
=(y^2+m)^2-\bigl((2m-p)y^2-qy+m^2-r\bigr).
\]
要使括号中的二次式成为 \(\bigl(sy-q/(2s)\bigr)^2\)，比较二次项与常数项就要求
\[
s^2=2m-p,\qquad m^2-r=\frac{q^2}{4s^2}.
\]
消去 \(s\)，得到关于 \(m\) 的三次方程
\[
4(2m-p)(m^2-r)=q^2
\]
取它的任一根 \(m\)，因为 \(q\neq0\)，有 \(2m-p\neq0\)，可选取非零 \(s\) 满足 \(s^2=2m-p\)。三次方程保证上述常数项条件也成立，于是
\begin{align*}
y^4+py^2+qy+r
&=(y^2+m)^2-\Bigl(sy-\frac q{2s}\Bigr)^2\\
&=\Bigl(y^2-sy+m+\frac q{2s}\Bigr)
  \Bigl(y^2+sy+m-\frac q{2s}\Bigr).
\end{align*}
后面只剩两个二次方程。这说明四次求解为何能够归结为三次求解及连续开平方。

\ProseParagraphBreak
上述配方给出了两个二次因子，还可以从根的配对解释辅助三次方程的来源。设
\[
f(X)=X^4+bX^3+cX^2+dX+e
=\prod_{i=1}^4(X-\alpha_i),
\]
把四个根的指标分成两组，每组两个；不区分组内顺序，也不区分两组的先后。这三种无序划分为 \(12\mid34\)、\(13\mid24\)、\(14\mid23\)，分别对应下面三个配对量：
\[
\begin{aligned}
\beta_1&=\alpha_1\alpha_2+\alpha_3\alpha_4,\\
\beta_2&=\alpha_1\alpha_3+\alpha_2\alpha_4,\\
\beta_3&=\alpha_1\alpha_4+\alpha_2\alpha_3.
\end{aligned}
\]
若 \(e_i\) 表示四个根的第 \(i\) 个初等对称多项式，则 \(e_1=-b\)、\(e_2=c\)、\(e_3=-d\)、\(e_4=e\)。要找一个以三个 \(\beta_j\) 为根的基域多项式，需要把它们的初等对称式改写为 \(e_i\)。其中两两乘积之和展开后，恰含十二个形如 \(\alpha_i^2\alpha_j\alpha_k\)、\(i,j,k\) 互异的单项式，各出现一次；\(e_1e_3\) 包含同样的十二项，并另有四项 \(e_4\)，所以两者相差 \(4e_4\)。三重乘积的八项则可分成两类：
\[
\beta_1\beta_2\beta_3
=e_4\sum_{i=1}^4\alpha_i^2
 +\sum_{1\le i<j<k\le4}(\alpha_i\alpha_j\alpha_k)^2.
\]
第一个和为 \(e_1^2-2e_2\)；第二个和为 \(e_3^2-2e_2e_4\)，因为四个三重乘积的六个两两乘积之和恰为 \(e_2e_4\)。于是得到
\begin{align*}
\beta_1+\beta_2+\beta_3&=e_2=c,\\
\beta_1\beta_2+\beta_1\beta_3+\beta_2\beta_3
&=e_1e_3-4e_4=bd-4e,\\
\beta_1\beta_2\beta_3&=e_1^2e_4+e_3^2-4e_2e_4
=b^2e+d^2-4ce.
\end{align*}
因此它们满足同一个三次辅助方程，其多项式为
\begin{equation}\label{b1:eq:quartic-cubic-resolvent}
\begin{aligned}
R(T)&\coloneqq\prod_{j=1}^3(T-\beta_j)\\
&=T^3-cT^2+(bd-4e)T+4ce-b^2e-d^2.
\end{aligned}
\end{equation}
置换四个根就置换这三个量，这正是\subsecref{b1:subsec:0404:03}中三种二元划分的作用。它的核为
\[
V_4=\{1,(12)(34),(13)(24),(14)(23)\}.
\]
这些置换保持每一种无序划分的整体，允许交换组内的根，也允许交换两组。若四根互异，由
\[
\beta_1-\beta_2=(\alpha_1-\alpha_4)(\alpha_2-\alpha_3)
\]
及另两式可知三个 \(\beta_j\) 也互异。因此保持全部三种划分等价于固定全部三个配对量，\(V_4\) 在这里有了具体的域论含义。仍在四根互异的条件下，原四次方程的 Galois 群 \(G\leq S_4\) 在这些量上的作用核就是 \(G\cap V_4\)。

\ProseParagraphBreak
现在将配对量与配方中的参数对应起来。对前面的 \(y^4+py^2+qy+r\)，有 \(b=0\)、\(c=p\)、\(d=q\)、\(e=r\)，故
\[
R(2m)=4(2m-p)(m^2-r)-q^2.
\]
辅助方程使 \(R(2m)=0\)，而 \(R\) 的根正是三个 \(\beta_j\)，故 \(2m=\beta_j\) 对某个 \(j\) 成立。重新标号，把该配对写为 \(\{\alpha_1,\alpha_2\}\mid\{\alpha_3,\alpha_4\}\)，于是 \(2m=\alpha_1\alpha_2+\alpha_3\alpha_4\)。暂令 \(s_0=\alpha_1+\alpha_2=-(\alpha_3+\alpha_4)\)。四个交叉乘积之和为 \((\alpha_1+\alpha_2)(\alpha_3+\alpha_4)=-s_0^2\)，故 Vieta 公式给出
\[
p=\alpha_1\alpha_2+\alpha_3\alpha_4
  +(\alpha_1+\alpha_2)(\alpha_3+\alpha_4)=2m-s_0^2.
\]
再比较两组根组成的两个二次因子的一次项系数，得到
\[
s_0^2=2m-p,\qquad
q=s_0(\alpha_1\alpha_2-\alpha_3\alpha_4),\qquad
\alpha_1\alpha_2+\alpha_3\alpha_4=2m.
\]
因此 \(s_0=\pm s\)；若 \(s_0=-s\)，交换两组根后便有 \(s_0=s\)。当 \(q\neq0\) 时，\(s\neq0\)，两组根的积于是分别为 \(m+q/(2s)\) 与 \(m-q/(2s)\)，恰好得到上面两个二次因子。辅助三次方程确定配对量，开平方确定两组根的和，最后各解一个二次方程；群作用与求根计算在此描述的是同一过程。

\subsection{一般五次方程不可根式求解}
\label{b1:subsec:2103:02}

下面的“一般”有一个精确含义：把系数当作代数独立的变量，而不是先给它们某组具体数值。

\begin{theorem}[一般方程]\label{b1:thm:generic-equation}
设 \(n\geq1\) 为整数，并设 \(u_1,\ldots,u_n\) 在 \(\mathbb Q\) 上代数独立。多项式
\[
X^n-u_1X^{n-1}+u_2X^{n-2}-\cdots+(-1)^nu_n
\]
在 \(\mathbb Q(u_1,\ldots,u_n)\) 上的 Galois 群是 \(S_n\)。因此当 \(n\geq5\) 时，一般 \(n\) 次方程不根式可解。
\end{theorem}
\begin{proof}
先取代数独立的 \(t_1,\ldots,t_n\)，令 \(E=\mathbb Q(t_1,\ldots,t_n)\)，\(e_j\) 为这些变量的第 \(j\) 个初等对称多项式，并令 \(F=\mathbb Q(e_1,\ldots,e_n)\)。于是
\[
\prod_{i=1}^n(X-t_i)=X^n-e_1X^{n-1}+\cdots+(-1)^ne_n.
\]
每个 \(t_i\) 都是右端多项式的根，因而在 \(F\) 上代数。依次加入 \(t_1,t_2,\ldots,t_n\)，在已知一个根后除去对应线性因子，下一根满足的多项式次数就降低一。所以各步次数分别至多为 \(n,n-1,\ldots,1\)，由\cref{b1:thm:tower-law}得到 \([E:F]\leq n!\)，特别地，扩张有限。各 \(t_i\) 互异，故 \(E/F\) 是右端可分多项式的分裂域，为 Galois 扩张。

变量的任意置换保持多项式环的加法与乘法，并延拓为分式域 \(E\) 的自同构。每个 \(e_j\) 都是对称多项式，所以这些置换固定全部 \(e_j\)，也就逐点固定它们生成的系数域 \(F\)。不同置换在某个 \(t_i\) 上取值不同，给出 \(n!\) 个不同的 \(F\) 自同构。\cref{b1:thm:separable-embeddings}的嵌入数上界于是给出
\[
n!\leq|\operatorname{Aut}_F(E)|\leq[E:F]\leq n!.
\]
全部等号成立，所有自同构恰为这些变量置换，因此 Galois 群就是 \(S_n\)。

还需证明这些初等对称多项式代数独立。若 \(H\in\mathbb Q[U_1,\ldots,U_n]\) 满足 \(H(e_1,\ldots,e_n)=0\)，因 \(t_1,\ldots,t_n\) 代数独立，这就是关于 \(t_i\) 的多项式恒等式。\cref{b1:thm:fundamental-symmetric-polynomials}的唯一性具体说：一个对称多项式写成初等对称多项式的多项式表达时，所用的表达式唯一。按 \(t_1>\cdots>t_n\) 的词典序，\(e_1^{a_1}\cdots e_n^{a_n}\) 的首项指数依次为 \(a_1+\cdots+a_n,a_2+\cdots+a_n,\ldots,a_n\)，依次作差便恢复各 \(a_i\)，所以不同单项式的首项也不同。现在零多项式同时由 \(H\) 与零表达，所以 \(H=0\)，证明了 \(e_1,\ldots,e_n\) 的代数独立性。由 \(u_i\mapsto e_i\) 得到系数域 \(\mathbb Q(u_1,\ldots,u_n)\cong F\) 的同构，并把题述多项式送到上面已经考察的多项式，因此两者的 Galois 群同构。

当 \(n\geq5\) 时，\cref{b1:cor:symmetric-nonsolvable}说明 \(S_n\) 不可解；\cref{b1:thm:radical-solvability}遂排除根式求解。
\end{proof}

这就是\term[Abel-Ruffini-dingli]{Abel--Ruffini 定理}[Abel--Ruffini theorem]的通用方程形式。这里的“一般”始终指系数域 \(\mathbb Q(u_1,\ldots,u_n)\) 上那一个以代数独立变量为系数的多项式，不是关于随机选取数值系数的概率陈述。它排除的是以这些独立系数为起点的通用根式表达，并不排除特殊方程。例如 \(X^5-1\) 的分裂域是分圆域，Galois 群 \((\mathbb Z/5\mathbb Z)^\times\cong C_4\) 可解，因而它的全部根都能用根式表示。


\subsection{一个不可根式求解的五次方程}
\label{b1:subsec:2103:04}

一般方程的结论不能直接判定某一个给定方程。现在考察序言中的 \(X^5-X-1\)，用\cref{b1:lem:reduction-cycle-type}把模素数分解得到的循环型与群的传递性结合起来，确定它的 Galois 群。


\begin{proposition}\label{b1:prop:quintic-x5-x-1}
多项式 \(f=X^5-X-1\) 在 \(\mathbb Q\) 上的 Galois 群为 \(S_5\)，因而不根式可解。
\end{proposition}
\begin{proof}
先证明不可约性。五次多项式若可约，就至少有两个不可约因子；若每个因子的次数都不小于三，总次数至少为六，矛盾。因此只须排除一次与不可约二次因子。模三时，\(f\) 在 \(0,1,2\) 处的值都为二，没有一次因子。\(\mathbb F_3\) 上的首一不可约二次多项式只有
\(X^2+1\)、\(X^2+X+2\)、\(X^2+2X+2\)：对常数项非零的六个首一二次式逐一检查是否在 \(0,1,2\) 处有根，即得此表。用它们分别除 \(\bar f\)，余式依次为
\[
2,\qquad X+2,\qquad X+2.
\]
所以也没有不可约二次因子，\(\bar f\) 不可约；\cref{b1:prop:reduction-irreducible}说明 \(f\) 在 \(\mathbb Q\) 上不可约。

模二时，直接相乘可核验
\[
\bar f=(X^2+X+1)(X^3+X^2+1).
\]
两个因子在 \(\mathbb F_2\) 中均无根，次数分别为二和三，因而均不可约且互异，所以约化无重因子。有限域上的不可约多项式都可分，故约化也无重根，满足\cref{b1:lem:reduction-cycle-type}的条件。该引理给出群中的一个 \((2)(3)\) 型置换 \(\tau=(a\,b)(d\,e\,h)\)。这两个循环不相交，彼此交换，故
\[
\tau^3=(a\,b)^3(d\,e\,h)^3=(a\,b),
\]
得到群中的一个换位。

在分裂域中，若 \(a,b\) 是同一不可约因子的两个根，\cref{b1:lem:simple-embedding-extension}对基域恒等映射给出 \(K(a)\to K(b)\)、\(a\mapsto b\) 的同构；\cref{b1:thm:embedding-extension}延拓这个同构，正规性使得到的嵌入成为分裂域的自同构。因此一个不可约因子的全部根恰组成一个轨道。现在 \(f\) 不可约，故五个根只有一个轨道，群的作用传递。再由\cref{b1:thm:orbit-stabilizer}，五整除群阶。\cref{b1:thm:cauchy}给出五阶元素，在五个点上它必为五轮换 \(c\)。

将五个点按 \(c\) 标为 \(\mathbb Z/5\mathbb Z\)，使 \(c\) 的作用为加一。设刚找到的换位交换 \(i,j\)，记非零差为 \(d=j-i\)。用 \(c\) 的各次幂共轭这个换位，就得到所有 \((k\,\,k+d)\)。由于五为素数，从任一点反复加 \(d\)，依次经过全部五个点；所以任意两点都可用这些相邻关系连成一条不重复顶点的路径。设路径为 \(v_0,\ldots,v_r\)，它的各相邻换位都在群中。由第一个换位开始，若已经得到 \((v_0\,v_j)\)，就有
\[
(v_0\,v_j)(v_j\,v_{j+1})(v_0\,v_j)=(v_0\,v_{j+1})
\qquad(1\le j<r).
\]
沿路径归纳，得到两端的换位 \((v_0\,v_r)\)。任意两点均可作两端，故群包含全部换位；换位生成 \(S_5\)，于是原 Galois 群等于 \(S_5\)。由\cref{b1:cor:symmetric-nonsolvable}及\cref{b1:thm:radical-solvability}，方程不根式可解。
\end{proof}
这个具体判定解释了序言中 \(X^5-1\) 与 \(X^5-X-1\) 的差别：前者有交换的分圆 Galois 群，后者的群为不可解的 \(S_5\)。次数同为五，并不意味着两者有相同的根式可解性。

\ProseParagraphBreak
对可约多项式，能求出一个根也不意味着全部根可用根式求出。例如 \((X-1)(X^5-X-1)\) 有有理根 \(1\)，但它的五次因子已由上面的命题证明不根式可解。

\subsection{可解但非交换的 Galois 群}
\label{b1:subsec:2103:03}

回到 \(X^3-2\)。令 \(a=\sqrt[3]{2}\)、\(\zeta=\zeta_3\)、\(L=\mathbb Q(a,\zeta)\)。\cref{b1:prop:cubic-splitting-field-correspondence}已确定 \([L:\mathbb Q]=6\)、\(\operatorname{Gal}(L/\mathbb Q)\cong S_3\) 及其全部中间域。这里利用这些结论，解释根式求解怎样反映在子群与域塔上。

\ProseParagraphBreak
由 Galois 对应，正规子群链和相应的域塔为
\[
S_3\trianglerighteq A_3\trianglerighteq1,
\qquad
\mathbb Q\subseteq\mathbb Q(\zeta)\subseteq L.
\]
两步扩张的 Galois 群分别是 \(S_3/A_3\cong C_2\) 与 \(A_3\cong C_3\)。第一步通过
\[
\zeta=\frac{-1+\sqrt{-3}}2,
\qquad \mathbb Q(\zeta)=\mathbb Q(\sqrt{-3})
\]
开平方完成，反向关系 \(\sqrt{-3}=2\zeta+1\) 也说明两域相等。这一步使用复平方根，已离开实域，仍属于允许的根式运算。第二步加入 \(a\)，其立方 \(a^3=2\) 已在基域中。于是这条对应于正规子群链的域塔本身就是一条根式塔。

\ProseParagraphBreak
先加入单位根，使第二步也能由 Galois 理论描述：在 \(\mathbb Q(\zeta)\) 上，加入 \(a\) 便同时包含 \(a,\zeta a,\zeta^2a\) 三个共轭根，循环自同构把 \(a\) 送到 \(\zeta a\)。若先走 \(\mathbb Q\subseteq\mathbb Q(a)\)，虽然已经通过开立方取得一个根，这一步却不正规，因为另外两个根不在实域 \(\mathbb Q(a)\) 中。

\ProseParagraphBreak
若 \(\rho(a)=\zeta a\)、\(\rho(\zeta)=\zeta\)，而 \(\kappa\) 为复共轭，则 \(\kappa\rho(a)=\zeta^2a\)、\(\rho\kappa(a)=\zeta a\)，两者不同，故全群 \(S_3\) 非交换。但上述两层商群 \(S_3/A_3\cong C_2\) 与 \(A_3\cong C_3\) 都是交换群，所以 \(S_3\) 仍然可解。交换群本身总可解，而这个例子说明可解群可以非交换；根式求解所需要的是逐层出现交换的商群，不要求整个 Galois 群交换。

% !TeX root = ../../Book1.tex
% 纲要标题：### 21.4 尺规作图
% 纲要位置：计划纲要.md，第 724 行。

\section{尺规作图}
\label{b1:sec:2104}

尺规作图限制了允许执行的操作：直尺只能连接已知点，圆规只能以已知点为圆心、以已知线段为半径画圆。把点写成坐标以后，直线与圆的交点只要求解一次或二次方程。因此，这种几何限制可以转化为扩张次数的限制。

% 纲要内容（仅作写作计划，尚非教材正文）：
%
% * 二次扩张塔与规矩数
% * 倍立方体、一般三等分角与正多边形
% * 化圆为方与 π 的超越性；超越性论证给出证明概要
%
% ---
%

\subsection{二次扩张塔与规矩数}
\label{b1:subsec:2104:01}

\begin{definition}\label{b1:def:constructible-number}
从平面上的 \((0,0)\)、\((1,0)\) 出发，允许有限次画过两个不同已知点的直线、以已知点为圆心且以两个已知点间距离为半径的圆，并把任意两条所画曲线的孤立交点加入已知点。重合曲线不由此指定新的点。如果一个点能经过有限次这种操作得到，则称它为\term[keguitu-dian]{尺规可作图点}[constructible point]。给定实数 \(a\)，如果 \((a,0)\) 是尺规可作图点，则称 \(a\) 为\term[guiju-shu]{规矩数}[constructible number]，也称可构造数、可作图数。\footnote{北京交通大学《近世代数》课程使用“规矩数”\cite[第 23 章第 1 节“几何作图简史、规矩数”]{BJTU2026ModernAlgebraCourse}；“可构造数”见席南华《基础代数》第一卷\cite[p. 125]{Xi2016BasicAlgebraI}，该书从复平面讨论可构造点。}\termidx[kegouzao-shu]{可构造数}\termidx[kezuotu-shu]{可作图数}
\end{definition}

标准的作垂线、平行线与转移线段长度均由上述操作实现。因此一个点可作图当且仅当它的两个坐标都是规矩数：从点向坐标轴作垂线得到坐标，反过来在两轴上标出坐标再作平行线，交点即为所需点。

\begin{theorem}\label{b1:thm:constructible-quadratic-tower}
实数 \(a\) 可作图，当且仅当存在全部位于 \(\mathbb R\) 中的域塔
\[
\mathbb Q=K_0\subseteq K_1\subseteq\cdots\subseteq K_r,
\qquad [K_i:K_{i-1}]=2,
\]
使 \(a\in K_r\)。允许 \(r=0\)。特别地，规矩数在 \(\mathbb Q\) 上的次数必为 \(2\) 的幂。
\end{theorem}
\begin{proof}
若当前全部已知坐标在实域 \(F\) 中，过已知点的直线方程系数属于 \(F\)。以已知点为圆心、已知线段为半径的圆方程是
\[
(x-u)^2+(y-v)^2=r^2,
\qquad u,v,r^2\in F.
\]
两直线的孤立交点由一次方程求出，坐标仍在 \(F\)。一条直线与一个圆联立，消去一个变量便得到至多二次方程；求出该变量后，另一个变量是它的 \(F\) 线性表达式。两圆相减得到一次方程，故也归结为前一种情况。重合的曲线不指定新的孤立交点；无交点时没有可执行步骤。因此每次新添交点至多需要一个实二次扩张。沿有限作图过程依次加入坐标，删去次数为 \(1\) 的步骤，就得到所需塔。

反过来，可作图的有向线段长度对加、减封闭，乘、除由相似三角形实现。例如连结 \((1,0)\) 与 \((0,b)\)，过 \((a,0)\) 作平行线，其纵轴截距为 \(ab\)。若 \(a\neq0\)，改为连结 \((a,0)\) 与 \((0,b)\)，过 \((1,0)\) 作平行线，纵轴截距便是 \(b/a\)。这些操作连同符号方向使规矩数构成一个实子域。

若 \(u>0\) 已可作图，在横轴上取 \(A=(0,0)\)、\(C=(1,0)\)、\(B=(1+u,0)\)，画以 \(AB\) 为直径的半圆，并过 \(C\) 作垂线，与上半圆交于 \(P=(1,h)\)，如\cref{b1:fig:square-root-semicircle}所示。圆心为 \(AB\) 的中点，半径为 \(AB\) 的一半，故
\[
\left(1-\frac{1+u}{2}\right)^2+h^2
=\left(\frac{1+u}{2}\right)^2,
\qquad h^2=u.
\]
取上半圆使 \(h>0\)，便作出了 \(\sqrt u\)。

% !TeX root = ../Book1.tex
\begin{figure}[htbp]
\centering
\begin{tikzpicture}[scale=1.35]
  % The proportions illustrate u=9/4; the construction works for every u>0.
  \pgfmathsetmacro{\diagramu}{2.25}
  \pgfmathsetmacro{\diagramr}{(1+\diagramu)/2}
  \pgfmathsetmacro{\diagramh}{sqrt(\diagramu)}
  \coordinate (A) at (0,0);
  \coordinate (C) at (1,0);
  \coordinate (B) at ({1+\diagramu},0);
  \coordinate (O) at (\diagramr,0);
  \coordinate (P) at (1,\diagramh);
  \draw[thick] (A)--(B);
  \draw[thick,color=AlgebraBlue] (A) arc[start angle=180,end angle=0,radius=\diagramr];
  \draw[thick,color=AlgebraBlue] (C)--(P);
  \draw[gray,dashed] (O)--(P);
  \draw (1,0.14)--(1.14,0.14)--(1.14,0);
  \foreach \point in {A,B,C,O,P} {\fill (\point) circle (1.2pt);}
  \node[below left] at (A) {\(A\)};
  \node[below] at (C) {\(C\)};
  \node[below right] at (B) {\(B\)};
  \node[below right] at (O) {\(O\)};
  \node[above left] at (P) {\(P\)};
  \node[below] at (0.5,-0.29) {\(1\)};
  \node[below] at ({1+\diagramu/2},-0.29) {\(u\)};
  \node[left,color=AlgebraBlue] at (0.92,{\diagramh/2}) {\(h\)};
\end{tikzpicture}
\caption[半圆开平方作图]{半圆开平方作图。\(AC=1\)、\(CB=u>0\)，\(O\) 为 \(AB\) 的中点，过 \(C\) 的垂线与上半圆交于 \(P\)。高度 \(CP=h\) 满足 \(h^2=u\)。图形比例仅作示意。}
\label{b1:fig:square-root-semicircle}
\end{figure}


最后，对任意实二次扩张 \(F(\beta)/F\)，把 \(\beta\) 的首一最小多项式写成 \(X^2+bX+c\)，配方得到 \((2\beta+b)^2=b^2-4c\)。判别式属于 \(F\) 且为正，所以 \(\beta=(-b\pm\sqrt{b^2-4c})/2\) 可由域运算和开正平方根得到。对塔归纳，末端全部元素都可作图。由\cref{b1:thm:tower-law}，\([K_r:\mathbb Q]=2^r\)，而 \(\bigl[\mathbb Q(a):\mathbb Q\bigr]\) 整除它，故为 \(2\) 的幂。
\end{proof}

若从更多给定点出发，应把 \(\mathbb Q\) 换为这些点坐标生成的域，证明逐字适用。次数为 \(2\) 的幂只是一个必要条件；完整判据要求存在二次扩张塔，不能省掉“塔”而只检查单个次数。

\subsection{倍立方体、三等分角与正多边形}
\label{b1:subsec:2104:02}

倍立方体要求从单位棱长作出棱长 \(\sqrt[3]{2}\)。\cref{b1:thm:eisenstein}对素元 \(2\) 给出 \(X^3-2\) 在 \(\mathbb Q\) 上不可约，故该数次数为 \(3\)，由\cref{b1:thm:constructible-quadratic-tower}不可作图。

\ProseParagraphBreak
一般三等分角也不可能，因为已经可作出的 \(60^\circ\) 角不能被尺规三等分。若 \(20^\circ\) 可作图，则 \(x=2\cos20^\circ\) 可作图；三倍角恒等式给出 \(x^3-3x-1=0\)。由\cref{b1:prop:rational-root}，这个首一整系数三次多项式的有理根只能是 \(\pm1\)，代入都非零。三次多项式可约必有一次因子，因此它不可约，\(x\) 的次数为 \(3\)，矛盾。这里排除的是一种对任意角都适用的作图方法；例如 \(90^\circ\) 可以三等分，并不与结论冲突。

\ProseParagraphBreak
正多边形问题的完整回答是\term[Gauss-Wantzel-dingli]{Gauss--Wantzel 定理}[Gauss--Wantzel theorem]。\footnote{旺策尔（Pierre Laurent Wantzel，1814-1848），法国数学家，证明若干古典尺规作图问题不可能完成。}

\begin{theorem}[Gauss--Wantzel]\label{b1:thm:constructible-polygon}
设 \(n\geq3\) 为整数。正 \(n\) 边形可尺规作图，当且仅当 \(\varphi(n)\) 为 \(2\) 的幂；等价地，
\[
n=2^a p_1\cdots p_s,
\]
其中 \(a,s\geq0\) 为整数，各 \(p_i\) 是两两不同的奇素数，并且每个 \(p_i\) 都可写成 \(2^{2^{r_i}}+1\)，其中 \(r_i\geq0\) 为整数；这样的素数称为\term[Fermat-sushu]{Fermat 素数}[Fermat prime]。\footnote{费马（Pierre de Fermat，1601？-1665），法国数学家，研究数论与解析几何；其出生年份有不同记载。}
\end{theorem}
\begin{proof}
正 \(n\) 边形可作图等价于 \(\zeta_n\) 的实部和虚部可作图。若可作图，把这两个坐标所在的实二次塔合并，再加入 \(i\)，得到一个次数为 \(2\) 的幂的域，其中含有 \(\zeta_n\)。由\cref{b1:thm:cyclotomic-irreducible}及塔公式，\(\varphi(n)=\bigl[\mathbb Q(\zeta_n):\mathbb Q\bigr]\) 也是 \(2\) 的幂。

反过来，若 \(\varphi(n)\) 为 \(2\) 的幂，\cref{b1:thm:cyclotomic-galois-group}说明分圆域的 Galois 群是有限交换 \(2\) 群。复共轭子群的固定域是实子域 \(F=\mathbb Q(\zeta_n+\zeta_n^{-1})\)：这里 \(n\geq3\)，故 \(\zeta_n\) 不实，且它在 \(F\) 上满足 \(X^2-(\zeta_n+\zeta_n^{-1})X+1=0\)。复共轭子群在交换群中正规，所以\cref{b1:thm:galois-quotient}说明 \(F/\mathbb Q\) 也是 Galois 扩张，其群是有限交换 \(2\) 群。这样的群有一条各步指数均为 \(2\) 的子群列：对非平凡有限交换 \(2\) 群取极大真子群，商为交换单群，由\cref{b1:prop:abelian-simple}为 \(C_2\)，再对阶较小的子群归纳。由\cref{b1:thm:finite-galois-correspondence}，\(F\) 因而位于一条实二次扩张塔的末端。于是 \(\cos(2\pi/n)\) 可作图，\(\sin(2\pi/n)\) 也可由 \(\sqrt{1-\cos^2(2\pi/n)}\) 作出，正 \(n\) 边形便可作图。

最后设 \(n=2^a\prod p^{e_p}\)，其中 \(p\) 为奇素数。由模素数幂计数及中国剩余定理，
\[
\varphi(n)=\varphi(2^a)\prod_{p\mid n,\ p\text{ 奇}}p^{e_p-1}(p-1).
\]
它为 \(2\) 的幂当且仅当每个奇素数的指数 \(e_p=1\)，且 \(p-1\) 是 \(2\) 的幂。若 \(p=2^m+1\) 是素数，写 \(m=2^r u\)，其中 \(u\) 为奇数。若 \(u>1\)，则 \(z+1\) 整除 \(z^u+1\)，取 \(z=2^{2^r}\) 会给出 \(p\) 的非平凡因子。因此 \(u=1\)，即 \(m=2^r\)。反向代入上式立即得到 \(\varphi(n)\) 为 \(2\) 的幂。
\end{proof}

例如正 \(3,5,17\) 边形及正 \(15\) 边形可作图，正 \(7\) 边形与正 \(9\) 边形不可作图，因为相应次数为 \(6\)。每个给定的 \(n\) 都可以按其素因数分解单独检查。

\subsection{化圆为方与圆周率的超越性}
\label{b1:subsec:2104:03}

化圆为方要求作出与单位圆等面积的正方形，其边长必须为 \(\sqrt\pi\)。为排除这种长度，需要研究指数函数在代数数处的取值。下述结果是\term[Lindemann-Weierstrass-dingli]{Lindemann--Weierstrass 定理}[Lindemann--Weierstrass theorem]的单变量推论\footnote{林德曼（Carl Louis Ferdinand von Lindemann，1852-1939），德国数学家，证明圆周率的超越性，从而解决化圆为方的不可能性。}；其证明方法将整数的离散性与指数积分的估计结合起来\footnote{魏尔斯特拉斯（Karl Theodor Wilhelm Weierstrass，1815-1897），德国数学家，发展复函数论与分析的严格基础。}。

\begin{theorem}[Lindemann--Weierstrass 的单变量推论]\label{b1:thm:lindemann-weierstrass-special-case}
若 \(\alpha\) 是非零代数数，则 \(e^\alpha\) 是超越数。
\end{theorem}
\begin{proofsketch}
假设 \(e^\alpha=\beta\) 为代数数。取非零多项式 \(P(T)\in\mathbb Z[T]\)，使 \(P(\beta)=0\) 且 \(P(0)\ne0\)，并记 \(\alpha_1=\alpha,\ldots,\alpha_r\) 为 \(\alpha\) 的全部共轭。令
\[
\Lambda=\mathbb Z\alpha_1+\cdots+\mathbb Z\alpha_r\subseteq(\mathbb C,+).
\]
这是有限生成无挠交换群，\cref{b1:thm:fga-structure}保证它有一组整数基 \(\lambda_1,\ldots,\lambda_s\)：每个 \(\gamma\in\Lambda\) 唯一写成 \(\sum_j a_j\lambda_j\)，其中 \(a_j\in\mathbb Z\)。其整群代数 \(\mathbb Z[\Lambda]\) 由形式有限和 \(\sum_\gamma b_\gamma E^\gamma\) 组成，乘法规定为 \(E^\gamma E^\eta=E^{\gamma+\eta}\)。整数基给出同构
\[
\mathbb Z[\Lambda]\cong\mathbb Z[T_1^{\pm1},\ldots,T_s^{\pm1}],
\qquad E^{\sum_j a_j\lambda_j}\longmapsto\prod_j T_j^{a_j}.
\]
右端是 Laurent 多项式环，因而是整环。在这个形式环中构造
\[
Q=\prod_{i=1}^r P(E^{\alpha_i}).
\]
每个 \(\alpha_i\ne0\)，故 \(P(E^{\alpha_i})\) 中不同次数对应不同形式指数，是非零元素；整环性质于是给出 \(Q\ne0\)。可写 \(Q=\sum_\gamma b_\gamma E^\gamma\)，其中 \(b_\gamma\in\mathbb Z\)。代数共轭置换保留 \(Q\)。现在使用求值同态
\[
\operatorname{ev}:\mathbb Z[\Lambda]\longrightarrow\mathbb C,
\qquad E^\gamma\longmapsto e^\gamma,
\]
在 Laurent 多项式坐标中即 \(T_j\mapsto e^{\lambda_j}\)。由 \(P(e^\alpha)=0\)，\(Q\) 的求值为零。这一步并未假定求值同态单射：形式环中的非零元素是否可能求值为零，正是这里要排除的情形。

在形式环中令 \(Q^-=\sum_\gamma b_\gamma E^{-\gamma}\)，即把 \(Q\) 的各指数取负，系数不变。上标 \(-\) 表示这个操作，既不表示乘法逆，也不表示复共轭。乘积 \(QQ^-\) 的常数项为 \(\sum_\gamma b_\gamma^2>0\)，因为一项的指数与另一项的负指数相消，恰好要求原来的两个指数相等。\(QQ^-\) 仍在代数共轭下保持，求值仍为零。合并相同指数，得到
\[
c_0+\sum_{j=1}^n c_j e^{z_j}=0,
\]
其中 \(c_j\) 为整数，\(c_0>0\)，各 \(z_j\) 为互异的非零代数数；指数集合在共轭下保持，且共轭置换保留相应系数。这里共轭只作用于形式表达式的指数；取通常的指数值是随后单独进行的代入。

余下的目标是构造一族不全为零的代数数 \(S_{m,i}\)，使它们乘以一个可控制的整数因子后成为代数整数，同时使全部代数共轭都足够小。非零代数整数的范数是非零整数，绝对值至少为 \(1\)；若所有共轭的乘积反而小于 \(1\)，便得到矛盾。这里约束的是范数，不是某一个共轭的绝对值。

令 \(z_0=0\)、\(f(z)=\prod_{j=0}^n(z-z_j)\)。在各 \(z_j\) 处安排高重零点，可以使许多低阶导数消失；除以 \((m-1)!\) 后，剩余导数中的阶乘比仍为整数，便于控制整性。为把这些导数与一个小积分联系起来，取全部导数之和，使相邻导数相减时消项。具体地，对正整数 \(m\) 及 \(0\le i\le n\)，记
\[
P_{m,i}(z)=\frac{f(z)^m}{z-z_i},\qquad
d_m=(n+1)m-1,\qquad
H_{m,i}(z)=\frac1{(m-1)!}\sum_{r=0}^{d_m}P_{m,i}^{(r)}(z).
\]
这里 \(P_{m,i}\) 是次数为 \(d_m\) 的多项式；它在 \(z_i\) 处有 \(m-1\) 重零，在其余 \(z_k\) 处有 \(m\) 重零。除去一个因子使这些多项式的次数小于 \((n+1)m\)，同时保留后面证明非奇异性所需的重零点。定义端点矩阵 \(A_m=(H_{m,i}(z_k))_{0\le i,k\le n}\)，并沿从 \(0\) 到 \(z_k\) 的线段取积分
\[
\varepsilon_{m,i}(z_k)
=\frac1{(m-1)!}\int_0^{z_k}e^{-z}P_{m,i}(z)\,\mathrm{d}z.
\]
这里的线段积分可用实参数直接表示：将 \(z=tz_k\)、\(0\le t\le1\) 代入，有
\[
\int_0^{z_k}F(z)\,\mathrm{d}z
=z_k\int_0^1F(tz_k)\,\mathrm{d}t,
\]
右端按实部、虚部分别作普通实积分；\(z_k=0\) 时两端均为零。因此积分的绝对值至多为 \(\lvert z_k\rvert\max_{0\le t\le1}\lvert F(tz_k)\rvert\)。
由于 \(H_{m,i}-H_{m,i}'=P_{m,i}/(m-1)!\)，分部积分给出可逐项核对的端点等式
\[
\varepsilon_{m,i}(z_k)
=H_{m,i}(0)-e^{-z_k}H_{m,i}(z_k).
\]

矩阵 \(A_m\) 非奇异，这一点不需要预设其行列式的精确公式。设 \(\sum_i\lambda_iH_{m,i}(z_k)=0\) 对全部 \(k\) 成立，令 \(H=\sum_i\lambda_iH_{m,i}\) 及 \(G=H-H'\)。由定义，\(G\) 被 \(f^{m-1}\) 整除，故在每个 \(z_k\) 处至少有 \(m-1\) 重零。利用 \(H'=H-G\) 逐次求导，并从 \(H(z_k)=0\) 出发归纳，可知 \(H\) 在每个 \(z_k\) 处至少有 \(m\) 重零。可是 \(\deg H\le d_m<(n+1)m\)，所以 \(H=0\)。再由
\[
G(z)=\frac{f(z)^{m-1}}{(m-1)!}
\sum_{i=0}^n\lambda_i\prod_{j\ne i}(z-z_j)
\]
可将括号内的多项式依次在 \(z_i\) 处取值：由于各 \(z_j\) 不同，每次只有第 \(i\) 项不为零，故全部 \(\lambda_i=0\)。

令 \(c=(c_0,\ldots,c_n)^{\mathsf T}\)，并设 \(S_{m,i}=(A_mc)_i=\sum_k c_kH_{m,i}(z_k)\)。由于 \(c_0>0\) 且 \(A_m\) 非奇异，这些代数数不全为零。把端点等式乘以 \(c_ke^{z_k}\) 后求和，由 \(\sum_k c_ke^{z_k}=0\) 得
\[
S_{m,i}=-\sum_{k=0}^n c_ke^{z_k}\varepsilon_{m,i}(z_k).
\]
线段只有有限条；其上 \(e^{-z}\)、\(f(z)\) 与各多项式 \(f(z)/(z-z_i)\) 有共同的有限上界。写 \(P_{m,i}=f^{m-1}\bigl(f/(z-z_i)\bigr)\)，并使用上面的实参数积分界，便得到与 \(m,i\) 无关的 \(C,B>0\)，使 \(\lvert S_{m,i}\rvert\le CB^m/(m-1)!\)。

取正整数 \(q\)，使每个 \(qz_j\) 都是代数整数。\(P_{m,i}\) 在 \(z_k\) 处至少有 \(m-1\) 重零；展开 \(P_{m,i}(z_k+t)\) 后，\(t^r\) 的系数是含 \(d_m-r\) 个差 \(z_k-z_j\) 的整数系数多项式，乘以 \(q^{d_m-r}\) 即成为代数整数。对非零导数项，\(P_{m,i}^{(r)}(z_k)/(m-1)!\) 中的阶乘比 \(r!/(m-1)!\) 是整数。因此取固定的 \(D=q^{n+1}\)，则 \(D^mH_{m,i}(z_k)\)，从而 \(D^mS_{m,i}\)，均为代数整数。

取包含全部 \(z_j\) 的固定 Galois 数域 \(M\)。指数 \(z_j\) 及对应系数 \(c_j\) 在代数共轭下同步置换；若 \(\sigma(z_i)=z_j\)，则 \(\sigma\) 把 \(P_{m,i}\)、\(H_{m,i}\) 分别变为 \(P_{m,j}\)、\(H_{m,j}\)，从而 \(\sigma(S_{m,i})=S_{m,j}\)。故上面的积分界同时控制全部共轭。若 \(D^mS_{m,i}\ne0\)，其范数 \(N_{M/\mathbb Q}(D^mS_{m,i})\) 是非零整数；但其绝对值至多为 \((CD^mB^m/(m-1)!)^{[M:\mathbb Q]}\)，右边随 \(m\) 趋于零。故充分大的 \(m\) 下全部 \(S_{m,i}=0\)。这与它们不全为零矛盾。

原文\cite[\S 1, pp. 214--216]{Lindemann1882Pi}使用 Hermite 积分\footnote{埃尔米特（Charles Hermite，1822-1901），法国数学家，证明自然常数 \(e\) 的超越性，并研究椭圆函数。}及另行构造的系数矩阵；这里直接证明端点矩阵的非奇异性，不沿用原文系数矩阵的精确行列式等式。共轭对称与积分估计在原文\cite[\S\S 2--4, pp. 216--225]{Lindemann1882Pi}中用于推出单变量结论。这个证明概要使用复指数函数、线段上的积分估计和代数整数范数，不能仅由 Galois 对应替代。
\end{proofsketch}

一般定理断言，互异代数数的指数在代数数域上线性无关。取指数 \(0,\alpha\)，若 \(\beta=e^\alpha\) 为代数数，就有代数系数关系 \(-\beta e^0+e^\alpha=0\)，因此上述结论确是它的单变量推论。

\begin{corollary}\label{b1:cor:pi-transcendental}
圆周率 \(\pi\) 是超越数，因而不能用尺规将圆化为等面积的正方形。
\end{corollary}
\begin{proof}
若 \(\pi\) 代数，则非零数 \(i\pi\) 也代数，由\cref{b1:thm:lindemann-weierstrass-special-case}，\(e^{i\pi}\) 应当超越。这与 Euler 恒等式 \(e^{i\pi}=-1\) 矛盾，故 \(\pi\) 超越。

若 \(\sqrt\pi\) 是规矩数，由\cref{b1:thm:constructible-quadratic-tower}，它在 \(\mathbb Q\) 上代数；平方后 \(\pi\) 也应代数，仍然矛盾。因此化圆为方不可能。
\end{proof}
在化圆为方问题中，作图判据先给出可作图数必须代数这一必要条件；超越性论证另使用复指数函数及积分估计。Euler 恒等式也属于复指数函数的性质，因此化圆为方的不可作图性同时包含代数与分析的论证。


\begin{chaptersummary}
在特征零基域中添入所需单位根后，素数阶循环扩张可以由一个根式生成。在 \(\mathbb Q\) 上，分圆多项式的不可约性给出分圆域的次数与完整 Galois 群；在一般基域上，单位根扩张的群是相应单位群的子群。特征零下，求根问题的障碍由可解群刻画：一方面用正规包络把根式塔置于可解 Galois 扩张中，另一方面沿素数阶循环商逐层加入根式。\(S_3\)、\(S_4\) 可解而 \(S_n\)、\(n\geq5\)，不可解，解释了通用公式在五次处发生的变化。

\ProseParagraphBreak
尺规作图只允许实二次步骤，所以比一般根式表达更受限制。作图判据的核心是一条二次扩张塔；可作图实代数数的最小多项式还必须有 \(2\)-群作为分裂域的 Galois 群，单看该数的次数为 \(2\) 的幂并不够。分圆域的交换性又使正多边形的判据简化为 \(\varphi(n)\) 是否为 \(2\) 的幂。对化圆为方，域论只能说明可作图长度必须代数；排除 \(\sqrt\pi\) 还要使用独立的超越性结果。
\end{chaptersummary}

\BookExercises

\begin{exercise}\label{b1:ex:21-1}
计算 \(\Phi_9(X)\)、\(\Phi_{12}(X)\)，并核对它们的次数。
\end{exercise}
\begin{solution}
由\eqref{b1:eq:cyclotomic-factorization}，
\[
\Phi_9(X)=\frac{X^9-1}{X^3-1}=X^6+X^3+1.
\]
又因 \(X^{12}-1=(X^6-1)(X^6+1)\)，在除去阶整除 \(6\) 的根后，\(X^6+1\) 收集阶为 \(4\) 或 \(12\) 的根，所以
\[
\Phi_{12}(X)=\frac{X^6+1}{X^2+1}=X^4-X^2+1.
\]
相应次数为 \(6=\varphi(9)\)、\(4=\varphi(12)\)，两多项式由\cref{b1:thm:cyclotomic-irreducible}均不可约。
\end{solution}

\begin{exercise}\label{b1:ex:21-2}
设 \(n>1\) 为奇数。证明 \(\Phi_{2n}(X)=\Phi_n(-X)\) 及 \(\mathbb Q(\zeta_{2n})=\mathbb Q(\zeta_n)\)。解释为何必须排除 \(n=1\) 的多项式等式。
\end{exercise}
\begin{solution}
对奇数 \(n\)，本原 \(n\) 次根 \(\alpha\) 的负数具有阶 \(2n\)：其 \(n\) 次幂为 \(-1\)，平方具有阶 \(n\)。这个操作为两类本原根建立双射。又因 \(n>1\) 为奇数，互素剩余类由 \(a\) 与 \(-a\) 成对，\(\varphi(n)\) 为偶数，所以 \(\Phi_n(-X)\) 仍首一，得到多项式等式。域方面，不妨选取兼容的本原根 \(\zeta_n=\zeta_{2n}^2\)；此时 \(\zeta_{2n}=-\zeta_n^{(n+1)/2}\)，故两边包含。\(n=1\) 时 \(\Phi_1(-X)=-X-1=-\Phi_2(X)\)，符号不同。
\end{solution}

\begin{exercise}\label{b1:ex:21-3}
证明 \(\mathbb Q(\zeta_8)=\mathbb Q(i,\sqrt2)\)，列出其全部二次子域。
\end{exercise}
\begin{solution}
由 \(\zeta_8=(1+i)/\sqrt2\)，得到左域包含于右域。另一方面 \(i=\zeta_8^2\)，且 \(\sqrt2=\zeta_8+\zeta_8^{-1}\)，给出反向包含。群 \((\mathbb Z/8\mathbb Z)^\times\) 有四个元素，非单位元阶均为 \(2\)，故为 \(V_4\)，恰有三个二阶子群。相应三个二次子域为 \(\mathbb Q(i)\)、\(\mathbb Q(\sqrt2)\)、\(\mathbb Q(\sqrt{-2})\)；它们分别由三种非零平方类生成，彼此不同，已经穷尽。
\end{solution}

\begin{exercise}\label{b1:ex:21-4}
令 \(t=\zeta_5+\zeta_5^{-1}=2\cos(2\pi/5)\)。求 \(t\) 的最小多项式及 \(\cos(2\pi/5)\) 的根式表达。
\end{exercise}
\begin{solution}
将 \(1+\zeta_5+\cdots+\zeta_5^4=0\) 除以 \(\zeta_5^2\)，得 \((t^2-2)+t+1=0\)，即 \(t^2+t-1=0\)。它没有有理根，所以不可约。因为 \(2\pi/5\) 位于第一象限，\(t>0\)，故 \(t=(-1+\sqrt5)/2\)，于是 \(\cos(2\pi/5)=(\sqrt5-1)/4\)。
\end{solution}

\begin{exercise}\label{b1:ex:21-5}
分别取 \(K=\mathbb Q(\sqrt5),\ n=5\) 和 \(K=\mathbb Q(i),\ n=8\)。求 \(K(\mu_n)/K\) 的次数及其 Galois 群，并写出它在 \((\mathbb Z/n\mathbb Z)^\times\) 中对应的子群。
\end{exercise}
\begin{solution}
对第一组数据，\cref{b1:ex:21-4}中的计算给出 \(\sqrt5=2(\zeta_5+\zeta_5^{-1})+1\)，故 \(K\subseteq\mathbb Q(\zeta_5)\)，而 \(K(\mu_5)=\mathbb Q(\zeta_5)\)。两域在 \(\mathbb Q\) 上的次数分别为 \(2\)、\(4\)，所以相对次数为 \(2\)。复共轭固定 \(\sqrt5\)，送 \(\zeta_5\) 到 \(\zeta_5^{-1}\)；它与恒等自同构已穷尽相对 Galois 群。因此\cref{b1:prop:cyclotomic-general-base}的嵌入将这个 \(C_2\) 识别为 \((\mathbb Z/5\mathbb Z)^\times\) 中的 \(\{\overline1,\overline4\}\)。

对第二组数据，\(i=\zeta_8^2\)，故 \(K(\mu_8)=\mathbb Q(\zeta_8)\)，相对次数也是 \(4/2=2\)。自同构 \(\zeta_8\mapsto\zeta_8^a\) 固定 \(i\) 当且仅当 \(i^a=i\)，即 \(a\equiv1\pmod4\)。相应子群为 \(\{\overline1,\overline5\}\)，也同构于 \(C_2\)。这两个例子中的嵌入均非满射：扩大基域后，必须保留新基域元素满足的条件，允许的单位根置换随之减少。
\end{solution}

\begin{exercise}\label{b1:ex:21-6}
用三次求根的配对规则求 \(y^3-3y-2=0\) 的全部根。若独立选择 \(u^3=v^3=1\)，为何不一定得到解？
\end{exercise}
\begin{solution}
此时 \(\Delta=0\)，\(U=V=1\)，但还必须有 \(uv=1\)。选 \(u=v=1\)，配对所得根为 \(2\)、\(\zeta_3+\zeta_3^2=-1\)、\(-1\)，与 \((y-2)(y+1)^2\) 一致。若错误地选 \(u=1,v=\zeta_3\)，则 \(uv=\zeta_3\neq1\)，展开式留下 \((3uv-3)(u+v)\neq0\)，所得和不是原方程的根。
\end{solution}

\begin{exercise}\label{b1:ex:21-7}
对 \(f(X)=X^4+4X-1\)，找出一个二次扩张 \(E/\mathbb Q\)，使 \(f\) 在 \(E[X]\) 中分解为两个二次因子，并写出分解。再说明 \(E/\mathbb Q\) 的非平凡自同构如何作用于这两个因子。
\end{exercise}
\begin{solution}
此处一般四次式的系数为 \(b=c=0\)、\(d=4\)、\(e=-1\)。因三次项已经为零，不需要平移；写成 \(X^4+pX^2+qX+r\) 时，系数就是 \(p=0,q=4,r=-1\)。故辅助多项式为
\[
R(T)=T^3+4T-16=(T-2)(T^2+2T+8).
\]
取辅助根 \(T=2\)，对应的配方参数为 \(m=1\)。由 \(s^2=2m-p=2\)，可取 \(s=\sqrt2\)，而 \(q/(2s)=\sqrt2\)。因此
\[
X^4+4X-1
=\bigl(X^2-\sqrt2X+1+\sqrt2\bigr)
 \bigl(X^2+\sqrt2X+1-\sqrt2\bigr).
\]
这里可以取 \(E=\mathbb Q(\sqrt2)\)。其非平凡自同构把 \(\sqrt2\) 变为 \(-\sqrt2\)，因而交换这两个二次因子。
\end{solution}

\begin{exercise}\label{b1:ex:21-8}
求 \(X^4-2\) 在 \(\mathbb Q\) 上的分裂域及其群，说明它根式可解而 Galois 群非交换。
\end{exercise}
\begin{solution}
令 \(\alpha=\sqrt[4]{2}>0\)，分裂域为 \(L=\mathbb Q(\alpha,i)\)。Eisenstein 判别给出 \(\bigl[\mathbb Q(\alpha):\mathbb Q\bigr]=4\)；该域为实域，加入 \(i\) 再扩张两倍，故 \([L:\mathbb Q]=8\)。在 \(\mathbb Q(i)\) 上，\(\alpha\) 的最小多项式仍为四次，故映射 \(r(\alpha)=i\alpha,r(i)=i\) 定义四阶自同构。复共轭 \(s\) 满足 \(s(\alpha)=\alpha,s(i)=-i\)，并有 \(s^2=1\)、\(srs=r^{-1}\)。八个映射 \(r^j,r^js\) 两两不同，所以群为二面体群 \(D_8\)，本书的下标表示其阶为 \(8\)。塔 \(\mathbb Q\subseteq\mathbb Q(\alpha)\subseteq L\) 直接给出根式表达，而 \(sr\neq rs\)。
\end{solution}

\begin{exercise}\label{b1:ex:21-9}
设特征零域 \(K\) 上的两个多项式 \(f,g\) 都根式可解。证明 \(fg\) 也根式可解，并用 Galois 群给出另一说明。
\end{exercise}
\begin{solution}
先取包含 \(f\) 全部根的根式塔，再把一条为 \(g\) 选定的根式塔逐步与当前末端取合成域。每一步新生成元的某次幂仍在前一层中，所以连接后仍是根式塔，包含两多项式的全部根。群论方面，设分裂域为 \(L_f,L_g\)，则合成域 \(L_fL_g\) 的 Galois 群通过两次限制嵌入 \(\operatorname{Gal}(L_f/K)\times\operatorname{Gal}(L_g/K)\)。这里限制映射的核是平凡群：若一个自同构在 \(L_f\)、\(L_g\) 上都恒等，它便固定二者生成的整个合成域。两个因子可解，有限直积及其子群可解，故\cref{b1:thm:radical-solvability}再次给出结论。
\end{solution}

\begin{exercise}\label{b1:ex:21-10}
求 \(X^5-2\) 在 \(\mathbb Q\) 上的 Galois 群，并与一般五次方程比较。
\end{exercise}
\begin{solution}
令 \(\alpha=\sqrt[5]{2}\)、\(F=\mathbb Q(\zeta_5)\)。Eisenstein 判别给出 \(\bigl[\mathbb Q(\alpha):\mathbb Q\bigr]=5\)，而 \([F:\mathbb Q]=4\)。若 \(\alpha\in F\)，则 \(\mathbb Q(\alpha)\subseteq F\)，塔公式会给出 \(5=[\mathbb Q(\alpha):\mathbb Q]\mid[F:\mathbb Q]=4\)，矛盾。因此 \(\alpha\notin F\)。因 \(F\) 含 \(\mu_5\)，\cref{b1:prop:single-radical-cyclic}说明 \(\bigl[F(\alpha):F\bigr]\) 为 \(1\) 或 \(5\)，只能为 \(5\)。分裂域 \(L=F(\alpha)\) 因而在 \(\mathbb Q\) 上次数为 \(20\)。

定义 \(\sigma(\alpha)=\zeta_5\alpha\)、\(\sigma(\zeta_5)=\zeta_5\)，得到五阶自同构。\(F\) 上的自同构 \(\tau_0\) 由 \(\tau_0(\zeta_5)=\zeta_5^2\) 给定。因 \(X^5-2\) 在 \(F[X]\) 中不可约，映射
\[
F[X]/(X^5-2)\xrightarrow{\ \sim\ }L,\qquad \overline X\longmapsto\alpha
\]
是域同构。通过这一同构，将系数上的 \(\tau_0\) 延拓为 \(L\) 的自同构
\[
\tau\!\left(\sum_{j=0}^4c_j\alpha^j\right)
=\sum_{j=0}^4\tau_0(c_j)\alpha^j\qquad(c_j\in F).
\]
基表示的唯一性保证映射有定义；\(\tau_0\) 固定 \(2\)，所以关系 \(\alpha^5=2\) 得到保持，\(\tau\) 也保持乘法，并且固定 \(\alpha\)。它的阶为 \(4\)。计算得 \(\tau\sigma\tau^{-1}=\sigma^2\)，二者生成全部 \(20\) 个自同构，所以群为非平凡半直积 \(C_5\rtimes C_4\)。它有循环正规子群和循环商，因而可解。这一具体五次方程可解，并不改变独立系数的一般五次方程具有不可解群 \(S_5\) 的结论。
\end{solution}

\begin{exercise}\label{b1:ex:21-11}
为 \(\sqrt2+\sqrt3\) 与 \(\sqrt{1+\sqrt2}\) 分别写出实二次扩张塔，并求它们在 \(\mathbb Q\) 上的次数。
\end{exercise}
\begin{solution}
第一条塔为 \(\mathbb Q\subseteq\mathbb Q(\sqrt2)\subseteq\mathbb Q(\sqrt2,\sqrt3)\)，两个平方类独立，末端次数为 \(4\)。若 \(\alpha=\sqrt3+\sqrt2\)，则 \(\alpha^{-1}=\sqrt3-\sqrt2\)；取和差恢复 \(\sqrt2,\sqrt3\)，所以该和生成末端，次数为 \(4\)。

第二条塔为 \(\mathbb Q\subseteq\mathbb Q(\sqrt2)\subseteq\mathbb Q(\sqrt{1+\sqrt2})\)。\(1+\sqrt2\) 不是 \(\mathbb Q(\sqrt2)\) 中的平方，因为其范数为 \(-1\)，而平方的范数是一个有理数的平方，不可能为负。故第二步也是二次，且 \(\sqrt2=\alpha^2-1\) 属于 \(\mathbb Q(\alpha)\)，其中 \(\alpha=\sqrt{1+\sqrt2}\)。于是 \(\alpha\) 的次数为 \(4\)。两数均由\cref{b1:thm:constructible-quadratic-tower}可作图，因为已经写出了包含它们的实二次塔。这与\cref{b1:prop:constructible-splitting-two-group}后的四次反例形成对照：次数同为 \(4\)，是否可作图仍取决于能否组成这样的塔。
\end{solution}

\begin{exercise}\label{b1:ex:21-12}
已给定角 \(\theta\)，令 \(K=\mathbb Q(\cos\theta,\sin\theta)\)，即其单位圆终边点坐标生成的实域。证明这个角能够尺规三等分，当且仅当 \(X^3-3X-2\cos\theta\) 在 \(K[X]\) 中可约。
\end{exercise}
\begin{solution}
该三次多项式的三个根为 \(2\cos\bigl((\theta+2j\pi)/3\bigr)\)、\(j=0,1,2\)。若不可约，每个根在 \(K\) 上次数为 \(3\)，由相对于 \(K\) 的\cref{b1:thm:constructible-quadratic-tower}，任何一个都不可作出，因此不能三等分。若可约，三次多项式有一个根在 \(K\) 中，除去该线性因子后其余两个根仍是上面列出的实余弦值，所以剩余二次因子的判别式非负，两根在至多二次的实扩张中。于是全部三个余弦可作出，相应正弦由开平方得到，能从所得有限个终边中选出位于给定角内的三等分线。重根时同样适用，不需要三根互异。
\end{solution}

\begin{exercise}\label{b1:ex:21-13}
判断正 \(9,21,34,60,85\) 边形能否尺规作图，并逐一说明理由。
\end{exercise}
\begin{solution}
\(9=3^2\) 含奇素数的重复因子，\(\varphi(9)=6\)，不可作图。\(21=3\cdot7\)，其中 \(7-1=6\) 不是 \(2\) 的幂，\(\varphi(21)=12\)，也不可作图。其余分解为 \(34=2\cdot17\)、\(60=2^2\cdot3\cdot5\)、\(85=5\cdot17\)，奇素因子均为不同的 Fermat 素数。相应 Euler 函数值为 \(16,16,64\)，由\cref{b1:thm:constructible-polygon}都可作图。
\end{solution}

\begin{exercise}\label{b1:ex:21-14}
设 \(n\geq3\)。证明 \(\Bigl[\mathbb Q\bigl(2\cos(2\pi/n)\bigr):\mathbb Q\Bigr]=\varphi(n)/2\)，并据此判断 \(\cos(2\pi/7)\) 能否作图。
\end{exercise}
\begin{solution}
记 \(t=\zeta_n+\zeta_n^{-1}\)。\(\zeta_n\) 满足 \(X^2-tX+1=0\)，且非实，而 \(\mathbb Q(t)\subseteq\mathbb R\)，所以 \(\bigl[\mathbb Q(\zeta_n):\mathbb Q(t)\bigr]=2\)。分圆域次数为 \(\varphi(n)\)，由塔公式得所需等式。对 \(n=7\)，\(t\) 的次数为 \(3\)，乘以有理常数 \(1/2\) 不改变生成域，因此 \(\cos(2\pi/7)\) 不可作图。
\end{solution}

\begin{exercise}\label{b1:ex:21-15}
从单位线段出发，能否作出一条长度等于单位圆周长的线段？这与化圆为方所用的障碍是否相同？
\end{exercise}
\begin{solution}
单位圆周长为 \(2\pi\)。若能作出该线段，规矩数域对除以 \(2\) 封闭，所以 \(\pi\) 也可作图，因而代数。这与由 Lindemann--Weierstrass 定理推出的\cref{b1:cor:pi-transcendental}矛盾，故不能。化圆为方要求 \(\sqrt\pi\) 可作图，平方后同样会迫使 \(\pi\) 代数；两者都依赖圆周率的超越性，不能仅由某个有限扩张次数的计算代替。
\end{solution}


\begin{exercise}\label{b1:ex:21-reduction-quartic}
令 \(g(X)=X^4-6X-2\)，\(L\) 为它在 \(\mathbb Q\) 上的分裂域。
\begin{enumerate}
\item 证明 \(g\) 不可约，核验分解
\[
\begin{aligned}
g(X)&\equiv(X-1)(X^3+X^2+X+2)\pmod7,\\
g(X)&\equiv(X-6)(X-8)(X^2-3X+12)\pmod{17},
\end{aligned}
\]
并据此求 \(\operatorname{Gal}(L/\mathbb Q)\)。
\item 证明 \(g\) 在 \((1,2)\) 中有实根，判断这个根能否用根式表达、能否尺规作出。
\item 已知 \(\operatorname{disc}(g)=-37040\)，求 \(L\) 中的二次子域，并证明它是唯一的。
\end{enumerate}
\end{exercise}
\begin{solution}
1. 多项式 \(g\) 对素数 \(2\) 满足 Eisenstein 条件，因而在 \(\mathbb Q[X]\) 中不可约。展开乘积可核验题中的两处分解。模 \(7\) 的三次因子在 \(0,1,\ldots,6\) 处的值依次为 \(2,5,2,6,2,3,1\)，没有根，故不可约。模 \(17\) 的二次因子的判别式为 \(9-48\equiv12\)，而非零平方剩余为 \(1,2,4,8,9,13,15,16\)，所以它不可约；两个一次根 \(6,8\) 互异，也不是二次因子的根。两次约化均无重根。

由\cref{b1:lem:reduction-cycle-type}，群 \(G=\operatorname{Gal}(L/\mathbb Q)\leq S_4\) 含三轮换和换位；不可约性又使它在四根上传递。因此 \(4\bigm\vert|G|\) 且 \(3\bigm\vert|G|\)，只能有 \(|G|=12\) 或 \(24\)。若阶为十二，指数为二使 \(G\) 在 \(S_4\) 中正规，因而包含所含换位的全部共轭，也就是所有换位；但换位生成 \(S_4\)，矛盾。故 \(G=S_4\)。只用三轮换不能得到这一结论，因为传递子群 \(A_4\) 也包含三轮换。

2. 因 \(g(1)=-7\)、\(g(2)=2\)，连续性给出一个实根 \(\alpha\in(1,2)\)。群 \(S_4\) 有正规列 \(1\subset V_4\subset A_4\subset S_4\)，各因子为交换群，故由\cref{b1:thm:radical-solvability}，\(\alpha\) 能用根式表达。但 \(S_4\) 的阶为 \(24\)，不是 \(2\) 的幂；\cref{b1:prop:constructible-splitting-two-group}排除了尺规作图。不可约性给出 \([\mathbb Q(\alpha):\mathbb Q]=4\)，所以这个计算再次区分了代数次数的条件与作图条件。

3. 将四根记为 \(a_1,\ldots,a_4\)，置 \(\delta=\prod_{i<j}(a_j-a_i)\)。\cref{b1:prop:discriminant-alternating-group}的证明给出 \(\sigma(\delta)=\operatorname{sgn}(\sigma)\delta\)，故 \(\delta\) 的稳定子群恰为 \(A_4\)。于是
\[
L^{A_4}=\mathbb Q(\delta)
=\mathbb Q(\sqrt{-37040})
=\mathbb Q(\sqrt{-2315}).
\]
这里 \(-37040=16\cdot(-2315)\)，且负有理数不是有理数的平方，故这个子域确为二次。若还有二次子域，Galois 对应会给出 \(S_4\) 的另一个指数为二的子群，也就是某个满同态 \(S_4\longrightarrow C_2\) 的核。所有换位共轭，在该同态下有相同的像；换位生成 \(S_4\)，满射性迫使这个像为 \(C_2\) 的非单位元。因此该同态就是符号同态，其核只能是 \(A_4\)。这证明了唯一性。
\end{solution}

% !TeX root = ../../Book1.tex
% 纲要标题：## 第22章　无限 Galois 理论
% 纲要位置：计划纲要.md，第 767 行。

\chapter{无限 Galois 理论}
\label{b1:ch:22}
\BookChapterNumber{b1:ch:22}

\begin{chapterabstract}
无限代数扩张可由有限子扩张拼合，其自同构须在各有限层上相容。本章把无限 Galois 群表示为有限群的逆极限，建立 Krull 拓扑与紧性，证明中间域同闭子群之间的无限 Galois 对应。
\end{chapterabstract}

\begin{learninggoals}
\begin{itemize}
\item 把无限Galois群写成有限商的逆极限，区分相容族与任意坐标选择。
\item 理解Krull拓扑的基本开集，掌握紧性证明中的有限交性质与选择原理。
\item 运用闭子群对应、开子群判别和正规商群结论，识别删除闭性后的反例。
\item 计算有限域及实数域的绝对Galois群，利用有限层判断连续性和表示的相容性。
\end{itemize}
\end{learninggoals}

依次把四次、八次、十六次单位根添入有理数域，便得到不断增大的分圆域。每一层都有有限 Galois 群；若想在所有层的并上定义一个自同构，在较大层上的取值必须限制成较小层已有的取值。独立任选各层的自同构，通常无法拼成同一个映射。无限 Galois 理论把这种相容性放进群的拓扑中，再问无限扩张的中间域究竟对应哪些子群。有限对应仍在每一层起作用，但直接照搬“任意子群对应中间域”会遗漏闭性条件。

\ProseParagraphBreak

本章的对应定理主要使用有限 Galois 对应、一般代数扩张的嵌入延拓、正规性与可分性，其中有限对应所需内容见\secref{b1:sec:1901}和\secref{b1:sec:1902}。有限域的 Frobenius 计算用于\secref{b1:sec:2204}的典型例子。拓扑术语在使用前定义；任意指标下的紧性证明使用 Zorn 引理。实数域的例子使用第12章引用的\cref{b1:thm:fundamental-theorem-algebra}；其证明位于\secref{b1:sec:E02}，调用实分析的中间值定理、Sylow 理论与有限 Galois 对应。末尾简述连续群上同调与无限 Galois 群的联系。


\ProseParagraphBreak
本章的无限 Galois 扩张可与 Lang 的《Algebra》\cite[Chapter VI, \S 14]{Lang2002}对照。Gille 与 Szamuely 的《Central Simple Algebras and Galois Cohomology》\cite[Sections 4.1--4.2]{GilleSzamuely2017CentralSimpleAlgebras}进一步讨论绝对 Galois 群及连续上同调，使用本章的拓扑群知识，也需要系数模和群上同调。有限群的线性表示与特征标则见 Serre 的《Linear Representations of Finite Groups》\cite[Part I, Chapters 1--2, pp. 3--24]{Serre1977FiniteRepresentations}。该书这部分主要讨论复系数表示；本章的有限域系数表示若经过阶数可被系数域特征整除的有限群，完全可约性和普通迹函数的判别能力就不能照搬复系数情形。

% !TeX root = ../../Book1.tex
% 纲要标题：### 22.1 无限代数扩张
% 纲要位置：计划纲要.md，第 769 行。

\section{无限代数扩张}
\label{b1:sec:2201}

无限代数扩张中的每一个元素，仍落在某个有限子扩张中。困难不在单个元素，而在于不同有限层上的选择必须彼此相容。本节先把自同构写成一族有限层上的自同构，再说明为什么任意选择各层自同构通常不能拼成一个整体。

% 纲要内容（仅作写作计划，尚非教材正文）：
%
% * 有限 Galois 子扩张的有向系统；说明有限层须被全局自同构保持，并区分过渡映射与整体限制映射的满射性
% * 代数闭包的自同构群；解释可分闭包与纯不可分部分，以及自同构群同构为何不能代替扩张的可分性
% * 逐层相容的自同构；证明拼合与所选有限层无关
%

\subsection{有限子扩张的有向系统}
\label{b1:subsec:2201:01}
本章固定代数扩张 \(L/K\)。若它正规且可分，仍称 \(L/K\) 为 Galois 扩张，但不再要求次数有限；这时把它的自同构群记为 \(G=\operatorname{Gal}(L/K)\)。有限层上的群阶仍等于相应的扩张次数，整个无限扩张的域次数与自同构群的普通基数却一般不相等。记录各个有限 Galois 层上的自同构及其限制映射，能够保留有限理论中可以逐层计算的信息；这些有限层还将决定群上的拓扑。

\begin{proposition}\label{b1:prop:finite-galois-layers}
设 \(L/K\) 为代数 Galois 扩张。\(L\) 的每个有限子集都包含在某个有限 Galois 子扩张 \(N/K\) 中。所有这样的 \(N\) 按包含关系构成有向集合，且 \(L=\bigcup_N N\)。
\end{proposition}
\begin{proof}
取有限集 \(a_1,\ldots,a_r\in L\)。每个 \(a_i\) 的最小多项式可分，且由正规性在 \(L\) 中分裂。只有有限多个这样的多项式，每个又只有有限多个根，所以它们的全部根组成有限集。将这些根加入 \(K\)，所得域 \(N\) 仍在 \(L\) 内，由有限多个代数元生成，故 \(N/K\) 有限。把相同的首一最小多项式只保留一次，所得乘积可分，因为不同不可约多项式没有公共根；\(N\) 正是该乘积的分裂域，因此正规且可分，并含原有限集。空集可取 \(N=K\)。

若 \(N_1,N_2\) 是两个有限 Galois 子扩张，各取有限生成元，将它们的并按上一段处理，就有同时包含两域的有限 Galois 子扩张。这就是按包含关系有向的含义；有限次重复，任意有限多个这样的域都能放进一个共同有限 Galois 层中。每个 \(a\in L\) 所成的单点集也被某层包含，故它们的并为 \(L\)。
\end{proof}
有向集合的抽象定义见\cref{b1:def:directed-system}；这里的共同上层已由上述证明构造。记这个有向集合为 \(\mathcal N\)，并记 \(G_N=\operatorname{Gal}(N/K)\)。选取有限层时要求 \(N/K\) 正规，是为了能在它上面记录整体自同构的限制：由\cref{b1:thm:normal-embedding-criterion}，对任意 \(\sigma\in G\)，有 \(\sigma(N)=N\)。若只取任意有限子扩张 \(E/K\)，\(\sigma(E)\) 可能是另一个共轭子域，\(\sigma|_E\) 就不能视为 \(E\) 的自同构。同理，对 \(N\subseteq M\)，\(M\) 的每个 \(K\)-自同构都保持 \(N\)，所以限制给出群同态
\[
\operatorname{res}_{M,N}:G_M\longrightarrow G_N.
\]
这个有限层间的限制映射是否满射，以及 \(G\) 能否实现有限层上的每个自同构，都还需要嵌入延拓定理保证。

\begin{lemma}\label{b1:lem:galois-restriction-surjective}
设 \(L/K\) 为代数 Galois 扩张，\(G=\operatorname{Gal}(L/K)\)。若 \(N\subseteq M\subseteq L\) 都是 \(L/K\) 的有限 Galois 子扩张，并记 \(G_N=\operatorname{Gal}(N/K)\)、\(G_M=\operatorname{Gal}(M/K)\)，则限制同态
\[
G\longrightarrow G_N,
\qquad
\operatorname{res}_{M,N}:G_M\longrightarrow G_N
\]
都满射。
\end{lemma}
\begin{proof}
取包含 \(L\) 的 \(K\) 的代数闭包 \(\Omega\)。给定 \(\tau\in G_N\)，把它看作 \(N\rightarrow\Omega\) 的嵌入。\cref{b1:thm:embedding-extension}应用于代数扩张 \(L/N\)，将它延拓为 \(\sigma:L\rightarrow\Omega\)。因为 \(\sigma|_N=\tau\)，它固定 \(K\)；\cref{b1:thm:normal-embedding-criterion}及 \(L/K\) 的正规性给出 \(\sigma(L)=L\)，故 \(\sigma\in G\)，且它在 \(N\) 上恰为 \(\tau\)。这证明 \(G\rightarrow G_N\) 满射，即有限层的每个自同构都能在整体扩张上实现。

把同一论证中的 \(L\) 换成 \(M\)，\(M/N\) 仍代数而 \(M/K\) 正规，就得到 \(\tau\) 到 \(M\) 的自同构延拓。因此 \(G_M\rightarrow G_N\) 也满射，即较小有限层上的自同构能提升到较大有限 Galois 层。第一种满射控制整体群在每个有限坐标上的像，第二种控制有限层之间的过渡映射。对于三层域，逐次限制等于直接限制。
\end{proof}
这两种满射性来自延拓与正规性，并不是把各有限群写成坐标以后自动得到的性质。

\subsection{代数闭包的自同构群}
\label{b1:subsec:2201:02}
不完美域的代数闭包中，可能有基域之外的元素被每个基域自同构固定。例如 \(K=\mathbb F_p(t)\)，在代数闭包中取 \(u=t^{1/p}\)。多项式 \(X^p-t\) 在 \(K\) 上不可约，且只有 \(u\) 这一个不同根，所以 \(u\notin K\)，每个 \(K\)-自同构却都固定 \(u\)。因此 \(\operatorname{Aut}_K(\overline K)\) 在整个代数闭包上的不动域可以严格大于 \(K\)。要使自同构的作用能够检测扩张中的不同共轭根，需要考察可分部分。

\ProseParagraphBreak
固定 \(K\) 的代数闭包 \(\overline K\)。其中所有在 \(K\) 上可分的元素组成子域，记为 \(K^{\mathrm{sep}}\)，称为 \(K\) 的\term[kefen-bibao]{可分闭包}[separable closure]。它是子域的依据为\cref{b1:prop:separable-subfield}；任一可分元的全部共轭根仍然可分，且已在 \(\overline K\) 中，所以 \(K^{\mathrm{sep}}/K\) 正规而可分。
\symidx{\(K^{\mathrm{sep}}\)：所选代数闭包中的可分闭包}

\begin{definition}\label{b1:def:absolute-galois-group}
设 \(K\) 为域，固定一个代数闭包 \(\overline K\)，并令 \(K^{\mathrm{sep}}\subseteq\overline K\) 为所有在 \(K\) 上可分的元素组成的可分闭包。群
\[
G_K\coloneqq\operatorname{Gal}(K^{\mathrm{sep}}/K)
\]
称为域 \(K\) 的\term[juedui-galois-qun]{绝对 Galois 群}[absolute Galois group]。换一份可分闭包会给出同构的群，但同构一般依赖于闭包之间同构的选择。
\end{definition}

\symidx{\(G_K\)：域的绝对 Galois 群}

\begin{proposition}\label{b1:prop:closure-automorphisms}
设 \(K\) 为域，\(\overline K\) 为 \(K\) 的代数闭包，\(K^{\mathrm{sep}}\subseteq\overline K\) 为其中所有在 \(K\) 上可分的元素组成的子域，并令 \(G_K=\operatorname{Gal}(K^{\mathrm{sep}}/K)\)。则限制映射 \(\operatorname{Aut}_K(\overline K)\rightarrow G_K\) 为群同构。若 \(K\) 完美，则 \(K^{\mathrm{sep}}=\overline K\)；一般情形中不能把 \(\overline K/K\) 直接视为 Galois 扩张。
\end{proposition}
\begin{proof}
\(K\)-自同构保持最小多项式及其可分性，故保持可分闭包，限制映射有定义。若特征为零，所有代数元可分，限制映射就是恒等识别。以下设特征为 \(p>0\)。对任意 \(a\in\overline K\)，\cref{b1:prop:inseparable-polynomial-shape}把它的最小多项式写为
\[
m_{a,K}(X)=g(X^{p^r}),
\]
其中 \(g\) 首一、不可约且可分。代入 \(a\) 得 \(g(a^{p^r})=0\)，故 \(g\) 正是 \(a^{p^r}\) 的最小多项式。因此 \(a^{p^r}\) 在 \(K\) 上可分，属于 \(K^{\mathrm{sep}}\)。这说明每个 \(a\) 都在 \(K^{\mathrm{sep}}\) 上纯不可分。

若两个 \(K\)-自同构 \(\sigma,\tau\) 在 \(K^{\mathrm{sep}}\) 上有相同限制 \(\nu\)，对给定 \(a\) 取上述的 \(r\)，则
\[
\sigma(a)^{p^r}=\nu(a^{p^r})=\tau(a)^{p^r}.
\]
两像都是方程 \(X^{p^r}-\nu(a^{p^r})=0\) 的根。这个方程至多有一个不同根，因为两根之差的 \(p^r\) 次幂为零，而域没有非零幂零元。因此 \(\sigma(a)=\tau(a)\)；对每个 \(a\) 都如此，限制映射便单射。

反过来，给定 \(\nu\in G_K\)，把它看作 \(K^{\mathrm{sep}}\rightarrow\overline K\) 的嵌入。\cref{b1:thm:embedding-extension}应用于代数扩张 \(\overline K/K^{\mathrm{sep}}\)，将它延拓为 \(\sigma:\overline K\rightarrow\overline K\)。这个嵌入固定 \(K\)，代数闭包在 \(K\) 上正规，\cref{b1:thm:normal-embedding-criterion}保证 \(\sigma(\overline K)=\overline K\)，所以它是自同构，得到满射性。完美域的所有代数元可分，最后的等式由定义得到。
\end{proof}
纯不可分部分上的延拓唯一，因而没有增加自同构的选择，这正是两个群同构的原因。但这个同构不会使原本不可分的元素变得可分；这些不能被移动的元素仍在代数闭包中，使该群在那里作用的不动域可以严格大于 \(K\)。所以，还须区分群本身与它所作用的域，不能由上述群同构把 \(\overline K/K\) 直接当作 Galois 扩张。

\subsection{相容自同构}
\label{b1:subsec:2201:03}
从有限层上的自同构恢复整体映射，自然要对 \(a\in L\) 规定 \(\sigma(a)=\sigma_N(a)\)，其中 \(N\) 是任意包含 \(a\) 的有限 Galois 层。首先须解决的是选择无关性：同一个元素可能属于许多层，只有这些层给出的像一致，才能在 \(L\) 上定义一个映射。

\begin{proposition}\label{b1:prop:galois-compatible-family}
设 \(L/K\) 为代数 Galois 扩张，\(G=\operatorname{Gal}(L/K)\)，并令 \(\mathcal N\) 为 \(L/K\) 的全部有限 Galois 子扩张组成的有向集合。对每个 \(N\in\mathcal N\)，记 \(G_N=\operatorname{Gal}(N/K)\)，并给定 \(\sigma_N\in G_N\)。则存在唯一 \(\sigma\in G\) 在每个 \(N\) 上等于 \(\sigma_N\)，当且仅当
\begin{equation}\label{b1:eq:galois-compatibility}
\sigma_M|_N=\sigma_N\qquad(N\subseteq M).
\end{equation}
\end{proposition}
\begin{proof}
限制一个自同构当然满足相容条件。反过来，对 \(a\in L\)，选含 \(a\) 的有限 Galois 层 \(N\)，规定 \(\sigma(a)=\sigma_N(a)\)。若还选了 \(M\)，由\cref{b1:prop:finite-galois-layers}取同时包含两者的一层 \(P\)，相容性给出 \(\sigma_N(a)=\sigma_P(a)=\sigma_M(a)\)，所以定义无关选择。

对任意两个元素，取包含它们的同一有限层，就可在该层验证保持和、积与单位，以及固定 \(K\)。还须保证所得映射可逆。若 \(N\subseteq M\)，相容性给出 \(\sigma_M|_N=\sigma_N\)，而 \(\sigma_N\) 是 \(N\) 的自同构；所以 \(\sigma_M^{-1}\) 在 \(N\) 上的限制正是 \(\sigma_N^{-1}\)。因此逆族 \((\sigma_N^{-1})_N\) 也相容，同样拼成一个映射，在每层上都是 \(\sigma\) 的双边逆，因而在 \(L\) 上也是双边逆。所得 \(\sigma\) 是自同构；\(L\) 是各层的并，全部限制也就唯一确定了它。
\end{proof}
在每层能够延拓，不意味着可以把所有层上任意选定的自同构拼在一起。例如在有限域的塔 \(\mathbb F_q\subset\mathbb F_{q^2}\subset\mathbb F_{q^4}\) 中，取四次层上的 Frobenius \(F(x)=x^q\)。它在二次层上的限制仍为 \(x\mapsto x^q\)；由\cref{b1:thm:finite-field-automorphisms}，这个限制的阶为 \(2\)，是二次扩张的非平凡自同构。所以若预先在二次层选择恒等，它就不能与四次层上选定的 \(F\) 相容，违反\eqref{b1:eq:galois-compatibility}。各层的选择分别合法，还须满足限制相容条件才能拼成整体。下一节把相容族组成的群称为逆极限，并在其上建立适合这些有限信息的拓扑。

% !TeX root = ../../Book1.tex
% 纲要标题：### 22.2 Krull 拓扑
% 纲要位置：计划纲要.md，第 775 行。

\section{Krull 拓扑}
\label{b1:sec:2202}

一个自同构在有限层上的取值，只提供有限的信息。把“在某个有限层上相同”作为接近的标准，就得到 Krull 拓扑。它不是额外附加的距离，而是把有限 Galois 群之间的相容关系变成可用的开集。所需拓扑概念及紧性证明均在本节给出。

% 纲要内容（仅作写作计划，尚非教材正文）：
%
% * 有限商与逆极限；区分一般有向系统与可数序列，坐标投影的满射性须另行证明
% * Galois 群的拓扑；把基本开陪集解释为在某一有限层上作用相同，并证明有限元素观察给出同一拓扑
% * 就地介绍离散拓扑、积拓扑、有限交性质及闭子空间，完整证明有限离散空间积的紧性；展开生成滤子、极大滤子的二择一性质与有限划分取值，并说明选择原理的位置
% * 把逆极限识别为相容条件定义的闭子空间，连续证明 Galois 群的紧性、全不连通性与开子群性质
% * 附录 A 仅汇总和补充这段拓扑，不承担本章缺失的必需引理
%

\subsection{有限商与逆极限}
\label{b1:subsec:2202:01}
有限 Galois 层按域包含关系排列：\(N\subseteq M\) 时，\(M\) 记录更多元素上的取值，而限制映射却从较大层的群 \(G_M\) 指向较小层的群 \(G_N\)。下面的逆系统沿用这个方向，指标向后增大，坐标信息则限制到前层。

\begin{definition}\label{b1:def:inverse-limit}
设 \((I,\leq)\) 为有向集。对每个 \(i\in I\) 给定群 \(G_i\)，并对 \(i\leq j\) 给定同态 \(r_{ji}:G_j\rightarrow G_i\)，满足 \(r_{ii}=\operatorname{id}_{G_i}\)，且对所有 \(i\leq j\leq k\) 都有 \(r_{ki}=r_{ji}r_{kj}\)。在直积 \(\prod_iG_i\) 中取相容子群
\[
\varprojlim_iG_i
=\Bigl\{\,(g_i)_i\in\prod_iG_i\mid
r_{ji}(g_j)=g_i\text{ 每当 }i\leq j\,\Bigr\},
\]
并逐坐标定义乘法。称这个群为所给系统的\term[nijixian]{逆极限}[inverse limit]。
\end{definition}
\symidx{\(\varprojlim_iG_i\)：群的逆极限}
同态保持乘法及逆元，故相容族的乘积与逆仍相容；各坐标单位组成单位族。因此这确实是直积的子群。投影到第 \(i\) 个坐标记为 \(\pi_i\)。逆极限把相互一致的各层信息看作一个整体：若群 \(H\) 上给出同态 \(f_i:H\rightarrow G_i\)，且 \(r_{ji}f_j=f_i\) 对每个 \(i\le j\) 成立，则映射 \(h\mapsto\bigl(f_i(h)\bigr)_i\) 的像落在逆极限中。它保持群运算，并且是唯一满足 \(\pi_i f=f_i\) 的同态 \(f:H\rightarrow\varprojlim_iG_i\)，因为每个坐标已经指定。在本章取 \(H=G\)、\(f_N\) 为限制映射，正是把整体自同构的有限层信息放进逆极限。其底集合构造与唯一性可对照\cref{b1:prop:inverse-limit-set-universal}。

\ProseParagraphBreak
例如，通常的模数约化给出逆系统
\[
\mathbb Z/2\mathbb Z\longleftarrow
\mathbb Z/4\mathbb Z\longleftarrow
\mathbb Z/8\mathbb Z\longleftarrow\cdots.
\]
模 \(2^r\) 的剩余类 \([a]_{2^r}\) 在下一层有两个提升：\([a]_{2^{r+1}}\) 与 \([a+2^r]_{2^{r+1}}\)。逐层相容地选定剩余类，就得到逆极限 \(\mathbb Z_2=\varprojlim_r\mathbb Z/2^r\mathbb Z\) 的一个元素。这个例子是一条可数链；一般有向逆系统未必能化为这样的可数塔，也未必存在一个“下一层”。在 Galois 系统中，任意两个已经考虑的有限层都能进入某个共同上层，各个限制必须在较小层上一致；重复取共同上层，可以同时比较有限多个坐标。

\ProseParagraphBreak
对上一节的有限 Galois 层应用\cref{b1:prop:galois-compatible-family}，得到群同构
\begin{equation}\label{b1:eq:galois-inverse-limit}
G\longrightarrow\varprojlim_{N\in\mathcal N}G_N,
\qquad\sigma\longmapsto(\sigma|_N)_N.
\end{equation}
每个坐标投影满射的依据是\cref{b1:lem:galois-restriction-surjective}：给定 \(G_N\) 中的自同构，嵌入延拓与 \(L/K\) 的正规性使它延拓为 \(G\) 中的自同构，因而在逆极限中得到一个具有指定第 \(N\) 坐标的相容族。一般逆极限的定义并不保证坐标投影满射；它只收集能够扩成完整相容族的坐标，没有保证某一层的每个取值都能如此扩充。今后把两侧识别，不仅记录抽象群运算，还保留全部限制映射。

\subsection{Galois 群的拓扑}
\label{b1:subsec:2202:02}
设 \(X\) 为集合，\(\mathcal T\) 是 \(X\) 的一族子集。若 \(\mathcal T\) 包含 \(\varnothing,X\)，并且对任意并和有限交封闭，则称 \(\mathcal T\) 为 \(X\) 上的一个\term[tuopu]{拓扑}[topology]；\(\mathcal T\) 中的集合称为开集，其补集称为闭集。若 \(\mathcal T=\mathcal P(X)\)，即 \(X\) 的每个子集都开，则称这个拓扑为\term[lisan-tuopu]{离散拓扑}[discrete topology]。设 \(X,Y\) 为拓扑空间，\(f:X\rightarrow Y\) 为映射。若 \(Y\) 中每个开集的原像都是 \(X\) 中的开集，则称 \(f\) 为连续映射；若 \(f\) 还是双射，并且 \(f\) 与 \(f^{-1}\) 都连续，则称 \(f\) 为同胚。本节的每个有限群 \(G_N\) 都采用离散拓扑。

\ProseParagraphBreak
设 \((X_i)_{i\in I}\) 为拓扑空间族，\(P=\prod_{i\in I}X_i\)，并以 \(p_i:P\rightarrow X_i\) 表示坐标投影。由所有集合
\[
\bigcap_{i\in F}p_i^{-1}(U_i),
\qquad F\subseteq I\text{ 有限},\quad U_i\subseteq X_i\text{ 开},
\]
的任意并组成的拓扑称为 \(P\) 上的\term[ji-tuopu]{积拓扑}[product topology]，其中 \(F=\varnothing\) 时交集约定为 \(P\)；所列有限交称为基本开集。对任意子集 \(Y\subseteq P\)，若把 \(Y\cap U\)，其中 \(U\) 遍历 \(P\) 的开集，规定为 \(Y\) 中的开集，则称所得拓扑为 \(Y\) 的子空间拓扑。应用于有限离散群族 \((G_N)\) 时，基本开集只规定有限多个坐标的取值，其余坐标任意。

\begin{samepage}
\begin{definition}\label{b1:def:krull-topology}
设 \(L/K\) 为代数 Galois 扩张，\(G=\operatorname{Gal}(L/K)\)，并令 \(\mathcal N\) 为 \(L/K\) 的全部有限 Galois 子扩张组成的有向集合。对 \(N\in\mathcal N\)，记 \(G_N=\operatorname{Gal}(N/K)\)，并令 \(\pi_N:G\rightarrow G_N\) 为限制映射。通过同构 \(G\cong\varprojlim_{N\in\mathcal N}G_N\)，把 \(G\) 看作有限离散群直积 \(\prod_NG_N\) 的子空间。由该子空间拓扑得到的 \(G\) 上的拓扑称为\term[krull-tuopu]{Krull 拓扑}[Krull topology]。若
\[
U_N\coloneqq\operatorname{Gal}(L/N)=\ker\pi_N,
\]
则对每个 \(\sigma\in G\)，其基本开邻域为 \(\sigma U_N\)，其中 \(N\) 遍历有限 Galois 子扩张。
\end{definition}
\end{samepage}
\symidx{\(U_N=\operatorname{Gal}(L/N)\)：Krull 拓扑的基本开正规子群}
对 \(\sigma,\tau\in G\)，陪集的具体含义是
\[
\tau\in\sigma U_N
\iff\pi_N(\tau)=\pi_N(\sigma)
\iff\tau|_N=\sigma|_N.
\]
也就是说，在有限层 \(N\) 上无法区分 \(\tau\) 与 \(\sigma\)。一个基本开集原本可能限制有限多个坐标 \(N_1,\ldots,N_r\)，但\cref{b1:prop:finite-galois-layers}给出同时包含这些域的有限 Galois 层 \(M\)。于是 \(\sigma U_M\subseteq\bigcap_j\sigma U_{N_j}\)，只规定第 \(M\) 个坐标就得到包含 \(\sigma\) 且落在原基本开集内的邻域。因此所列陪集确实构成邻域基。它们也都是闭集：其补集是同一有限商中其余陪集的并。

\begin{proposition}\label{b1:prop:krull-group-operations}
设 \(L/K\) 为代数 Galois 扩张，\(G=\operatorname{Gal}(L/K)\)，并在 \(G\) 上赋予 Krull 拓扑。则乘法 \(G\times G\rightarrow G\) 与取逆 \(G\rightarrow G\) 都连续，所以 \(G\) 为拓扑群。
\end{proposition}
\begin{proof}
若要确定 \(\sigma\tau\) 在 \(N\) 上的取值，只需分别确定 \(\sigma,\tau\) 在 \(N\) 上的取值，因为限制保持乘法。具体地，\(\sigma'\in\sigma U_N\)、\(\tau'\in\tau U_N\) 使 \(\pi_N(\sigma'\tau')=\pi_N(\sigma\tau)\)，故它们的乘积落在 \(\sigma\tau U_N\) 中。这在每个点给出乘法连续性。限制也保持逆元，因此确定 \(\sigma|_N\) 就确定 \(\sigma^{-1}|_N\)，得到取逆连续。
\end{proof}
于是左右平移均为同胚。若只要求 \(\tau(a_j)=\sigma(a_j)\) 对有限多个 \(a_1,\ldots,a_r\in L\) 成立，取包含它们的有限 Galois 层 \(N\)，则 \(\sigma U_N\) 包含于这个有限观察集合；该集合的每个点都有这样的陪集邻域，所以它开。反过来，取 \(N/K\) 的一组有限基 \(b_1,\ldots,b_s\)。自同构固定 \(K\) 并保持线性组合，故在这些基元素上取值相同就等价于在整个 \(N\) 上取值相同，所得集合正是 \(\sigma U_N\)。因此“在有限多个元素上取值相同”与“在某个有限 Galois 层上取值相同”给出相同拓扑；前者表示有限观察，后者便于使用有限群的结构。

\subsection{有限离散空间积的紧性}
\label{b1:subsec:2202:04}
相容数据的条件可能有无限多个，我们只能一次检查有限批条件。对于闭条件，紧性将把“每一有限批条件都有共同解”加强为“全部条件有共同解”，这也是后面处理闭子群与商群拓扑时需要的性质。在有限离散空间的任意直积中，各坐标只有有限个取值，这将使上述性质能够由有限划分的选择得到。

\ProseParagraphBreak
设 \(X\) 为拓扑空间。若 \(X\) 的每个开覆盖都有有限子覆盖，则称 \(X\) 为\term[nijin-kongjian]{拟紧空间}[quasi-compact space]。若 \(X\) 中任意两个不同的点都有互不相交的开邻域，则称 \(X\) 具有 Hausdorff 性\footnote{豪斯多夫（Felix Hausdorff，1868-1942），德国数学家，研究集合论与拓扑学，建立拓扑空间的系统理论。}；若 \(X\) 同时拟紧并具有 Hausdorff 性，则称 \(X\) 为\term[jin-kongjian]{紧空间}[compact space]。

\ProseParagraphBreak
设 \(X\) 为集合，\(\mathcal C\subseteq\mathcal P(X)\) 为一族子集。若 \(\mathcal C\) 中任意有限多个集合的交都非空，其中空族的交约定为 \(X\)，则称 \(\mathcal C\) 具有\term[youxian-jiao-xingzhi]{有限交性质}[finite intersection property]。对坐标条件而言，每个集合收集满足某一条件的点，例如满足 \(r_{ji}(g_j)=g_i\) 或 \(g_i=a_i\) 的点；有限交性质就是任意有限批要求都能同时满足。

\ProseParagraphBreak
有限子覆盖性质等价于每个具有有限交性质的闭集族都有非空总交。事实上，若这种闭集族的总交为空，其开补集就覆盖 \(X\)；有限子覆盖会使相应闭集的有限交为空，矛盾。反过来，若一个开覆盖没有有限子覆盖，则开集的闭补集具有有限交性质，总交却为空，也与闭集族的条件矛盾。在 Hausdorff 空间中，这个闭集判据便等价于紧性。有限交判据及闭子空间的结论另在\cref{b1:prop:compact-fip}中统一陈述并证明。

\ProseParagraphBreak
设 \(X\) 为集合。若 \(\mathcal F\subseteq\mathcal P(X)\) 包含 \(X\) 而不包含空集，对有限交封闭，并且包含其每个成员的所有超集，则称 \(\mathcal F\) 为\term[zhen-lvzi]{真滤子}[proper filter]。若真滤子在包含关系下极大，则称为\term[chaolvzi]{超滤子}[ultrafilter]。这里的极大是指不能再严格扩大而仍保持为真滤子，并不是包含所有其他真滤子的最大者；尤其不能把空集加入其中。

\ProseParagraphBreak
每个有限离散空间本身都是紧空间，因为有限多个点就能从任意开覆盖中选出有限子覆盖，单点开集又可分离不同点。对任意指标的乘积，我们将把有限条件组成的滤子扩大为超滤子，由它在每个有限坐标中确定一个值，再验证所得点满足全部原条件。

\begin{lemma}\label{b1:lem:finite-discrete-product-compact}
任意一族有限离散空间的直积，在积拓扑下是紧空间。
\end{lemma}
\begin{proof}
记直积为 \(X=\prod_{i\in I}X_i\)。两个不同点必有某个坐标不同，规定该坐标值的两个基本开集将它们分离，故 \(X\) 为 Hausdorff 空间。若 \(X\) 为空，紧性直接成立。以下给定具有有限交性质的闭集族 \((C_\alpha)_\alpha\)，证明总交非空。

把原闭集族的有限交及其超集组成的族记为
\[
\mathcal F=
\left\{\,A\subseteq X\,\middle|\,
C_{\alpha_1}\cap\cdots\cap C_{\alpha_r}\subseteq A
\text{ 对某个有限指标列表成立}\,\right\},
\]
其中允许空列表，其交为 \(X\)。它包含 \(X\)，对取超集封闭；两个成员各包含一个有限交，它们的交便包含这两批闭集的共同有限交，故仍属于 \(\mathcal F\)。有限交性质又保证这些交都非空，所以 \(\varnothing\notin\mathcal F\)。这证明 \(\mathcal F\) 为真滤子，并包含每个 \(C_\alpha\)。

对包含 \(\mathcal F\) 的全体真滤子按包含关系应用\cref{b1:thm:A-zorn}。这些滤子组成 \(\mathcal P(\mathcal P(X))\) 的子集，因而是一个集合；它非空，因为 \(\mathcal F\) 本身在其中。若一条非空链中的滤子取并，该并仍包含 \(X\)，不包含空集，并且对超集封闭。有限多个属于这个并的集合，分别属于链中的有限多个滤子；链全序，故其中有一个滤子同时包含它们，有限交仍属于该滤子，也属于链的并。因此链的并是一个包含 \(\mathcal F\) 的真滤子，作为上界；空链则以 \(\mathcal F\) 为上界。Zorn 引理给出一个极大者 \(\mathcal U\)。任何严格包含它的真滤子也会包含 \(\mathcal F\)，所以这里的极大者确为超滤子。

对任何 \(A\subseteq X\)，\(A\) 与 \(X\setminus A\) 恰有一个属于 \(\mathcal U\)。它们不能同时属于，因为交为空。现设 \(A\notin\mathcal U\)。若有某个 \(B\in\mathcal U\) 与 \(A\) 不相交，则 \(B\subseteq X\setminus A\)，向上封闭性给出 \(X\setminus A\in\mathcal U\)。若反而每个 \(B\in\mathcal U\) 都与 \(A\) 相交，则加入 \(A\) 后的任意有限交为 \(B_1\cap\cdots\cap B_r\cap A\)；前 \(r\) 个集合的交仍属于 \(\mathcal U\)，按假设与 \(A\) 相交，所以这个有限交非空。更具体地，族
\[
\mathcal V=\{\,D\subseteq X\mid B\cap A\subseteq D
\text{ 对某个 }B\in\mathcal U\text{ 成立}\,\}
\]
包含 \(\mathcal U\) 和 \(A\)，对有限交与取超集封闭，且不包含空集，因而是严格包含 \(\mathcal U\) 的真滤子，违背极大性。这就证明了二择一性质。

对于 \(X\) 的任意有限划分，恰有一块属于 \(\mathcal U\)。不同块的交为空，故至多一块属于；若一块也没有，二择一性质使全部块的补集都属于 \(\mathcal U\)，有限求交却为空，矛盾。现在对每个坐标 \(i\)，纤维 \(p_i^{-1}(\{t\})\)、\(t\in X_i\)，恰好给出一个有限划分，因而唯一确定其中属于 \(\mathcal U\) 的一块。将相应的坐标值记为 \(a_i\)，得到点 \(a=(a_i)_i\in X\)。对于每个有限集 \(F\subseteq I\)，柱状集合
\[
B_F(a)=\bigcap_{i\in F}p_i^{-1}(\{a_i\})
\]
属于 \(\mathcal U\)；这些集合构成 \(a\) 的基本开邻域，因为每个 \(X_i\) 都离散。

若 \(a\notin C_\alpha\)，其开补集包含 \(a\)。积拓扑的定义保证它包含一个只限制有限坐标的基本开邻域；各坐标空间离散，再把所限制的坐标收紧到单点，就有某个有限集 \(F\subseteq I\) 满足 \(a\in B_F(a)\subseteq X\setminus C_\alpha\)。可是 \(B_F(a),C_\alpha\) 都属于 \(\mathcal U\)，它们的交却为空，矛盾。所以 \(a\) 属于全部 \(C_\alpha\)，证明总交非空。
\end{proof}
这里借助超滤子同时作出与有限交条件相容的坐标选择。闭子空间的紧性也随即得到：闭子空间中的闭集仍是整个空间中的闭集，有限交性质可以原样用于整个紧空间，分离邻域与子空间相交又保证 Hausdorff 性。

\ProseParagraphBreak
以后还会用到两个基本事实。连续像保留有限子覆盖性质：将像上的开覆盖拉回，取有限子覆盖再映回即可；目标为Hausdorff空间时，像因而紧。Hausdorff 空间中的紧子集闭：给紧子集外一点 \(x\)，对每个子集内点 \(y\) 取分离 \(x,y\) 的开邻域，覆盖紧子集的那些邻域可选有限个；对应的 \(x\) 邻域有限求交，就与整个紧子集不交。因而从紧空间到 Hausdorff 空间的连续双射为同胚：闭集的像紧而闭，使逆映射连续。这些事实汇总为\cref{b1:prop:compact-maps}，将在比较 Galois 商群的拓扑时使用。

\ProseParagraphBreak
上述紧性证明中，选择原理明确用在把有限相容条件生成的真滤子扩充为超滤子的 Zorn 步骤。极大性使它对每个子集与其补集作出二择一决定，进而在每个有限划分中决定唯一一块。每个坐标的值一旦由 \(\mathcal U\) 决定，就是唯一的，不再另作任意选择。代数闭包与嵌入延拓中的 Zorn 引理用于取得极大的部分代数扩张或部分嵌入；这里扩大的是记录相容条件的集合族，两处取得的极大对象不同。

\subsection{紧性、全不连通性与开子群}
\label{b1:subsec:2202:03}
设 \(X\) 为拓扑空间，\(A\subseteq X\) 采用子空间拓扑。若 \(A\) 不能写成两个不相交非空开集的并，则称 \(A\) 为连通子空间。若 \(X\) 的每个非空连通子空间都只有一个点，则称 \(X\) \term[quan-buliantong]{全不连通}[totally disconnected]。

\begin{theorem}\label{b1:thm:krull-compact}
设 \(L/K\) 为代数 Galois 扩张，\(G=\operatorname{Gal}(L/K)\)，并在 \(G\) 上赋予 Krull 拓扑。对 \(L/K\) 的每个有限 Galois 子扩张 \(N/K\)，记 \(U_N=\operatorname{Gal}(L/N)\)。则 \(G\) 是紧、Hausdorff、全不连通的拓扑群。子群 \(H\leq G\) 开，当且仅当它包含某个 \(U_N\)；此时 \(H\) 同时闭且指数有限。
\end{theorem}
\begin{proof}
在有限群直积 \(X=\prod_NG_N\) 中，对 \(N\subseteq M\) 考察映射
\[
\Phi_{M,N}:X\longrightarrow G_N\times G_N,\qquad
(g_P)_P\longmapsto\bigl(\operatorname{res}_{M,N}(g_M),g_N\bigr).
\]
它只依赖两个坐标，目标采用有限离散拓扑，所以连续。相容条件 \(\operatorname{res}_{M,N}(g_M)=g_N\) 定义的集合正是对角集 \(\Delta_N=\{\,(h,h)\mid h\in G_N\,\}\) 的原像；有限离散空间中的对角集闭，故每个相容条件定义闭集。逆极限是这些闭集的交，因此为闭子空间。\cref{b1:lem:finite-discrete-product-compact}及闭子空间的紧性给出 \(G\) 紧。

若 \(\sigma\ne\tau\)，存在 \(a\in L\) 使 \(\sigma(a)\ne\tau(a)\)。由\cref{b1:prop:finite-galois-layers}，\(L\) 是有限 Galois 层之并，可取包含 \(a\) 的有限层 \(N\)，于是 \(\pi_N(\sigma)\ne\pi_N(\tau)\)。两个陪集 \(\sigma U_N,\tau U_N\) 不相交，故 \(G\) 为 Hausdorff 空间。又因 \(\sigma U_N\) 同时开闭，对任何包含 \(\sigma,\tau\) 的子空间 \(A\)，两部分 \(A\cap\sigma U_N\) 与 \(A\setminus\sigma U_N\) 都是非空相对开集，并且不交、合起来为 \(A\)。因此连通子空间不能含两个不同点，得到全不连通性。

若 \(H\) 是开子群，它作为单位元的邻域包含某个 \(U_N\)。反过来，若包含 \(U_N\)，则 \(H\) 为若干个 \(U_N\) 陪集的并，所以开。此时 \(H\) 的指数不超过有限数 \([G:U_N]=|G_N|\)。其补集是其他 \(H\) 陪集的并，而平移保持开集，所以 \(H\) 同时闭。拓扑群运算已由\cref{b1:prop:krull-group-operations}验证。
\end{proof}

全不连通只禁止含有两个不同点的连通子空间，并不要求每个单点都是开集。特别地，若这里的 \(G\) 无限，就不可能离散：否则全体单点开集覆盖 \(G\)，却没有有限子覆盖，与紧性矛盾。因此有限层陪集能够分离不同点，不表示每个点都已经孤立。

\ProseParagraphBreak
设 \(H\) 为拓扑群。若 \(H\) 与某个有限离散群逆系统的逆极限拓扑同构，则称 \(H\) 为\term[toushe-youxian-qun]{投射有限群}[profinite group]；这里拓扑同构指同时为群同构和同胚。\eqref{b1:eq:galois-inverse-limit}说明 \(G\) 就是这样的群。“投射有限”并不意味着整个群只有有限多个元素；有限的是各层群，每个实际出现的坐标像都是 \(H\) 的有限商群。拓扑也是定义的一部分，不能只保留抽象群同构而忘记开邻域；连续性与稠密性都由这份拓扑决定。

% !TeX root = ../../Book1.tex
% 纲要标题：### 22.3 拓扑 Galois 对应
% 纲要位置：计划纲要.md，第 748 行。

\section{拓扑 Galois 对应}
\label{b1:sec:2203}

有限 Galois 对应将全部子群与全部中间域配对。无限情形中，固定一个元素只是在某个有限层上的条件，因此不动域只能区分子群的有限层信息。正确的对应要在群的一侧取闭子群；这个限制既有直接的证明，也有有限域给出的明确反例。

% 纲要内容（仅作写作计划，尚非教材正文）：
%
% * 中间域与闭子群的对应
% * 闭性不可删除的例子
% * 有限子扩张、开子群与正规子群
%

\subsection{中间域与闭子群的对应}
\label{b1:subsec:2203:01}
对中间域 \(K\subseteq E\subseteq L\)，记 \(G_E=\operatorname{Gal}(L/E)\)。对任意子群 \(H\leq G\)，记
\[
L^H=\{\,a\in L\mid\sigma(a)=a\text{ 对每个 }\sigma\in H\,\}.
\]
不动元素对加减乘及非零取逆封闭，故构成中间域。域扩大时固定它的自同构减少，子群扩大时不动域缩小。

\begin{theorem}[无限 Galois 对应]\label{b1:thm:infinite-galois-correspondence}
设 \(L/K\) 为代数 Galois 扩张，\(G=\operatorname{Gal}(L/K)\)，并在 \(G\) 上赋予 Krull 拓扑。映射
\[
E\longmapsto\operatorname{Gal}(L/E),\qquad
H\longmapsto L^H
\]
给出全部中间域与 \(G\) 的全部闭子群之间的反向一一对应。
\end{theorem}
\begin{proof}
先证明 \(G_E\) 闭。对每个 \(a\in E\)，选含 \(a\) 的有限 Galois 层 \(N\)。固定 \(a\) 的自同构，是有限群 \(G_N\) 中固定该元素的那些自同构的原像，因此是有限个 \(U_N\) 陪集的并，既开又闭。\(G_E\) 是对全部 \(a\in E\) 取这些闭子集的交，故闭。

接着证明 \(L^{G_E}=E\)。包含 \(E\subseteq L^{G_E}\) 由定义成立。若 \(a\in L\setminus E\)，则 \(a\) 在 \(E\) 上的最小多项式次数大于一。它整除 \(a\) 在 \(K\) 上的最小多项式，故可分，并由 \(L/K\) 的正规性在 \(L\) 中分裂。取另一个根 \(b\ne a\)，代入最小多项式给出 \(E(a)\rightarrow L\) 的 \(E\)-嵌入 \(a\mapsto b\)。由\cref{b1:thm:embedding-extension}延拓到 \(L\)。此延拓固定 \(K\)，由正规嵌入判据其像等于 \(L\)，所以属于 \(G_E\)，却不固定 \(a\)。这说明不动域中不可能多出这样的 \(a\)。同一论证也表明，对任何中间域 \(E\)，\(L/E\) 仍正规且可分。

最后取闭子群 \(H\)。显然 \(H\subseteq\operatorname{Gal}(L/L^H)\)。若 \(\sigma\notin H\)，开集 \(G\setminus H\) 含 \(\sigma\)，故有有限 Galois 层 \(N\)，使 \(\sigma U_N\cap H=\varnothing\)。于是 \(\pi_N(\sigma)\notin H_N\coloneqq\pi_N(H)\)。对有限 Galois 扩张 \(N/K\) 应用\cref{b1:thm:finite-galois-correspondence}，有 \(H_N=\operatorname{Gal}(N/N^{H_N})\)。因此存在 \(a\in N^{H_N}\) 不被 \(\pi_N(\sigma)\) 固定。每个 \(h\in H\) 在 \(N\) 上的限制属于 \(H_N\)，所以 \(a\in L^H\)，而 \(\sigma(a)\ne a\)。故 \(\sigma\notin\operatorname{Gal}(L/L^H)\)，得到反向包含。两个操作于是互逆。
\end{proof}
闭性在最后一段有明确用途：它使一个不在 \(H\) 中的自同构能够在某个有限商里与 \(H\) 区分开。没有闭性时，这种有限层分离可能失败。

\subsection{闭性条件与反例}
\label{b1:subsec:2203:02}
子集 \(H\) 的闭包 \(\overline H\) 是包含它的全部闭集之交；一个点属于闭包，当且仅当它的每个开邻域都与 \(H\) 相交。后一个刻画来自取开补集。子群的闭包仍为子群：若 \(x,y\) 的每个有限层邻域都遇到 \(H\)，则可取 \(h,k\in H\) 在同一有限层上分别与 \(x,y\) 相同，\(hk^{-1}\) 就在该层上与 \(xy^{-1}\) 相同，故 \(xy^{-1}\in\overline H\)。

\begin{corollary}\label{b1:cor:fixed-field-closure}
设 \(L/K\) 为代数 Galois 扩张，\(G=\operatorname{Gal}(L/K)\)，并在 \(G\) 上赋予 Krull 拓扑。对任意子群 \(H\leq G\)，记 \(\overline H\) 为 \(H\) 在 \(G\) 中的闭包。则
\[
L^H=L^{\overline H},\qquad
\operatorname{Gal}(L/L^H)=\overline H.
\]
\end{corollary}
\begin{proof}
元素 \(a\) 的稳定子群在上一证明中已经说明是闭子群。因此若它包含 \(H\)，就包含 \(\overline H\)，得到 \(L^H\subseteq L^{\overline H}\)；反向由 \(H\subseteq\overline H\) 得到。对闭子群 \(\overline H\) 应用\cref{b1:thm:infinite-galois-correspondence}，即得第二式。
\end{proof}
下面给出两个不同子群具有相同不动域的例子。取 \(L=\overline{\mathbb F}_q\)、\(K=\mathbb F_q\)，并令 \(F(a)=a^q\)。\cref{b1:thm:finite-field-subfields}及\cref{b1:thm:finite-field-automorphisms}说明有限层为 \(\mathbb F_{q^n}\)，其 Galois 群由 \(F\) 的限制生成，阶为 \(n\)。于是循环子群 \(H=\langle F\rangle\) 在每个有限商中都有满像，故每个基本开陪集都与 \(H\) 相交，即 \(H\) 稠密。

\ProseParagraphBreak
这个循环子群却不等于整个 \(G\)。将 \(n\) 写为 \(n=2^r m\)、\(m\) 为奇数，用中国剩余定理选唯一剩余类 \(e_n\pmod n\)，使
\[
e_n\equiv0\pmod{2^r},\qquad e_n\equiv1\pmod m.
\]
若 \(d\mid n\)，\(e_n\) 模 \(d\) 后满足定义 \(e_d\) 的同一组条件，所以 \(e_n\equiv e_d\pmod d\)。因此有限层上的 \(F^{e_n}\) 相容，\cref{b1:prop:galois-compatible-family}给出一个整体自同构 \(\sigma\)。若它等于某个整数幂 \(F^a\)，则对所有 \(r\) 都有 \(a\equiv0\pmod{2^r}\)，迫使整数 \(a=0\)；但 \(e_3\equiv1\pmod3\)，又排除了 \(a=0\)。故 \(H\ne G\)。

\ProseParagraphBreak
\(F\) 固定的元素满足 \(x^q=x\)，恰为 \(\mathbb F_q\)。所以真子群 \(H\) 与整个 \(G\) 的不动域都是 \(\mathbb F_q\)，说明闭性不能从无限 Galois 对应中删除。

\subsection{有限子扩张、开子群与正规子群}
\label{b1:subsec:2203:03}
\begin{proposition}\label{b1:prop:finite-open-correspondence}
设 \(L/K\) 为代数 Galois 扩张，\(G=\operatorname{Gal}(L/K)\)，并在 \(G\) 上赋予 Krull 拓扑。对中间域 \(K\subseteq E\subseteq L\)，记 \(G_E=\operatorname{Gal}(L/E)\)。则 \([E:K]<\infty\) 当且仅当 \(G_E\) 是 \(G\) 的开子群；此时
\[
[E:K]=[G:G_E].
\]
\end{proposition}
\begin{proof}
若 \(E/K\) 有限，取有限个域生成元并包含在有限 Galois 层 \(N\) 中，则 \(U_N\subseteq G_E\)，所以 \(G_E\) 开。反过来，若 \(G_E\) 开，\cref{b1:thm:krull-compact}给出 \(U_N\subseteq G_E\)。取不动域并用\cref{b1:thm:infinite-galois-correspondence}，得 \(E\subseteq L^{U_N}=N\)，故扩张有限。

选择这样的 \(N\)，限制映射给出陪集集合之间的双射
\(G/G_E\rightarrow G_N/\operatorname{Gal}(N/E)\)：满射来自限制的满射性，两个限制落在同一陪集当且仅当两个原自同构之商固定 \(E\)，所以单射。有限 Galois 对应给出右端的大小为 \([E:K]\)，完成次数公式。
\end{proof}
任意开子群都指数有限。反向若再假设子群闭，则其有限多个陪集都闭，补集也闭，因而子群开。对中间域所对应的子群，\cref{b1:thm:infinite-galois-correspondence}保证闭性，所以有限指数与开性等价。

\begin{theorem}\label{b1:thm:infinite-galois-quotient}
设 \(L/K\) 为代数 Galois 扩张，\(G=\operatorname{Gal}(L/K)\)，并在 \(G\) 上赋予 Krull 拓扑。设 \(H\leq G\) 为闭子群，并令 \(E=L^H\)。则 \(H\) 正规当且仅当 \(E/K\) 正规；在这两个等价条件成立时，\(E/K\) 为 Galois 扩张，限制映射诱导群同构
\[
G/H\cong\operatorname{Gal}(E/K).
\]
\end{theorem}
\begin{proof}
若 \(E/K\) 正规，任何 \(\sigma\in G\) 都保持 \(E\)，故限制给出同态 \(G\rightarrow\operatorname{Gal}(E/K)\)，核为 \(G_E=H\)，所以 \(H\) 正规。任取目标群中的自同构，由\cref{b1:thm:embedding-extension}延拓到 \(L\) 的嵌入，再用 \(L/K\) 正规性得到 \(L\) 的自同构，证明限制满射。群的第一同构定理给出所列同构。

反过来，若 \(H\) 正规，对 \(a\in E,h\in H,\sigma\in G\)，有
\[
h\bigl(\sigma(a)\bigr)
=\sigma\bigl((\sigma^{-1}h\sigma)(a)\bigr)=\sigma(a),
\]
因为 \(\sigma^{-1}h\sigma\in H\)。所以 \(\sigma(E)\subseteq E\)，再用 \(\sigma^{-1}\) 得到相等。任意 \(K\)-嵌入 \(E\rightarrow\overline L\) 可延拓到 \(L\)，延拓属于 \(G\)，因此将 \(E\) 映为自身。\cref{b1:thm:normal-embedding-criterion}说明 \(E/K\) 正规；可分性则从 \(L/K\) 继承。
\end{proof}
这个同构还保持自然拓扑。为了把商拓扑的含义与有限商检测放在一起，我们将在\subsecref{b1:subsec:2205:04}给出其拓扑方面的证明；本小节已完成群与域的代数对应。

% !TeX root = ../../Book1.tex
% 纲要标题：### 22.4 典型绝对 Galois 群
% 纲要位置：计划纲要.md，第 789 行。

\section{典型绝对 Galois 群}
\label{b1:sec:2204}

绝对 Galois 群汇集了基域全部有限可分扩张的信息，但一般很难完整描述。这里有限 Galois 扩张与一般有限可分扩张所对应的群论对象有所不同：有限 Galois 扩张 \(N/K\) 给出有限商群 \(G_K/\operatorname{Gal}(K^{\mathrm{sep}}/N)\cong\operatorname{Gal}(N/K)\)；一般有限可分扩张 \(E/K\) 对应开子群 \(\operatorname{Gal}(K^{\mathrm{sep}}/E)\)，它未必正规，因而未必给出商群。若要用有限商群记录 \(E\)，可先取它的有限正规闭包。上述区别来自\cref{b1:prop:finite-open-correspondence,b1:thm:infinite-galois-quotient}。有限域给出可显式写成逆极限的模型，实数域给出有限模型。它们也说明，连续到有限群的映射必须能由某个有限层决定。

% 纲要内容（仅作写作计划，尚非教材正文）：
%
% * 有限域的绝对 Galois 群；区分全部模数的投射有限整数与单一素数的进整数，以及拓扑生成与抽象生成
% * 实数域与实闭域的联系
% * 连续表示与离散集上的连续作用；正式证明到有限离散群的连续同态由整数 1 的像唯一决定，再解释任意可逆矩阵均可指定为 Frobenius 的像；说明连续群上同调的后续位置
%

\subsection{有限域的绝对 Galois 群}
\label{b1:subsec:2204:01}
把正整数按整除关系排列；当 \(n\mid m\) 时，以模 \(n\) 的约化
\(\mathbb Z/m\mathbb Z\rightarrow\mathbb Z/n\mathbb Z\) 为过渡映射。所得拓扑环
\[
\widehat{\vphantom{Z}\mathbb Z}=\varprojlim_{n\geq1}\mathbb Z/n\mathbb Z,
\]
称为\term[toushe-youxian-zhengshu]{投射有限整数环}[profinite integers]。其元素是相容剩余类族 \((a_n)_n\)，加法与乘法均逐坐标进行；拓扑是有限离散环直积中的子空间拓扑。
\symidx{\(\widehat{\vphantom{Z}\mathbb Z}\)：投射有限整数环}

\begin{theorem}\label{b1:thm:finite-field-absolute-galois}
设 \(q\) 为素数幂，\(G_{\mathbb F_q}=\operatorname{Gal}(\overline{\mathbb F}_q/\mathbb F_q)\) 为有限域 \(\mathbb F_q\) 的绝对 Galois 群，并在其上赋予 Krull 拓扑。令 \(\widehat{\vphantom{Z}\mathbb Z}=\varprojlim_{n\geq1}\mathbb Z/n\mathbb Z\) 按整除关系及自然约化构成，并取逆极限拓扑。则存在拓扑群同构
\[
\widehat{\vphantom{Z}\mathbb Z}\longrightarrow G_{\mathbb F_q},
\qquad(a_n)_n\longmapsto\sigma,
\qquad \sigma|_{\mathbb F_{q^n}}=F^{a_n},
\]
其中 \(F(x)=x^q\)。整数 \(1\) 对应 Frobenius 自同构。
\end{theorem}
\begin{proof}
有限域完美，其可分闭包就是代数闭包。每个代数元 \(a\in\overline{\mathbb F}_q\) 都生成有限扩张 \(\mathbb F_q(a)/\mathbb F_q\)；若其次数为 \(n\)，这个域便有 \(q^n\) 个元素，因此是当前闭包中由 \(X^{q^n}-X\) 的全部根组成的唯一子域 \(\mathbb F_{q^n}\)。由\cref{b1:thm:finite-field-subfields}，这些有限子域的包含关系对应 \(n\mid m\)。故它们已经覆盖整个代数闭包，并给出全部有限 Galois 层。

由\cref{b1:thm:finite-field-automorphisms}，第 \(n\) 层的自同构群是阶 \(n\) 的循环群，以 \(F\) 的限制为生成元。把它与加法群 \(\mathbb Z/n\mathbb Z\) 识别以后，\(F^a\) 从第 \(m\) 层限制到第 \(n\) 层，只需把指数 \(a\) 模 \(n\) 约化。因此\cref{b1:prop:galois-compatible-family}给出所述群同构。在逆极限一侧，规定第 \(n\) 个坐标的剩余类，对应于在 Galois 群一侧规定向 \(\mathbb F_{q^n}\) 的限制；有限多个坐标条件又可合并到一个共同倍数所对应的坐标。这两类条件分别构成两侧拓扑的邻域基，故该群同构及其逆都连续。
\end{proof}
整数到 \(\widehat{\vphantom{Z}\mathbb Z}\) 的自然映射为 \(a\mapsto(a\bmod n)_n\)，它单射，因为被所有正整数整除的整数只能为零。其像稠密：有限多个模数的条件可由一个共同倍数的坐标统一规定，再选该剩余类的一个整数代表。\subsecref{b1:subsec:2203:02}中构造的 \((e_n)_n\) 表明这个像不等于全部投射有限整数。因此整数幂组成的普通循环子群 \(\langle F\rangle=\{F^a:a\in\mathbb Z\}\cong\mathbb Z\) 是真子群，却满足
\[
\overline{\langle F\rangle}=G_{\mathbb F_q}.
\]
称 \(F\) 为拓扑生成元，是说它的整数幂稠密，而不是说每个群元素都是某个整数幂。

\ProseParagraphBreak
\(\widehat{\vphantom{Z}\mathbb Z}\) 同时记录所有正整数模数的相容信息。给定素数 \(p\)，改为只记录模数为 \(p^r\) 的坐标，所得环
\(\mathbb Z_p=\varprojlim_r\mathbb Z/p^r\mathbb Z\) 称为 \(p\) 进整数环；从 \(\widehat{\vphantom{Z}\mathbb Z}\) 到 \(\mathbb Z_p\) 的自然映射就是舍去其余坐标。两种逆极限所保留的有限信息不同。以后讨论 \(p\) 进表示时，正是这些模 \(p^r\) 的有限信息决定连续性。
\symidx{\(\mathbb Z_p\)：\(p\) 进整数环}

\subsection{实数域与实闭域}
\label{b1:subsec:2204:02}
\begin{proposition}\label{b1:prop:real-absolute-galois}
实数域的绝对 Galois 群为 \(G_{\mathbb R}=\operatorname{Gal}(\mathbb C/\mathbb R)\cong C_2\)，其中非单位元为复共轭；其 Krull 拓扑为离散拓扑。
\end{proposition}
\begin{proof}
由 \(\mathbb C=\mathbb R(i)\)、\(i^2=-1\)，以及 \(X^2+1\) 在 \(\mathbb R\) 上不可约，得 \([\mathbb C:\mathbb R]=2\)，所以 \(\mathbb C/\mathbb R\) 是代数扩张。\cref{b1:thm:fundamental-theorem-algebra}又保证 \(\mathbb C\) 代数闭，故它确为 \(\mathbb R\) 的一个代数闭包；该定理的证明见\secref{b1:sec:E02}。特征零保证所有代数元可分，所以这个代数闭包也是可分闭包。

\(\mathbb C\) 是 \(X^2+1\) 的分裂域。自同构必须把 \(i\) 送到 \(i\) 或 \(-i\)，两种选择分别给出恒等与复共轭，所以群为 \(C_2\)。整个闭包已经是有限 Galois 层，因而 \(\operatorname{Gal}(\mathbb C/\mathbb C)=\{1\}\) 是基本开子群；它的每个陪集都是单点开集，故拓扑离散。
\end{proof}
这个例子的代数闭性来自实闭域刻画与实分析的中间值定理；自同构群及其拓扑则由二次扩张直接算出。

\ProseParagraphBreak
相同计算对特征零、满足“\(-1\) 不是平方且添加 \(i\) 后得到代数闭域”的域 \(F\) 也适用：\(F(i)/F\) 为二次可分闭包，故 \(G_F\cong C_2\)。这些条件是实闭域的一种代数刻画，具体见\cref{b1:thm:real-closed-characterizations}；\appref{b1:app:E}从有序域与实闭包出发加以说明。此处只作这一条件下的计算，不把任意有序域或任意实数子域都视为实闭域。例如 \(\mathbb Q\subseteq\mathbb R\)，但 \(\mathbb Q\) 还有许多不同次数的有限 Galois 扩张。

\subsection{连续表示与离散集上的连续作用}
\label{b1:subsec:2204:03}
\begin{definition}\label{b1:def:continuous-finite-representation}
设 \(K\) 为域，在一个固定代数闭包中取可分闭包 \(K^{\mathrm{sep}}\)，并令 \(G_K=\operatorname{Gal}(K^{\mathrm{sep}}/K)\) 为赋予 Krull 拓扑的绝对 Galois 群；再设 \(k\) 为有限域，\(d\geq1\) 为整数。如果群同态
\[
\rho:G_K\longrightarrow\operatorname{GL}_d(k),
\]
在目标群 \(\operatorname{GL}_d(k)\) 取离散拓扑时连续，则称 \(\rho\) 为 \(G_K\) 在 \(k^d\) 上的 \(d\) 维\term[lianxu-biaoshi]{连续表示}[continuous representation]。这里 \(\operatorname{GL}_d(k)\) 表示可逆矩阵群，\(\rho\) 规定群元素在线性空间 \(k^d\) 上怎样作用。
\end{definition}
\symidx{\(\operatorname{GL}_d(k)\)：域上的一般线性群}
连续性使 \(\rho^{-1}(\{I_d\})\) 成为单位元的开邻域，所以它包含某个 \(U_N=\operatorname{Gal}(K^{\mathrm{sep}}/N)\)，其中 \(N/K\) 是有限 Galois 扩张。这意味着两个自同构只要在 \(N\) 上限制相同，它们作用在 \(k^d\) 上的矩阵就相同：每个连续表示的全部作用信息已经由某个有限 Galois 商决定。这个判别及其逆命题将在\cref{b1:prop:finite-continuous-hom}中完整证明。对于有限域，Frobenius 的一个像就能确定这样的表示，但这里还需要连续性。

\begin{proposition}\label{b1:prop:profinite-integer-finite-hom}
设 \(D\) 为赋予离散拓扑的有限群，并令 \(\widehat{\vphantom{Z}\mathbb Z}=\varprojlim_{n\geq1}\mathbb Z/n\mathbb Z\) 按整除关系及自然约化构成，赋予逆极限拓扑。把 \(\widehat{\vphantom{Z}\mathbb Z}\) 看成加法群，并用 \(1\) 表示整数 \(1\) 的自然像。则连续群同态
\[
\rho:\widehat{\vphantom{Z}\mathbb Z}\longrightarrow D
\]
通过 \(\rho\mapsto\rho(1)\) 与 \(D\) 的元素建立一一对应。若 \(d\in D\) 的阶为 \(m\)，对应同态可先约化到 \(\mathbb Z/m\mathbb Z\)，再令 \(\bar a\mapsto d^a\)。
\end{proposition}
\begin{proof}
给定 \(d\in D\)，它的阶 \(m\) 有限。因为 \(d^m=1_D\)，映射 \(\mathbb Z/m\mathbb Z\rightarrow D\)、\(\bar a\mapsto d^a\) 定义良好且为群同态。把它与第 \(m\) 个坐标投影复合，便得到所需同态。坐标投影连续，有限离散群之间的映射也连续，所以这个复合同态连续。

若两个连续同态 \(\rho,\eta\) 在 \(1\) 上取值相同，则由同态性，它们在所有整数的自然像上取值相同。这个整数子群稠密，而等值集合
\[
\{x:\rho(x)=\eta(x)\}
=\bigcup_{c\in D}\bigl(\rho^{-1}(\{c\})\cap\eta^{-1}(\{c\})\bigr)
\]
是闭集：目标离散，单点闭，两映射连续，并且这里是有限个闭集的并。等值集合包含稠密子群，故为整个 \(\widehat{\vphantom{Z}\mathbb Z}\)，这就证明唯一性。
\end{proof}
这里的唯一性来自整数子群的稠密性与同态的连续性，不能用“\(1\) 生成整个抽象群”来解释。由\cref{b1:thm:finite-field-absolute-galois}，对任一 \(A\in\operatorname{GL}_d(k)\)，恰有一个连续表示 \(\rho:G_{\mathbb F_q}\rightarrow\operatorname{GL}_d(k)\) 满足 \(\rho(F)=A\)；它经过模 \(\operatorname{ord}(A)\) 的有限循环商。这是有限域绝对 Galois 群的特殊性质，一般绝对 Galois 群的元素及其关系会限制可指定的矩阵像。

\ProseParagraphBreak
例如设 \(\ell\) 为素数，取系数域 \(k=\mathbb F_\ell\)。矩阵
\[
A=\begin{pmatrix}1&1\\0&1\end{pmatrix}
\]
满足 \(A^r=\begin{psmallmatrix}1&r\\0&1\end{psmallmatrix}\)，因此其阶为 \(\ell\)。令 Frobenius 元 \(F\) 的像为 \(A\)，便得到经过 \(\mathbb Z/\ell\mathbb Z\) 的非平凡二维连续表示。每个群元素的像都是某个 \(A^r\)，其迹均为 \(2\cdot1_k\)，与二维平凡表示的迹相等；两表示却不同构，因为 \(F\) 在其中分别作用为 \(A\) 与单位矩阵，而 \(A\ne I_2\) 不可能与 \(I_2\) 共轭。

\ProseParagraphBreak
这个表示中，由第一个标准基向量张成的直线是平凡子表示，商空间上的作用也平凡，但整个表示不是两个平凡表示的直和。迹只给出这两个一维作用的迹之和，不能辨认它们在这里怎样组合成二维作用。系数域的特征 \(\ell\) 正好整除有限像群的阶 \(\ell\)，这一例说明在这种特征下，迹函数可能无法区分不同构的表示。

\ProseParagraphBreak
设拓扑群 \(G_K\) 作用在离散集合 \(M\) 上。如果作用映射 \(G_K\times M\rightarrow M\) 连续，则称这个作用连续。在 \(M\) 离散时，这等价于每个元素的稳定子群开：连续性在 \((1,m)\) 附近给出一个固定 \(m\) 的开邻域；反向则在每个 \((g,m)\) 用稳定子群的陪集构造保持作用值的开邻域。第四卷附录C将定义连续离散模的群上同调，并介绍相应的上链与系数条件；本章的拓扑群知识是学习这些内容的基础。

% !TeX root = ../../Book1.tex
% 纲要标题：### 22.5 有限层检测、闭包与连续性
% 纲要位置：计划纲要.md，第 760 行。

\section{有限层检测、闭包与连续性}
\label{b1:sec:2205}

闭子群可以通过全部有限商中的像辨认。到有限离散群的同态，其连续性可以通过核是否包含某个基本开子群判定；到赋有逆极限拓扑的群的同态，则可逐个有限坐标检验连续性。本节将这些判别写成公式，并说明商群的拓扑为何仍须保留闭性要求。

% 纲要内容（仅作写作计划，尚非教材正文）：
%
% * 闭子群由其在各有限 Galois 商中的像确定；写出子群闭包的有限商检测公式
% * 以有限域的 Frobenius 生成的稠密循环子群比较 \(\mathbb Z\) 与 \(\widehat{\mathbb Z}\)，说明具有同一不动域的子群为何必须取闭包
% * 连续到有限离散群的同态具有开核；有限层上的表示与相容的连续表示
% * 无限 Galois 对应中的商群结论要求闭正规子群；说明与第四卷连续上同调的联系，不调用其定理
%
% ---
%

\subsection{闭子群与闭包的有限商检测}
\label{b1:subsec:2205:01}
\begin{proposition}\label{b1:prop:closure-finite-quotients}
设 \(L/K\) 为代数 Galois 扩张，\(G=\operatorname{Gal}(L/K)\)，并在 \(G\) 上赋予 Krull 拓扑。令 \(\mathcal N\) 为 \(L/K\) 的全部有限 Galois 子扩张组成的有向集合；对 \(N\in\mathcal N\)，记 \(\pi_N:G\rightarrow\operatorname{Gal}(N/K)\) 为限制映射，\(U_N=\ker\pi_N=\operatorname{Gal}(L/N)\)。则对任意子群 \(H\leq G\)，有
\begin{equation}\label{b1:eq:closure-finite-quotients}
\overline H=\bigcap_{N\in\mathcal N}\pi_N^{-1}\bigl(\pi_N(H)\bigr)
=\bigcap_{N\in\mathcal N}HU_N.
\end{equation}
其中 \(\overline H\) 表示 \(H\) 在 \(G\) 中的闭包。所以两个闭子群相等，当且仅当它们在每个有限 Galois 商中的像相等。
\end{proposition}
\begin{proof}
点 \(\sigma\) 属于 \(\overline H\)，当且仅当每个基本开邻域 \(\sigma U_N\) 与 \(H\) 相交。相交意味着有 \(h\in H\) 满足 \(\pi_N(h)=\pi_N(\sigma)\)，这等价于 \(\sigma\in\pi_N^{-1}\bigl(\pi_N(H)\bigr)\)，也等价于 \(\sigma\in HU_N\)。逐个 \(N\) 求交得到公式。若两个闭子群的有限像相同，两者分别等于公式中的同一个交；反向由取像直接成立。
\end{proof}
把这个公式与\cref{b1:cor:fixed-field-closure}合用，得到一个完整的辨别准则：任意两个子群具有相同不动域，当且仅当它们的闭包相同，也当且仅当它们在全部有限商中的像相同。对子群本身而言，有限信息只确定它的闭包。

\subsection{Frobenius 的幂、闭包与有限子域}
\label{b1:subsec:2205:02}
设 \(q\) 为素数幂，\(F(x)=x^q\)。在 \(G_{\mathbb F_q}=\operatorname{Gal}(\overline{\mathbb F}_q/\mathbb F_q)\cong\widehat{\vphantom{Z}\mathbb Z}\) 中，\(\langle F\rangle\) 对应整数的自然像 \(\mathbb Z\)。\subsecref{b1:subsec:2203:02}已证明这一像稠密而非全群；现在考虑 \(F\) 的幂所生成的子群。

\ProseParagraphBreak
固定正整数 \(d\)，有
\begin{equation}\label{b1:eq:profinite-multiple-kernel}
d\widehat{\vphantom{Z}\mathbb Z}
=\ker\left(\widehat{\vphantom{Z}\mathbb Z}\longrightarrow\mathbb Z/d\mathbb Z\right).
\end{equation}
一个方向由乘 \(d\) 直接得到。反过来，设相容族 \(a\) 在模 \(d\) 下为零。对每个 \(n\)，选择 \(a_{dn}\) 的一个被 \(d\) 整除的整数代表，除以 \(d\) 后取模 \(n\)，记为 \(b_n\)。改变代表只改变 \(n\) 的倍数，故定义不变；原族的相容性说明 \((b_n)_n\) 相容，并满足 \(db=a\)。这证明另一方向。

\ProseParagraphBreak
\(\langle F^d\rangle\) 在第 \(n\) 个商 \(\mathbb Z/n\mathbb Z\) 中的像为
\[
d(\mathbb Z/n\mathbb Z)=\{[dk]_n:k\in\mathbb Z\}
=\frac{d\mathbb Z+n\mathbb Z}{n\mathbb Z}.
\]
这些有限像也正是 \(d\widehat{\vphantom{Z}\mathbb Z}\) 的有限像；后者由\eqref{b1:eq:profinite-multiple-kernel}知是闭子群。因此\eqref{b1:eq:closure-finite-quotients}给出
\(\overline{\langle F^d\rangle}=d\widehat{\vphantom{Z}\mathbb Z}\)。它恰为固定 \(\mathbb F_{q^d}\) 的自同构组成的子群，所以 \(\langle F^d\rangle\) 与其闭包的不动域同为 \(\mathbb F_{q^d}\)。

下表比较这几个子群的拓扑性质与不动域，其中 \(d\geq2\)，\(\mathbb Z\) 表示整数在 \(\widehat{\vphantom{Z}\mathbb Z}\) 中的自然像。\(d\widehat{\vphantom{Z}\mathbb Z}\) 由\eqref{b1:eq:profinite-multiple-kernel}知既闭又开；\(\mathbb Z\) 稠密而非全群，因而不闭，也不可能开。零子群闭；若它开，紧群 \(\widehat{\vphantom{Z}\mathbb Z}\) 就成为有限离散空间，与它无限相矛盾。特别地，\(\mathbb Z\) 与全群虽不相等，却有同一个不动域。
\begin{center}
\begin{tabular}{@{}lccc@{}}
\multicolumn{4}{c}{Frobenius 子群的闭性、开性与不动域}\\
\addlinespace[0.4em]
\toprule
子群 \(H\) & 闭 & 开 & \(\overline{\mathbb F}_q^{\,H}\) \\
\midrule
\(\widehat{\vphantom{Z}\mathbb Z}\) & 是 & 是 & \(\mathbb F_q\) \\
\(d\widehat{\vphantom{Z}\mathbb Z}\) & 是 & 是 & \(\mathbb F_{q^d}\) \\
\(\mathbb Z\) & 否 & 否 & \(\mathbb F_q\) \\
\(\{0\}\) & 是 & 否 & \(\overline{\mathbb F}_q\) \\
\bottomrule
\end{tabular}
\end{center}

\subsection{开核与有限层表示}
\label{b1:subsec:2205:03}
\begin{proposition}\label{b1:prop:finite-continuous-hom}
设 \(L/K\) 为代数 Galois 扩张，\(G=\operatorname{Gal}(L/K)\)，并在 \(G\) 上赋予 Krull 拓扑。对每个有限 Galois 子扩张 \(N/K\)，记 \(G_N=\operatorname{Gal}(N/K)\)，\(\pi_N:G\rightarrow G_N\) 为限制映射。再设 \(D\) 为有限离散群，\(\rho:G\rightarrow D\) 为群同态。以下条件等价：\(\rho\) 连续；\(\ker\rho\) 开；存在某个这样的 \(N\) 及同态 \(\rho_N:G_N\rightarrow D\)，使 \(\rho=\rho_N\pi_N\)。
\end{proposition}
\begin{proof}
连续性使单点 \(\{1\}\) 的原像开。若核开，取 \(U_N\subseteq\ker\rho\)，规定 \(\rho_N(\sigma|_N)=\rho(\sigma)\)。两个提升之商在 \(U_N\) 中，故像相同；限制满射保证公式在全部 \(G_N\) 上有定义，并保持群乘法。反过来，有限层投影连续，有限离散集合之间的任意映射连续，所以经过有限商的同态连续。
\end{proof}
有限层不一定唯一。若某表示经过 \(G_N\)，也经过任何包含 \(N\) 的更大有限 Galois 层；应比较其限制相容性，而不是把选择层的不同误当作表示不同。

\ProseParagraphBreak
取素数 \(\ell\) 和整数 \(d\geq1\)。相容的有限层表示可合成以 \(\ell\) 进整数环为系数的表示。给 \(M_d(\mathbb Z_\ell)\cong\mathbb Z_\ell^{d^2}\) 赋予积拓扑，并给 \(\operatorname{GL}_d(\mathbb Z_\ell)\) 赋予子空间拓扑；这里的矩阵群不取离散拓扑。作为拓扑群，\(\operatorname{GL}_d(\mathbb Z_\ell)\) 可识别为
\[
\varprojlim_r\operatorname{GL}_d(\mathbb Z/\ell^r\mathbb Z).
\]
相容矩阵逐项给出 \(\mathbb Z_\ell\) 中的矩阵，有限层的逆矩阵也相容，从而给出整体逆矩阵。两侧拓扑均由模 \(\ell^r\) 的坐标条件生成，故这一识别还是同胚。具体地，若
\[
V_r=\ker\!\left(\operatorname{GL}_d(\mathbb Z_\ell)
\longrightarrow\operatorname{GL}_d(\mathbb Z/\ell^r\mathbb Z)\right),
\]
则 \(AV_r\)（\(r\geq1\)）构成矩阵 \(A\in\operatorname{GL}_d(\mathbb Z_\ell)\) 的邻域基。若 \(L/K\) 为代数 Galois 扩张，\(G=\operatorname{Gal}(L/K)\) 取 Krull 拓扑，且 \(\rho_r:G\rightarrow\operatorname{GL}_d(\mathbb Z/\ell^r\mathbb Z)\) 均连续、模约化后相容，就得到唯一同态 \(\rho:G\rightarrow\operatorname{GL}_d(\mathbb Z_\ell)\)。目标的基本开集只限制有限个坐标，其原像为有限多个开集的交，故 \(\rho\) 连续；反向将连续的 \(\rho\) 与各坐标投影复合即可。每个 \(\rho_r\) 都可由\cref{b1:prop:finite-continuous-hom}经过某个有限 Galois 商 \(\operatorname{Gal}(N_r/K)\)。这里 \(r\) 记录矩阵系数的精度，\(N_r\) 记录决定这一级表示的域扩张，两者是不同的层次；随着 \(r\) 增大，所需的 \(N_r\) 也可以增大。逐个精度能经过有限商，并不要求存在一个固定有限层同时决定全部精度，因而整体 \(\rho\) 不必有开核或有限像。

\ProseParagraphBreak
以有限域为例，由\cref{b1:thm:finite-field-absolute-galois}，可以把 \(G_{\mathbb F_q}\) 识别为加法群 \(\widehat{\vphantom{Z}\mathbb Z}\)。对相容族 \(a=(a_n)_n\)，只取素数幂模数的坐标，得到
\[
a_\ell=(a_{\ell^r})_{r\geq1}\in\mathbb Z_\ell.
\]
定义二维表示
\[
\rho(a)=\begin{pmatrix}1&a_\ell\\0&1\end{pmatrix}
\in\operatorname{GL}_2(\mathbb Z_\ell).
\]
矩阵乘法将右上角相加，故这是群同态。其模 \(\ell^r\) 约化只依赖 \(a_{\ell^r}\)，因此连续，并经过有限 Galois 层 \(N_r=\mathbb F_{q^{\ell^r}}\)；全部约化相容，所以 \(\rho\) 连续。普通整数 \(k\) 的像是 \(\begin{psmallmatrix}1&k\\0&1\end{psmallmatrix}\)，不同整数在 \(\mathbb Z_\ell\) 中仍不同，故整体像无限。若所有约化都经过同一个有限商，那么两元素在该商中相等时，其矩阵在每个模 \(\ell^r\) 下都相等，从而整体矩阵相等；这会使 \(\rho\) 本身经过该有限商，像也有限，矛盾。因此这一例中没有统一的有限层。其核也不开：若核开，紧性使其陪集只有有限个，而这些陪集与表示的像一一对应，又会迫使像有限。

\subsection{闭正规子群与 Galois 商群}
\label{b1:subsec:2205:04}
设 \(G\) 为拓扑群，\(H\trianglelefteq G\) 为正规子群，\(q:G\rightarrow G/H\) 为自然映射。在抽象商群 \(G/H\) 上定义\term[shang-tuopu]{商拓扑}[quotient topology]：集合 \(V\subseteq G/H\) 开，当且仅当 \(q^{-1}(V)\) 在 \(G\) 中开。本节应用时，\(G=\operatorname{Gal}(L/K)\)，\(H\) 还是闭子群，并令 \(E=L^H\)。

\begin{proposition}\label{b1:prop:galois-quotient-homeomorphism}
设 \(L/K\) 为代数 Galois 扩张，\(G=\operatorname{Gal}(L/K)\)，并在 \(G\) 上赋予 Krull 拓扑。设 \(H\leq G\) 为闭正规子群，\(E=L^H\)。在 \(G/H\) 上赋予商拓扑，在 \(\operatorname{Gal}(E/K)\) 上赋予 Krull 拓扑。则限制映射诱导的群同构
\(G/H\rightarrow\operatorname{Gal}(E/K)\) 是拓扑群同构。
\end{proposition}
\begin{proof}
先看限制同态 \(r:G\rightarrow\operatorname{Gal}(E/K)\)。目标中一个基本邻域规定某个有限 Galois 子扩张 \(F/K\)、\(F\subseteq E\) 上的作用，其原像在 \(G\) 中也是同样规定 \(F\) 上作用的基本开集，所以 \(r\) 连续。写成 \(r=\bar r q\)，由商拓扑定义可知诱导双射 \(\bar r\) 连续。

若 \(C\subseteq G/H\) 闭，则 \(q^{-1}(C)\) 是紧群 \(G\) 的闭子集，因而紧。由\cref{b1:prop:compact-maps}，其连续像 \(r\bigl(q^{-1}(C)\bigr)=\bar r(C)\) 在 Hausdorff 群 \(\operatorname{Gal}(E/K)\) 中紧，故闭。因此 \(\bar r\) 还是闭映射，逆映射连续，所列同构为同胚。
\end{proof}
若 \(H\) 是非闭正规子群，\(G/H\) 的商拓扑就不可能是 Hausdorff：否则单位元单点 \(\{H\}\) 闭，由商映射 \(q\) 的连续性，\(H=q^{-1}(\{H\})\) 也闭，矛盾。以 \(G=\widehat{\vphantom{Z}\mathbb Z}\)、\(H=\mathbb Z\) 的稠密像为例，商群非平凡，商拓扑却只有空集和全空间两个开集。事实上，若商中有非空开集，其原像 \(V\) 是非空开集且满足 \(V+H=V\)。对任意 \(g\in G\)，开集 \(g-V\) 与稠密的 \(H\) 相交，于是 \(g\in V+H=V\)，故 \(V=G\)。

\ProseParagraphBreak
无限 Galois 对应将中间域与闭子群配对，有限中间扩张对应开子群，正规中间扩张对应闭正规子群；相应中间扩张的 Galois 群是原群对该闭正规子群的拓扑商群。更进一步的连续群上同调和算术表示，还需指定系数对象及相应的连续性条件。


\begin{chaptersummary}
有限层上的限制将一个 Galois 自同构变为相容坐标族，Krull 拓扑则规定有限个坐标所能观察到的信息。相容条件把逆极限刻画为有限离散群直积中的闭子群，由此得到紧性、Hausdorff 性和全不连通性。证明中的选择原理用于极大真滤子的存在。

\ProseParagraphBreak
闭子群恰好能从其不动域恢复。任意子群与其闭包有相同不动域，有限域中Frobenius生成的稠密真子群给出这一现象的明确例子。有限子扩张对应开子群，正规子扩张对应闭正规子群；商群同构还须带上自然的商拓扑。

\ProseParagraphBreak
这些结论使无限对象仍可通过有限层计算。连续到有限离散群的同态经过某个有限Galois商，相容的有限层表示则可以合成到逆极限系数中的连续表示。有限与无限Galois理论由此有了统一的描述；无限情形须要求子群闭、同态连续。进一步的算术与上同调应用将以这些结论为基础。
\end{chaptersummary}

\BookExercises

\begin{exercise}\label{b1:ex:22-1}
固定 \(\overline{\mathbb F}_q\)，令 \(E_r=\mathbb F_{q^{r!}}\)。证明 \(\overline{\mathbb F}_q=\bigcup_{r\geq1}E_r\)。若给定 \(\tau_r\in\operatorname{Gal}(E_r/\mathbb F_q)\)，且 \(\tau_{r+1}|_{E_r}=\tau_r\)，证明它们唯一确定 \(\overline{\mathbb F}_q\) 的一个 \(\mathbb F_q\)-自同构，并说明只用阶乘层的限制条件仍给出相同的 Krull 拓扑。
\end{exercise}
\begin{solution}
每个代数元属于某个 \(\mathbb F_{q^n}\)。若 \(r\geq n\)，则 \(n\mid r!\)，由\cref{b1:thm:finite-field-subfields}，\(\mathbb F_{q^n}\subseteq E_r\)，故这些域的并是整个闭包。这里并不要求每个有限域都等于某个 \(E_r\)；关键是每个有限层都包含在某个阶乘层中。给定题中的相容族，对任意有限层 \(\mathbb F_{q^n}\) 选取 \(r\) 使 \(n\mid r!\)，用 \(\tau_r\) 的限制定义其作用。换另一 \(s\) 时再取更大的阶乘层，相容性保证两次限制相同。这样得到全部有限层的相容族，且每个限制都固定 \(\mathbb F_q\)，由\cref{b1:prop:galois-compatible-family}唯一拼成整体 \(\mathbb F_q\)-自同构。阶乘层给出的开陪集本来就是 Krull 基本开集，而任意有限层给出的基本开集又能由更大的阶乘层细化，所以只保留这些层也给出相同的 Krull 拓扑。
\end{solution}

\begin{exercise}\label{b1:ex:22-2}
设 \(K=\mathbb F_p(t)\)，在固定代数闭包中令 \(L=\bigcup_{r\geq0}K(t^{1/p^r})\)。证明 \(L/K\) 无限，且 \(\operatorname{Aut}_K(L)=1\)。为什么不能在此使用无限 Galois 对应？
\end{exercise}
\begin{solution}
在 \(\mathbb F_p[t][X]\) 中，多项式 \(X^{p^r}-t\) 对素元 \(t\) 满足\cref{b1:thm:eisenstein}：常数项被 \(t\) 整除但不被 \(t^2\) 整除，其他非首项系数为零。因此 \(K(t^{1/p^r})/K\) 的次数为 \(p^r\)，这些次数无界，故 \(L/K\) 无限。每个 \(t^{1/p^r}\) 是 \(X^{p^r}-t\) 在闭包中的唯一根，任何 \(K\) 自同构都必须固定它，所以固定整个并域。其自同构群的不动域为 \(L\)，并非 \(K\)。扩张是纯不可分的，不满足\cref{b1:thm:infinite-galois-correspondence}的可分性假设。
\end{solution}

\begin{exercise}\label{b1:ex:22-3}
在 \(\overline{\mathbb F}_q\) 中令 \(F(x)=x^q\)。能否同时在 \(\mathbb F_{q^4}\) 上指定 \(F\)，在 \(\mathbb F_{q^6}\) 上指定 \(F^2\)，再延拓为一个整体自同构？若把后一指定改成 \(F^3\)，答案如何？
\end{exercise}
\begin{solution}
两域的交为 \(\mathbb F_{q^2}\)。第一种指定在交域上的限制分别为非恒等的 \(F\) 与恒等的 \(F^2\)，违反\eqref{b1:eq:galois-compatibility}，所以不可能同时延拓。第二种指定在交域上相同。相应的指数须满足
\[
a\equiv1\pmod4,\qquad a\equiv3\pmod6.
\]
这里模数不互素；可解性要求两个指定指数模 \(\gcd(4,6)=2\) 相同，恰好就是交域上的相容性。写 \(a=1+4u\)，第二个同余式化为 \(2u\equiv1\pmod3\)，即 \(u\equiv2\pmod3\)，所以 \(a\equiv9\pmod{12}\)。整体自同构 \(F^9\) 因而同时实现这两个指定。分别可以延拓不等于所选延拓彼此相容，交域上的限制是必须检查的条件。
\end{solution}

\begin{exercise}\label{b1:ex:22-4}
对每个 \(r\geq1\) 令 \(A_r=\mathbb Z/2\mathbb Z\)，过渡映射 \(A_{r+1}\rightarrow A_r\) 全为零映射。求 \(\varprojlim_r A_r\)，并说明逆极限的坐标投影是否必为满射。
\end{exercise}
\begin{solution}
相容族 \((a_r)_r\) 必须满足 \(a_r=0(a_{r+1})=0\)，故只有零族，逆极限为平凡群。每个坐标投影的像都只有零，而 \(A_r\) 有两个元素，所以不满射。\cref{b1:lem:galois-restriction-surjective}为 Galois 系统提供了额外的满射性，不能把它误认为任意逆极限定义的直接后果。
\end{solution}

\begin{exercise}\label{b1:ex:22-5}
设 \((A_i,r_{ji})\) 为有向指标上的有限群逆系统，且所有过渡映射均满射。证明逆极限到每个 \(A_i\) 的坐标投影满射，明确指出紧性在哪一步使用。
\end{exercise}
\begin{solution}
固定 \(i\) 与 \(a\in A_i\)。在非空积空间 \(X=\prod_j A_j\) 中，要求第 \(i\) 坐标为 \(a\)，并要求所有相容等式成立。这些要求分别定义闭集。任取有限多个要求，选取同时大于固定指标 \(i\) 及这一有限组要求中全部指标的 \(k\)。由 \(r_{ki}\) 满射，取 \(b\in A_k\) 使其在第 \(i\) 层的像为 \(a\)；把所有涉及坐标指定为 \(b\) 的相应像，未涉及坐标取单位元，就满足这一有限组要求。因此闭集族有有限交性质。

\cref{b1:lem:finite-discrete-product-compact}使总交非空，其中任一点都是第 \(i\) 坐标为 \(a\) 的相容族。此处从全部有限组可解推出同时可解，正是紧性步骤；正文给出的紧性证明在极大真滤子的构造处使用 Zorn 引理。
\end{solution}

\begin{exercise}\label{b1:ex:22-6}
令 \(\mathbb Z_2=\varprojlim_r\mathbb Z/2^r\mathbb Z\)。证明乘以 \(2\) 的映射 \(\mathbb Z_2\rightarrow\mathbb Z_2\) 单射，其像恰为模 \(2\) 投影的核，因而是开且闭的指数 \(2\) 子群。
\end{exercise}
\begin{solution}
若 \(2a=0\)，则在模 \(2^{r+1}\) 的坐标上，\(2a_{r+1}=0\)，故 \(a_{r+1}\) 被 \(2^r\) 整除。约化到第 \(r\) 层便得 \(a_r=0\)。对所有 \(r\) 都成立，所以 \(a=0\)。像当然落在模 \(2\) 的核中。

反过来，若 \(a_1=0\)，则每个 \(a_{r+1}\) 有偶数代表。将这个代表除以 \(2\) 再模 \(2^r\)，记为 \(b_r\)；换代表改变的是 \(2^r\) 的倍数，所以良定义。相容性保证 \((b_r)_r\in\mathbb Z_2\)，并有 \(2b=a\)。于是像就是模 \(2\) 的核。该投影满射到离散的二元群，其核开且闭，指数为 \(2\)。
\end{solution}

\begin{exercise}\label{b1:ex:22-7}
设 \(L/K\) 为代数 Galois 扩张，\(G=\operatorname{Gal}(L/K)\) 取 Krull 拓扑。对有限 Galois 中间域 \(N/K\)，记 \(G_N=\operatorname{Gal}(N/K)\)，并记限制映射为 \(\pi_N:G\to G_N\)。证明每个既开又闭的子集 \(C\subseteq G\) 都可写成 \(C=\pi_N^{-1}(D)\)，其中 \(N/K\) 为某个有限 Galois 中间域，\(D\subseteq G_N\)。这里不要求 \(C\) 是子群。
\end{exercise}
\begin{solution}
\(C\) 闭，所以由\cref{b1:thm:krull-compact}及闭子空间的紧性，\(C\) 紧。\(C\) 开，故其每个点有一个包含在 \(C\) 中的基本开陪集。这些陪集覆盖 \(C\)，紧性给出有限子覆盖 \(\sigma_jU_{N_j}\)。取同时包含全部 \(N_j\) 的有限 Galois 层 \(N\)。每个 \(U_{N_j}\) 包含 \(U_N\)，所以每个所选陪集以及它们的并，都是若干个 \(U_N\) 陪集的并。令 \(D=\pi_N(C)\)，便得等式。空集可直接取 \(D=\varnothing\)。
\end{solution}

\begin{exercise}\label{b1:ex:22-8}
设 \(L/K\) 为代数 Galois 扩张，\(G=\operatorname{Gal}(L/K)\) 取 Krull 拓扑，\(D\) 为离散群，不要求有限。证明每个连续同态 \(\rho:G\rightarrow D\) 的像有限，并且 \(\rho\) 经过某个有限 Galois 商 \(\operatorname{Gal}(N/K)\)。
\end{exercise}
\begin{solution}
由\cref{b1:thm:krull-compact}，\(G\) 紧，连续像 \(\rho(G)\) 也紧。离散空间的全部单点构成开覆盖，若其为无限集就没有有限子覆盖，所以 \(\rho(G)\) 必有限。将值域缩到这个有限离散群，\cref{b1:prop:finite-continuous-hom}给出一个有限 Galois 层 \(N\)，使 \(\rho\) 经过 \(G_N\)。还可直接从 \(\ker\rho=\rho^{-1}(1)\) 开并包含某个 \(U_N\) 得到同一结论。
\end{solution}

\begin{exercise}\label{b1:ex:22-9}
设 \(L/K\) 为代数 Galois 扩张，\(G=\operatorname{Gal}(L/K)\) 取 Krull 拓扑。若 \(H\leq G\) 为闭且指数有限的子群，证明 \(H\) 开。再取积拓扑群 \(V=\prod_{n\geq1}\mathbb F_2\) 与有限支集子群 \(W=\bigoplus_{n\geq1}\mathbb F_2\)，构造一个包含 \(W\) 的指数 \(2\) 真子群，并证明它不开。
\end{exercise}
\begin{solution}
左右平移为同胚，所以每个 \(H\) 陪集都闭。有限指数使 \(G\setminus H\) 只由有限多个这样的闭陪集组成，因此补集闭，\(H\) 开。对应于中间域的子群由\cref{b1:thm:infinite-galois-correspondence}保证闭，才可将此结论用于域论。

令 \(v=(1,1,\ldots)\in V\)，则 \(v\notin W\)。在 \(W\oplus\mathbb F_2v\) 上定义 \(\lambda_0(w+cv)=c\)，再把该子空间的一组基扩充为 \(V\) 的基，并规定 \(\lambda\) 在新增基向量上取零。这一步使用向量空间基的扩充，所得映射依赖所选的代数基，并不是由某个有限坐标规则给出的。它给出线性映射 \(\lambda:V\to\mathbb F_2\)。子群 \(H'=\ker\lambda\) 包含 \(W\)，不包含 \(v\)，且指数为 \(2\)。任意有限个坐标上的取值都可由 \(W\) 中的元素实现，故 \(W\) 稠密，\(H'\) 也稠密。若 \(H'\) 开，它的其他陪集也开，因而 \(H'\) 同时闭，与它稠密而非整个 \(V\) 矛盾。所以有限指数本身不能保证子群开。也可直接看出 \(\lambda\) 不连续：若连续，\(\lambda^{-1}(0)\) 应闭，包含稠密的 \(W\) 就必须等于 \(V\)，这与 \(\lambda(v)=1\) 矛盾。
\end{solution}

\begin{exercise}\label{b1:ex:22-10}
设 \(L/K\) 为代数 Galois 扩张，\(G=\operatorname{Gal}(L/K)\) 取 Krull 拓扑，\(K\subseteq E\subseteq L\) 且 \([E:K]<\infty\)。令 \(H=\operatorname{Gal}(L/E)\)。证明 \(C=\bigcap_{g\in G}gHg^{-1}\) 是开正规子群，其不动域是 \(E/K\) 在 \(L\) 内的正规闭包。
\end{exercise}
\begin{solution}
由\cref{b1:prop:finite-open-correspondence}，\(H\) 开且指数有限。群在有限集合 \(G/H\) 上作左乘，其核正是 \(C\)。这就是\cref{b1:prop:coset-action-core}中的正规核：它是包含在 \(H\) 中的最大 \(G\)-正规子群，因为若 \(J\trianglelefteq G\) 且 \(J\leq H\)，则 \(J=gJg^{-1}\leq gHg^{-1}\) 对所有 \(g\) 成立，故 \(J\leq C\)。只需选有限个陪集代表即可求这个交：若 \(g'=gh\)、\(h\in H\)，则 \(g'Hg'^{-1}=gHg^{-1}\)。因此 \(C\) 是有限个开子群的交，开且正规。

对任意 \(g\)，一个元素被 \(gHg^{-1}\) 固定当且仅当它属于 \(g(E)\)：把固定条件用 \(g^{-1}\) 共轭，应用\cref{b1:thm:infinite-galois-correspondence}即可。因此 \(C\) 恰为固定全部 \(g(E)\) 的自同构，\(L^C\) 就是这些域的合成域。它有限且在 \(K\) 上正规，因为 \(G\) 置换全部共轭域；任何包含 \(E\) 的正规中间扩张都包含这些共轭，所以这恰为最小的正规扩张，即正规闭包。子群一侧的“包含在 \(H\) 中且最大”，经包含反向的对应，成为域一侧的“包含 \(E\) 且最小”。
\end{solution}

\begin{exercise}\label{b1:ex:22-11}
令 \(G_{\mathbb F_q}=\operatorname{Gal}(\overline{\mathbb F}_q/\mathbb F_q)\cong\widehat{\vphantom{Z}\mathbb Z}\)，\(F(x)=x^q\) 为 Frobenius 自同构。取正整数 \(a,b\)，求 \(\langle F^a,F^b\rangle\) 的闭包，并证明
\[
a\widehat{\vphantom{Z}\mathbb Z}\cap b\widehat{\vphantom{Z}\mathbb Z}
=\operatorname{lcm}(a,b)\widehat{\vphantom{Z}\mathbb Z}.
\]
写出对应不动域的交与合成域。
\end{exercise}
\begin{solution}
在抽象循环子群 \(\langle F\rangle\) 中，Bézout 等式给出 \(\langle F^a,F^b\rangle=\langle F^d\rangle\)，其中 \(d=\gcd(a,b)\)。由\eqref{b1:eq:profinite-multiple-kernel}及其后的闭包计算，闭包为 \(d\widehat{\vphantom{Z}\mathbb Z}\)。

令 \(m=\operatorname{lcm}(a,b)\)。同一公式把左侧描述为模 \(a\)、模 \(b\) 同时为零的相容族。看其模 \(m\) 的坐标，一个整数剩余类同时被 \(a,b\) 整除当且仅当被 \(m\) 整除，所以左侧为模 \(m\) 的核，等于右侧。固定域反向对应遂给出
\[
\mathbb F_{q^a}\cap\mathbb F_{q^b}=\mathbb F_{q^d},
\qquad\mathbb F_{q^a}\mathbb F_{q^b}=\mathbb F_{q^m}.
\]
第一式对应生成子群的闭包，第二式对应闭子群的交。
\end{solution}

\begin{exercise}\label{b1:ex:22-12}
固定素数 \(\ell\)，令 \(E=\bigcup_{r\geq0}\mathbb F_{q^{\ell^r}}\subseteq\overline{\mathbb F}_q\)。证明 \(E/\mathbb F_q\) 是 Galois 扩张且 \(\operatorname{Gal}(E/\mathbb F_q)\cong\mathbb Z_\ell\)。在自然识别 \(\operatorname{Gal}(\overline{\mathbb F}_q/\mathbb F_q)\cong\widehat{\vphantom{Z}\mathbb Z}\) 下，求到 \(\operatorname{Gal}(E/\mathbb F_q)\) 的限制映射之核，并判断该核是否为零、是否开。
\end{exercise}
\begin{solution}
任一元素位于有限 Galois 层，其全部共轭都在该层中，故并域正规且可分。\cref{b1:prop:galois-compatible-family}把其 Galois 群识别为 \(\varprojlim_r\mathbb Z/\ell^r\mathbb Z\)，且有限坐标的基本开集一致，故为拓扑同构。由\cref{b1:thm:infinite-galois-quotient}，整个代数闭包上的限制满射。核为
\[
H=\bigcap_{r\geq1}\ell^r\widehat{\vphantom{Z}\mathbb Z}.
\]
每一项闭，所以 \(H\) 闭。这个交并不是零子群。对每个正整数 \(n\)，写 \(n=\ell^s m\)、\(\gcd(m,\ell)=1\)，用中国剩余定理选取唯一的剩余类
\[
e_n\equiv0\pmod{\ell^s},\qquad e_n\equiv1\pmod m.
\]
若 \(n\mid n'\)，则 \(e_{n'}\) 约化后仍满足这两个同余式，故 \((e_n)_n\) 是相容族。它在每个模 \(\ell^r\) 坐标为零，因而属于 \(H\)；但取任一不同于 \(\ell\) 的素数 \(p\)，其模 \(p\) 坐标为一，所以它非零。若 \(H\) 开，由\cref{b1:prop:finite-open-correspondence}，其固定域 \(E\) 应在 \(\mathbb F_q\) 上有限，与层次数 \(\ell^r\) 无界矛盾。因此该核不开。
\end{solution}

\begin{exercise}\label{b1:ex:22-13}
设 \(\ell\) 为素数，\(G_{\mathbb R}=\operatorname{Gal}(\mathbb C/\mathbb R)\) 取 Krull 拓扑。按基变换同构分类全部连续同态 \(G_{\mathbb R}\to\operatorname{GL}_2(\mathbb F_\ell)\)，其中目标群取离散拓扑。
\end{exercise}
\begin{solution}
由\cref{b1:prop:real-absolute-galois}，源群是离散的 \(C_2\)，所以每个表示都连续，由复共轭的像 \(A\) 决定，且 \(A^2=I\)。若 \(\ell\neq2\)，多项式 \(X^2-1\) 有不同根 \(1,-1\)，故空间分解为两个特征空间之和。按负特征空间的维数，恰有三个同构类，代表为 \(I\)、\(\operatorname{diag}(1,-1)\)、\(-I\)。

若 \(\ell=2\)，写 \(A=I+N\)，则 \(N^2=0\)。\(N=0\) 给出平凡表示；若 \(N\neq0\)，取 \(v\) 使 \(Nv\neq0\)，则 \(Nv,v\) 线性无关，否则 \(v\) 是 \(Nv\) 的倍数将推出 \(Nv=0\)。在这组基下 \(N\) 将第二个基向量 \(v\) 送到第一个基向量 \(Nv\)，并将 \(Nv\) 送到零，故 \(A\) 的代表为 \(\begin{psmallmatrix}1&1\\0&1\end{psmallmatrix}\)。因此特征 \(2\) 时恰有两个同构类。
\end{solution}

\begin{exercise}\label{b1:ex:22-14}
设 \(m,n\) 为正整数，\(\widehat{\vphantom{Z}\mathbb Z}\) 取逆极限拓扑，\(\mathbb Z/m\mathbb Z\) 取离散拓扑。求经过模 \(n\) 投影的连续同态 \(\widehat{\vphantom{Z}\mathbb Z}\to\mathbb Z/m\mathbb Z\) 的数目，其中有多少个是满射？再设 \(\ell\) 为素数，\(\rho:G_{\mathbb F_q}\to\operatorname{GL}_2(\mathbb F_\ell)\) 是把 Frobenius 自同构 \(F(x)=x^q\) 送到 \(A=\begin{psmallmatrix}1&1\\0&1\end{psmallmatrix}\) 的连续表示。对正整数 \(s\)，求其在 \(G_{\mathbb F_{q^s}}\) 上的限制，并求使该限制平凡的最小 \(s\)。所有绝对 Galois 群均在同一个 \(\overline{\mathbb F}_q\) 内取定。
\end{exercise}
\begin{solution}
由\cref{b1:prop:profinite-integer-finite-hom}，这样的连续同态由 \(1\) 的像 \(c\in\mathbb Z/m\mathbb Z\) 唯一决定。它经过模 \(n\) 投影，当且仅当 \(nc=0\)：这是 \(\mathbb Z/n\mathbb Z\) 的关系在目标中必须成立的条件。若 \(g=\gcd(m,n)\)，写 \(m=gm'\)、\(n=gn'\)，则 \(\gcd(m',n')=1\)。于是 \(m\mid nc\) 等价于 \(m'\mid n'c\)，再由互素性等价于 \(m'\mid c\)。模 \(m\) 的这些类为 \(0,m',\ldots,(g-1)m'\)，共有 \(g\) 个。满射还要求 \(c\) 的阶为 \(m\)，故当 \(m\mid n\) 时有 \(\varphi(m)\) 个满射，否则没有。

扩大有限基域以后，\(G_{\mathbb F_{q^s}}\) 是 \(G_{\mathbb F_q}\) 中的闭子群 \(s\widehat{\vphantom{Z}\mathbb Z}\)，它的拓扑生成元是 \(F^s\)。限制表示把这一生成元送到
\[
A^s=\begin{pmatrix}1&\bar s\\0&1\end{pmatrix},\qquad \bar s\in\mathbb F_\ell.
\]
连续性使这个取值唯一决定限制表示；具体地，它仍经过模 \(\ell\) 坐标，按右上角相加计算。若 \(\ell\mid s\)，则限制平凡；否则 \(A^s\) 的阶为 \(\ell\)，限制的像仍是由 \(A\) 生成的 \(\ell\) 阶群。因此使限制平凡的最小正整数是 \(s=\ell\)，相应基域为 \(\mathbb F_{q^\ell}\)。
\end{solution}

\begin{exercise}\label{b1:ex:22-15}
固定素数 \(\ell\)。对每个 \(r\geq1\)，记 \(a_{\ell^r}\) 为 \(a\in\widehat{\vphantom{Z}\mathbb Z}\) 的模 \(\ell^r\) 坐标，并定义
\[
\rho_r:\widehat{\vphantom{Z}\mathbb Z}\longrightarrow
\operatorname{GL}_2(\mathbb Z/\ell^r\mathbb Z),\qquad
a\longmapsto\begin{pmatrix}1&a_{\ell^r}\\0&1\end{pmatrix}.
\]
证明这些表示相容且连续，并求它们给出的 \(\mathbb Z_\ell\) 值矩阵表示的像与核。该表示是否经过单个有限商？
\end{exercise}
\begin{solution}
两个所列矩阵相乘时，右上角相加，其余位置不变，故 \(\rho_r\) 为同态；它经过有限坐标投影，所以连续。约化矩阵正好约化右上角，因而各层相容。根据正文在\subsecref{b1:subsec:2205:03}给出的逆极限构造，得到连续表示 \(\rho\)，其像为
\[
\Biggl\{\,\begin{pmatrix}1&b\\0&1\end{pmatrix}\ \Bigg|\ b\in\mathbb Z_\ell\,\Biggr\}.
\]
自然投影 \(\widehat{\vphantom{Z}\mathbb Z}\rightarrow\mathbb Z_\ell\) 在每一有限层都满射，因而由有限离散积的紧性及有限交性质可知它满射。故 \(\rho\) 的核为 \(\bigcap_r\ell^r\widehat{\vphantom{Z}\mathbb Z}\)。像是无限群，例如所有整数给出不同矩阵，因此不能经过单个有限商。每个有限坐标都经过有限层，并不意味着整族共有一个有限决定层。
\end{solution}

\begin{exercise}\label{b1:ex:22-16}
给 \(\overline{\mathbb F}_q\) 赋予离散拓扑，给 \(G_{\mathbb F_q}=\operatorname{Gal}(\overline{\mathbb F}_q/\mathbb F_q)\) 赋予 Krull 拓扑。证明 \(G_{\mathbb F_q}\) 对 \(\overline{\mathbb F}_q\) 的自然作用连续，每个元素的稳定子群都开，但整个作用的核不开。
\end{exercise}
\begin{solution}
若元素 \(a\) 在 \(\mathbb F_q\) 上次数为 \(d\)，则 \(\mathbb F_q(a)=\mathbb F_{q^d}\)，其稳定子群是模 \(d\) 投影的核 \(d\widehat{\vphantom{Z}\mathbb Z}\)，由\eqref{b1:eq:profinite-multiple-kernel}为开子群。对任意 \((g,a)\)，集合 \(g\operatorname{Stab}(a)\times\{a\}\) 是一个开邻域，整个邻域上的作用值均为 \(g(a)\)，因而作用连续。

整体核必须固定全部代数元，所以只有恒等自同构。若它开，则由\cref{b1:prop:finite-open-correspondence}，其固定域 \(\overline{\mathbb F}_q\) 应是有限扩张，矛盾。这里对每个元素分别找到开稳定子群，不能换成“存在一个开子群同时固定所有元素”。
\end{solution}

\begin{exercise}\label{b1:ex:22-17}
设 \(L/K\) 为代数 Galois 扩张，\(G=\operatorname{Gal}(L/K)\) 取 Krull 拓扑，\(H,J\leq G\) 为闭子群。若 \(N/K\) 是 \(L/K\) 的有限 Galois 中间扩张，\(\pi_N:G\to\operatorname{Gal}(N/K)\) 为限制映射，是否必有 \(\pi_N(H\cap J)=\pi_N(H)\cap\pi_N(J)\)？用一个有限 Galois 扩张给出反例。
\end{exercise}
\begin{solution}
取 \(L=\mathbb Q(\sqrt2,\sqrt3)\)。\cref{b1:prop:biquadratic-example}给出 \([L:\mathbb Q]=4\)；它又是可分多项式 \((X^2-2)(X^2-3)\) 的分裂域，所以\cref{b1:prop:galois-automorphism-count}给出 \(|G|=4\)。每个自同构由两个平方根的像决定，像分别只能为原平方根或其负数，故取符号给出单射 \(G\rightarrow C_2^2\)。两边同有四个元素，因而是同构。以这两个符号为坐标，令 \(H=\langle(1,1)\rangle\)、\(J=\langle(1,0)\rangle\)，并令 \(\pi\) 为限制到 \(\mathbb Q(\sqrt2)\)，即取第一坐标。

这时 \(H\cap J=0\)，故 \(\pi(H\cap J)=0\)；但两个子群各自都满射到第一个 \(C_2\)，所以 \(\pi(H)\cap\pi(J)=C_2\)。所有子群在有限离散群中都闭，反例仍成立。\cref{b1:prop:closure-finite-quotients}使用的是全部有限层的相容信息，并未断言单个投影保持子群的交。
\end{solution}

\begin{exercise}\label{b1:ex:22-18}
将 \(\mathbb Z\) 视为它在 \(\widehat{\vphantom{Z}\mathbb Z}\) 中的自然像，令 \(Q=\widehat{\vphantom{Z}\mathbb Z}/\mathbb Z\) 带商拓扑。证明任何从 \(Q\) 到 Hausdorff 空间的连续映射都是常值映射，并解释为何这不说明抽象群 \(Q\) 平凡。
\end{exercise}
\begin{solution}
先直接核对商拓扑。非空开集的原像 \(V\subseteq\widehat{\vphantom{Z}\mathbb Z}\) 是非空开集且满足 \(V+\mathbb Z=V\)。给任意 \(g\)，开集 \(g-V\) 与稠密的 \(\mathbb Z\) 相交，故 \(g\in V+\mathbb Z=V\)。所以商中只有空集与全空间两个开集。

若连续映射取两个不同值，Hausdorff 性给出分别包含这两个值的不相交开邻域，其原像是两个非空不相交开集，与上述描述矛盾，故映射常值。正文在\subsecref{b1:subsec:2203:02}构造的相容族，在所有模 \(2^r\) 坐标为零、在模 \(3\) 坐标为一，不能来自整数。因此 \(\mathbb Z\) 是真子群，抽象商非平凡。失去的是用 Hausdorff 拓扑区分这些陪集的能力。
\end{solution}

\begin{exercise}\label{b1:ex:22-19}
设 \(L/K\) 为代数 Galois 扩张，\(G=\operatorname{Gal}(L/K)\) 取 Krull 拓扑，\(H\) 为 \(G\) 的闭正规子群。对 \(L/K\) 的每个有限 Galois 中间域 \(N/K\)，记 \(G_N=\operatorname{Gal}(N/K)\)，并记限制映射为 \(\pi_N:G\to G_N\)。证明有限群 \(G_N/\pi_N(H)\) 构成逆系统，并证明自然映射
\[
G/H\longrightarrow\varprojlim_N G_N/\pi_N(H)
\]
是拓扑群同构。
\end{exercise}
\begin{solution}
\(\pi_N(H)\) 正规；限制映射将 \(\pi_M(H)\) 满射到 \(\pi_N(H)\)，所以诱导的商映射良定义且相容。先定义连续同态
\[
\Phi:G\longrightarrow\varprojlim_NG_N/\pi_N(H),
\qquad g\longmapsto\bigl(\pi_N(g)\pi_N(H)\bigr)_N.
\]
各坐标相容，故 \(\Phi\) 确实取值于逆极限；每个坐标连续，故 \(\Phi\) 连续。由\cref{b1:prop:closure-finite-quotients}，其核为
\[
\ker\Phi=\bigcap_N\pi_N^{-1}\bigl(\pi_N(H)\bigr)=\overline H=H.
\]
因此 \(\Phi\) 诱导单同态 \(\overline\Phi:G/H\to\varprojlim_NG_N/\pi_N(H)\)。由商拓扑的定义，\(\overline\Phi\) 连续。

给定右侧的相容陪集族，对每个 \(N\) 令 \(C_N\subseteq G\) 为这一陪集在 \(G\) 中的原像。限制满射使 \(C_N\) 非空，它又开且闭。任取有限个 \(N\)，选共同上层 \(M\)；相容性保证 \(C_M\) 包含于它们的交。所以 \((C_N)_N\) 有有限交性质，由\cref{b1:thm:krull-compact}总交非空。这证明满射。

\(G\) 紧，故其商空间 \(G/H\) 也紧；目标是有限离散群直积的子空间，具有 Hausdorff 性。\cref{b1:prop:compact-maps}说明连续双射 \(\overline\Phi\) 为同胚，得到所需拓扑同构。闭性在核由 \(\overline H\) 变为 \(H\) 的步骤不可省略。
\end{solution}

\begin{exercise}\label{b1:ex:22-20}
设 \(T\) 是有限零一字组成的集合，对取前缀封闭，且每个长度都有 \(T\) 中的字。证明存在无限零一序列，其每个有限前缀都属于 \(T\)。先用有限离散积的紧性证明，再给出按次序选 \(0\) 或 \(1\) 的直接构造，说明此特殊情形无需为每一步作无规则的选择。
\end{exercise}
\begin{solution}
在 \(\{0,1\}^{\mathbb N}\) 中，令 \(C_r\) 为前 \(r\) 位组成的字属于 \(T\) 的序列。它是有限多个基本开闭集的并，故闭且非空；非空性可将任一长度为 \(r\) 的字后补零。前缀封闭使 \(C_{r+1}\subseteq C_r\)，所以具有有限交性质。\cref{b1:lem:finite-discrete-product-compact}给出非空总交。

直接构造时，如果一个字在 \(T\) 中具有任意长的延长字，就称它可以无限延长。空字有此性质。若一个字有此性质，它的两个一步延长至少一个仍有此性质；否则两个分支的长度各有上界，取较大上界便产生矛盾。规定能选 \(0\) 分支时就选 \(0\)，否则选 \(1\)，递归产生所需无限序列。这个规则在每一步唯一确定所选的分支。
\end{solution}

